% Paper-series v2.381: complete algebraic exceptional torsion atlas and all-cocycle twist-stability; original claims retained.
% Paper-series v2.380: unified bibliography, exact exceptional double/length/gauge mathematics in I-B; previous theorem labels retained.
% Paper-series v2.379: full equivariant Dirac/Hodge/eta/torsion atlas in Paper I-B; paper-root references synchronized.
% Paper-series v2.378: 2026 exceptional literature consolidation and reference synchronization. The mathematical additions are located in Paper I-B.
% Paper-series v2.377: reference-only rebind to current Paper I-B v1.1; original mathematics retained except new Paper I-B comparison chapter.
% Framework v2.370: derived exceptional-root/Hurwitz comparison in I-A/I-C; other papers reference-only synchronized.
% Framework v2.368: preserve prior II-B convention/unital fixes and verified II-D corrections; current references synchronized.
% Framework v2.358: current 30-root source reference synchronization; selectively updated III-E, III-F, VI-B, VI-C.
% Framework v2.356: bibliography/source-registry synchronization; mathematical body unchanged.
% Framework v2.349: integrated exact conditional mathematical closures; G0 and GP1--GP3 retained as axioms.
% Framework v2.347: exact operational test-resolution tiers, one-body response iff criterion, terminating Majorana hierarchy, odd-transfer witness, cutoff cumulants, basis-free regular-CCR covariance. Prior domains retained.
% Framework v2.320: signed Petersen quantum sectors, mirror-transfer CP no-go, conditional interference and source-visible response; exact finite certificates. Preserves inherited original-source theorem bodies.
% Framework v2.275: exact positive dense-twin word cubature, finite repeatability and eighth-moment classification, nested sampling frames, joint action and POVM-jet exports; preserves v2.274.
% Framework v2.274: two-scalar positive convolution certificates, first-leaked exceptional coset detection, branch-symmetric twin mixing, and finite endpoint moments; all v2.273 branches preserved.
% Framework v2.273: restored missing Problem 3 theorems from v2.268; exact exceptional frames, twin geometry, convolution repeatability and projector/occurrence firewalls; merged on v2.272 without overwriting independent branches.
% Framework v2.272: merged v2.270 covariant native-readout mathematics with v2.271 Problem 4 finite-action Hopf hypercharge/holonomy corrections; no universal bulk-gauge uniqueness claim.
% Independent branches: v2.270 readout/guard exports and v2.270 K16 action-consumed charge closure; all prior scoped assumptions retained.
% Framework v2.267: exact Paper II-A converse/observable certificates, III-F action tests, V-B right orbit and V-C native gauge scope; prior v2.266 preserved.
% Framework v2.266: exact Paper II-C outcome-count/readout converse results and scoped downstream imports; 30-root references synchronized.
% Framework v2.265: II-C coordinate-stage constitution and exact first-symbol host inheritance; prior mathematical domains retained.
% Framework v2.264: II-B reverse and action-bridge integration; preserves prior mathematical domains; active 30-root references synchronized.
% Framework v2.263: Paper I-D adds bridge-controlled root-carry parity and source-relative old-cap cancellation; other root mathematical bodies retained; active 30-root references synchronized.
% Framework v2.262: merges Paper II-F three-gap closure from v2.260 into v2.261, preserving newer II-B/II-D repairs; 30-root references synchronized.
% Framework v2.261: merged Paper II-B GP2-level/geometric repairs; preserves v2.260 II-D repairs; active references synchronized.
% Framework v2.260: Paper II-D signed RS, finite-symbol typing, and conditional curvature/Jordan repairs; references synchronized.
% Framework v2.259: conditional native SU(2) atlas and explicit later discharge; v2.258 mathematical bodies retained elsewhere.
% Framework v2.258: preserves parallel II-A/dynamics update and adds complete readout/frequency closures.
% Framework v2.257: exact minimal complete position readout and arbitrary-preparation frequency law.
% Framework v2.256: exact native inertia and sharp bounded ordered-source closure; preserves concurrent v2.255 work.
% Framework v2.255: preserves parallel detector integration and adds exact finite readout closures.
% Framework v2.254: shared bibliography/reference sync; substantive Pasterski detector-Ward additions only IV-B/IV-C/VII-C.
% Framework v2.253: typed raw coordinate defect in I-A and synchronized VI-A import; active references synchronized.
% Framework v2.252: optional shell-frame source-return calculation in III-G; active references synchronized.
% Framework v2.251: forward common-action execution, joint source jets and scoped coupled control.
% Framework v2.251: exact full-heat variations, fermion source jets and explicit physical completion frontier.
% Framework v2.251: certified interacting source jets, pole stability and scoped completion proofs.
% Framework v2.251: interacting finite-source evaluation and coupled local vacuum classification.
% Framework v2.251: exact finite source execution and native causal common-parent bridge.
% Framework v2.251: constructive parent and scoped source-incidence audit; see release report.
% Framework v2.251: common-action two-jet, corrected source/response closure and synchronized references.
% Framework v2.244: scoped rooted source-order Schur diagnostics and synchronized paper references.
% Framework v2.242: signed rooted phase lift, complete transverse return and exact family cubic.
% Framework v2.241: proved rooted conormal geometry, conditional cap realization and scoped action tests.
% Framework v2.239: coherent 30-paper reference rebinding; selective new proofs on established owners.
% Framework v2.238: Paper IV restructured; inherited mathematics retained.
\documentclass[12pt,a4paper]{article}
\usepackage{mathtools}
\usepackage{float}
\usepackage{caption}
\usepackage{tikz}
\usepackage{tikz-cd}
\usetikzlibrary{arrows,positioning,calc,decorations.pathmorphing}
\usepackage{microtype}
\usepackage{xltabular}
% Paper-series v2.381: new I-B torsion/modular results, bibliography and current source roots.
% Paper-series v2.380: new exceptional bibliography and synchronized paper references.
% Paper-series v2.379 current 31-paper macro/aux registry.
% Paper-series v2.378: current 31-paper AUX/bibliography registry updated for Paper I-B v1.2 literature additions.
% Framework v2.371: current-source reference registry refreshed; old theorem labels retained.
% Framework v2.370: current 30-paper source/AUX root rebinding; native mathematical macros unchanged.
% Framework v2.368: synchronized merged 30-root registry; mathematical notation unchanged.
% Framework v2.358: active 30-paper source registry synchronized; existing mathematical macros unchanged.
% Framework v2.357: selective new source-owned certificates; active reference roots updated for five revised papers.
% Framework v2.356: active 30-root registry synchronized to the 30 papers in this release.
% Framework v2.275: exact positive dense-twin word cubature, finite repeatability and eighth-moment classification, nested sampling frames, joint action and POVM-jet exports; preserves v2.274.
% Framework v2.274: filename and root-source synchronization.
% Framework v2.273: restored missing Problem 3 theorems from v2.268; exact exceptional frames, twin geometry, convolution repeatability and projector/occurrence firewalls; merged on v2.272 without overwriting independent branches.
% Framework v2.272: merged v2.270 covariant native-readout mathematics with v2.271 Problem 4 finite-action Hopf hypercharge/holonomy corrections; no universal bulk-gauge uniqueness claim.
% Independent branches: v2.270 readout/guard exports and v2.270 K16 action-consumed charge closure; all prior scoped assumptions retained.
% Framework v2.269: synchronized 30 active AUX roots; mathematics and bibliography preserved from Macros-v2_161.
% Framework v2.266: Paper II-C converse-derived outcome resources; II-E, III-C, IV-A/B, V-A and VI-B downstream mathematical closures; 30-root registry synchronized.
% Framework v2.265: II-C coordinate-stage constitution and exact first-symbol host inheritance; prior mathematical domains retained.
% Framework v2.264: paper II-B reverse/audit and IV-A observable-minimal exchange; synchronized external references to 30 active roots.
% Framework v2.263: Paper I-C adds bridge-controlled root-carry parity and source-relative old-cap cancellation; other root mathematical bodies retained; active 30-root references synchronized.
% Framework v2.262: integrates Paper II-F radius, Casimir and readout results; active 30-root references synchronized.
% Framework v2.261: merged Paper II-B GP2-level/geometric repairs; preserves v2.260 II-D repairs; active references synchronized.
% Framework v2.260: Paper II-D signed RS, finite-symbol typing, and conditional curvature/Jordan repairs; references synchronized.
% Framework v2.259: conditional native SU(2) atlas and explicit later discharge; v2.258 mathematical bodies retained elsewhere.
% Framework v2.258: preserves parallel II-A/dynamics update and adds complete readout/frequency closures.
% Framework v2.257: exact minimal complete position readout and arbitrary-preparation frequency law.
% Framework v2.256: exact native inertia and sharp bounded ordered-source closure; preserves concurrent v2.255 work.
% Framework v2.255: preserves parallel detector integration and adds exact finite readout closures.
% Framework v2.254: Pasterski bibliography and synchronized 30-paper external-reference registry.
% Framework v2.253: typed raw coordinate defect in I-A and synchronized VI-A import; active references synchronized.
% Framework v2.252: synchronized source registry; optional III-G shell-frame calculation.
% Framework v2.251: forward common-action execution, joint source jets and scoped coupled control.
% Framework v2.251: exact full-heat variations, fermion source jets and explicit physical completion frontier.
% Framework v2.251: certified interacting source jets, pole stability and scoped completion proofs.
% Framework v2.251: synchronized active source registry; mathematical notation unchanged.
% Framework v2.251: synchronized active source registry; mathematical notation unchanged.
% Framework v2.251: synchronized active source registry; mathematical notation unchanged.
% Framework v2.251: synchronized source registry; no mathematical notation changes.
% Framework v2.244: synchronized 30-paper AUX registry and rooted phase-source action-order tests.
% Framework v2.242: signed rooted phase lift, complete transverse return and exact family cubic.
% Framework v2.241: proved rooted conormal geometry, conditional cap realization and scoped action tests.
% Framework v2.239: all 30 active source roots synchronized; no mathematical macro changes.
% Framework v2.238: Paper IV A/B/C split and 30-root reference registry.
% Framework v2.237: Paper VII-D technical split and 28-root reference registry.
% Framework v2.236: Paper VI-C technical-volume reference added; all current root bindings advanced.
% Framework v2.235: 26-root AUX and bibliography bindings synchronized; mathematical macros unchanged.
% BEGIN INLINE SOURCE: Paper00-TeXBuildCompatibility-v1_0.tex
% Rendering-only repairs for inherited long appendices and bitmap text fonts.
% Integer counters through 26 retain the original labels; later values use AA,
% AB, ... . This file changes no action coefficient, state or theorem premise.
\RequirePackage{alphalph}
\makeatletter
\let\BHPrimitiveAlph\@Alph
\newalphalph{\BHExtendedAlph}{\BHPrimitiveAlph}{26}
\let\@Alph\BHExtendedAlph
\AtBeginDocument{%
  \@ifpackageloaded{microtype}{\microtypesetup{expansion=false}}{}%
}
\makeatother

% END INLINE SOURCE: Paper00-TeXBuildCompatibility-v1_0.tex

% Framework v2.164: current reference/bibliography profile; mathematical definitions retained.
% Framework v2.162: finite metric, functional Hilbert and scoped Ward integration.
% Framework v2.161: inherits Macros-v2_57.tex exactly except the current external-reference binding.
% Macros v2.57 / framework v2.156: current reference registry only; mathematical definitions unchanged.
% Framework v2.143 / Macros v2.55: current external-reference bindings; mathematical notation unchanged.
% Framework v2.110 / Macros v2.54: current external-reference registry and bibliography; mathematical notation inherited.
% Framework v2.67 / Macros-v2_53: adds Born-action-cell and joint-EH readout symbols; absolute unit realization is metrological; MATCH_{J/scal}^{EH} remains the sole G24 normalization bridge.
% Framework v2.35 / Macros-v2_52: grouped paper numbering, active AUX registry, and Born-v6_20 bibliography binding synchronized; notation unchanged.
% Framework v2.26: active-source AUX and bibliography bindings synchronized; notation unchanged.
% Framework v2.09: reviewed early-reciprocity integration; scoped physical imports.
% Framework v2.08: Born--Euler mathematics and conditional complete-action potential criterion.
% Framework v2.07: active AUX registry includes Paper I-B v1.2; shared bibliography Born-v6_18.
% Macros-v2_46 / framework v2.04: mathematical atlas and downstream realization correction.
% Macros-v2_46 / framework v2.03: scoped structure integration and current dependency bindings.
% Macros-v2_46 / framework v2.01: shared source bindings; prior revision comments are historical.
% v2.42 / framework v1.99: Synchronized RG/AS owner and bibliography bindings.
% v2.41 / framework v1.99: active manuscript AUX bindings and geometric-closure bibliography synchronized.
% Release-v1.80 centralized source synchronization for scalar-circle/comparison-band integration.
% Release v1.79: Hurwitz hierarchy and kinetic-normalized action response integration; physical mass matching remains conditional.
% Release v1.78: approved M8 closure-batch integration; physical and analytic domains retained.
% Current v1.78 registry: all active TeX roots are synchronized below.
% Earlier release comments are historical scope records, not current status.
% Release-v1.77 synchronization: native same-owner Noether closure, generator--flow/history typing, analytic-certification firewalls, scalaron/K16, LF/QCD, quantum-Ward, and RAR boundary-owner transport are registered; shared mathematical notation is unchanged.
% Release-v1.76 synchronization: registry points to the v1.76 active paper roots.
% Release-v1.76 synchronization: registers the local A+Phi+g Noether/constraint closure and dynamics-frontier propagation; shared mathematical notation is unchanged.
% Release-v1.75 synchronization: registry points to the v1.75 active paper roots.
% Release-v1.75 synchronization: integrates the joint-Dirac Einstein bulk real form and the local odd-scalar heat-coefficient action on explicit analytic first-jet/interior domains; v1.74 Yang--Mills closure is retained. Fermionic reality, pure-metric curvature-squared/exact-nonlocal terms, noncanonical boundaries, physical vacuum selection, and actual solution-history tokens remain open.
% Release-v1.74 synchronization: registry points to the v1.74 active paper roots and retains the v1.73 theorem labels.
% Release-v1.74 synchronization: canonical Macros-v2_34.tex; generated local nonlinear Yang--Mills action, constrained connection-field phase space, source-free Gauss momentum map, and Legendre closure are integrated on their proved domains. The remaining non-gauge ACT_cl sectors and actual solution-history selection remain open.
% Release-v1.71 synchronization: closes the former JOINT_x/JOINT_5
% connection-occurrence gap by canonical structure-group extension of the
% nonflat Born connection, faithful associated representation, and the local
% cotangent--Clifford operator constructor.  Actual fluctuation members,
% actions, vacua, solutions, calibrations, and analytic QFT domains remain
% separately typed.  The active-AUX registry is advanced once.
% Release-v1.71 canonical macro and active-AUX registry.
% ======================================================================
% Macros-v2_39.tex -- canonical publication macro file
% Release v1.71 retains the v1.70 null-ray/QCD boundaries and additionally
% removes the duplicate structural connection actuality premise.  It does not
% claim global bundle topology, select a gauge/Higgs member, or promote the
% finite construction to an interacting continuum theory.
% Release v1.68 constructs the finite light-front total-fraction operators and
% parton number/momentum operator measures, proves their simplex and Hausdorff-limit
% theorems, and separates positive cut measures from scheme-renormalized
% light-ray distributions.  It also separates finite real Schur poles, stable
% norm-resolvent/Riesz pole-residue limits, and continued resonance
% pole/residue limits.  The Paper-VII-C scattering package is claim-indexed:
% local QFT and spectral claims import no S-matrix or infrared premise, while
% asymptotic amplitudes, LSZ, S-matrix poles, completeness, and cross sections
% import only their named massive or massless/long-range clauses.
% Release v1.67 factorizes BND48_H into a same-action Hessian/Schur
% descendant and still-open typed boundary-trace, critical-lift, and
% refinement-limit conditions.  It proves the finite-collar Dirichlet-to-
% Neumann law sqrt(y)tanh(L sqrt(y)), its nonuniform critical limit, and the
% finite-return ultra-infrared turnover.  The Paper-I-A Hardy constant mode is
% an exact geometric near-hit, not yet the physical metric boundary mode.
% The current local f(R) Hessian remains massive and therefore still refutes
% a RAR prediction on the minimal action branch.
% Release v1.66 derives the transverse depth-six normalizer bridge from the
% already generated C6 Reynolds image and Paper-I-A binary-dihedral normalizer.
% The bridge is the unique least 2D3-stable scalar-containing realization of
% the faithful quotient 2D3/C6=C2, has trace defect six and hence K=48, but is
% transverse to the nested active carrier.  Exact K=44 and K=52 contraction,
% semigroup, and energy-midpoint identities are registered as structural
% readouts rather than projector events.  The active E48 carrier gap remains
% unchanged.  A same-action boundary-symbol/Calderon intertwiner is isolated
% as the necessary and still unproved owner bridge to the static DtN sort;
% representation equivalence alone neither derives RAR nor creates a K=48
% action junction, pole, occupation, or cosmological transition.
% Release v1.65 derives the native E_48-centered carrier mass gap from the
% GP1--GP3-generated active spectrum, distinguishes its carrier-free interval
% from a particle located at E_48, and withdraws the unsupported promotion to
% a positive continuum Yang--Mills mass gap.  The latter remains a valid
% conditional transfer theorem on an enlarged continuum branch, but it is not
% an active framework prediction or an unfinished native-gap proof.
% Release v1.64 synchronizes Derivation Batches 14--22.  It registers the
% GP1--GP3 heat-profile orbit, same-Hessian electroweak response and pole
% ownership, K16/K48 action-response corrections, perturbation state and
% non-Gaussianity typing, finite-vacuum/gap/QCD limit repairs, factorized
% position-area ownership, the actual static metric response, and the ordered
% Wilson/string-tension limits.  Strict continuum convergence, positive
% physical gap, surviving higher cumulants, and positive tension remain open.
% Release v1.61 closes TIMEWIDTH16 as an independent direct-production datum.
% JUNC_{2|7} continues to own the finite intrinsic microscopic profile/width,
% while the actual same-action background/production history determines its
% physical-time pullback, support, and duration.  The current Paper-I-A--VII-C AUX
% registry is advanced once; shared mathematical notation is unchanged.
% Release v1.60 derives the final-octet primitive-remnant depletion, canonical
% Hopf-current representative, exact Hopf descent to dOmega/(4pi), and
% closed-FLRW a^{-3} neutral-current shape.  It retypes K=48 as the onset
% of the final one-unit geometric transfer and records cDM/RAR as distinct
% physical projections of that common neutral parent while preserving
% occurrence, Hilbert-owner, stress-matching, and static-kernel premises.
% Release v1.59 transports the v1.58 exceptional-shell/K16 closure through
% Papers III-F and VI-B--VII-C.  It replaces Paper VI-B's obsolete mixed-Hessian
% output typing by the third-variation pair vertex and Bogoliubov occupation
% support, derives the post-seed Higgs-pair feedback vertex and its exact
% alpha--delta kernel factors, propagates the unique K16/K48 action-junction
% split, identifies the QCD delta count with the upstream 2O shell atom only
% as a common-scalar numerical compatibility, and narrows THERM16 to dynamical
% approach/beta selection because the finite Gibbs/KMS target is already
% unique on a supplied simple observable algebra once its Hamiltonian and
% inverse temperature are supplied; direct-sum central weights remain typed.
% Release v1.58 registers the exceptional nearest-shell closure, the derived
% transverse K=16 action event, the unique interior action junction, the
% type-correct ordered pair vertex, the thin-layer squeeze law, and the
% radiation-like/matter-like kinematic split.  REG_0, transfer magnitude,
% physical layer width, and thermalization remain explicitly conditional.
% Release v1.57 registers the active-depth-six obstruction, the native
% K=48 carrier desert, the typed 7|9 owner reduction, and the dynamical
% reheating handoff across Papers III-B--VI-B.
% Release v1.56 registers the least-support n_2 inflation dressing of the
% persistent n_7 Laurent response, the analytic global-convexity and positive
% endpoint-measure proofs, the derived inflation core and K=16 coefficient
% freeze, and the no-particle dark-sector split into static RAR plus the
% conditional homogeneous effective-curvature response at reflected K=48.
% There is no n_6 action, scalaron-DM mass, or second Einstein--Hilbert term.
% Release v1.55 registers the protected-octet resolved/residual reflection
% q_sp+Z_W=8, its j<->8-j and K<->64-K scale covariance, and the derived
% rank-one-exit quadratic grade C_1^2/U=C_2 in Papers III-C, III-E, and VIII.
% The reflection is a typed trace/scale theorem, not a time reversal or a
% physical event-occurrence theorem; reheating transfer remains conditional.
% Release v1.54 registers the exact positive-part Laurent indices, the
% active-domain/terminal-endpoint action split, the no-second-EH firewall, and
% the q_DM=6/K_DM=48 scalaron readout typing across Papers III-D/E/F, VI,
% VIII, and IX.
% Release v1.53 registers the common nonnegative Laurent obstacle projection
% for n_1^+ and n_7^+, discharges the three former exit/freeze assumptions,
% and advances the Paper-III-D/E/F, VIII, and IX consumers.  Shared notation
% and bibliography ownership are unchanged.
% Release v1.51 registers the PQTI derivation batch and its updated
% Paper-III-A/C/D/E/F, VIII, and IX consumers.  Shared notation and
% bibliography ownership are unchanged.
% Release v1.50 advances the Paper-III-B AUX target to v1.40 after the
% rigorous audit repair; shared notation and bibliography ownership are unchanged.
% Release v1.49 advances the Paper-II-E/F AUX target to v3.14 after the
% minimal audit repair; shared notation and bibliography ownership are unchanged.
% Release v1.46 synchronizes Paper II-A's minimal frozen constructor-origin
% theorem and its five dependent AUX targets.  Shared notation is unchanged.
% Release v1.45 updates the Paper-II-B AUX target to v2.23 after the
% GP3-rigidity and appendix-placement repair; shared notation is unchanged.
% Release v1.44 advances the centralized Paper-I-A--VII-C AUX registry after the
% Paper-I-A v5.16 protected-completion scope repair.  Every current
% paper imports this one file; no paper-level opt-in or \externaldocument list
% is permitted.  The bibliography remains centralized here through Born-v6_13.
% Shared bibliography, theorem infrastructure, colors, notation alignment,
% and current Paper-I-A--VII-C external-reference registration.
% This file centralizes common imports, bibliography sorting, colors, theorem
% infrastructure, and the notation shared by the complete series.
% v2.9 (release v1.43) made the Paper-I-A--VII-C AUX registry automatic and
% centralized the bibliography and all external-reference ownership.
% v2.8 release-registry refresh (v1.33) registers the rooted generated-parent
% cascade and the POL/REG/free-law/11+4/Born-group closure revisions.  No
% shared mathematical notation or Paper-II-A fundamental principle is changed.
% v2.8 release-registry refresh (v1.32) registers the universal generated-
% continuation revisions in Papers II-A, III-A, III-C, IV, and VIII.  No
% shared mathematical notation or Paper-II-A fundamental principle is changed.
% v2.8 release-registry refresh (v1.30) registers the CRP4/projective-state
% source revisions and their direct registry-only consumers.  No shared
% mathematical notation or Paper-II-A axiom is changed.
% v2.7 release-registry refresh (v1.27) updates only the current external
% Paper-I-A--VII-C basenames after the placement/assumption synchronization; no
% shared mathematical notation or Paper-II-A axiom is changed.
% v2.7 synchronizes the geometric defect-repair revision: Paper-II-A v3.42,
% Paper-III-C v1.26, Paper-III-D v1.27, Paper-III-E v1.32, and Paper-III-F
% v1.34.  It registers the common trace-defect/spectral-filtration theorem,
% its owner-typed response pullback, and the reduced GRED and TERM interfaces.
% No shared mathematical notation is changed.
% v2.6 synchronizes Paper-III-E v1.31's cone-type repair, minimal-response
% classification, and explicit classical action-role boundary, together with
% Paper-III-F v1.32's downstream response-atlas synchronization.  No shared
% mathematical notation is changed.
% v2.5 synchronizes the Paper-III-D v1.26 faithful-trace-duality proof,
% exact ordered-owner joint locus, and typed alternative-presentation
% bookkeeping repair, together with its Paper-III-E and closure-frontier
% references.  No shared mathematical notation is changed.
% v2.4 synchronizes the Paper-III-C v1.25 orbit/process/owner repair, its
% downstream Papers III-B--F, V, VI, VIII, and IX revisions, and Born-v6_13.
% v2.3 synchronizes the Paper-II-A v3.41 witness/support/principal inheritance
% layer, the Paper-III-A v1.24 joint-realizability/readout firewall, all dependent
% host-rate revisions, and the current Paper-I-A--VII-B source registry.
% v2.2 synchronizes the generated active-realization, determinant-depth,
% dimension-weighted hierarchy, and typed cosmological-event interfaces.
% v2.1 synchronizes the exceptional-support interface revisions: Paper I-A
% v5.9, Paper II-A v3.30, Paper II-E/F v3.7, and Paper IV v4.16.
% v2.0 introduces the stage-safe G0 occurrence notation, synchronizes the
% active Paper-I-A--IV source registry (current rigorous Paper-III closure:
% III-A v1.17, III-B v1.22, III-C--E v1.17, III-F v1.14, IV v4.14), and makes
% external-reference self-skip depend on logical paper role rather than a
% versioned job name.
% ======================================================================

% ======================================================================
% macros.tex — Shared macros for the Born Geometry series
% ======================================================================

% ======================================================================
% 1. PACKAGES, BIBLIOGRAPHY, COLORS, AND UNIT DEFINITIONS
% ======================================================================
% Core packages shared across the complete Paper-I-A--VII-C series.  Loading them here makes
% the macro file a genuine series preamble while remaining compatible with existing
% paper preambles that load the same packages before \input.
\usepackage{amsmath}
\usepackage{amssymb}
\usepackage{amsthm}
\usepackage{mathrsfs}
\usepackage{array}
\usepackage{booktabs}
\usepackage{xcolor}
\makeatletter
\@ifpackageloaded{geometry}{}{\usepackage{geometry}\geometry{margin=1in}}
\makeatother
\usepackage{mdframed}

% Shared support packages required by the macros below and by most papers.
\usepackage{ragged2e}
\usepackage{longtable}
\usepackage{tabularx}
\usepackage{enumitem}
\usepackage[most]{tcolorbox}

% Series-wide numeric bibliography convention using the portable BibTeX path.
% All Papers I-A--VII-C use the ordinary \cite interface.
\usepackage[numbers,sort&compress]{natbib}
\providecommand{\printbibliography}{%
  \bibliographystyle{plainnat}%
  \bibliography{Born-v6_59}%
}

% Shared colors.  \providecolor avoids clobbering a paper-specific override.
\providecolor{postulatecolor}{RGB}{230,242,255}
\providecolor{compare}{RGB}{230,245,230}
\providecolor{definitioncolor}{RGB}{230,242,255}

\providecolor{openriskcolor}{RGB}{255,235,235}
\providecolor{bridgecolor}{RGB}{240,240,255}
\providecolor{gold}{RGB}{255,215,0}

% ======================================================================
% 2. CUSTOM ENVIRONMENTS
% ======================================================================
% Common table column types used in all papers.  Guarded so that a paper may
% define its own versions before loading this file.
\makeatletter
\@ifundefined{NC@find@L}{%
  \newcolumntype{L}[1]{>{\RaggedRight\arraybackslash}p{#1}}%
}{}
\@ifundefined{NC@find@C}{%
  \newcolumntype{C}[1]{>{\Centering\arraybackslash}p{#1}}%
}{}
\@ifundefined{NC@find@Y}{%
  \newcolumntype{Y}{>{\RaggedRight\arraybackslash}X}%
}{}
\makeatother

% ======================================================================
% 3. THEOREM STYLES AND ENVIRONMENTS
% ======================================================================
% Custom underlined theorem style (useful for remarks or notes)
\newtheoremstyle{underlined}
  {3pt}   % Space above
  {3pt}   % Space below
  {}      % Body font (empty means normal/roman text, typical for remarks)
  {}      % Indent amount
  {}      % Theorem head font (can be \itshape, \bfseries, etc.)
  {.}     % Punctuation after theorem head
  { }     % Space after theorem head
  {\underline{\thmname{#1}\thmnumber{ #2}}\thmnote{ (#3)}} % Theorem head formatting

% Standard plain theorem style environments.  Each environment is guarded so
% the shared macro file can be loaded after a paper-specific theorem setup.
\theoremstyle{plain}
\makeatletter
\@ifundefined{theorem}{\newtheorem{theorem}{Theorem}[section]}{}
\@ifundefined{lemma}{\newtheorem{lemma}{Lemma}[section]}{}
\@ifundefined{proposition}{\newtheorem{proposition}{Proposition}[section]}{}
\@ifundefined{corollary}{\newtheorem{corollary}{Corollary}[section]}{}
\makeatother

% Definition style environments (typically non-italicized body)
\theoremstyle{definition}
\makeatletter
\@ifundefined{definition}{\newtheorem{definition}{Definition}[section]}{}
\@ifundefined{postulate}{\newtheorem{postulate}{Postulate}[section]}{}
\@ifundefined{closure}{\newtheorem{closure}{Closure}[section]}{}
\makeatother

% Applying the custom underlined style
\theoremstyle{underlined}
\makeatletter
\@ifundefined{remark}{\newtheorem{remark}{Remark}[section]}{}
\@ifundefined{convention}{\newtheorem{convention}{Convention}[section]}{}
\@ifundefined{principle}{\newtheorem{principle}{Principle}[section]}{}
\makeatother

% Framed proof environment.
% Every \begin{proof} ... \end{proof} is automatically framed/shaded.
% This avoids mdframed/trivlist nesting, which can drop the top border after
% theorem-like environments.
\makeatletter
\renewenvironment{proof}[1][\proofname]{%
  \par\addvspace{6pt}%
  \noindent
  \begin{tcolorbox}[%
    enhanced,
    breakable,
    colback=black!4,
    colframe=black!35,
    boxrule=0.35pt,
    arc=1pt,
    left=6pt,
    right=6pt,
    top=4pt,
    bottom=4pt,
    boxsep=0pt,
    before skip=0pt,
    after skip=6pt,
    parbox=false
  ]%
  \pushQED{\qed}%
  \normalfont\itshape #1\@addpunct{.}\enspace
  \normalfont\ignorespaces
}{%
  \unskip\nobreak\hfill\popQED
  \end{tcolorbox}%
  \par\addvspace{2pt}%
}
\makeatother

\mdfdefinestyle{postulatebox}{
  backgroundcolor=yellow!10,
  linecolor=black,
  linewidth=0.4pt,
  skipabove=\baselineskip,
  skipbelow=\baselineskip,
  nobreak=true
}

% ======================================================================
% PROOF-READABILITY HELPERS
% ======================================================================

\newcommand{\ProofPlan}[1]{%
  \par\smallskip
  \noindent\textbf{Proof plan. }#1%
  \par\smallskip
}

\newlist{proofsteps}{enumerate}{2}
\setlist[proofsteps,1]{
  label=\textbf{(\alph*)},
  ref=\alph*,
  leftmargin=2.2em,
  itemsep=0.45\baselineskip,
  topsep=0.35\baselineskip
}
\setlist[proofsteps,2]{
  label=\textbf{(\roman*)},
  ref=\roman*,
  leftmargin=2.0em,
  itemsep=0.25\baselineskip,
  topsep=0.2\baselineskip
}


% ======================================================================
% 4. NUMBER SYSTEMS AND ALPHABETS
% ======================================================================
\providecommand{\RR}{\mathbb{R}} % Real numbers
\providecommand{\PP}{\mathbb{P}} % Projective space / Probability
\providecommand{\CC}{\mathbb{C}} % Complex numbers
\providecommand{\ZZ}{\mathbb{Z}} % Integers
\providecommand{\NN}{\mathbb{N}} % Natural numbers
\providecommand{\HH}{\mathbb{H}} % Quaternions / Hyperbolic space
\providecommand{\LFRayGeom}{\mathsf{LF\!\mbox{-}RAY\!\mbox{-}GEOM}}
\providecommand{\LFRep}{\mathsf{LF\!\mbox{-}REP}}
\providecommand{\LFState}{\mathsf{LF\!\mbox{-}STATE}}
\providecommand{\LFLim}{\mathsf{LF\!\mbox{-}LIM}}
\providecommand{\LFScat}{\mathsf{LF\!\mbox{-}SCAT}}
% Number systems / projective spaces
\providecommand{\CP}{\mathbb{CP}}
\providecommand{\RP}{\mathbb{RP}}

% ======================================================================
% 5. GROUPS AND LIE THEORY
% ======================================================================
\providecommand{\SO}{\ensuremath{\mathrm{SO}}}     % Special Orthogonal group
\providecommand{\SU}{\ensuremath{\mathrm{SU}}}     % Special Unitary group
\providecommand{\Spin}{\ensuremath{\mathrm{Spin}}} % Spin group
\providecommand{\CS}{\ensuremath{\mathrm{CS}}}     % Chern-Simons (capitalized)
\providecommand{\cs}{\ensuremath{\mathrm{cs}}}     % Chern-Simons (lowercase)

% ======================================================================
% 6. COMMON MANIFOLDS AND TOPOLOGY
% ======================================================================
\providecommand{\Sone}{S^1}                 % Standard circle

% ======================================================================
% 7. MATH OPERATORS (Declared via \DeclareMathOperator)
% ======================================================================
\DeclareMathOperator{\End}{End}   % Endomorphisms
\DeclareMathOperator{\Aut}{Aut}   % Automorphisms
\DeclareMathOperator{\spec}{spec} % Lowercase spectrum
\DeclareMathOperator{\tr}{tr}     % Lowercase trace (if you already have \Tr)
\DeclareMathOperator{\supp}{supp} % Support
\DeclareMathOperator{\TrA}{Tr_{A}} % Partial trace over subsystem A
\DeclareMathOperator{\TrB}{Tr_{B}} % Partial trace over subsystem B
% Lie/algebra operators
\DeclareMathOperator{\Ad}{Ad}
\DeclareMathOperator{\Sym}{Sym}
\DeclareMathOperator{\su}{su}
% ======================================================================
% 8. STANDARD OPERATOR-LIKE SYMBOLS (Safe definitions via \providecommand)
% ======================================================================
\providecommand{\vol}{\operatorname{vol}}   % Volume form
\providecommand{\Vol}{\operatorname{Vol}}   % Total volume
\providecommand{\Ric}{\operatorname{Ric}}   % Ricci tensor
\providecommand{\Rm}{\operatorname{Rm}}     % Riemann tensor
\providecommand{\Tr}{\operatorname{Tr}}     % Uppercase trace
\providecommand{\Spec}{\operatorname{Spec}} % Uppercase spectrum
\providecommand{\diag}{\operatorname{diag}} % Diagonal matrix
\providecommand{\rank}{\operatorname{rank}} % Rank
\providecommand{\scal}{\mathcal R} % Scalar curvature
\providecommand{\Cl}{\operatorname{Cl}}     % Clifford algebra
\providecommand{\diam}{\operatorname{diam}} % Diameter

% ======================================================================
% 9. FUNCTIONS, TRIGONOMETRY, AND CALCULUS
% ======================================================================

% Codifferential naming conventions
% (Avoid \ddag: commonly defined by physics.sty and symbol packages)
\providecommand{\codiff}{\mathrm{d}^{\dagger}} % Codifferential operator

% Quantum mechanics bra-ket notation (needed for basis tables)
\providecommand{\ket}[1]{\lvert #1 \rangle}
\providecommand{\bra}[1]{\langle #1 \rvert}
\providecommand{\braket}[2]{\langle #1 \mid #2 \rangle}

% ======================================================================
% 10. PHYSICAL SCALES AND CONSTANTS (Paper 1c / Born Geometry)
% ======================================================================
\providecommand{\Rp}{\ensuremath{R_{\mathrm P}}}     % matched common EH curvature scale (not defined from observed G)
\providecommand{\hBorn}{\ensuremath{h_{\mathrm B}}} % physical representative 2\pi\hbar of one primitive Born scalar action cell
\providecommand{\UEHjoint}{\ensuremath{U_{J}}} % joint-Dirac EH curvature-coordinate scale 8 kappa_J
\providecommand{\Gread}{\ensuremath{G_{\rm scal}}} % scalar-action coordinate c_op^3/(h_B U); physical matching explicit
\providecommand{\GreadJoint}{\ensuremath{G_{J,\mathrm{raw}}}} % raw ordinary-trace coordinate = G_scal/96
\providecommand{\Rds}{R_{\mathrm{dS}}}    % de Sitter Curvature
\providecommand{\Mpl}{M_{\mathrm{Pl}}}    % Planck mass
\providecommand{\lp}{\ell_{\mathrm{Pl}}}  % Planck length
\providecommand{\Pmax}{P_{\mathrm{max}}}  % Maximum momentum
\providecommand{\Mplbar}{\bar{M}_{\mathrm{Pl}}} % Reduced Planck mass
\providecommand{\eff}{\mathrm{eff}}       % Subscript for effective quantities
\providecommand{\Cosmo}{\ensuremath{\Lambda}}                % Cosmological constant (use \Cosmo in source)
\providecommand{\LCDM}{\ensuremath{\Cosmo\mathrm{CDM}}}  % \LambdaCDM (typesets with \Cosmo)
\providecommand{\Lcut}{\ensuremath{\Lambda_{\mathrm{cut}}}}  % Cutoff scale (brace when used as script)
\providecommand{\Rcs}{\ensuremath{R_{\mathrm{cs}}}}      % Configuration-space radius
\providecommand{\CosmoEff}{\ensuremath{\Cosmo_{\eff}}}   % Effective cosmological constant

% Disambiguation: \Rmom is the momentum-space (S^3) radius. 
% Reserve plain R for the 4D curvature scalar in f(R).
\providecommand{\Rmom}{\ensuremath{R_{\mathrm{mom}}}} 
\providecommand{\Punit}{P_{0}} % Momentum unit set by \Rmom (= \hbar/\Rmom)
\providecommand{\Rvec}{\ensuremath{R_{\mathrm{vec}}}} % invariant vector-space radius in the reduced round sector
% Specific numeric values and expressions
\providecommand{\Vlogcap}{L_{\mathrm{Hopf}}^{(1)}}
\providecommand{\Hopf}{\operatorname{Hopf}}
% Scalar phase (S^1) sector + universal cover notation (Paper I-A)
\providecommand{\tper}{T_\star}      % Time period/scale
\providecommand{\Eper}{E_\star}      % Energy period/scale
\providecommand{\ttilde}{\tilde t}   % Universal cover time
\providecommand{\Etilde}{\tilde E}   % Universal cover energy

% ======================================================================
% 10A. SERIES-WIDE LOGICAL / NOTATIONAL DISAMBIGUATIONS
% ======================================================================
% Use these symbols to avoid collisions among: (i) the Paper--III
% no-insertion/geometric-closure class, (ii) Paper--I-A trimer diagonals,
% (iii) Paper--V-A/V-C family-space diagonals, (iv) scalar-index running,
% and (v) protected shell-clock labels.
\providecommand{\Rprim}{\ensuremath{\mathsf R_{\mathrm{prim}}^{\Sigma}}} % primitive finite structural readouts
\providecommand{\Rloop}{\ensuremath{\mathsf R_{\circ}^{\Sigma}}} % closed-loop readout role
\providecommand{\Rbdry}{\ensuremath{\mathsf R_{\partial}^{\Sigma}}} % boundary-conditioned readout role
\providecommand{\Rgen}{\ensuremath{\mathsf R_{\mathrm{fin}}^{\Sigma}}} % finite generated closure of primitive roles
\providecommand{\PhysOcc}{\operatorname{PhysOcc}}
\providecommand{\Actual}{\operatorname{Actual}}
\providecommand{\GeomLic}{\operatorname{GeomLic}}
\providecommand{\InvDist}{\operatorname{InvDist}}
\providecommand{\Dcl}{\ensuremath{\mathsf{D}_1}}      % unmarked positive Dirac-even closed spectral-readout form class
\providecommand{\Dheat}{\ensuremath{\mathsf D_{\mathrm{heat}}}} % one-parameter heat-semigroup subclass
\providecommand{\DtriI}{\ensuremath{D^{\mathrm{tri}}_1}} % Paper--I-A cyclic-trimer diagonal
\providecommand{\DtriII}{\ensuremath{D^{\mathrm{tri}}_2}}
\providecommand{\DfamI}{\ensuremath{D^{\mathrm{fam}}_1}} % Family-space diagonal selector
\providecommand{\DfamII}{\ensuremath{D^{\mathrm{fam}}_2}}
\providecommand{\alpharun}{\ensuremath{\alpha_{\mathrm{run}}}} % scalar spectral-index running
\providecommand{\Ksh}{\ensuremath{K_{\mathrm{sh}}}} % protected shell-clock coordinate
\providecommand{\ddef}{\ensuremath{\delta_{\mathrm{def}}}} % universal defect scalar after Paper--III realization

% ======================================================================
% 10B. SERIES-WIDE LOGICAL STATUS TAGS
% ======================================================================
% Upright small-caps wrapper for audit/status labels.  It prevents LaTeX from
% trying to use unavailable small-caps italic font shapes inside theorem bodies
% such as TU/lmr/m/scit.
\providecommand{\AuditLabel}[1]{}
\renewcommand{\AuditLabel}[1]{{\normalfont\scshape #1}}

% These tags are bookkeeping labels for consistency audits. They do not
% create new mathematical hypotheses; they record how a datum is being used
% in the proof chain.
\providecommand{\StatusPrimitive}{}
\renewcommand{\StatusPrimitive}{\AuditLabel{Primitive}}
\providecommand{\StatusDerived}{}
\renewcommand{\StatusDerived}{\AuditLabel{Derived}}
\providecommand{\StatusImported}{}
\renewcommand{\StatusImported}{\AuditLabel{Imported}}
\providecommand{\StatusConditional}{}
\renewcommand{\StatusConditional}{\AuditLabel{Conditional}}
\providecommand{\StatusBridge}{}
\renewcommand{\StatusBridge}{\AuditLabel{Bridge-derived}}
\providecommand{\StatusComparison}{}
\renewcommand{\StatusComparison}{\AuditLabel{Comparison-only}}

% ======================================================================
% 10C. SERIES-WIDE CLAIM-WORDING AUDIT LABELS
% ======================================================================
% These are textual audit labels used to keep modal language consistent
% across the series.  They introduce no new mathematical assumptions.
\providecommand{\ClaimForced}{}
\renewcommand{\ClaimForced}{\AuditLabel{forced}}
\providecommand{\ClaimUnique}{}
\renewcommand{\ClaimUnique}{\AuditLabel{unique}}
\providecommand{\ClaimCanonical}{}
\renewcommand{\ClaimCanonical}{\AuditLabel{canonical}}
\providecommand{\ClaimMinimal}{}
\renewcommand{\ClaimMinimal}{\AuditLabel{minimal}}
\providecommand{\ClaimClosed}{}
\renewcommand{\ClaimClosed}{\AuditLabel{closed}}
\providecommand{\ClaimPrediction}{}
\renewcommand{\ClaimPrediction}{\AuditLabel{prediction}}

% ======================================================================
% 10D. SERIES-WIDE THEOREM-CHAIN AUDIT LABELS
% ======================================================================
% These labels mark the direction of load-bearing theorem chains. They are
% audit annotations only; they introduce no new mathematical assumptions.
\providecommand{\ChainInput}{}
\renewcommand{\ChainInput}{\AuditLabel{Chain input}}
\providecommand{\ChainOutput}{}
\renewcommand{\ChainOutput}{\AuditLabel{Chain output}}
\providecommand{\ChainGuard}{}
\renewcommand{\ChainGuard}{\AuditLabel{Chain guard}}
\providecommand{\ChainNoBackflow}{}
\renewcommand{\ChainNoBackflow}{\AuditLabel{No backflow}}

% ======================================================================
% 10E. SERIES-WIDE LOCAL THEOREM-AUDIT LABELS
% ======================================================================
% These labels mark a statement-level dependency audit around individual
% theorem/definition nodes. They are bookkeeping labels only.
\providecommand{\LocalAudit}{}
\renewcommand{\LocalAudit}{\AuditLabel{Local audit}}
\providecommand{\UsesOnly}{}
\renewcommand{\UsesOnly}{\AuditLabel{Uses only}}
\providecommand{\Exports}{}
\renewcommand{\Exports}{\AuditLabel{Exports}}
\providecommand{\DoesNotUse}{}
\renewcommand{\DoesNotUse}{\AuditLabel{Does not use}}

% ======================================================================
% 10F. SERIES-WIDE NUMERICAL-ROW AUDIT LABELS
% ======================================================================
% These labels classify individual rows in numerical, benchmark, status,
% and comparison tables.  They are audit annotations only and do not change
% any theorem, value, formula, or comparison number.
\providecommand{\RowAudit}{}
\renewcommand{\RowAudit}{\AuditLabel{Numerical-row audit}}
\providecommand{\RowTheorem}{}
\renewcommand{\RowTheorem}{\AuditLabel{Theorem-fixed row}}
\providecommand{\RowBridge}{}
\renewcommand{\RowBridge}{\AuditLabel{Bridge-fixed row}}
\providecommand{\RowRepresentative}{}
\renewcommand{\RowRepresentative}{\AuditLabel{Representative row}}
\providecommand{\RowComparison}{}
\renewcommand{\RowComparison}{\AuditLabel{Comparison row}}
\providecommand{\RowCalibrated}{}
\renewcommand{\RowCalibrated}{\AuditLabel{Calibrated row}}
\providecommand{\RowEstimate}{}
\renewcommand{\RowEstimate}{\AuditLabel{Estimate row}}
\providecommand{\RowOpen}{}
\renewcommand{\RowOpen}{\AuditLabel{Open row}}
\providecommand{\RowNoSelector}{}
\renewcommand{\RowNoSelector}{\AuditLabel{No selector}}

% ======================================================================
% 10G. TABLE HELPERS
% ======================================================================
% Table helpers used by Paper V-C.  The unused readout-label preset has been
% removed; only active table helpers remain.
\providecommand{\gpsection}[1]{%
  \addlinespace[0.55em]
  \multicolumn{3}{@{}l}{\textsc{#1}}\\[-0.25em]
  \midrule
}
\providecommand{\mathcell}[1]{%
  \(\begin{gathered}#1\end{gathered}\)%
}

\newcommand{\DerBy}[1]{%
  \begin{tabular}[t]{@{}l@{}}#1\end{tabular}%
}

\newcommand{\TableNotes}[1]{%
  \par\medskip
  \begin{minipage}{0.94\linewidth}
  \footnotesize\emph{Notes.} #1
  \end{minipage}\par
}
% ======================================================================
% 10H. PAPER VI-A / QCD NOTATION
% ======================================================================
% Shared symbols used by Paper VI-A and downstream QCD/Yang--Mills handoffs.
\providecommand{\Pqcd}{\mathfrak P_{\rm VIII}^{\rm QCD}}
\providecommand{\Iqcd}{\mathfrak I_{\rm QCD}^{(7)}}
\providecommand{\muQCD}{\mu_{\rm QCD}^{\rm int}}
\providecommand{\LambdaQCD}{\Lambda_{\rm QCD}}
\providecommand{\LambdaBH}{\Lambda_{\rm BH}^{\rm QCD}}
\providecommand{\HQCD}{\mathcal H_{\rm QCD}^{\rm BH}}
\providecommand{\HQCDfin}{H_{\rm QCD}^{\rm fin}}
\providecommand{\Hraw}{\mathcal H_{\rm raw}}
\providecommand{\HphysQCD}{\mathcal H_{\rm phys}^{\rm QCD}}
\providecommand{\Vcol}{V_C}
\providecommand{\Ps}{P_{\rm sing}}
\providecommand{\PiOct}{\Pi_8}
\providecommand{\MSbar}{\overline{\rm MS}}
\providecommand{\YM}{\mathrm{YM}}
\providecommand{\Ktri}{K_{\triangle}}
\providecommand{\KYM}{K_{\YM}^{\rm BH}}
\providecommand{\aBH}{a_{\rm BH}}
\providecommand{\epsBH}{\varepsilon_{\rm BH}}
\newcommand{\RowConstructive}{\textnormal{\textsc{Constructive}}}
\newcommand{\RowObligation}{\textnormal{\textsc{Bridge obligation}}}
\newcommand{\RowTarget}{\textnormal{\textsc{Target row}}}
\newcommand{\RowValidation}{\textnormal{\textsc{Validation row}}}
\providecommand{\slashed}[1]{\not{#1}}

% ======================================================================
% 10I. PAPER VI-B / CONTINUOUS-QFT NOTATION
% ======================================================================
% Shared symbols used by Paper VI-B and any downstream continuous-QFT or
% Yang--Mills handoff paper.  All definitions are guarded so individual
% papers may override them locally before loading this file if needed.
\providecommand{\A}{\mathfrak A}
\providecommand{\Hphys}{\mathcal H_{\rm phys}}
\providecommand{\QCD}{\mathrm{QCD}}
\providecommand{\BH}{\mathrm{BH}}
\providecommand{\loc}{\mathrm{loc}}

\providecommand{\RBH}{\mathcal R_{\rm BH}}
\providecommand{\DG}{\mathfrak D_G}
\providecommand{\PW}{\mathcal{P\!W}} % Peter--Weyl spectral-symbol subspace
\providecommand{\Toe}{\operatorname{Toe}} % Toeplitz compression/representative
\providecommand{\Hol}{\operatorname{Hol}} % Holonomy

% ======================================================================
% 11. LATE-LOADED PACKAGES (Must be at the end of the preamble)
% ======================================================================
% Load late packages only if a paper has not already loaded them, avoiding
% option clashes in paper-specific preambles.
\makeatletter
\@ifpackageloaded{hyperref}{}{\usepackage{hyperref}}
\makeatother
% --- Series-wide disambiguations (Pass 3) ---
\providecommand{\diracdelta}{\delta_{\rm D}} % Dirac delta distribution (delta reserved for hierarchy)
\providecommand{\Hg}{H_{\mathrm g}} % Hodge harmonic projector (keep H for Hubble)
\providecommand{\Gr}{G_{\mathrm r}} % Green operator (keep G for Newton)
\providecommand{\FD}{\mathrm{D}} % functional/variational derivative symbol (delta reserved)
\providecommand{\Id}{\mathbf{1}} % identity/Kronecker tensor (avoid \delta_{ij})


% BEGIN INLINE SOURCE: Paper00-ActiveExternalReferences-v2_130.tex
% Framework v2.191: current reference binding; existing definitions retained.
% Framework v2.190: current22-root registry with native clock-evolution proofs.
% Framework v2.189: exact current22-root external registry, including retained v2.188 hardening owners.
% Framework v2.187: current registry synchronization only; source release v2.186 preserved.
% Framework v2.186: authoritative current 22-paper registry.
% Framework v2.182: all current 22 logical paper AUX bindings.
% Framework v2.170: current I-B v1.43 external AUX binding.
% Framework v2.167: exact current 22-root registry after CP--BL2 integration.
% Framework v2.166: exact current 22-root compact-release AUX registry.
% Framework v2.165: authoritative current root bindings.
% Framework v2.164: exact current 22 manuscript AUX bindings.
% Framework v2.162: finite metric, functional Hilbert and scoped Ward integration.
% Framework v2.161: current 22 manuscript AUX bindings; source only.
% Framework v2.156 / registry v2.113: synchronized current 22 manuscript roots; source-only registry.
% Framework v2.143 / registry v2.111: active manuscript AUX basenames; no generated AUX or PDF included.
% ======================================================================
% Central external-reference imports, numbers only, no PDF jumping
% ======================================================================

% xr-hyper is loaded centrally so that papers may retain their existing
% \usepackage{xr-hyper} lines without redefining \externaldocument after this
% registry is installed.  If plain xr was loaded earlier, retain it.
\makeatletter
\@ifpackageloaded{xr-hyper}{}{%
  \@ifpackageloaded{xr}{}{\usepackage{xr-hyper}}%
}
\makeatother

% Make external-document registration idempotent.  The registry identifies a
% logical paper by its Paper<number><optional letter> prefix.  Prefix-qualified
% imports remain independent.  Current paper roots contain no local
% \externaldocument calls; this guard also makes accidental duplicates harmless.
\ExplSyntaxOn
\prop_new:N \g_bh_external_refs_prop
\cs_new_eq:NN \BHOriginalExternalDocument \externaldocument

\cs_new_protected:Npn \bh_external_document:nnnn #1#2#3#4
 {
  \tl_set:Nn \l_tmpa_tl {#3}
  \regex_replace_once:nnN
   { \A (?: xref/ )? (Paper[0-9]+[A-F]?) .* \Z }
   { \1 }
   \l_tmpa_tl
  \tl_set:Nx \l_tmpb_tl { #1 | \tl_use:N \l_tmpa_tl }
  \prop_if_in:NVTF \g_bh_external_refs_prop \l_tmpb_tl
   {
    \PackageInfo{BornHopf}
     {Skipping~duplicate~external~refs~for~#3}
   }
   {
    \prop_gput:NVn \g_bh_external_refs_prop \l_tmpb_tl {#3}
    \IfNoValueTF{#4}
     { \BHOriginalExternalDocument[#1][#2]{#3} }
     { \BHOriginalExternalDocument[#1][#2]{#3}[#4] }
   }
 }

\RenewDocumentCommand{\externaldocument}{O{} O{} m o}
 { \bh_external_document:nnnn {#1}{#2}{#3}{#4} }

\NewDocumentCommand{\BHExternalRefs}{m}
 {
  % Do not import the aux file of the logical paper currently being compiled,
  % even when the job and registry use different version suffixes.
  \tl_set:Nx \l_tmpa_tl { \c_sys_jobname_str }
  \regex_replace_once:nnN
   { \A (Paper[0-9]+[A-F]?) .* \Z }
   { \1 }
   \l_tmpa_tl
  \tl_set:Nn \l_tmpb_tl { #1 }
  \regex_replace_once:nnN
   { \A (Paper[0-9]+[A-F]?) .* \Z }
   { \1 }
   \l_tmpb_tl
  \str_if_eq:eeTF { \tl_use:N \l_tmpa_tl } { \tl_use:N \l_tmpb_tl }
   {
    \PackageInfo{BornHopf}
     {Skipping~external~refs~for~#1~because~this~is~the~current~document}
   }
   {
    \IfFileExists{xref/#1.aux}
     { \externaldocument[][nocite]{xref/#1}[#1.pdf] }
     {
      \IfFileExists{#1.aux}
       { \externaldocument[][nocite]{#1} }
       {
        \PackageInfo{BornHopf}
         {External~aux~file~#1.aux~not~found;~no~labels~imported}
       }
     }
   }
 }
\ExplSyntaxOff

% External references are automatic.  This file is the sole owner of both the
% bibliography and the current Paper-I-A--VII-D AUX registry (30 roots including I-C).  Missing auxiliaries
% are informational on the first pass and are populated by the source-only
% two-pass series build.  Basenames deliberately omit temporary download/copy
% suffixes such as ``(1)'' or ``-1''.
  % Paper I-A: intrinsic S3 and exceptional group introduction.
  \BHExternalRefs{Paper1A-Geometry-v5_175}

  % Paper I-B: complete exceptional geometry and finite spectral structures.
  \BHExternalRefs{Paper1B-ExceptionalGeometry-v1_5}
  % Paper I-C: reciprocal product M8 geometry.
  \BHExternalRefs{Paper1D-TechnicalMaterial-v1_65}
  % Paper I-D: terminal technical proofs and calculations.
  \BHExternalRefs{Paper1C-M8Geometry-v1_137}

\ifdefined\BHMathOnly
  % Papers I-A and I-B import only the mathematical atlas labels.
\else
  % Paper II: five physical-mathematical main papers plus one technical proof volume.
  \BHExternalRefs{Paper2A-GeometricConstitution-v3_221}
  \BHExternalRefs{Paper2B-MomentumCarrier-v2_144}
  \BHExternalRefs{Paper2C-ReciprocalLocalization-v2_169}
  \BHExternalRefs{Paper2D-SymmetrySpectra-v2_141}
  % The former II-F export labels remain in II-E; current II-F is the technical volume.
  \BHExternalRefs{Paper2E-KinematicExports-v3_138}
  % Paper II-F: complete technical proof and calculation owner for II-A--II-E.
  \BHExternalRefs{Paper2F-TechnicalMaterial-v1_67}

  % Paper III: Dynamics and Intrinsic Geometric Constitution (six components).
  \BHExternalRefs{Paper3A-GeneratedDynamics-v1_176}
  \BHExternalRefs{Paper3B-SpectralAction-v1_187}
  \BHExternalRefs{Paper3C-ProtectedHierarchy-v1_187}
  \BHExternalRefs{Paper3D-RunningResponse-v1_165}
  \BHExternalRefs{Paper3E-GravityAction-v1_197}
  \BHExternalRefs{Paper3F-CompleteDynamics-v1_263}
  % Paper III-G: complete preserved technical proof owner.
  \BHExternalRefs{Paper3G-TechnicalMaterial-v1_53}

  % Paper IV-A: generated finite quantum mechanics.
  \BHExternalRefs{Paper4A-FiniteQuantumMechanics-v4_214}
  % Paper IV-B: complete finite quantum processes and physical response.
  \BHExternalRefs{Paper4B-QuantumResponse-v1_49}
  % Paper IV-C: extended proofs and finite technical calculations.
  \BHExternalRefs{Paper4C-TechnicalMaterial-v1_49}

  % Paper V-A: The Standard Model.
  \BHExternalRefs{Paper5A-FiniteMatterGeometry-v3_209}

  % Paper V-B: Finite Spectrum and Flavor Closure.
  \BHExternalRefs{Paper5B-TerminalSpectrum-v5_199}

  % Paper V-C: Finite Interaction Boundaries of the Standard Model.
  \BHExternalRefs{Paper5C-InteractionResponse-v2_177}
  % Paper V-D: preserved technical proof and calculation owner.
  \BHExternalRefs{Paper5D-TechnicalMaterial-v1_53}

  % Paper VI-A: Gravity and Closed Cosmology.
  \BHExternalRefs{Paper6A-GravityCosmology-v2_260}

  % Paper VI-B: Geometric Dark Sector and Cosmological Readouts.
  \BHExternalRefs{Paper6B-CosmologyReadouts-v1_206}

  % Paper VI-C: complete technical proof and comparison owner.
  \BHExternalRefs{Paper6C-TechnicalMaterial-v1_50}

  % Paper VII-A: Finite QCD Kernel and Nonperturbative Closure.
\BHExternalRefs{Paper7A-QCDKernel-v1_172}

  % Paper VII-B: Hadronic, Thermal, and Nuclear QCD Readouts.
  \BHExternalRefs{Paper7B-QCDReadouts-v1_157}

  % Paper VII-C: Continuous Quantum Field Theory.
  \BHExternalRefs{Paper7C-ContinuousQFT-v2_202}

  % Paper VII-D: complete technical material for VII-A through VII-C.
  \BHExternalRefs{Paper7D-TechnicalMaterial-v1_53}

\fi

% END INLINE SOURCE: Paper00-ActiveExternalReferences-v2_130.tex


\hypersetup{hidelinks}
\setlength{\emergencystretch}{3em}
\title{\textbf{Born--Hopf Geometry in a Closed Universe IV-C: Technical Material} \\[0.5em]\large Preserved Algebraic, Spectral, Euclidean, and Finite Source-Response Calculations}
\author{Peter Kugelmann\thanks{\texttt{pkugelmann@proton.me}}}
\date{Version 1.50 --- 19 November 2026 \\ Framework v2.379}
\begin{document}
\maketitle
\begin{abstract}
Paper IV-C contains the complete technical and optional mathematical calculations redistributed from the former unified Paper IV. Original proof, definition, theorem and numbered-equation bodies are preserved without altering their hypotheses. The calculations support the quantum structure of Paper IV-A and the observable response of Paper IV-B but introduce neither an independent physical Hamiltonian nor a continuum-field or detector-occurrence certificate. The v2.254 appendix proves finite charge covariance transfer, static spectral noise retention and conditional round-sphere spin-two inversion.
\end{abstract}
\setcounter{tocdepth}{2}
\tableofcontents
\clearpage
\noindent\textit{Original-source reference convention.} In unchanged inherited passages, Part~I refers to the finite quantum-core material now owned by Paper~IV-A. Former appendix references may resolve to the technical companion Paper~IV-C; their theorem labels and hypotheses remain unchanged.

\section[Finite algebra, group-word and angular-momentum calculations]{Finite algebra, group-word and\\ angular-momentum calculations}
\label{p4c:sec:v238-c1}
\noindent\textit{Source and claim boundary.} These are original extended calculations. Statements on physical sources, measurements or continuum limits retain the inherited antecedents of Papers IV-A/IV-B.

\subsection{Parent group-word closure versus the fixed quantum carrier}
\label{p4:sec:v188-group-word-carrier-firewall}

\begin{theorem}[Relation closure does not enlarge $m_p=1$]
\label{p4:thm:v188-relation-closure-not-carrier}
The parent operations retained at $M_8$---multiplication, inverse, Cartan
squaring and commutator---act on group roles and on the descendant
observable/relation algebra.  Their restriction to the finite momentum
supports does not by itself add vectors to the declared level-one carrier
$\mathcal H_1$ and therefore does not alter
\[
 \boxed{m_p=1.}
\]
On the Hopf circle the Cartan square sends the primitive character
$e^{i\alpha}$ to $e^{2i\alpha}$; on the tetrahedral abelian residue
$C_3\cong2T/Q_8$ it sends a generator to its inverse.  These are exact
quadratic parent relations.

The existing comparison theorem remains the proof that
\[
 \boxed{m_{\rm cmp}=2}
\]
is the unique first comparison/action band: the clock algebra requires
$m\ge2$ and exact quartic $2I$ evaluation requires $m\le2$.  The parent square
explains why degree two is the first intrinsic relation layer; it does not
replace that two-sided uniqueness proof and does not define a second momentum
truncation.
\end{theorem}

\begin{proof}
The imported momentum-host and carrier-rigidity theorems of Paper~II-B,
Theorems~\ref{p2:thm:momentum-host-H1}
and~\ref{p2b:thm:momentum-carrier-rigidity}, fix the carrier
$\mathcal H_p\cong\mathcal H_1$; $m_p=1$ is its derived constructor-length
label.  No new representation truncation is declared here.  Under the
inherited typed constructor grammar, restricting the parent words acts on
support elements, relations and observable products without adjoining
representation sectors.  The square formulas are direct.
The value $m_{\rm cmp}=2$ then follows from the previously proved independent
lower bound from the complete clock algebra and upper bound from exact $2I$
quartic evaluation.
\end{proof}

% BEGIN v2.60 DEGREE-TWO / FINITE-INTERFACE CLOSURE
\begin{theorem}[Degree-two fixed-carrier covariant saturation]
\label{p4:thm:v260-degree-two-covariant-saturation}
For every fixed $\Spin(4)$-equivariant linear map
\[
 \Psi:\mathcal P_x(m)\longrightarrow\End(\mathcal H_1),
\]
Paper~II-C implies
\[
 \Psi(E_k^x)=0\qquad(k\ge3),
 \qquad
 \operatorname{rank}\Psi\le14.
\]
The explicit unital completely positive covariant map $\Phi$ of Paper~II-C
attains this bound and satisfies, for every $m\ge2$,
\[
 \boxed{
 \operatorname{rank}\Phi|_{\mathcal P_x(m)}=14,
 \qquad
 \ker\Phi|_{\mathcal P_x(m)}=\bigoplus_{k=3}^{m}E_k^x.}
\]
Consequently
\[
 \boxed{
 \Phi(\mathcal P_x(m))=\Phi(\mathcal P_x(2))\qquad(m\ge2),}
\]
and the maximally visible fixed-carrier covariant fraction obeys
\[
 \boxed{
 \frac{14}{\dim\mathcal P_x(k)}
 =\frac{84}{(k+1)(k+2)(2k+3)}\longrightarrow0.}
\]
Thus growth of reciprocal-position resolution beyond degree two cannot add a
new scalar component to any fixed $\Spin(4)$-covariant linear comparison into
the frozen momentum carrier.
\end{theorem}

\begin{proof}
This is the representation-theoretic content of
Proposition~\ref{p2c:prop:no-faithful-allband-xp-map} together with the sharp
positive witness of Corollary~\ref{p2c:cor:optimal-povm-harmonic-kernel}.
Existence of $\Phi$ does not select a physical measurement and does not make
$\Phi$ multiplicative.
\end{proof}

\begin{corollary}[Fourfold characterization of the comparison degree]
\label{p4:cor:v260-degree-two-fourfold-characterization}
The value $m_{\rm cmp}=2$ is now characterized independently by all four of
the following facts:
\begin{enumerate}[label=(\roman*),leftmargin=2.1em]
\item $\operatorname{Sym}^2(E_1^x)\cong E_0^x\oplus E_2^x$;
\item it is the unique clock/quartic comparison band already proved above;
\item it is the highest complete scalar band visible faithfully to the frozen
      $\mathcal H_1$ carrier through a fixed covariant linear comparison;
\item on the sharp-frontier branch, $\tau_2$ is the first event at which this
      complete 14-dimensional band is physically retained.
\end{enumerate}
At $\tau_2$, Paper~I-A's map
\[
 \mathsf S_{2T,2}:\mathcal P_x(2)\hookrightarrow\mathbb C^{24}
\]
is an exact isometric sampling embedding.  Hence the whole retained
comparison band is finitely reconstructible on the inherited Hurwitz support
once $\tau_2$ is reached.  The support is a comparison/sample support, not a
microscopic position lattice.
\end{corollary}

\begin{proposition}[Finite $x/p$ interface no-go and shared-support taxonomy]
\label{p4:prop:v260-finite-xp-interface-taxonomy}
The current finite interface results impose the following simultaneous
constraints.
\begin{enumerate}[label=(\alph*),leftmargin=2.1em]
\item Although $\dim\mathcal H_1=\dim\mathcal E_{15}^{\rm BM}=15$, every
natural left-$2T$-equivariant linear map
$\mathcal H_1\to\mathcal E_{15}^{\rm BM}$ has
\[
 \boxed{\operatorname{rank}\le5,\qquad\dim\ker\ge10,}
\]
and, more sharply,
\[
 \boxed{V_3\subseteq\ker F,}
\]
where $V_3$ is the restriction of the left chiral adjoint block $V_{1,0}$.
Thus equal dimensions do not supply a natural invertible Born exchange, and
the natural scalar Hurwitz target is structurally blind to an entire chiral
adjoint sector.
\item A normalized positive classical first-moment realization of the
intrinsic coordinate tuple has least outcome radius
\[
 \boxed{r_{\min}=\sqrt3\,R_{\rm mom}.}
\]
On the undeformed radius-$R_{\rm mom}$ ball, exact scaled moments $tX_a^D$
exist iff $|t|\le1/\sqrt3$; exact same-radius moments are impossible.
\item The same 24-point Hurwitz support can nevertheless carry both the
faithful scalar sampling of $\mathcal P_x(2)$ and the exact intrinsic-coordinate
POVM at radius $\sqrt3R_{\rm mom}$.  Shared support does not identify the two
maps, carriers, traces, or radii.
\item The 14-dimensional scalar comparison band
$E_0\oplus E_1\oplus E_2$ is not the 16-dimensional quadratic coordinate-
operator image
$\mathbf1\oplus(1,1)\oplus\Lambda^2_{\mathbb C}$.  Their irreducible content
is different, so the word ``quadratic'' supplies no $\Spin(4)$-equivariant
identification.
\end{enumerate}
Any stronger $\mathrm{PAIR}_{xp}^{\rm op}$ proposal must therefore introduce
a differently typed target/action, a calibrated dilation or unsharpness, or
additional licensed structure; it cannot be inferred from dimension equality
or common Hurwitz support alone.
\end{proposition}

\begin{proof}
Part (a) is Paper~II-E, Proposition~\ref{p2e:prop:sharp-natural-2t-rank} and Corollary~\ref{p2e:cor:v261-natural-hurwitz-v3-blindness}.  Part (b) is Paper~II-C,
Theorem~\ref{p2c:thm:optimal-intrinsic-povm-radius} and
Proposition~\ref{p2c:prop:exact-intrinsic-classical-moment-no-go}.  Part (c)
combines Paper~I-A, Corollary~\ref{p1:cor:v185-2t-degree-two-sampling-isometry},
with Paper~II-C, Corollary~\ref{p2c:cor:optimal-hurwitz-povm}.  Part (d) is the
representation decomposition recorded in Paper~I-A,
Remark~\ref{p1:rem:symmetric-vs-clifford-quadratic}.
\end{proof}
\begin{proposition}[Intrinsic localization bounds and measurement-map transport]
\label{p4:prop:intrinsic-localization-moment-transport}
For every density representative \(\sigma\) on the complete generated
intrinsic fifteen-dimensional carrier,
\[
 \boxed{\sum_{a=1}^4\operatorname{Var}_\sigma(X_a^D)
 \ge\frac{1823}{4608}R_p^2.}
\]
The bound is sharp.  It concerns the sum of four variances and does not
give a positive lower bound for each coordinate separately.  On a protected
Clifford block \(X_a=c\Gamma_a\), the stronger complete symmetric
covariance identity is
\[
 V_{ab}:=\tfrac12\Tr\!\left(\sigma
 \{X_a-\mu_aI,X_b-\mu_bI\}\right)
 =c^2\delta_{ab}-\mu_a\mu_b,
 \qquad \mu_a=\Tr(\sigma X_a),\qquad |\mu|\le c.
\]
For the intrinsic protected class \(c=\sqrt3R_p/4\); this is a
protected-block specialization rather than the full-carrier bound.

Let \(\mathsf Q_\theta\) range over the complete sphere-valued
\(\Spin(4)\)-covariant optimal POVM family of Paper~II-C, at outcome
radius \(r=\sqrt3R_p\).  Its first moments
\(\mathsf Q_\theta(x_a)=X_a^D\) and the intrinsic products
\(X_a^DX_b^D\) are fixed, while individual classical quadratic readouts
\(\mathsf Q_\theta(x_ax_b)\) may depend on the endpoint parameters.
Their summed noise is nevertheless fixed:
\[
 \sum_a\left[\mathsf Q_\theta(x_a^2)-(X_a^D)^2\right]
 =3R_p^2I-\sum_a(X_a^D)^2.
\]
Consequently a downstream quadratic prediction using the classical
readout must specify its map or prove parameter independence.  The sharp
radius and the continuous fixed-map rank-fourteen comparison result are
unchanged throughout this family; no multiplicativity of
\(\mathsf Q_\theta\) is implied.
\end{proposition}

\begin{proof}
The first statement is Paper~II-C,
Theorem~\ref{p2c:thm:intrinsic-total-variance}; the covariance identity is
Corollary~\ref{p2c:cor:protected-coordinate-covariance} and follows from
\(\{\Gamma_a,\Gamma_b\}=2\delta_{ab}I\).
The family and its quadratic qualification are
Theorem~\ref{p2c:thm:complete-intrinsic-endpoint-seeds} and
Proposition~\ref{p2c:prop:endpoint-fixed-map-qualification}.  Since each
outcome lies on the same radius-\(\sqrt3R_p\) sphere, unitality gives
\(\mathsf Q_\theta(\sum_a x_a^2)=3R_p^2I\), yielding the noise identity
of Corollary~\ref{p2c:cor:optimal-endpoint-total-noise}.
\end{proof}
% END v2.60 DEGREE-TWO / FINITE-INTERFACE CLOSURE

% BEGIN v2.61 RECIPROCITY-BOTTLENECK / DEGREE-TWO STRENGTHENING
\begin{corollary}[Further independent characterizations of the degree-two comparison event]
\label{p4:cor:v261-degree-two-strengthening}
On the sharp-frontier branch, the comparison event \(\tau_2\) has the
following additional independent characterizations.
\begin{enumerate}[label=(\roman*),leftmargin=2.1em]
\item It is the first event at which nonzero protected transverse
      Hodge--Dirac response is mathematically possible; before it, every
      reciprocal-odd nonconstant profile lies in one scalar Laplace shell and
      the protected second-order transverse term vanishes.
\item It is the first event at which the complete position Hopf-charge algebra
      contains all five momentum readout projectors:
      \[
       C^*(I,Q_H|_{\mathcal P_x(2)})\cong
       C^*(I,Q_H|_{\mathcal H_1})\cong\mathbb C^5.
      \]
      The inherited normalized trace weights remain different,
      \[
       \tau_{x,2}:\frac1{14}(3,2,4,2,3),\qquad
       \tau_p:\frac1{15}(1,3,7,3,1),
      \]
      so this algebraic coincidence is not a carrier identification.
\item The common right-diagonal \(SU(2)\)-singlet sector of
      \(\mathcal P_x(k)\otimes\mathcal H_1\) reaches its terminal dimension
      \(15\) from \(k=2\) onward.
\item Within the exceptional family \(2T,2O,2I\), simultaneous demand for the
      one-handed \(\mathcal H_1\) quotient and exact uniform integration of
      all quartic scalar-band products uniquely gives
      \[
       \boxed{\Gamma=2I,\qquad m_x=2.}
      \]
\end{enumerate}
These facts strengthen the comparison-degree justification without identifying
the 14-dimensional scalar band with the 15-dimensional momentum carrier or
actualizing an exceptional-group occurrence.
\end{corollary}

\begin{theorem}[Fixed-\(m_p=1\) reciprocity bottleneck]
\label{p4:thm:v261-fixed-momentum-reciprocity-bottleneck}
Let the physical momentum carrier remain \(\mathcal H_1\) while the complete
position band grows to \(\mathcal P_x(k)\), \(k\ge1\).  Then:
\begin{enumerate}[label=(\alph*),leftmargin=2.1em]
\item no complete position band is Hopf-\(U(1)\)-equivalent to the frozen
      momentum carrier,
      \[
       \boxed{\mathcal P_x(k)\not\cong_{U(1)_H}\mathcal H_1
       \quad\text{for every }k\ge1;}
      \]
\item the Hopf-diagonal matched dimension grows only quadratically,
      \[
       M_k:=\dim(\mathcal P_x(k)\otimes\mathcal H_1)^{U(1)_{H,\rm diag}}
       =O(k^2),\qquad
       \frac{M_k}{15\,\dim\mathcal P_x(k)}\sim\frac{3}{4k};
      \]
\item the common right-diagonal singlet dimension is
      \[
       \dim(\mathcal P_x(k)\otimes\mathcal H_1)^{SU(2)_{R,\rm diag}}
       =\begin{cases}12,&k=1,\\15,&k\ge2;\end{cases}
      \]
\item the full diagonal-\(\Spin(4)\) singlet dimension is exactly
      \[
       \boxed{\dim(\mathcal P_x(k)\otimes\mathcal H_1)^{\Spin(4)_{\rm diag}}=2;}
      \]
\item literal unequal-level Born exchange has the fixed common part
      \[
       \boxed{
       (\mathcal P_x(k)\otimes\mathcal H_1)
       \cap(\mathcal H_1\otimes\mathcal P_x(k))
       =\mathcal P_x(1)\otimes\mathcal P_x(1),
       \qquad\dim=25.}
      \]
\end{enumerate}
Thus the unrestricted mixed carrier grows as \(O(k^3)\), but increasingly
strong reciprocal/symmetry matching leaves \(O(k^2)\), then 15, then 2
channels, while literal exchange overlap stays 25-dimensional.  Increasing
\(m_x\) therefore adds predominantly one-sided reciprocal-position resolution;
it does not generate a growing supply of strongly reciprocal channels.
\end{theorem}

\begin{proof}
These are the all-level representation consequences of Paper~I-C,
Theorem~\ref{p1b:thm:all-level-chiral},
Proposition~\ref{p1b:prop:all-level-hopf},
Theorem~\ref{p1b:thm:mixed-diagonal}, and
Theorem~\ref{p1b:thm:unequal-exchange}.  The Hopf-support mismatch is already
visible at \(k=1\); at \(k=2\) the supports coincide but the multiplicities do
not, and for \(k\ge3\) the position band carries additional Hopf charges.
\end{proof}
% END v2.61 RECIPROCITY-BOTTLENECK / DEGREE-TWO STRENGTHENING



\begingroup
\small
\setlength{\LTleft}{0pt}
\setlength{\LTright}{0pt}
\setlength{\tabcolsep}{5pt}
\renewcommand{\arraystretch}{1.16}
\emergencystretch=2em
\hbadness=10000
\captionsetup{hypcap=false}

\captionof{table}[Downstream QFT support map]{Downstream QFT support map.}
\label{p4:tab:qft-support-map}

\begin{longtable}{@{}
  >{\raggedright\arraybackslash}p{0.17\linewidth}
  >{\raggedright\arraybackslash}p{\dimexpr0.83\linewidth-2\tabcolsep\relax}
@{}}

\toprule
\textbf{Paper}
& \textbf{Supported structure and status supplied by Paper~IV} \\
\midrule
\endfirsthead

\toprule
\textbf{Paper}
& \textbf{Supported structure and status supplied by Paper~IV} \\
\midrule
\endhead

\midrule
\multicolumn{2}{r}{\emph{Continued on next page}}\\
\bottomrule
\endfoot

\bottomrule
\endlastfoot

Paper~I-A \cite{Kugelmann2026a}
&
\textit{Supported structure:} Hopf-to-ray comparison, Reeb charge blocks,
protected finite cell, and, on the actualized or conditionally generated
\(\mathsf{Supp}_{2T}\) domain, Hurwitz-link internal supports interpreted
quantum mechanically.
\textit{Paper~IV status:} ordinary ray phase is projective, while the
degree-one associated character of
Theorem~\ref{p4:thm:hopf-ray-phase} compares it with the oriented Hopf line;
Reeb weights become exact or effective
superselection sectors by Corollary~\ref{p4:cor:hopf-reeb-superselection}; the
Paper-I-A Hurwitz-link octet, reduced septet, and trimer pair are admissible
finite internal supports for the Paper-IV Hilbert/Fock calculus only after
\(\mathsf{Supp}_{2T}\) and the relevant statistics/occupation selectors
have been supplied. \\

Paper~II \cite{Kugelmann2026b}
&
\textit{Supported structure:} kinematic constitutive map, finite resolution,
and reciprocal large-window language.
\textit{Paper~IV status:} finite carriers, tensor composition, spectral
evolution, and large-window canonical interfaces explain why standard
laboratory variables can be used without making them primitive. \\

Paper~III \cite{Kugelmann2026c}
&
\textit{Supported structure:} on an occurring closed-process domain, the
loop-derived \(\Dcl\) Euclidean spectral form; unconditionally, scalar-cover time,
the GP1-calibrated Lorentzian cone/sign, the conditional local second-order
free law, and finite spectral thresholds.
\textit{Paper~IV status:} finite Euclidean kernels reconstruct positive
Hamiltonians on their cyclic OS spaces, with same-carrier identification
requiring the semigroup-factorization criterion; the upstream dispersion
then becomes a finite one-particle Hamiltonian and threshold theorem. \\

Paper~V-A
&
\textit{Supported structure:} finite Standard--Model carrier, chiral
projectors, gauge-sector blocks, internal tensor factors, and field-language
interfaces.
\textit{Paper~IV status:} the generated finite-algebra calculus, the optional
\(\mathsf{FENV}/\mathsf{ENV}\) criteria, selector-subalgebra rule, tensor composition,
gauge projection, finite Fock calculus, and finite channels supply the quantum
foundation.  Chronologically, Paper~IV retains generic \(\mathsf{1P}\) and
\(\mathsf{STAT}\) families; on Paper~V-A's independently constituted visible-matter
branch the one-particle carrier
\((S_L\otimes F_L)\oplus(S_R\otimes F_R)\) is \textsc{derived-later}, and on
its stated Lorentzian positive-energy spin--statistics domain the Fermi branch
is likewise \textsc{derived-later}.  These later discharges do not retroactively
make the generic Paper-IV families unique. \\

Paper~V-B
&
\textit{Supported structure:} finite particle spectra, mass operators,
rank-three mixing blocks, Yukawa representatives, and heavy-neutral carriers.
\textit{Paper~IV status:} spectral calculus, projector-overlap mixing,
positive block elimination, finite resonance reduction, and density-matrix
transport supply the quantum machinery; Paper~V-B owns the spectrum and flavor
specializations. \\

Paper~V-C
&
\textit{Supported structure:} gauge and interaction closure, electroweak
breaking, heavy-neutral decay channels, reduced thermal kinetics, and the
finite-trace lepto--baryogenesis readout.
\textit{Paper~IV status:} finite thermal blocks, Dyson expansions, GKLS
reductions, resonance kernels, and the finite-temperature seesaw interface
license the interaction and kinetic language; Paper~V-C owns the sector
couplings, source selector, and baryogenesis readout. \\

Paper~VI-A
&
\textit{Supported structure:} gravitational realization, homogeneous
backgrounds, scalaron and perturbation modes, reheating, and curved-background
particle production.
\textit{Paper~IV status:} stable-quadratic, Bogoliubov, fixed-product
open-system, and reheating-interface theorems establish the quantum status
of the supplied cosmological mode blocks on their displayed domains;
Paper~VI-A owns the background, stability data, couplings, and
gravitational sector. \\

Paper~VI-B
&
\textit{Supported structure:} cosmological transfer kernels, determinant
objects, late-time coefficients, and observational readout maps.
\textit{Paper~IV status:} CPTP transfer, CP selected branches, operator
effects, response functionals, and determinant kernels are kept as distinct
typed objects; their well-defined outputs are not new microscopic constants.
Paper~VI-B owns their cosmological construction and propagation. \\

Paper~VII-A
&
\textit{Supported structure:} finite \(SU(3)_c\) singlet projection,
Hamiltonian/transfer blocks, reflection-compatible correlators, and finite
QCD kernel observables.
\textit{Paper~IV status:} gauge projection, finite Gibbs/state calculus,
positive transfer, spectral evolution, and Euclidean positivity provide the
quantum foundation; Paper~VII-A owns the QCD kernel. \\

Paper~VII-B
&
\textit{Supported structure:} hadronic, thermal, nuclear, and validation
readouts over the Paper-VII-A kernel.
\textit{Paper~IV status:} states, spectral measures, thermal traces, matrix
elements, channels, and transfer operators are inherited finite-quantum
operations; Paper~VII-B owns only the QCD readout and export layer. \\

Paper~VII-C
&
\textit{Supported structure:} controlled regulator towers, OS/Wightman
reconstruction, continuum local nets, covariance, local fields, and continuum
gap transfer.
\textit{Paper~IV status:} the finite algebra, state/Born calculus, selector
restrictions, gauge projection, thermal identities, and CCR boundary defect
are imported rather than reproved.  Paper~VII-C owns the additional compactness,
positivity, locality, covariance, spectral, and sector certificates required
for the continuum limit. \\

\end{longtable}
\TableNotes{This downstream map records what Paper~IV licenses later papers to use.  It supplies permission to use field-theoretic language without silently adding continuum QFT as a new foundation.}
\endgroup

The downstream consequence is therefore sharp: when a later paper uses particles, fields, propagators, thresholds, decay channels, running, finite-temperature states, or Bogoliubov particle production, Paper~IV supplies the finite theorem-level foundation.  When a later paper needs physical sector identifications, beta functions, measured couplings, branching fractions, Boltzmann/density-matrix evolution, or cosmological likelihood propagation, those are not closed by Paper~IV and must be supplied in the relevant sector paper.  In particular, the Paper~V-B mass/Yukawa representatives and the Paper~V-C source-amplitude, interaction, and finite-trace baryogenesis outputs are not Paper~IV assumptions; they are later sector-specific readouts inserted into the Paper~IV kinetic interface.



\paragraph{Eight framework-native QFT strengthening modules}
\label{p4:sec:eight-native-qft-modules}

This subsection expands the bridge dictionary into the eight framework-native QFT modules that produce the largest downstream gain.  Each module is deliberately finite-first.  The theorem-level part is an algebraic or spectral statement inside a finite carrier.  The continuum-QFT words attached to it--renormalization, interaction, thermal field theory, scattering, Euclidean path integral, gauge field, CCR, or Wick calculus--are representation or interface language, not new primitive axioms.

\paragraph{Module 1: finite resolution coarse-graining and threshold matching}
\label{p4:sec:module-finite-rg-thresholds}

Let \(A_N=\End(\mathcal H_N)\) be a resolved finite observable algebra.  A
nested-resolution interface supplies coherent isometries
\[
V_{NM}:\mathcal H_M\longrightarrow\mathcal H_N,
\qquad
V_{NL}=V_{NM}V_{ML}\quad(N\ge M\ge L).
\]
The canonical Heisenberg coarse-graining associated with these data is the
unital completely positive compression \(X\mapsto V_{NM}^*XV_{NM}\).
It is not, merely by being a compression, a conditional expectation onto a
subalgebra of \(A_N\); that stronger term requires a specified unital
subalgebra embedding and the bimodule property.  In a nested spectral
realization whose low/high split is supplied by the upstream resolution
theorem or by an explicitly declared interface domain, write
\(P_L=V_{NL}V_{NL}^*\) for the lower-resolution projector and
\(P_H=\mathbf 1-P_L\).  Neither \(P_L\) nor the block generator below selects
a physical sector; they express the already supplied resolution split.  For
a finite self-adjoint block
\[
H_N=\begin{pmatrix}
H_{LL} & H_{LH}\\
H_{HL} & H_{HH}
\end{pmatrix},
\]
the native threshold operation is not subtraction of an ultraviolet infinity.  It is finite compression, possibly corrected by a Schur complement when the eliminated sector couples virtually to the retained sector.


\begin{theorem}[Coherent finite compression and threshold matching]
\label{p4:thm:finite-resolution-rg-threshold-matching}
Let \(A_N=\End(\mathcal H_N)\) be a nested finite resolution family and let
the coherent isometries \(V_{NM}\) above be supplied by its interface domain.
Define
\[
\mathcal E_{M\leftarrow N}:A_N\longrightarrow A_M,
\qquad
\mathcal E_{M\leftarrow N}(X)=V_{NM}^*XV_{NM}.
\]
Then the maps are unital and completely positive and satisfy
\[
\mathcal E_{L\leftarrow M}\circ \mathcal E_{M\leftarrow N}=\mathcal E_{L\leftarrow N}.
\]
Hence the supplied coherent resolution system forms an exact finite
coarse-graining projective family without inserting a regulator profile.
If, additionally, unital embeddings \(j_{NM}:A_M\hookrightarrow A_N\) are
supplied and \(\mathcal E_{M\leftarrow N}\) is
\(j_{NM}(A_M)\)-bimodular, then \(j_{NM}\mathcal E_{M\leftarrow N}\) is a
conditional expectation onto that embedded subalgebra.  No such
conditional-expectation claim is made without these extra data.

If \(H_N\) is a self-adjoint element of the supplied finite dynamics domain,
\(z\notin\spec(H_{HH})\), and
\[
z-H_{LL}-H_{LH}(z-H_{HH})^{-1}H_{HL}
\]
is invertible, equivalently \(z\in\rho(H_N)\cap\rho(H_{HH})\), then the retained-sector resolvent is exactly
\begin{equation}
\label{p4:eq:finite-rg-schur-resolvent}
P_L(z-H_N)^{-1}P_L
=
\bigl(z-H_{LL}-\Sigma_H(z)\bigr)^{-1}P_L,
\qquad
\Sigma_H(z)=H_{LH}(z-H_{HH})^{-1}H_{HL}.
\end{equation}
Consequently, at a real matching point \(\mu\notin\spec(H_{HH})\), the finite effective retained operator is
\begin{equation}
\label{p4:eq:finite-rg-effective-hamiltonian}
H_{\eff}^{(L)}(\mu)
=
H_{LL}+H_{LH}(\mu-H_{HH})^{-1}H_{HL},
\end{equation}
with the sign convention fixed by \eqref{p4:eq:finite-rg-schur-resolvent}.  It is self-adjoint, and its resolvent interpretation is valid away from the poles and Schur-complement spectrum.  Expanding \(H_{\eff}^{(L)}(\mu)\) in any retained operator basis \(\{\mathcal O_a^{(L)}\}\),
\begin{equation}
\label{p4:eq:finite-rg-running-coordinates}
H_{\eff}^{(L)}(\mu)=\sum_a g_a^{(L)}(\mu)\mathcal O_a^{(L)},
\end{equation}
defines finite running representatives \(g_a^{(L)}\).  These representatives are coordinate data of the same finite operator under resolution change; their existence is theorem-level, while their numerical beta functions are model-specific.
\end{theorem}


\begin{proof}
Each map \(X\mapsto V_{NM}^*XV_{NM}\) is unital because
\(V_{NM}^*V_{NM}=\mathbf1_{\mathcal H_M}\), and it is completely positive
because it has one Kraus operator.  Coherence gives
\[
\mathcal E_{L\leftarrow M}(\mathcal E_{M\leftarrow N}(X))
=V_{ML}^*V_{NM}^*XV_{NM}V_{ML}
=V_{NL}^*XV_{NL},
\]
which is the projective composition law.  The stated bimodule condition is
exactly the additional criterion making the embedded map a conditional
expectation.

The resolvent identity is the finite block inverse formula for \(z-H_N\).
Since \(z-H_{HH}\) and the retained Schur complement are invertible,
eliminating the high block gives
\(z-H_{LL}-H_{LH}(z-H_{HH})^{-1}H_{HL}\).  The expansion in a finite
retained basis is ordinary finite-dimensional linear algebra.  No continuum
regulator or subtraction procedure is used.
\end{proof}

\paragraph{Module 2: finite interacting mode Hamiltonians and Dyson calculus}
\label{p4:sec:module-finite-interactions}

A framework-native interaction is a self-adjoint finite operator polynomial in resolved creation and annihilation maps, for example
\begin{equation}
\label{p4:eq:finite-interacting-mode-polynomial}
V_N
=
\sum_{ij}m_{ij}a_i^\dagger a_j
+\sum_{ijkl}\lambda_{ijkl}a_i^\dagger a_j^\dagger a_k a_l
+\sum_\alpha \bigl(y_\alpha\mathcal Y_\alpha+y_\alpha^*\mathcal Y_\alpha^*\bigr),
\end{equation}
where each displayed term is an operator on the finite Fock carrier and the coefficients obey the adjointness relations making \(V_N=V_N^*\).  The interaction is therefore not formal: it is a finite matrix.


\begin{theorem}[Finite-window interacting QFT calculus]
\label{p4:thm:finite-interacting-qft-calculus}
Let \(H_N=H_{0,N}+V_N\) be a finite self-adjoint Hamiltonian on a resolved Bose/Fermi Fock sector, with \(V_N\) a self-adjoint finite mode polynomial.  Then:
\begin{enumerate}
\item \(U_N(t)=e^{-itH_N}\) exists globally and is unitary.
\item For \(t\ge0\), the interaction-picture propagator is given by the
norm-convergent Dyson series
\begin{equation}
\label{p4:eq:finite-native-dyson-series}
U_{I,N}(t)
=
\mathbf 1+
\sum_{m\ge1}(-i)^m
\int_{0\le t_m\le\cdots\le t_1\le t}
V_{I,N}(t_1)\cdots V_{I,N}(t_m)\,dt_1\cdots dt_m .
\end{equation}
For \(t<0\), it is given by the reverse-ordered Picard simplex,
equivalently by the signed time-ordered exponential
\(\mathcal T\exp(-i\int_0^tV_{I,N}(\tau)\,d\tau)\).
\item Every finite correlation function
\[
\langle\Omega_N,A_1(t_1)\cdots A_k(t_k)\Omega_N\rangle
\]
has a well-defined finite-matrix perturbative expansion obtained by inserting the norm-convergent Dyson series \eqref{p4:eq:finite-native-dyson-series}.  At every fixed perturbative order the coefficient is a finite matrix expression, and at fixed \(N\) no ultraviolet divergence occurs; diagrammatic rules are bookkeeping for finite matrix elements.
\end{enumerate}
\end{theorem}


\begin{proof}
Finite-dimensional self-adjointness gives a global unitary exponential by the
spectral theorem.  The interaction-picture propagator solves
\[
  \dot U_{I,N}(t)=-iV_{I,N}(t)U_{I,N}(t)
\]
with bounded norm-continuous coefficient.  Picard iteration gives
\eqref{p4:eq:finite-native-dyson-series}.  The series is bounded by the
exponential of
\[
  \int_{\min(0,t)}^{\max(0,t)}\|V_{I,N}(s)\|\,ds .
\]
Inserting the expansion into finite matrix elements gives the stated
perturbative correlation calculus.
\end{proof}

\begin{remark}[What becomes theorem-level]
\label{p4:rem:finite-interaction-status}
The theorem-level object is the interacting finite operator theory.  The continuum names ``vertices'', ``loops'', ``counterterms'', and ``Feynman graphs'' are optional coordinates for finite sums until a continuum or thermodynamic limit is taken.  This is the native support for decay amplitudes, CP-violating interference terms, finite effective potentials, and the weak-coupling kernels used later in Papers~V-A--V-B.
\end{remark}


% BEGIN INLINE SOURCE: Paper4-FiniteSourceResponse-Main-v1_1.tex
% Framework v2.160: complete finite response and licensed reduced-source transport.
% Framework v2.156: positive-chain finite source-response theorem.
\paragraph{Exact finite interacting response to a licensed source history}
\label{p4:sec:v4104-finite-source-response}

The finite interacting calculus also determines the response to an already
constituted geometric or external source.  This consumes an energy-complete
Hamiltonian family, not only an observable-flow generator modulo its centre.
Metric-dependent central energies and the inherited vacuum allocation remain
part of that family.  An actual source family, state and kinetic transport are
antecedents; finite dimensionality does not manufacture them.

\begin{theorem}[Complete finite interacting source response]
\label{p4:thm:v4104-finite-interacting-source-response}
Let $I=[t_i,t_f]$ and let $H(t;\phi)$ be an admitted norm-continuous
self-adjoint Hamiltonian on a fixed-rank finite physical source algebra,
twice continuously differentiable in its licensed source coordinates.
Use a coherent orthonormal kinetic frame and the occurring initial state.
The finite domain means a finite CAR carrier or an actually finite bosonic
occupation truncation; finite one-particle rank alone does not bound bosonic
Fock occupations.  Retain the complete action pullback and define
\[
 D_a=\partial_aH,\qquad Q_{ab}=\partial_a\partial_bH.
\]
For a nonlinear parent map $q=f(\phi)$ this includes
\[
 Q_{ab}=H_{,AB}f^A_{,a}f^B_{,b}+H_{,A}f^A_{,ab}.
\]
The inherited unitary $U(t,s)$ obeys $i\hbar\partial_tU=\bar H(t)U$ and
$\bar\rho(t)=U(t,t_i)\rho_iU(t,t_i)^\dagger$.
For $v=\delta\phi$, the complete density tangent satisfies
\[
 \mathscr L\sigma=-\frac{i}{\hbar}[D_b,\bar\rho]v^b+q,
 \qquad \mathscr L=\partial_t+\frac{i}{\hbar}\operatorname{ad}_{\bar H}.
\]
Its unique zero-initial-data retarded inverse is
\[
 \boxed{(\mathscr L_R^{-1}f)(t)=
 \int_{t_i}^{t}U(t,s)f(s)U(t,s)^\dagger\,ds.}
\]
Here $q$ is an independently owned source-history tangent, not an inserted
collision term.  Let
\[
 \sigma_{\rm ind}=U(t,t_i)\sigma_iU(t,t_i)^\dagger+
                  \mathscr L_R^{-1}q.
\]
For $F_a=\operatorname{Tr}(\rho D_a)$ one has the exact linear response
\[
 \boxed{\delta F_a(t)=\bar Q_{ab}(t)v^b(t)
 +\int_{t_i}^{t}\Pi^R_{ab}(t,s)v^b(s)\,ds
 +\operatorname{Tr}(D_a(t)\sigma_{\rm ind}(t)),}
\]
where $\bar Q_{ab}=\operatorname{Tr}(\bar\rho Q_{ab})$ and
\[
 \boxed{\Pi^R_{ab}(t,s)=-\frac{i}{\hbar}\Theta(t-s)
 \operatorname{Tr}\!\left(\bar\rho(s)
 [U(t,s)^\dagger D_a(t)U(t,s),D_b(s)]\right).}
\]
All admitted interactions remain inside $\bar H$, $D$ and $Q$.
\end{theorem}
\begin{proof}
Differentiate $\dot\rho=-i[H,\rho]/\hbar$ to obtain the tangent equation.
Finite continuous matrix ODE existence gives $U$.  Differentiating the ordered
integral gives $\mathscr L\mathscr L_R^{-1}f=f$ and zero initial value;
the homogeneous solution is unitary transport, proving uniqueness.
Differentiate $F_a$ and substitute this integral.  Trace cyclicity gives
\[
 \operatorname{Tr}(D_aU[D_b,\bar\rho]U^\dagger)
 =\operatorname{Tr}(\bar\rho[U^\dagger D_aU,D_b]),
\]
which proves the kernel and its sign.  The complete second source jet supplies
the instantaneous term.  No interaction-strength expansion has been used.
\end{proof}

\begin{corollary}[Finite causal domain and complete source outputs]
\label{p4:cor:v4104-finite-retarded-source-domain}
In trace norm,
\[
 \|\mathscr L_R^{-1}f\|_{C(I;1)}\le |I|\|f\|_{C(I;1)},\qquad
 |\Pi^R_{ab}(t,s)|\le\frac{2}{\hbar}\|D_a(t)\|\|D_b(s)\|.
\]
No energy gap or stationary commutator inverse is needed.  For every already
generated source observable $T_A$,
\[
 \delta\langle T_A\rangle
 =\operatorname{Tr}(\bar\rho\,\delta T_A)
  +\operatorname{Tr}(T_A\sigma).
\]
Thus pressure, stress and current responses, when those operators are owned,
are responses of the same density tangent.  Metric-only unitary variation with
fixed initial state preserves $\delta\operatorname{Tr}(\rho^k)=0$.
These are finite-time results, not infinite-time relaxation or continuum
uniformity claims.
\end{corollary}
\begin{proof}
Unitary conjugation is an isometry in trace norm.  The kernel estimate follows
from $\|[A,B]\|\le2\|A\|\|B\|$ and a normalized positive state.
Differentiate the complete output pairing.  Unitary conjugation preserves the
density spectrum and hence its trace powers.
\end{proof}

The kinetic metric and connection remain inherited.  In a moving orthonormal
frame $V$, the generator is
$\widetilde H=VHV^\dagger+i\hbar\dot VV^\dagger$; in a nonorthonormal frame
the adjoint is $X^\sharp=G^{-1}X^\dagger G$.  If source dependence of the
frame creates derivative couplings, use the complete functional variation:
\[
 \sigma[v](t)=\sigma_{\rm ind}(t)-\frac{i}{\hbar}
 \int_{t_i}^{t}U(t,s)[\delta\widetilde H[v](s),\bar\rho(s)]
 U(t,s)^\dagger ds.
\]
No derivative of a constraint projector or source normalization is discarded.
The theorem does not produce missing Hilbert-stress operators, a quantum
diffeomorphism Ward certificate, a physical boundary, or a new environmental
projector.  It closes the finite interacting source-response step for every
owner in its stated domain.


% BEGIN INLINE SOURCE: Paper4-FiniteReducedSourceResponse-Main-v1_0.tex
% Framework v2.160: derived finite source-response reduction.
\begin{corollary}[Complete finite source response under a licensed reduction]
\label{p4:cor:finite-reduced-source-response}
Use the energy-complete joint Hamiltonian family, occurring state, coherent
kinetic frame and owned source tangents of
Theorem~\ref{p4:thm:v4104-finite-interacting-source-response} on a compact
interval $I$.  Let $\mathcal A_J$ and $\mathcal A_S$ be the admitted
fixed-rank finite joint and reduced physical algebras, and suppose an already
licensed CPTP state reduction
\[
 \mathcal R_\phi(t):\mathsf{States}(\mathcal A_J)
 \longrightarrow\mathsf{States}(\mathcal A_S)
\]
is supplied in those inherited frames.  Require this family to be
norm-continuous in time, with source derivatives continuous to every order
consumed below.  If second state variations are
consumed, require the complete joint state family to be $C^2$ in those
directions.  A $C^2$ finite Hamiltonian evolution with fixed or $C^2$
initial/source inputs supplies that regularity on its admitted compact
domain.  Put $\rho_S=\mathcal R_\phi\rho_J$.

For the complete joint density tangent $\sigma_J[v]$ and source direction
$v$, the reduced tangent is
\[
 \boxed{\delta\rho_S[v]
 =\mathcal R_\phi\sigma_J[v]
  +D\mathcal R_\phi[v]\bar\rho_J.}
\]
For two source directions $u,v$, the complete second variation is
\[
\begin{aligned}
 D^2\rho_S[u,v]={}&\mathcal R_\phi D^2\rho_J[u,v]
 +D\mathcal R_\phi[u]D\rho_J[v]
 +D\mathcal R_\phi[v]D\rho_J[u]\\
 &+D^2\mathcal R_\phi[u,v]\bar\rho_J .
\end{aligned}
\]
All entries are evaluated at the background; no source derivative of a
reduction, constraint projector, kinetic frame or normalization is discarded.
In particular, if $\mathcal R$ is source-independent, the instantaneous-source
chart of the cited theorem gives
\[
\begin{aligned}
 \delta\rho_S[v](t)={}&\mathcal R(t)\sigma_{J,\mathrm{ind}}(t)\\
 &-\frac{i}{\hbar}\mathcal R(t)\int_{t_i}^t
 U(t,s)[D_b(s),\bar\rho_J(s)]U(t,s)^\dagger v^b(s)\,ds .
\end{aligned}
\]
For derivative-dependent source couplings use the cited theorem's complete
functional variation $\delta\widetilde H[v]$ in this integral.  For every
already owned reduced observable $T$,
\[
 \delta\langle T\rangle
 =\operatorname{Tr}_S(\delta T\,\bar\rho_S)
  +\operatorname{Tr}_S(T\,\delta\rho_S).
\]
Thus the finite reduced response, including its contacts and memory, is a
generated composition of the complete joint source owner and reduction.
It has no independent response normalization or inserted collision law.
For a source-independent reduction and continuous Hermitian forcing $f$,
\[
 \|\mathcal R(\cdot)\mathscr L_R^{-1}f\|_{C(I;1)}
 \le |I|\,\|f\|_{C(I;1)}.
\]
\end{corollary}
\begin{proof}
The joint finite source theorem supplies $\sigma_J$ and its retarded
integral.  The product rule applied to the licensed linear reduction gives
the first variation, and applying it twice gives the two mixed derivatives
and the second derivative of $\mathcal R_\phi$.  Differentiating the full
observable pairing gives the output formula.  For source-independent
$\mathcal R$, substituting the joint integral gives the displayed response.
If $X=X_+-X_-$ is Hermitian, positivity and trace preservation imply
\[
 \|\mathcal R X\|_1
 \le\operatorname{Tr}(\mathcal R X_+)+\operatorname{Tr}(\mathcal R X_-)
 =\|X\|_1.
\]
Combining this contraction with
Corollary~\ref{p4:cor:v4104-finite-retarded-source-domain} proves the causal
bound.  Paper~II-A generated-image and co-consumption retain the complete
licensed composition, including all source jets, without a second physical
occurrence or coefficient.
\end{proof}

The physical subsystem reduction and complete joint source owner remain
antecedents.  A fixed product-environment assignment, or a legitimate CP
assignment on its admitted domain, yields a channel on the allowed system
inputs.  An arbitrary correlated initial state does not itself supply that
assignment.  Exact finite response can retain memory and recurrence;
irreversible GKLS/kinetic or unbounded/continuum claims still require their
own domains and limiting certificates.  The corollary supplies neither a
metric Hilbert-stress operator, a quantum Ward identity nor an environmental
projector.  These boundaries agree with
Theorems~\ref{p4:thm:fixed-finite-bath-semigroup-rigidity} and
\ref{p4:thm:microscopic-bath-gkls}.

% END INLINE SOURCE: Paper4-FiniteReducedSourceResponse-Main-v1_0.tex


% END INLINE SOURCE: Paper4-FiniteSourceResponse-Main-v1_1.tex



% BEGIN INLINE SOURCE: Paper4-StationarySourceEnergy-Main-v1_0.tex
% Framework v2.164: finite stationary specialization of the complete source owner.


\subsection{Fuzzy \texorpdfstring{$\CP^1\simeq S^2$}{CP1=S2} and noncommuting coordinates}
\label{p4:sec:part2-fuzzycp1}

Local two-state or spin-like sectors are represented by fuzzy $\CP^1\simeq S^2$ blocks inside the finite algebra.
Their coordinate observables are matrix-valued and therefore noncommuting.
This provides the framework source of effective spin and local noncommutative coordinate observables.




\subsection{Diagonal \texorpdfstring{$\SU(2)$}{SU(2)} and bounded angular momentum}
\label{p4:sec:part2-angular}

Paper~II identifies the finite kinematic carrier
\[
\mathcal H_m
=
\bigoplus_{j_L+j_R\le m} V_{j_L}\otimes V_{j_R},
\]
together with commuting dimensionless chiral \(\mathfrak{su}(2)\) actions generated by \(L_i^\pm\),
\[
[L_i^\pm,L_j^\pm]=i\epsilon_{ijk}L_k^\pm,
\qquad
[L_i^+,L_j^-]=0.
\]
The physical rotation generators are the diagonal combinations
\begin{equation}
\label{p4:eq:Ji}
\hat J_i:=\hbar(L_i^+ + L_i^-),
\qquad
[\hat J_i,\hat J_j]=i\hbar\,\epsilon_{ijk}\hat J_k.
\end{equation}
Thus angular momentum in the framework is the diagonal \(\SU(2)\) content of the finite \(\Spin(4)\) carrier.

\begin{theorem}[Framework angular-momentum cutoff]
\label{p4:thm:angular-cutoff}
On each irreducible \(\Spin(4)\)-sector \(V_{j_L}\otimes V_{j_R}\subset\mathcal H_m\),
restriction to the diagonal copy of \(\SU(2)\) gives the Clebsch--Gordan decomposition
\begin{equation}
\label{p4:eq:diag-cg}
V_{j_L}\otimes V_{j_R}
\cong
\bigoplus_{J=|j_L-j_R|}^{j_L+j_R} V_J.
\end{equation}
Consequently, only angular-momentum sectors satisfying
\begin{equation}
\label{p4:eq:Jmax-m}
J\le m
\end{equation}
occur anywhere in \(\mathcal H_m\).
Equivalently,
\begin{equation}
\label{p4:eq:J2-bound}
\hat J^2 \le \hbar^2 m(m+1)\,\mathbf 1,
\qquad
\|\hat J_i\|\le \hbar m.
\end{equation}
\end{theorem}

\begin{proof}
Equation~\eqref{p4:eq:diag-cg} is the standard decomposition for the diagonal \(\SU(2)\) action on
\(V_{j_L}\otimes V_{j_R}\); see for example \cite{Hall2015lie}.
Since the carrier contains only sectors with \(j_L+j_R\le m\), the largest diagonal spin allowed by any summand is \(J_{\max}=m\).
The spectral bound for \(\hat J^2\) follows from its eigenvalues \(\hbar^2J(J+1)\) on \(V_J\), and the component bound follows because each \(\hat J_i\) has spectrum contained in \([-\hbar J,\hbar J]\) on \(V_J\).
\end{proof}

\begin{corollary}[Finite ladder structure]
\label{p4:cor:finite-ladders}
With \(\hat J_\pm:=\hat J_1\pm i\hat J_2\), every angular-momentum multiplet in \(\mathcal H_m\) is finite:
\begin{equation}
\label{p4:eq:Jpm-action}
\hat J_\pm\ket{J,M}
=
\hbar\sqrt{J(J+1)-M(M\pm1)}\,\ket{J,M\pm1},
\end{equation}
with
\begin{equation}
\hat J_+\ket{J,J}=0,
\qquad
\hat J_-\ket{J,-J}=0,
\qquad
J\le m.
\end{equation}
Thus every finite carrier supports only finitely many diagonal multiplets.
On the physically derived/co-realized compact momentum sector one has \(m_p=1\), so only the minimal microscopic multiplets occur there.
The extended orbital tower belongs instead to the emergent position-side canonical regime described in Section~\ref{p4:sec:part2-momentum-position-split}.
\end{corollary}

\begin{proof}
The ladder formula is the standard finite-dimensional \(\mathfrak{su}(2)\) action on the irreducible multiplet \(V_J\).  The coefficient vanishes at \(M=J\) for \(\hat J_+\) and at \(M=-J\) for \(\hat J_-\), so the ladder terminates.  The inequality \(J\le m\) is exactly the cutoff from Theorem~\ref{p4:thm:angular-cutoff}.
\end{proof}


\begin{corollary}[Physical compact momentum sector]
\label{p4:cor:physical-m1}
For the microscopic compact momentum sector selected by \eqref{p4:eq:m-momentum-physical},
\begin{equation}
\label{p4:eq:Jmax-m1}
J\le 1,
\qquad
\hat J^2\le 2\hbar^2\,\mathbf 1.
\end{equation}
Hence the diagonal rotational content of the co-realized physical momentum
carrier is minimal and finite by construction.
Higher orbital towers belong to the reciprocal large-position presentation, not to the microscopic compact momentum sector itself.
\end{corollary}

\begin{proof}
The physical compact momentum branch has \(m_p=1\) by \eqref{p4:eq:m-momentum-physical}.  Substituting \(m=1\) into Theorem~\ref{p4:thm:angular-cutoff} gives \(J\le1\) and therefore \(\hat J^2\le\hbar^2 1(1+1)\mathbf1=2\hbar^2\mathbf1\).
\end{proof}


Paper~II also gives the momentum generators
\begin{equation}
\label{p4:eq:Pi-angular}
\hat P_i=\frac{\hbar}{\Rmom}\,(L_i^+-L_i^-),
\qquad
[\hat J_i,\hat P_j]=i\hbar\,\epsilon_{ijk}\hat P_k,
\end{equation}
so the momentum triple transforms as a vector under the same diagonal rotation algebra.
On a fixed \((j_L,j_R)\)-sector one has
\begin{equation}
\label{p4:eq:JP-sector-norm}
\|\hat P_i\|=\frac{\hbar}{\Rmom}(j_L+j_R),
\qquad
\|\hat J_i\|=\hbar(j_L+j_R),
\end{equation}
and hence
\begin{equation}
\label{p4:eq:JP-relation}
\|\hat J_i\|=\Rmom\,\|\hat P_i\|.
\end{equation}
Therefore the angular-momentum cutoff is not an independent postulate; it is the exact representation-theoretic image of the compact momentum cutoff.



\subsection{Vector operators and finite selection rules}
\label{p4:sec:part2-selection}

The diagonal \(\SU(2)\) action immediately constrains which transitions can be induced by observables transforming as tensors.
For a vector operator, write spherical components \(V_q\) with \(q\in\{-1,0,1\}\), satisfying
\begin{equation}
\label{p4:eq:vector-operator}
[\hat J_3,V_q]=\hbar q\,V_q,
\qquad
[\hat J_\pm,V_q]=\hbar\sqrt{(1\mp q)(1\pm q+1)}\,V_{q\pm1}.
\end{equation}
This is the finite version of a rank-one tensor operator under rotations.

\begin{theorem}[Finite vector selection rules]
\label{p4:thm:vector-selection}
Let \(V_q\) be a vector operator for the diagonal \(\SU(2)\) action on \(\mathcal H_m\).
If a matrix element
\begin{equation}
\label{p4:eq:vector-me}
\bra{J',M'}V_q\ket{J,M}
\end{equation}
is nonzero, then necessarily
\begin{equation}
\label{p4:eq:selection-M}
M'=M+q,
\end{equation}
and
\begin{equation}
\label{p4:eq:selection-J}
J'\in\{J-1,J,J+1\}.
\end{equation}
Because the carrier is truncated, one also has
\begin{equation}
\label{p4:eq:selection-trunc}
0\le J,J'\le m,
\end{equation}
so all channels that would leave the finite tower are absent.
In particular, from the top multiplet \(J=m\) there is no upward \(J\to m+1\) transition.
\end{theorem}

\begin{proof}
The commutator with \(\hat J_3\) gives
\[
\hat J_3(V_q\ket{J,M})=(M+q)\hbar\,V_q\ket{J,M},
\]
which proves \eqref{p4:eq:selection-M}.
Since a vector operator transforms in the spin-one representation, the state \(V_q\ket{J,M}\) lies in the tensor product \(V_J\otimes V_1\).
By the Clebsch--Gordan rule,
\[
V_J\otimes V_1\cong V_{J+1}\oplus V_J\oplus V_{J-1},
\]
with the obvious omission of the nonexistent \(V_{-1}\) when \(J=0\).
Hence only \(J'\in\{J-1,J,J+1\}\) can occur.
The truncation \(J'\le m\) follows from Theorem~\ref{p4:thm:angular-cutoff}, so a \(J=m+1\) channel is absent from \(\mathcal H_m\).
\end{proof}

\begin{corollary}[Momentum and coordinate selection rules]
\label{p4:cor:momentum-selection}
The momentum components \(\hat P_q\) obtained from \eqref{p4:eq:Pi-angular} obey the vector selection rules
\begin{equation}
\bra{J',M'}\hat P_q\ket{J,M}\neq 0
\quad\Longrightarrow\quad
M'=M+q,
\qquad
J'\in\{J-1,J,J+1\},
\end{equation}
with the same truncation \(J'\le m\).
On the physical momentum carrier \(m_p=1\), this
leaves only the channels internal to the minimal compact multiplet content.
Likewise, in a local fuzzy \(\CP^1\simeq S^2\) block the coordinate triple transforms as a vector, so dipole-type transitions are finite and terminate at the edge of the available tower.
\end{corollary}

\begin{proof}
The momentum triple transforms as a diagonal vector by \eqref{p4:eq:Pi-angular}, so Theorem~\ref{p4:thm:vector-selection} applies directly to its spherical components.  The local fuzzy coordinate triple has the same vector transformation type, so the same Clebsch--Gordan selection rule and finite-edge truncation apply.
\end{proof}


\begin{remark}[Local fuzzy-harmonic cutoff]
\label{p4:rem:fuzzy-harmonic-cutoff}
For a local spin-\(j\) fuzzy \(\CP^1\simeq S^2\) block, the observable algebra is
\begin{equation}
\label{p4:eq:endvj-decomp}
\End(V_j)
\cong
V_j\otimes V_j^*
\cong
\bigoplus_{\ell=0}^{2j}V_\ell.
\end{equation}
Hence local spherical-harmonic content terminates at \(\ell_{\max}=2j\).
The framework therefore contains two related finiteness statements: a global carrier cutoff \(J\le m\) and a local fuzzy-harmonic cutoff \(\ell\le 2j\).
\end{remark}


\begin{corollary}[Selected \(\mathcal H_1\) coordinate graph]
\label{p4:cor:h1-coordinate-graph-selection-rule}
On the series-selected finite cell, the projected coordinate representation
\(X=(\frac12,\frac12)\) on
\[
\mathcal H_1=(0,0)\oplus(\tfrac12,\tfrac12)\oplus(1,0)\oplus(0,1)
\oplus(\tfrac12,0)\oplus(0,\tfrac12)
\]
has exactly the retained channels
\[
(0,0)\leftrightarrow(\tfrac12,\tfrac12),
\qquad
(\tfrac12,\tfrac12)\leftrightarrow(1,0),
\qquad
(\tfrac12,\tfrac12)\leftrightarrow(0,1),
\]
and
\[
(\tfrac12,0)\leftrightarrow(0,\tfrac12).
\]
Hence, for projected coordinate coupling, the derived/co-realized level-one carrier splits as
\[
  11\text{-dimensional orbital block}
  \quad\oplus\quad
  4\text{-dimensional protected spinor block}.
\]
\end{corollary}


\begin{proof}
This is the quantum selection-rule form of Paper~I-A, Proposition~\ref{p1:prop:math-h1-coordinate-coupling-graph}, equivalently Paper~II, Proposition~\ref{p2:prop:h1-coordinate-coupling-graph}.
It is obtained by tensoring each Spin(4) summand with \((\frac12,\frac12)\) and projecting back to \(\mathcal H_1\).
The present corollary uses the result only as a finite matrix-element selection rule for coordinate-type observables.
\end{proof}

\begin{remark}[No physical leakage claim]
\label{p4:rem:h1-coordinate-graph-status}
Corollary~\ref{p4:cor:h1-coordinate-graph-selection-rule} is an operator-selection statement at the selected finite carrier.
It does not identify the protected block with any later physical gauge or
fermion sector, and it does not use Paper~IV to derive \(m_p=1\); that
structural derivation is supplied upstream by Paper~II-B.  No physical
momentum occurrence needs to be added beyond the upstream transfer-operation
root; a particular momentum test or readout remains a G0 token.
\end{remark}



\subsection{Addition of angular momentum on composite systems}
\label{p4:sec:part2-addition}

For independent subsystems, Theorem~\ref{p4:thm:tensor-product} gives the composite carrier
\begin{equation}
\label{p4:eq:Hab-spin}
\mathcal H_{AB}=\mathcal H_A\otimes\mathcal H_B.
\end{equation}
If each subsystem carries a diagonal rotation algebra generated by \(\hat J_i^{(A)}\) and \(\hat J_i^{(B)}\), the total rotation generators on the composite are
\begin{equation}
\label{p4:eq:Jtot-AB}
\hat J_i^{(AB)}:=\hat J_i^{(A)}\otimes \mathbf 1_B + \mathbf 1_A\otimes \hat J_i^{(B)}.
\end{equation}
These again satisfy
\[
[\hat J_i^{(AB)},\hat J_j^{(AB)}]=i\hbar\,\epsilon_{ijk}\hat J_k^{(AB)}.
\]
Thus angular-momentum addition is not appended separately; it follows from tensor-product composition together with the diagonal \(\SU(2)\) action.

\begin{theorem}[Finite spin addition]
\label{p4:thm:finite-addition}
\leavevmode\par\noindent
Suppose subsystem multiplets \(V_{J_A}\subset\mathcal H_A\) and
\(V_{J_B}\subset\mathcal H_B\) occur in the two carriers.  Then under the total
generators \eqref{p4:eq:Jtot-AB},
\begin{equation}
\label{p4:eq:cg-composite}
V_{J_A}\otimes V_{J_B}
\cong
\bigoplus_{J=|J_A-J_B|}^{J_A+J_B} V_J.
\end{equation}
Equivalently, there is a coupled basis
\begin{equation}
\label{p4:eq:coupled-basis}
\ket{J,M;J_A,J_B}
=
\sum_{M_A+M_B=M}
C^{J,M}_{J_A,M_A;J_B,M_B}
\ket{J_A,M_A}_A\otimes\ket{J_B,M_B}_B,
\end{equation}
with Clebsch--Gordan coefficients \(C^{J,M}_{J_A,M_A;J_B,M_B}\), where
\begin{equation}
\label{p4:eq:composite-J-range}
|J_A-J_B|\le J\le J_A+J_B,
\qquad
M=-J,-J+1,\dots,J.
\end{equation}
If the subsystem towers are bounded by \(J_A\le J_A^{\max}\) and \(J_B\le J_B^{\max}\), then the composite tower is bounded by
\begin{equation}
\label{p4:eq:JAB-bound}
J\le J_A^{\max}+J_B^{\max}.
\end{equation}
In particular, if the two subsystem carriers arise from truncation levels \(m_A\) and \(m_B\), then
\begin{equation}
\label{p4:eq:mAB-bound}
J\le m_A+m_B.
\end{equation}
\end{theorem}

\begin{proof}
Equation~\eqref{p4:eq:cg-composite} is the usual Clebsch--Gordan decomposition for the diagonal action of \(\SU(2)\) on the tensor product of two irreducible multiplets; see for example \cite{Hall2015lie}.
The coupled basis \eqref{p4:eq:coupled-basis} is the associated change of basis from uncoupled to total-spin eigenstates.
The upper bound \eqref{p4:eq:JAB-bound} follows because the largest value in \eqref{p4:eq:composite-J-range} is \(J_A+J_B\), and Theorem~\ref{p4:thm:angular-cutoff} gives \(J_A\le m_A\) and \(J_B\le m_B\) for subsystem carriers obtained from the framework truncations.
\end{proof}

\begin{corollary}[Finite singlet/triplet and higher coupled sectors]
\label{p4:cor:singlet-triplet}
Two spin-\(\tfrac12\) sectors satisfy
\begin{equation}
\label{p4:eq:half-half}
V_{1/2}\otimes V_{1/2}
\cong
V_1\oplus V_0,
\end{equation}
so the triplet and singlet sectors arise exactly.
More generally, every coupled multiplet remains finite, and there is no channel beyond the maximal value allowed by \eqref{p4:eq:JAB-bound}.
\end{corollary}

\begin{proof}
For \(J_A=J_B=\tfrac12\), the Clebsch--Gordan range \(|J_A-J_B|\le J\le J_A+J_B\) gives precisely \(J=0,1\), hence \(V_{1/2}\otimes V_{1/2}\cong V_0\oplus V_1\).  The general finiteness statement follows from the same finite Clebsch--Gordan range and the bound \eqref{p4:eq:JAB-bound}.
\end{proof}




\subsection{Finite recoupling for three or more subsystems}
\label{p4:sec:part2-recoupling}

For three independent subsystems, tensorial composition gives
\(
\mathcal H_{ABC}=\mathcal H_A\otimes\mathcal H_B\otimes\mathcal H_C
\)
and the total rotation generators are
\begin{equation}
\label{p4:eq:Jtot-ABC}
\hat J_i^{(ABC)}
=
\hat J_i^{(A)}\otimes\mathbf 1\otimes\mathbf 1
+
\mathbf 1\otimes\hat J_i^{(B)}\otimes\mathbf 1
+
\mathbf 1\otimes\mathbf 1\otimes\hat J_i^{(C)}.
\end{equation}
One may first couple \(A\) and \(B\), or first couple \(B\) and \(C\).
Both procedures span the same fixed total-spin sector and are related by a finite unitary recoupling matrix.

\begin{theorem}[Finite recoupling]
\label{p4:thm:finite-recoupling}
Fix subsystem spins \(J_A,J_B,J_C\) and a total multiplet \(V_J\subset V_{J_A}\otimes V_{J_B}\otimes V_{J_C}\).
The two coupled bases
\begin{equation}
\label{p4:eq:ab-coupling}
\ket{((J_AJ_B)J_{AB},J_C)JM}
\end{equation}
and
\begin{equation}
\label{p4:eq:bc-coupling}
\ket{(J_A,(J_BJ_C)J_{BC})JM}
\end{equation}
are related by a finite unitary transformation
\begin{equation}
\label{p4:eq:6j-recoupling}
\ket{(J_A,(J_BJ_C)J_{BC})JM}
=
\sum_{J_{AB}} U^J_{J_{AB}J_{BC}}
\ket{((J_AJ_B)J_{AB},J_C)JM},
\end{equation}
where the sum runs only over the finitely many values compatible with the triangle inequalities.
In standard normalization one may write
\begin{equation}
\label{p4:eq:6j-symbol}
U^J_{J_{AB}J_{BC}}
=
(-1)^{J_A+J_B+J_C+J}
\sqrt{(2J_{AB}+1)(2J_{BC}+1)}
\begin{Bmatrix}
J_A & J_B & J_{AB}\\
J_C & J   & J_{BC}
\end{Bmatrix}.
\end{equation}
All intermediate spins are bounded by the subsystem cutoffs, so every recoupling sum in the framework is finite.
\end{theorem}

\begin{proof}
This is the associativity of repeated Clebsch--Gordan decomposition.
The two bases arise from the two bracketings
\((V_{J_A}\otimes V_{J_B})\otimes V_{J_C}\) and \(V_{J_A}\otimes(V_{J_B}\otimes V_{J_C})\), which are canonically isomorphic because tensor products are associative.
Within a fixed total-spin sector the corresponding basis change is unitary.
Since each pairwise coupling admits only finitely many intermediate spins by \eqref{p4:eq:composite-J-range}, the recoupling matrix is finite-dimensional.
The explicit form \eqref{p4:eq:6j-symbol} is the standard Wigner--Racah parametrization.
\end{proof}

\begin{remark}[Framework significance of recoupling finiteness]
\label{p4:rem:finite-recoupling}
In ordinary continuum presentations, recoupling of angular momentum is already
discrete.  The framework sharpens this further: the allowed subsystem
multiplets themselves lie in finite towers.  Hence not only each individual
coupling channel, but the entire network of many-body recouplings, is cut off
by construction.  This bounded recoupling structure is the natural setting for
composite spin systems, finite molecular or nuclear multiplets, and later
internal multiplet couplings in the Standard Model sector.
\end{remark}



\section{CAR/CCR ordering words and quadratic-line verification}
\label{p4c:sec:v238-c2}
\noindent\textit{Source and claim boundary.} These are original extended calculations. Statements on physical sources, measurements or continuum limits retain the inherited antecedents of Papers IV-A/IV-B.

\subsection{A graded coordinate certifies its supplied occupation boundary}
\begin{proposition}[Graded field-port boundary extractor]
\label{p4:prop:v2174-graded-port-boundary}
On the standard complete total-occupation sublevel through $N\ge0$,
retain its constituted Bose grading $N_B$ and charged linear coordinate
$X=\chi(a_-^\dagger+a_+)$ with $\chi\ne0$. Then
\begin{equation}
 R_-:=\tfrac12(X+[N_B,X])=\chi a_-^\dagger,\qquad
 L_+:=\tfrac12(X-[N_B,X])=\chi a_+,\qquad
 \nu(R_-)=N+1,
 \label{p4:eq:v2174-graded-port-boundary}
\end{equation}
where $\nu(R)=\min\{s\ge1:R^s=0\}$. Thus $N=\nu(R_-)-1$ is
determined by this already generated marked matrix.
\end{proposition}
\begin{proof}
The degree equations $[N_B,a_-^\dagger]=a_-^\dagger$ and
$[N_B,a_+]=-a_+$ remain exact at the boundary, where the forbidden raising
map is zero on both sides. They separate the two summands of $X$.
Every $N+1$ successive raises vanish, whereas
$R_-^N|0_B\rangle=\chi^N\sqrt{N!}|0,N\rangle\ne0$.
This includes $N=0$, with $R_-^0=I$ and $\nu(0)=1$.
\end{proof}
The complete-initial-sublevel hypothesis is essential: on an arbitrary
retained set, the longest raising chain need not measure the largest degree.
This certificate determines a boundary from an already constituted operator;
it supplies no physical operator or boundary by itself. The Fock unit used
to test the matrix is not identified with the full interacting Gibbs ground.
The auxiliary geometric pair boundary is kept separate in
Appendix~\ref{p4:app:v2174-occupation-comparisons}.

% END INLINE SOURCE: Paper4-GradedOccupationBoundary-Main-v1_0.tex


% BEGIN INLINE SOURCE: Paper4-QuadraticOccupationOrdering-Main-v1_0.tex
% v2.174: exact weighted finite boundary; no physical ordering is chosen.


\subsection{The exact top edge of the quadratic ordering word}
\begin{lemma}[Products of compressed quadratures retain a weighted boundary]
\label{p4:lem:v2174-weighted-ordering-boundary}
On an already supplied complete total-occupation sublevel through $N$,
let $P_N$ be its top projector and diagonalize the actual $\Omega>0$.
Using the compressed quadratures of
Theorem~\ref{p4:thm:v245-finite-free-vacuum-covariance}, multiplication
after compression and compression of the parent quadratic product give
respectively
\begin{align}
 H_{\rm prod}^{(2)}={}&d\Gamma(\hbar\Omega)
   +\tfrac\hbar2\operatorname{tr}\Omega\,I
   -\tfrac\hbar2P_N(d\Gamma(\Omega)+\operatorname{tr}\Omega\,I)P_N,
   \nonumber\\
 H_{\rm comp}^{(2)}={}&d\Gamma(\hbar\Omega)
   +\tfrac\hbar2\operatorname{tr}\Omega\,I.
 \label{p4:eq:v2174-weighted-ordering-boundary}
\end{align}
\end{lemma}
\begin{proof}
Below the top degree the standard ladder products give the ordinary delta;
on the top degree the first creation is blocked. Hence, on the whole carrier,
\[
 [a_i,a_j^\dagger]=\delta_{ij}I
    -P_N(\delta_{ij}I+a_j^\dagger a_i)P_N.
\]
For each eigenfrequency $\omega_i$, the compressed quadratic word is
$\hbar\omega_i(a_ia_i^\dagger+a_i^\dagger a_i)/2$.
Summing the displayed commutator with weights $\omega_i$ proves the first
formula. The parent oscillator identity, compressed only after forming
its products, proves the second. Transport the spectral basis and the
occupation projector together to obtain the intrinsic statement.
\end{proof}
The actual complete action supplies its own ordering and central allocation.
Neither formula is chosen for an unevaluated physical member, and no
assumption $P_NP_0=0$ is made about its Gibbs support. The same calculation
with all frequency weights equal to one gives the corresponding action-moment
lifts $\hbar(N_B+dI/2)-\hbar P_N(N_B+dI)P_N/2$ and
$\hbar(N_B+dI/2)$.

% END INLINE SOURCE: Paper4-QuadraticOccupationOrdering-Main-v1_0.tex

% END C01 V2174 GRADED BOUNDARY AND ORDERING


\begin{proposition}[Finite thermal boundary correction for one truncated bosonic mode]
\label{p4:prop:v245-finite-thermal-boundary-correction}
For one oscillator of frequency \(\omega>0\) with occupation cap \(N\), let
\(x=e^{-\beta\hbar\omega}\).  Then
\[
 Z_N=\frac{1-x^{N+1}}{1-x},\qquad
 \bar n_N=\frac{x}{1-x}-\frac{(N+1)x^{N+1}}{1-x^{N+1}},
\]
and the top-level probability is
\[
 p_N=\frac{(1-x)x^N}{1-x^{N+1}}.
\]
With the finite ladder commutator
\([a,a^\dagger]=I-(N+1)P_N\), the exact thermal variances are
\[
 \boxed{\langle q^2\rangle_{\beta,N}
 =\frac{\hbar}{2\omega}
 \left(2\bar n_N+1-(N+1)p_N\right),}
\]
\[
 \boxed{\langle p^2\rangle_{\beta,N}
 =\frac{\hbar\omega}{2}
 \left(2\bar n_N+1-(N+1)p_N\right).}
\]
They reduce to the exact vacuum result as \(\beta\to\infty\), and to the
regular-CCR \(\coth(\beta\hbar\omega/2)\) factor only in the controlled
\(N\to\infty\) limit.  Hence no infinite-CCR thermal covariance is silently
used at fixed finite occupation capacity.
\end{proposition}

\begin{proof}
The finite geometric sum gives \(Z_N\); differentiating
\(-\partial_{\beta\hbar\omega}\log Z_N\) gives \(\bar n_N\).
The truncated ladder has
\(aa^\dagger=a^\dagger a+I-(N+1)P_N\).  Substitution into the standard
quadrature definitions gives the displayed factor.  The two limits follow
directly from \(x\to0\) and then from \(x^{N+1}\to0\), respectively.
\end{proof}


So the framework distinguishes sharply between fermionic and bosonic mode quantization:
fermionic CAR are exact on finite sectors, whereas bosonic CCR are recovered only away from the occupation boundary or in controlled large-cutoff regimes.
That distinction is exactly what a finite foundational theory should predict.



\subsection{Second-quantized observables and many-body Hamiltonians}
\label{p4:sec:part2-second-quantized}

Every one-body operator $A\in\End(\mathcal K)$ lifts canonically to the number-preserving many-body operator
\begin{equation}
\label{p4:eq:dGamma}
d\Gamma(A)
=
\sum_{i,j=1}^d A_{ij}\,\mathfrak a_i^\dagger \mathfrak a_j,
\end{equation}
where $\mathfrak a_i$ denotes either $c_i$ in the fermionic case or the truncated bosonic annihilation map in the bosonic case.
If $A=A^*$, then $d\Gamma(A)$ is a many-body observable. For general $A$, $d\Gamma(A)$ is a number-preserving operator; physical Hamiltonians built from quadratic or quartic expressions require the usual finite-dimensional self-adjointness conditions.
Likewise the total number operator is
\begin{equation}
\label{p4:eq:number-op}
N
=
\sum_{i=1}^d \mathfrak a_i^\dagger \mathfrak a_i.
\end{equation}
Thus the one-particle structural algebra lifts functorially to a many-particle operator algebra.

More general candidate quadratic and quartic generators have the form
\begin{equation}
\label{p4:eq:many-body-hamiltonian}
H_{\mathrm{mb}}
=
\sum_{i,j} h_{ij}\,\mathfrak a_i^\dagger \mathfrak a_j
+
\frac12\sum_{i,j} \bigl(\Delta_{ij}\,\mathfrak a_i^\dagger \mathfrak a_j^\dagger + \overline{\Delta}_{ij}\,\mathfrak a_j \mathfrak a_i\bigr)
+
\frac14\sum_{i,j,k,l} V_{ijkl}\,\mathfrak a_i^\dagger \mathfrak a_j^\dagger \mathfrak a_l \mathfrak a_k.
\end{equation}
The first term is the second quantization of a one-body generator, the pairing term describes effective Bogoliubov or condensate-type mixing, and the quartic term encodes two-body interactions.
All of them are legitimate self-adjoint finite elements on the mode carrier,
provided the coefficients satisfy the usual Hermiticity conditions.  They
represent actual Hamiltonians only on a supplied interaction/\(\mathsf{DYN}\)
domain.  They are structural only when generated or fixed by the constituted
carrier; otherwise they are interaction-domain data used by the theorem and
carry no upstream selecting power.


\begin{proposition}[One-particle data do not select exchange statistics]
\label{p4:prop:one-particle-does-not-select-statistics}
Let \(\mathcal K\) be a complex one-particle carrier with
\(d:=\dim\mathcal K\ge2\).  Both
\[
 \operatorname{Sym}^2\mathcal K
 \qquad\text{and}\qquad
 \Lambda^2\mathcal K
\]
are nonzero functorial constructions from the same \(\mathcal K\), have
dimensions \(d(d+1)/2\) and \(d(d-1)/2\), and carry respectively the
\(+1\) and \(-1\) characters of the factor transposition.  Therefore the
one-particle space, its dimension, and its complete unitary action do not
select a generic \(\mathsf{STAT}\) member.
\end{proposition}

\begin{proof}
The flip \(\tau(x\otimes y)=y\otimes x\) commutes with
\(U\otimes U\) for every \(U\in U(\mathcal K)\).  Its spectral projections
\((I\pm\tau)/2\) are consequently natural, and their ranks are the two
displayed positive dimensions.  Thus the same one-particle datum admits two
inequivalent nonzero exchange representations.  Selecting the realized one
requires additional geometry, such as a constituted exchange
connection/holonomy, or the full locality, covariance, positivity, and field
hypotheses of a spin--statistics theorem; choosing a branch for its convenient
cutoff behavior is not a derivation.
\end{proof}


\begin{theorem}[Finite second-quantization after one-particle and statistics selection]
\label{p4:thm:finite-second-quantization}
Let \(\mathcal K=\operatorname{Ran}P\) be an intrinsic finite mode block
generated by the constituted finite spectral, charge, chirality, grading, or
symmetry projectors of the Born framework, and supply a one-particle
assignment identifying \(\mathcal K\) with the \(n=1\) mode carrier.
Supply either the fermionic
sign-sector selector or the bosonic trivial-sector selector.  In the bosonic
case GP1 already forces finite occupation support; for the standard initial
sublevel used here additionally supply or derive
\(\mathsf P_{\rm gbd}(\mathrm{occ},N_{\max})\).
Then:
\begin{enumerate}[label=(\roman*)]
\item the many-particle carrier is a finite direct sum of statistics-selected tensor powers of $\mathcal K$;
\item every one-body observable $A\in\End(\mathcal K)$ induces a canonical many-body observable $d\Gamma(A)$;
\item every self-adjoint quadratic or quartic polynomial in the mode
operators defines a candidate self-adjoint many-body generator on that
carrier; it represents an actually realized Hamiltonian only on the
displayed interaction/\(\mathsf{DYN}\) domain;
\item if $\mathcal K\cong \mathcal K_{\mathrm{ext}}\otimes \mathcal K_{\mathrm{int}}$, then internal multiplets, family labels, and their unitary mixings lift modewise to the many-body theory.
\end{enumerate}
Therefore, after the named one-particle assignment, statistics selector,
and, on the bosonic initial-sublevel branch, the generated occupation-boundary
data have been supplied or derived, the passage to finite-mode field language
requires no further
algebraic quantization axiom.  It is a repackaging of the supplied or
later-derived one-particle/statistics/occupation structures together with the
already generated tensor and symmetry calculus.
\end{theorem}


\begin{proof}
Item~(i) is the reorganization of tensor powers by the supplied trivial- or
sign-sector selector.  Theorem~\ref{p4:thm:permutation-sector-decomposition}
shows that all nonzero Young sectors are algebraically available; Bose, Fermi,
and any other physical Young-sector realization require a selector of the same
logical kind.
Item~(ii) follows from the definition of $d\Gamma(A)$ in equation~\eqref{p4:eq:dGamma}.
Item~(iii) follows because any self-adjoint operator on a finite Hilbert
space mathematically generates a unitary group, and quadratic/quartic mode
polynomials are particular self-adjoint operators on the chosen carrier.
The \(\mathsf{DYN}\) role, rather than self-adjointness alone, assigns that
group actual dynamics.
For item~(iv), choose product bases on $\mathcal K_{\mathrm{ext}}$ and $\mathcal K_{\mathrm{int}}$; then the creation operators carry composite labels $\mathfrak a_{iA}^\dagger$.
Any internal unitary $U_{\mathrm{int}}$ acts by
\[
\mathfrak a_{iA}^\dagger
\mapsto
\sum_B (U_{\mathrm{int}})_{BA}\,\mathfrak a_{iB}^\dagger,
\]
which is exactly the lifted mixing action on the many-particle carrier.
No further algebraic postulate beyond the stated one-particle assignment,
statistics selector, and a bosonic generated-boundary datum is required for the standard initial-sublevel truncation.
\end{proof}




\subsection{Quadratic line reality and its faithful completion}

\begin{proposition}[Quadratic line reality]
\label{p4:prop:quadratic-line-reality}
Let \(N\) be a Hermitian complex line, with \(h\) antilinear in
its first argument, and let \(0\ne m\in(N^*)^{\otimes2}\).
Writing \(\sigma=\|m\|\), define
\[
 j_m=\sigma^{-1}R_h^{-1}m^\flat,\qquad
 m^\flat(v)(w)=m(v,w).
\]
Then \(j_m\) is antiunitary, \(j_m^2=1\), and
\[
 m(v,w)=\sigma h(j_mv,w).
\]
This is a natural equivalence between a nonzero symmetric line
form and a positive norm together with an antiunitary involution.
\end{proposition}
\begin{proof}
In a unit frame, write \(m(v,w)=\mu vw\).
Then \(j_m(z)=(\bar\mu/|\mu|)\bar z\), which is antiunitary,
squares to one and satisfies the displayed formula.  The
Riesz definition proves independence from the frame.  Conversely
the displayed formula reconstructs \(m\) and its norm.
\end{proof}

\begin{proposition}[Faithful conjugate-line completion]
\label{p4:prop:faithful-conjugate-line-completion}
Use the linear Hermitian identification \(N^*\simeq\bar N\)
to regard \(m^\flat\) as \(M_m:N\to\bar N\).  On
\(N\oplus\bar N\), set
\[
 D_m=\begin{pmatrix}0&M_m^\dagger\\M_m&0\end{pmatrix},
 \quad \Gamma=\begin{pmatrix}1&0\\0&-1\end{pmatrix},
 \quad J_0(v,\bar w)=(w,\bar v).
\]
Then \(D_m^\dagger=D_m\), \(\Gamma D_m=-D_m\Gamma\),
\(J_0^2=1\), \(J_0\Gamma=-\Gamma J_0\), and \(J_0D_m=D_mJ_0\).
Recovering \(m\) from the lower off-diagonal block is natural;
the completion has the same stabilizer as \(m\) under
\(g\oplus\bar g\).
\end{proposition}
\begin{proof}
In a unit frame the off-diagonal entries of \(D_m\) are
\(\bar\mu,\mu\), and \(J_0\) is conjugate swap.
Matrix multiplication gives all the identities.  The lower
block determines \(\mu\), proving recovery and equality of
stabilizers under the specified action.
\end{proof}

This completion applies in particular to the mathematical line
of Paper I-A when that line and its quadratic are the declared
inputs.  It is not a physical \(F\oplus\bar F\) construction or
a global \(J_F\).  If \(N\otimes G\) includes a family carrier,
\(m^\flat\otimes1_G\) lands in \(N^*\otimes G\), not
\((N\otimes G)^*=N^*\otimes G^*\).  A family bilinear is required
for the latter target; the line theorem supplies none.
% END M8 v1.78 FER02




\begin{corollary}[Generated gauge reduction supplies the finite Yang--Mills real-form domain]
\label{p4:cor:v174-yang-mills-reduction-real-form-domain}
On the local stationary-product domain of Paper~III-F,
Theorem~\ref{p3f:thm:v174-generated-joint-gauge-phase-space-action}, the
Gauss moment-map zero locus and its typed quotient supply the separately valid
gauge reduction required in
Theorem~\ref{p4:thm:stationary-quadratic-real-form-pair}.  Consequently every
stable finite physical-polarization block of that generated Yang--Mills action
inherits the complex-time Euclidean/Lorentzian quadratic pair and reciprocal
state/test dynamics without another gauge-reduction or Wick premise.  This
corollary does not construct the finite regulator action, path measure,
interacting continuum limit, state, or history.
\end{corollary}


\begin{remark}[Domain of the free dispersion theorem]
\label{p4:rem:dispersion-domain}
The affine relation in
Theorem~\ref{p4:thm:finite-one-particle-dispersion} is not a universal
selection rule for all interacting or effective Hamiltonians.  It applies
only to intrinsic finite blocks in the imported local second-order free
domain.  Interactions, higher-order effective operators, dissipative
coarse-graining, subluminal media, and sector-specific profiles belong to
their own domains and do not back-propagate into the GP1 cone theorem.
\end{remark}

\begin{corollary}[Finite spectral gaps as finite-volume energy thresholds]
\label{p4:cor:spectral-gap-threshold}
Let \(\mathcal K=\operatorname{Ran}P\) be an intrinsic finite spectral block satisfying the stability and isolation properties of Definition~\ref{p4:def:intrinsic-finite-spectral-block} for the reconstructed generator \(H_1\).  If \(\lambda_0\ge0\) is the lowest eigenvalue of \(A|_{\mathcal K}\), then the lowest one-particle energy in that block is
\begin{equation}
\label{p4:eq:threshold-from-gap}
  E_{\rm thr}
  =
  \sqrt{\hbar^2c^2\lambda_0+\mu^2c^4}.
\end{equation}
Within the stated local relativistic one-particle interpretation, one may
package this finite-volume threshold as
\begin{equation}
\label{p4:eq:effective-threshold-mass}
  m_{\rm thr}
  =
  \frac{E_{\rm thr}}{c^2}.
\end{equation}
In a sector with \(\mu=0\), any nonzero compact spatial spectral gap
\(\lambda_0>0\) therefore appears as a positive finite-volume energy
threshold, with the corresponding effective threshold parameter
\begin{equation}
\label{p4:eq:massless-gap-threshold}
  m_{\rm thr}
  =
  \frac{\hbar}{c}\sqrt{\lambda_0}.
\end{equation}
\end{corollary}

\begin{proof}
Theorem~\ref{p4:thm:finite-one-particle-dispersion} gives \(E_\lambda^2=\hbar^2c^2\lambda+\mu^2c^4\) for every eigenvalue of the positive spatial operator on the intrinsic block.  Stability gives invariance of the block under the free generator, and isolation makes its bottom energy a genuine separated threshold rather than a transient basis component.  Taking the lowest eigenvalue gives \eqref{p4:eq:threshold-from-gap}--\eqref{p4:eq:massless-gap-threshold}.
\end{proof}

\begin{remark}[One-way dependency from GP1 to finite dispersion]
\label{p4:rem:no-circularity-p3-dispersion}
\label{p4:rem:lorentzian-noncircular}
The dependency order is
\[
\begin{aligned}
&\text{derived clocked-Hopf doctrine and calibrated radius/time representatives}\\
&\qquad\Longrightarrow
\text{Paper-III operational cone and Lorentzian sign}\\
&\qquad\Longrightarrow
\text{Paper-III local second-order free law}\\
&\qquad\Longrightarrow
\text{Paper-IV positive finite Hamiltonian and thresholds}.
\end{aligned}
\]
The first two arrows require no quantum one-particle premise.  The third is
conditional on realizing the local second-order scalar sector.  Paper~IV
supplies the finite Born/Fock interpretation, intrinsic stable block, and
positive-frequency square root; it does not close an upstream sign theorem.
\end{remark}

\begin{corollary}[Downstream finite threshold export on the Lorentzian branch]
\label{p4:cor:closes-p3-lorentzian-input}
For every stable isolated intrinsic one-particle block satisfying
Theorem~\ref{p4:thm:finite-one-particle-dispersion}, the pair
\[
\left(
H_1=\sqrt{\hbar^2c^2A+\mu^2c^4\mathbf1},
\quad
E_{\rm thr}
=\sqrt{\hbar^2c^2\lambda_0+\mu^2c^4}
\right)
\]
is an exact finite quantum output on the upstream Lorentzian/local-free
branch.  Later papers may import this Hamiltonian and threshold only with
that branch tag.  They may not use this corollary as a premise for deriving
the GP1 ceiling, the causal cone, or the Lorentzian sign.
\end{corollary}

\begin{proof}
The Hamiltonian formula is
Theorem~\ref{p4:thm:finite-one-particle-dispersion}; taking the lowest
eigenvalue is Corollary~\ref{p4:cor:spectral-gap-threshold}.  Both theorems
state the same upstream domain, so their ordered pair is an exact export on
that domain.  The proof contains no reverse implication.
\end{proof}

\begin{remark}[Ownership relative to Paper~III]
\label{p4:rem:reciprocal-p3-lorentzian-closure}
Paper~III owns the GP1-calibrated cone, its Lorentzian conformal sign, and the
local second-order scalar mode relation.  Paper~IV owns the quantum-event
interpretation and the positive finite Hamiltonian realization on a stable
intrinsic block.  General interacting, medium-dependent, higher-order, Wick,
and continuum reconstructions remain outside this finite theorem unless
their separate domains are stated.
\end{remark}



\subsection[Field operators, correlation functions, and the particle/field dictionary]
{Field operators and the particle--field dictionary}
\label{p4:sec:part2-field-dictionary}

For any mode vector $f=\sum_i f_i e_i\in\mathcal K$, define
\begin{equation}
\label{p4:eq:mode-operators}
\mathfrak a^\dagger(f)
=
\sum_i f_i\,\mathfrak a_i^\dagger,
\qquad
\mathfrak a(f)
=
\sum_i \overline{f_i}\,\mathfrak a_i.
\end{equation}
The associated hermitian field operator is
\begin{equation}
\label{p4:eq:field-operator}
\Phi(f)=\mathfrak a(f)+\mathfrak a^\dagger(f).
\end{equation}
In a local presentation basis this is the finite-mode version of a quantum field smeared against the test function or mode profile $f$.
The particle/field dictionary is therefore basis-dependent but completely explicit: particles correspond to occupancy in distinguished one-particle sectors, while fields are the operator-valued mode superpositions that create, annihilate, and mix those occupancies.

Given a many-body Hamiltonian, the Heisenberg-evolved mode operators
\begin{equation}
\label{p4:eq:heis-mode}
\mathfrak a_i(t)=e^{itH_{\mathrm{mb}}}\mathfrak a_i e^{-itH_{\mathrm{mb}}}
\end{equation}
produce finite correlation functions such as
\begin{equation}
\label{p4:eq:two-point}
G_{ij}(t)
=
\Tr\!\bigl(\rho\,\mathfrak a_i(t)\mathfrak a_j^\dagger(0)\bigr).
\end{equation}
These are the finite analogues of propagators and Green functions.
Hence the usual QFT observables already have a precise finite-resolution form before any continuum limit is taken.



\section{Euclidean reflection, Gaussian boundary and finite kernel estimates}
\label{p4c:sec:v238-c3}
\noindent\textit{Source and claim boundary.} These are original extended calculations. Statements on physical sources, measurements or continuum limits retain the inherited antecedents of Papers IV-A/IV-B.

\subsection{Finite Euclidean kernels and Hamiltonian reconstruction}
\label{p4:sec:finite-euclidean-hamiltonian-reconstruction}

On an occurring closed-process domain, Paper~III uses a
\(\Dcl\)-form elliptic spectral loop, while Paper~III and later papers also
need real-time Hamiltonian language.  The finite bridge between a supplied
positive Euclidean kernel and real-time spectral language does not require
assuming continuum Lorentzian QFT.  It is a finite spectral reconstruction
statement and does not itself select the closed process or a Hamiltonian.

\begin{remark}[Claim-specific Wick and reconstruction ledger]
\label{p4:rem:wick-claim-ledger}
Paper~IV uses the hierarchy introduced in Paper~III-E, now with two logically
separate finite routes:
\begin{center}
\begin{tabular}{c|p{0.72\linewidth}}
\((\mathrm{WICK}_{2{\rm cyc}})\)
 & Kernel-first: finite positive-spectral cyclic continuation.  On the
 forward Hamiltonian-generated branch of
 Theorem~\ref{p4:thm:hamiltonian-generated-euclidean-branch}, the same finite
 certificate is automatic.\\
\((\mathrm{WICK}_{\rm same})\)
 & For an arbitrary supplied kernel: identification with the original carrier
 via the whitened same-carrier test of
 Lemma~\ref{p4:lem:positive-frequency-generator-from-euclidean-block}.  On
 the forward Hamiltonian-generated branch the criterion holds identically.\\
\((\mathrm{WICK}_{\rm OS,fin})\)
 & Kernel-first bosonic quasi-free or complete graded-fermionic OS
 certificate below; additionally, the ordinary/even Hamiltonian-generated
 transfer branch is reflection positive by
 Theorem~\ref{p4:thm:hamiltonian-transfer-reflection-positivity}.\\
\((\mathrm{WICK}_{\rm int,\infty})\)
 & Interacting infinite-volume reconstruction; not implied by any finite
 certificate below.
\end{tabular}
\end{center}
Thus the first three are not globally retired labels: their required
certificates are automatic or theorem-level only on the explicitly stated
finite branches.  The free Lorentzian mode equation imported from Paper~III
was derived without any of these tags.
\end{remark}


\begin{theorem}[Forward Hamiltonian-generated finite Euclidean branch]
\label{p4:thm:hamiltonian-generated-euclidean-branch}
Let \(\mathcal K\) be a finite Hilbert space, let
\(H=H^*\ge0\) on \(\mathcal K\), and let \(C_0>0\) be a positive invertible
operator on the same carrier.  Define, for \(\tau\ge0\),
\begin{equation}
\label{p4:eq:forward-hamiltonian-euclidean-kernel}
 C_H^+(\tau)
 :=C_0^{1/2}e^{-\tau H/\hbar}C_0^{1/2}.
\end{equation}
If
\[
 H=\sum_{a=1}^N E_aP_a,
 \qquad E_a\ge0,
\]
is the finite spectral resolution over distinct energies, then
\begin{equation}
\label{p4:eq:forward-hamiltonian-positive-spectral}
 C_H^+(\tau)
 =\sum_{a=1}^N Z^{(a)}e^{-E_a\tau/\hbar},
 \qquad
 Z^{(a)}:=C_0^{1/2}P_aC_0^{1/2}\succeq0.
\end{equation}
For every \(u\in\mathcal K\), the even scalar extension
\[
 c_u(\tau):=\langle u,C_H^+(|\tau|)u\rangle,
 \qquad \tau\in\mathbb R,
\]
therefore has the positive spectral form of
Theorem~\ref{p4:thm:finite-euclidean-hamiltonian-reconstruction} and has a
finite cyclic real-time continuation.

Moreover, whitening is exact:
\begin{equation}
\label{p4:eq:forward-hamiltonian-whitened-semigroup}
 C_0^{-1/2}C_H^+(\tau)C_0^{-1/2}
 =e^{-\tau H/\hbar}.
\end{equation}
Hence the original carrier and generator satisfy the same-carrier criterion,
and the recovered generator is uniquely \(H\).  The positive-time branch has
the entire operator-valued extension
\[
 C_H^+(z)=C_0^{1/2}e^{-zH/\hbar}C_0^{1/2},
 \qquad z\in\mathbb C,
\]
whose imaginary-axis value is
\[
 C_{H,\rm rt}(t)
 =C_0^{1/2}e^{-itH/\hbar}C_0^{1/2}.
\]
Thus on this explicitly forward-constructed finite branch the certificates
\((\mathrm{WICK}_{2{\rm cyc}})\) and \((\mathrm{WICK}_{\rm same})\) are
automatic consequences of finite spectral calculus.  The reflected two-sided
kernel obtained by replacing \(\tau\) with \(|\tau|\) is not claimed to be
analytic through \(\tau=0\); analyticity refers to the ordered positive-time
branch displayed above.
\end{theorem}

\begin{proof}
Finite spectral calculus gives
\[
 e^{-\tau H/\hbar}=\sum_a e^{-E_a\tau/\hbar}P_a,
\]
and multiplication by \(C_0^{1/2}\) on both sides yields
\eqref{p4:eq:forward-hamiltonian-positive-spectral}.  Positivity of every
\(Z^{(a)}\) follows from
\[
 \langle u,Z^{(a)}u\rangle
 =\|P_aC_0^{1/2}u\|^2\ge0.
\]
Scalar compression therefore gives nonnegative spectral weights, so
Theorem~\ref{p4:thm:finite-euclidean-hamiltonian-reconstruction} applies to
each such two-point channel.  Since \(C_0>0\), multiplying
\eqref{p4:eq:forward-hamiltonian-euclidean-kernel} by \(C_0^{-1/2}\) on both
sides proves \eqref{p4:eq:forward-hamiltonian-whitened-semigroup}.  Its
norm-continuous positive contraction semigroup has generator
\[
 -\hbar\left.\frac{d}{d\tau}e^{-\tau H/\hbar}\right|_{\tau=0}=H,
\]
so the generator is the same operator on the same carrier; equivalently it
satisfies Lemma~\ref{p4:lem:positive-frequency-generator-from-euclidean-block}.
Finally, the matrix exponential is entire in finite dimension, and setting
\(z=it\) gives the stated unitary real-time branch.  No Lorentzian sign,
causal cone, or continuum field limit is inferred by this construction.
\end{proof}



\begin{theorem}[Complex-time carrier dynamics and reciprocal evaluation square]
\label{p4:thm:complex-time-reciprocal-dynamics-square}
Under the hypotheses of
Theorem~\ref{p4:thm:hamiltonian-generated-euclidean-branch}, define
\begin{equation}
\label{p4:eq:complex-time-carrier-parent}
 \mathcal V_H(z):=e^{-zH/\hbar},
 \qquad z\in\mathbb C,
\end{equation}
and, on the finite observable algebra \(\mathcal B(\mathcal K)\), the entire
two-parameter operation family
\begin{equation}
\label{p4:eq:complex-time-two-sided-parent}
 \mathcal L_H(z,w)(A)
 :=e^{-zH/\hbar}A e^{-wH/\hbar}.
\end{equation}
Its positive-real and imaginary-axis carrier restrictions are
\begin{equation}
\label{p4:eq:complex-time-carrier-children}
 T_\tau:=\mathcal V_H(\tau)=e^{-\tau H/\hbar},
 \qquad
 U_t:=\mathcal V_H(it)=e^{-itH/\hbar},
\end{equation}
while the corresponding operation children are
\begin{align}
\label{p4:eq:complex-time-operation-children}
 \mathcal E_\tau(A)
 &:=\mathcal L_H(\tau/2,\tau/2)(A)
   =e^{-\tau H/(2\hbar)}A e^{-\tau H/(2\hbar)},\\
 \alpha_t(A)
 &:=\mathcal L_H(-it,it)(A)
   =e^{itH/\hbar}A e^{-itH/\hbar}.
\end{align}
For \(\tau\ge0\), \(\mathcal E_\tau\) is completely positive and
contractive; for \(t\in\mathbb R\), \(\alpha_t\) is a one-parameter
\(*\)-automorphism group.  Their trace-dual positive-functional operations are uniquely
\begin{align}
\label{p4:eq:complex-time-state-children}
 \mathcal E_{\tau,*}(\rho)
 &=e^{-\tau H/(2\hbar)}\rho e^{-\tau H/(2\hbar)},\\
 \alpha_{t,*}(\rho)
 &=e^{-itH/\hbar}\rho e^{itH/\hbar},
\end{align}
They preserve the positive cone and are trace-nonincreasing on the Euclidean
branch; normalization of a nonzero Euclidean output would be a separate
nonlinear readout.  They satisfy the evaluation identities
\begin{equation}
\label{p4:eq:complex-time-evaluation-square}
 \Tr\!\bigl(\mathcal E_{\tau,*}(\rho)A\bigr)
 =\Tr\!\bigl(\rho\,\mathcal E_\tau(A)\bigr),
 \qquad
 \Tr\!\bigl(\alpha_{t,*}(\rho)A\bigr)
 =\Tr\!\bigl(\rho\,\alpha_t(A)\bigr).
\end{equation}
Thus Euclidean contraction, Lorentzian unitary evolution, and their
state/test-dual descriptions are restrictions of one generated complex-time
operation diagram.  If \(H\) and the functional-calculus graph are generated
from one occurring root, then
Theorems~\ref{p2a:thm:evaluation-dual-corealization} and
\ref{p2a:thm:universal-corealization-bridge-normal-form} co-realize the whole
diagram without duplicate Euclidean, Lorentzian, Schr\"odinger, or Heisenberg
occurrence premises.

This result does not actualize a \(\mathsf{DYN}\) token, state history, or
trajectory.  It also does not identify the spectral resolution parameter
\(\Lambda^{-2}\) in a heat trace with the elapsed Euclidean time \(\tau\),
and it does not establish a nonlinear Lorentzian action or interacting
infinite-volume Wick reconstruction.
\end{theorem}

\begin{proof}
Finite-dimensional spectral calculus makes \(\mathcal V_H\) and
\(\mathcal L_H\) entire.  Since \(H\ge0\),
\(\|e^{-\tau H/(2\hbar)}\|\le1\), so
\(A\mapsto BAB\), with \(B=e^{-\tau H/(2\hbar)}\), is completely positive
and contractive.  For real \(t\), \(e^{-itH/\hbar}\) is unitary, giving the
stated automorphism group.  Cyclicity of the finite trace proves
\eqref{p4:eq:complex-time-evaluation-square} and uniqueness of the trace-dual
maps.  The operation and positive-functional children are obtained from the same
two-sided parent by restriction to the positive-real diagonal and imaginary
anti-diagonal, respectively.  The final co-realization statement is then the
evaluation-dual and rooted-incidence conclusion of the cited Paper-II-A
theorems.  None of those constructors supplies an actual process token or
identifies an independent cutoff coordinate with complex time.
\end{proof}

\begin{theorem}[Finite Euclidean-kernel to cyclic Hamiltonian reconstruction]
\label{p4:thm:finite-euclidean-hamiltonian-reconstruction}
Let \(C(\tau)\) be a finite Euclidean two-point kernel on a resolved block, defined for \(\tau\in\mathbb R\), with spectral form
\[
C(\tau)=\sum_{a=1}^N Z_a e^{-E_a|\tau|/\hbar},
\qquad
Z_a\ge0,
\qquad
E_a\ge0.
\]
Terms with equal \(E_a\) are first combined, so the retained energies are
pairwise distinct and their weights are nonzero.
Then there exists a finite cyclic realization
\((\mathcal K_C,H_C,\Omega_C,B_C)\), with \(H_C\ge0\) self-adjoint, such that
\[
C(\tau)=
\langle\Omega_C,
B_C e^{-|\tau|H_C/\hbar}B_C^*\Omega_C\rangle.
\]
The associated real-time correlator is the finite spectral continuation
\[
C_{\rm rt}(t)=\sum_{a=1}^N Z_a e^{-iE_a t/\hbar}
=
\langle\Omega_C,
B_C e^{-itH_C/\hbar}B_C^*\Omega_C\rangle.
\]
The reconstruction space \(\mathcal K_C\) is generated by the nonzero
spectral-weight channels and need not equal the finite carrier block on
which the kernel was recorded.  Thus positive Euclidean spectral data
determine a positive cyclic Hamiltonian realization at finite resolution;
a Hamiltonian on the original block requires the additional semigroup
criterion of
Lemma~\ref{p4:lem:positive-frequency-generator-from-euclidean-block}.
This theorem is the finite certificate for
\((\mathrm{WICK}_{2{\rm cyc}})\); it does not establish same-carrier,
full OS, or interacting infinite-volume reconstruction.
\end{theorem}


\begin{proof}
After combining equal energies, take an orthonormal vector \(e_a\) for each
retained nonzero weight \(Z_a\), let
\(H_Ce_a=E_ae_a\), and add a vacuum vector \(\Omega_C\) with
\(H_C\Omega_C=0\).  Define
\(B_C^*\Omega_C=\sum_a\sqrt{Z_a}\,e_a\) and
\(B_Ce_a=\sqrt{Z_a}\Omega_C\), extending \(B_C\) by zero on any remaining
orthogonal complement.  Then
\[
\langle\Omega_C,
B_Ce^{-|\tau|H_C/\hbar}B_C^*\Omega_C\rangle
=\sum_a Z_a e^{-E_a|\tau|/\hbar}.
\]
Because \(H_C\) is a finite diagonal matrix with nonnegative eigenvalues, it
is positive and self-adjoint.  Replacing the contraction semigroup
\(e^{-|\tau|H_C/\hbar}\) by the unitary group \(e^{-itH_C/\hbar}\) gives the
stated real-time correlator.  Since the retained energies are distinct,
Lagrange interpolation supplies each spectral projection of \(H_C\) as a
polynomial in \(H_C\); applying those projections to
\(B_C^*\Omega_C\) recovers every \(e_a\).  Thus \(\Omega_C\) is cyclic for
the finite algebra generated by \(H_C,B_C,B_C^*\).  This construction identifies no vector of
\(\mathcal K_C\) with a vector of the original carrier block.  No causal
cone, Lorentz group, or continuum field algebra enters the reconstruction.
\end{proof}


\begin{theorem}[Finite bosonic OS positivity and complete graded-fermionic criterion]
\label{p4:thm:finite-os-positivity}
Let \(C_{\alpha\beta}(\tau)\) be a finite Euclidean two-point kernel on a resolved block with positive spectral representation
\begin{equation}
\label{p4:eq:finite-os-spectral-representation}
C_{\alpha\beta}(\tau)
=
\sum_{a=1}^N
Z^{(a)}_{\alpha\beta}e^{-E_a|\tau|/\hbar},
\qquad
E_a\ge0,
\qquad
Z^{(a)}\ge0
\quad\text{as a matrix.}
\end{equation}
In the bosonic branch let \(\Theta\) be an antilinear, involutive,
\(*\)-compatible Euclidean reflection.  For neutral fields,
\((\Theta\Phi_\alpha)(\tau)=\Phi_\alpha(-\tau)\), with coefficients
complex-conjugated; for charged fields include the supplied adjoint/charge
conjugation.  Let \(F\) be a finite polynomial in fields smeared at positive
Euclidean times.  In the bosonic case, and hence on its even subalgebra, if the higher correlations are
the finite quasi-free/Wick extension of
\eqref{p4:eq:finite-os-spectral-representation}, then
\begin{equation}
\label{p4:eq:finite-os-positivity}
\langle \Theta F\,F\rangle_E\ge0.
\end{equation}

For a finite graded-fermionic sector, additionally supply a
\(\mathbb Z_2\)-graded factorization
\(\mathcal A=\mathcal A_-\widehat\otimes\mathcal A_+\), an antilinear
parity-preserving reflection
\(\Theta:\mathcal A_+\to\mathcal A_-\), and a fixed graded OS product
\(\circ\).  Equivalently, one may use the convention
\[
X_-\circ Y_+:=i^{|X_-|}X_-Y_+
\]
on homogeneous matching-parity elements, with any equivalent Klein phase
absorbed into \(\Theta\).  Choose a homogeneous vector-space basis
\(\{C_I\}_{I=1}^M\) of the finite algebra \(\mathcal A_+\), compatible with
its parity decomposition, and define the full reflection matrix
\[
R_{IJ}:=
\omega\!\left(\Theta(C_I)\circ C_J\right).
\]
Then
\[
\boxed{\quad
\omega\!\left(\Theta(F)\circ F\right)\ge0
\ \text{for every }F\in\mathcal A_+
\quad\Longleftrightarrow\quad
R\succeq0 .
\quad}
\]
For a gauge-invariant, number-conserving quasi-free state, every entry of
\(R\) is computed from the complete covariance by the finite determinant
Wick rule.  The reflected one-particle matrix
\[
q(f,g):=\omega\!\left(\Theta(c(f))\circ c(g)\right)
\]
is only the degree-one principal corner of \(R\).  Consequently,
\(R\succeq0\) implies \(q\succeq0\), but \(q\succeq0\) alone does not imply
OS positivity on all of \(\mathcal A_+\).  A non-number-conserving
quasi-free state is treated in a self-dual/Majorana basis, where the entries
of \(R\) are the corresponding Pfaffians.  Reflection-positive Majorana
Hamiltonian families also admit coupling-matrix criteria
\cite[Theorem~4.2]{JaffeJanssens2016Reflection}; these are computable sector tests, not a
universal consequence of finite CAR.

Positive matrices \(Z^{(a)}\) in an ungraded scalar spectral representation
do not by themselves imply the fermionic condition \(R\succeq0\).  Both the
bosonic result and the fermionic classification are finite and algebraic;
neither constructs an arbitrary interacting continuum Wightman theory.
They are the finite certificates for
\((\mathrm{WICK}_{\rm OS,fin})\).
\end{theorem}


\begin{proof}
For a linear positive-time observable
\[
F=\sum_{j,\alpha}c_{j\alpha}\Phi_\alpha(\tau_j),
\qquad \tau_j>0,
\]
reflection positivity is the positivity of
\begin{align}
\langle \Theta F\,F\rangle_E
&=
\sum_{i,j,\alpha,\beta}
\overline{c_{i\alpha}}c_{j\beta}
C_{\alpha\beta}(\tau_i+\tau_j) \\
&=
\sum_a (v^{(a)})^*Z^{(a)}v^{(a)}
\ge0,
\end{align}
where
\[
v^{(a)}_\alpha
:=
\sum_j c_{j\alpha}e^{-E_a\tau_j/\hbar}.
\]
This is nonnegative because each \(Z^{(a)}\) is positive semidefinite.  For bosonic
polynomial observables, Wick expansion rewrites
\(\langle\Theta F F\rangle_E\) as a symmetric-Fock norm square.

In the fermionic case, write
\[
F=\sum_{I=1}^M f_I C_I.
\]
Antilinearity of \(\Theta\) and bilinearity of the fixed graded product give
\[
\omega\!\left(\Theta(F)\circ F\right)
=
\sum_{I,J}\overline{f_I}R_{IJ}f_J
=f^*Rf.
\]
This is nonnegative for every coefficient vector \(f\) if and only if the
finite matrix \(R\) is positive semidefinite.  Changing the homogeneous basis
acts by congruence, so the criterion is basis independent.  For a quasi-free
state, the finite Wick theorem computes all entries of \(R\) from the full
covariance: determinants in the gauge-invariant number-conserving
presentation and Pfaffians in the self-dual/Majorana presentation.  The
degree-one form \(q\) is therefore a principal submatrix and is necessary,
but positivity of one principal submatrix cannot establish positivity of the
whole reflection matrix.  Without the graded reflection and twist, \(R\) is
not defined.
\end{proof}

\begin{theorem}[Hamiltonian transfer-semigroup reflection positivity]
\label{p4:thm:hamiltonian-transfer-reflection-positivity}
Let \(\mathcal A\subseteq\End(\mathcal K)\) be a finite unital
\(^*\)-algebra, let \(H=H^*\ge0\), and let \(\Omega\in\mathcal K\) be a
normalized invariant ground vector,
\[
 H\Omega=0.
\]
For an ordered time-labelled word
\[
 W=((B_1,s_1),\ldots,(B_n,s_n)),
 \qquad s_1<\cdots<s_n,
\]
define its Hamiltonian-generated Euclidean functional by
\begin{equation}
\label{p4:eq:transfer-euclidean-word-functional}
 \mathcal E_H(W)
 :=\left\langle\Omega,
 B_1e^{-(s_2-s_1)H/\hbar}B_2\cdots
 e^{-(s_n-s_{n-1})H/\hbar}B_n\Omega\right\rangle .
\end{equation}
For a positive-time word
\[
 F=((A_1,\tau_1),\ldots,(A_n,\tau_n)),
 \qquad 0<\tau_1<\cdots<\tau_n,
\]
define reflection by
\[
 \Theta F
 :=((A_n^*,-\tau_n),\ldots,(A_1^*,-\tau_1)).
\]
Then for any two positive-time words \(F,G\),
\begin{equation}
\label{p4:eq:transfer-reflection-gram}
 \mathcal E_H(\Theta F,G)
 =\langle V(F),V(G)\rangle,
\end{equation}
where
\begin{equation}
\label{p4:eq:transfer-halfspace-vector}
 V(F)
 :=e^{-\tau_1H/\hbar}A_1
 e^{-(\tau_2-\tau_1)H/\hbar}A_2\cdots
 e^{-(\tau_n-\tau_{n-1})H/\hbar}A_n\Omega.
\end{equation}
Extend the reflected pairing sesquilinearly to finite linear combinations
of positive-time words.  Consequently, for
\(X=\sum_r c_rF_r\),
\begin{equation}
\label{p4:eq:transfer-os-norm-square}
 \boxed{
 \mathcal E_H(\Theta X,X)
 =\left\|\sum_r c_rV(F_r)\right\|^2\ge0 .}
\end{equation}
The corresponding finite OS quotient is canonically isometric to
\(\operatorname{span}\{V(F)\}\subseteq\mathcal K\); if that span is the
chosen cyclic carrier, the OS reconstruction returns that carrier.  Thus
\((\mathrm{WICK}_{\rm OS,fin})\) is theorem-level on this
Hamiltonian-generated ordinary/ungraded branch and on any even observable
subalgebra to which the ordinary product applies.  This theorem does not
replace the full graded-fermionic reflection-matrix criterion of
Theorem~\ref{p4:thm:finite-os-positivity}.
\end{theorem}

\begin{proof}
Let
\(F=((A_1,\tau_1),\ldots,(A_n,\tau_n))\) and
\(G=((B_1,\sigma_1),\ldots,(B_m,\sigma_m))\).  The concatenated ordered word
\((\Theta F,G)\) has times
\[
 -\tau_n<\cdots<-\tau_1<\sigma_1<\cdots<\sigma_m.
\]
Using \eqref{p4:eq:transfer-euclidean-word-functional} therefore gives
\[
\begin{aligned}
 \mathcal E_H(\Theta F,G)
 =\langle\Omega,&
 A_n^*e^{-(\tau_n-\tau_{n-1})H/\hbar}\cdots
 A_1^*e^{-(\tau_1+\sigma_1)H/\hbar}B_1\\
 &\qquad\cdot e^{-(\sigma_2-\sigma_1)H/\hbar}\cdots B_m\Omega\rangle.
\end{aligned}
\]
Since every heat factor is self-adjoint and
\(e^{-\tau_1H/\hbar}e^{-\sigma_1H/\hbar}
=e^{-(\tau_1+\sigma_1)H/\hbar}\), this is exactly
\(\langle V(F),V(G)\rangle\).  Sesquilinearity then yields
\eqref{p4:eq:transfer-os-norm-square}.  Quotienting the null vectors of the
reflection form identifies \([F]\mapsto V(F)\) isometrically with the stated
subspace, proving the finite reconstruction statement.

If a physical Hamiltonian has a lowest eigenvalue \(E_0\) rather than zero,
replace it in this transfer representation by
\(H-E_0\mathbf1\ge0\) and choose a corresponding ground vector.  The scalar
shift does not change the adjoint real-time dynamics
\(A\mapsto e^{itH/\hbar}Ae^{-itH/\hbar}\).  No quasi-free/Wick factorization
is used in the proof.  Odd graded fields require the separate reflected
product and Klein/twist data of Theorem~\ref{p4:thm:finite-os-positivity}.
\end{proof}


\begin{remark}[Euclidean--Hamiltonian bridge and non-identifications]
\label{p4:rem:euclidean-hamiltonian-noncircular}
\leavevmode\par\noindent
Theorem~\ref{p4:thm:finite-euclidean-hamiltonian-reconstruction} is not a
Wick-rotation postulate for a pre-existing relativistic QFT.  It is finite
spectral theory: positivity of a supplied finite kernel gives a positive
Hamiltonian spectrum and Euclidean/real-time correlators on its cyclic
reconstruction space.  Identification of that space and generator with an
already constituted carrier block requires the separate whitened semigroup
criterion in
Lemma~\ref{p4:lem:positive-frequency-generator-from-euclidean-block}, unless
the kernel is instead defined forward from that Hamiltonian as in
Theorem~\ref{p4:thm:hamiltonian-generated-euclidean-branch}.  The transfer
construction of
Theorem~\ref{p4:thm:hamiltonian-transfer-reflection-positivity} similarly
proves finite ordinary/even reflection positivity because the Hamiltonian and
state are already part of its theorem domain.

No theorem in this subsection identifies Paper~III-E's gravitational heat
readout
\(\Tr e^{-\Delta_+^{\rm scal}/\Lambda_g^2}\) with a quantum two-point
covariance.  The later free-mode vacuum theorem derives, on its distinct
stationary oscillator domain, the convention-safe normalization
\(C_{QQ}=\hbar(2\Omega)^{-1}=\hbar^2(2H)^{-1}\) with \(H=\hbar\Omega\);
that result is not an identification of the gravitational heat trace with a
quantum covariance.  A bare mathematical functional-calculus map
\(H\mapsto e^{-\tau H/\hbar}\) is not by itself an actual DYN role.  Later
Theorem~\ref{p4:thm:pure-scalar-factor-representation} derives such an
exponential only as the representation of the already typed pure scalar-factor
translation on a stationary positive-frequency free-law block; actual process
selection remains under \(\mathsf{DYN}\).  The GP1 cone and Lorentzian sign are
already Paper-III inputs.  A general interacting
infinite-volume Euclidean reconstruction continues to require the explicit
\(({\rm WICK}_{\rm int,\infty})\) domain.
\end{remark}




\subsection{Gaussian boundary amplitudes and stationary-state distinctions}
\label{p4:app:v2109-gaussian-boundary}
\begin{theorem}[Finite Gaussian boundary constructor with its absolute source]
\label{p4:thm:v2109-gaussian-boundary-source}
On a supplied real positive Euclidean configuration branch, take a fixed
reference measure \(dx\), fixed domain \(\mathbb R^m\), and
\[
 S_E(q,x;g)=\tfrac12q^TA(g)q+q^TB(g)x+\tfrac12x^TC(g)x+s(g),
 \qquad C>0.
\]
The boundary weight and Schur coefficient are
\[
 \Psi(q;g)=(2\pi\hbar)^{m/2}(\det C)^{-1/2}e^{-s/\hbar}
 e^{-q^TKq/(2\hbar)},\qquad K=A-BC^{-1}B^T.
\]
Where a licensed boundary-amplitude map interprets this weight as a
wavefunction and \(K>0\), its normalized state is
\[
 \psi_K(q)=\left(\frac{\det K}{(\pi\hbar)^n}\right)^{1/4}
 e^{-q^TKq/(2\hbar)}.
\]
The absolute effective boundary action is
\[
 \Gamma_\partial=\tfrac12q^TKq+s+\frac\hbar2\log\det C
                  -\frac{m\hbar}{2}\log(2\pi\hbar),
\]
so at fixed boundary coordinates, measure and domain,
\[
 \delta\Gamma_\partial=\tfrac12q^T(\delta K)q+\delta s
 +\frac\hbar2\operatorname{Tr}(C^{-1}\delta C),
\]
\[
 \delta K=\delta A-(\delta B)C^{-1}B^T-BC^{-1}(\delta B)^T
            +BC^{-1}(\delta C)C^{-1}B^T.
\]
\end{theorem}
\begin{proof}
Complete the square by translating the interior coordinate to
\(x+C^{-1}B^Tq\), and evaluate the positive Gaussian integral.  Taking
\(-\hbar\log\) and differentiating the matrix inverse and determinant
proves the source formula.
\end{proof}
The normalization determinant contributes to gravitational/source variation;
it cannot be dropped because a normalized state is unchanged by a scalar
factor.  A commutative probability density is not automatically this
wavefunction amplitude: its square-root amplitude would be different.

\begin{proposition}[Ground-compatible width and collar gluing]
\label{p4:prop:v2109-ground-width-collar}
For \(H=p^Tp/(2\mu)+\mu\omega^2q^Tq/2\), \(\mu,\omega>0\), a scalar
positive Gaussian width \(K\) gives the oscillator ground precisely when
\(K=\mu\omega\).  Its annihilator is
\((p-iKq)\psi_K=0\).  Euclidean propagation through a collar of duration
\(\tau>0\) sends the width to
\[
 K_\tau=\mu\omega\,
 \frac{K+\mu\omega\tanh(\omega\tau)}
 {\mu\omega+K\tanh(\omega\tau)}.
\]
All positive widths obey the collar subdivision composition law, whereas
\(K=\mu\omega\) is the unique positive absorbing fixed width.
\end{proposition}
\begin{proof}
Apply \(p=-i\hbar\partial_q\) to the Gaussian and then apply \(H\).
The quadratic term in \(H\psi_K/\psi_K\) vanishes exactly at
\(K^2=\mu^2\omega^2\).  Positive width fixes the sign.  Schur
elimination of the oscillator Euclidean collar kernel gives the displayed
fractional-linear map.  The hyperbolic-tangent addition formula proves
subdivision.  Its positive fixed points solve \(K^2=\mu^2\omega^2\).
\end{proof}
The annihilator convention here is \(p-iKq\), which is proportional to
the usual annihilation operator; it does not impose both \(q\psi=0\)
and \(p\psi=0\).

\begin{proposition}[Three inequivalent stationarity conditions]
\label{p4:prop:v2109-stationarity-distinction}
For a finite \(E\succeq0\), real-time stationary densities are precisely
those commuting with \(E\).  A density fixed by the normalized Euclidean
map \(\rho\mapsto e^{-\tau E/2}\rho e^{-\tau E/2}/
\operatorname{Tr}(\rho e^{-\tau E})\), for one \(\tau>0\), is supported
in a single energy eigenspace, which may be excited.  Preserving the
unnormalized weight instead requires support in the zero-energy eigenspace.
\end{proposition}
\begin{proof}
Real-time stationarity is its generator commutator equation.  In an energy
basis, every nonzero diagonal entry of a normalized Euclidean fixed state
satisfies \(e^{-\tau e_i}=Z\), with the same normalization \(Z\).
Positivity removes matrix elements outside its diagonal support; injectivity
of the exponential gives a single energy.  The converse follows by direct
substitution.  The last statement is
Theorem~\ref{p4:thm:v2109-euclidean-weight-support}.
\end{proof}
Hence subdivision, normalized stationarity, zero means, scalar relative
support and classical zero homogeneous cotangent do not by themselves
supply quantum zero excitation.  This is the precise scope boundary, not a
new vacuum axiom.  A reset to a daughter vacuum also does not prove parent
exhaustion, so QROOT and VACCAP remain distinct physical incidences.

% END INLINE SOURCE: Paper4-GaussianBoundaryAppendix-v1_0.tex


% Bibliography moved to end of reorganized volume.
% BEGIN BOUNDARY-THERMAL-INTERFACE v4.96


\section{Field-inductive systems and formal continuum-presentation estimates}
\label{p4c:sec:v238-c4}
\noindent\textit{Source and claim boundary.} These are original extended calculations. Statements on physical sources, measurements or continuum limits retain the inherited antecedents of Papers IV-A/IV-B.

\subsection{Same-word transported-embedding covariance}
\label{p4:sec:v4118-transported-embedding-covariance}
Use the already licensed first-order parent and finite word of
Theorem~\ref{p3f:thm:v1148-finite-metric-constructor}, including its
field/dual slots, actual ordering, occupation and constraint projectors,
central allocation and consumed source jets. All transformations below
are changes of presentation of that word. They select no physical cutoff,
state or interaction. A spatial pullback is included only where it is a
licensed action on the parent domain, spin/gauge lift and complete marks.
The clock and foliation are fixed in this subsection. Here $x$ denotes
the numerical source coordinates (or supplied source history), whereas
$y$ denotes spatial position. The parent $T$ may act on fields as a
spatial pullback/operator in $y$, on its licensed common domain. In
contrast, $C(t,x)$ below is an invertible finite coefficient matrix
depending on numerical sources; it is not a pointwise matrix $C(t,y)$
acting on distinct field values. Dots denote total derivatives along
the actual source history.

\begin{theorem}[Exact parent-presentation and coefficient-frame covariance]
\label{p4:thm:v4118-transported-embedding-covariance}
Let $\psi'=T\psi$ for a differentiable invertible licensed parent
presentation $T(t,x)$; write $A=\dot T T^{-1}$. Transform the parent
first-order forms and the actual field embedding by
\begin{align*}
 G_p'&=T^{-\dagger}G_pT^{-1},&
 K_p'&=T^{-\dagger}K_pT^{-1},\\
 B_p'&=T^{-\dagger}B_pT^{-1}
                  +\tfrac12(A^\dagger G_p'-G_p'A),& J'&=TJ.
\end{align*}
Every inherited interaction tensor transforms contragrediently in all
its field and dual-field slots. The actual other-sector, occupation,
constraint and ordering maps are transported as maps of the same word.
Then the finite constructor returns exactly
\[
 G'=G,\qquad K'=K,\qquad B'=B,\qquad
 e'=e,\quad\gamma'=\gamma,\quad\widehat H'=\widehat H.
\]
The last equality includes every inherited vertex and central energy.
It holds as an identity of the complete source family, so all existing
first, mixed second and functional jets agree, not only one matrix value.

More generally take $c'=C^{-1}c$, $J'=TJC$ with an invertible coefficient
frame $C(t,x)$. Set $S'=(C^\dagger GC)^{1/2}$, $W=S^{-1}$ and
\[
 F=S'C^{-1}W.
\]
Then $F$ is unitary, the normalized one-particle state satisfies $q'=Fq$,
and
\[
 e'=FeF^\dagger,\qquad
 \gamma'=F\gamma F^\dagger-\dot F F^\dagger,\qquad
 Z'F=TZ.
\]
On the full finite CAR word use the exterior lift $\mathbb U=\Gamma(F)$
and the declared state convention $u'=\mathbb Uu$. The complete finite
Hamiltonian obeys
\[
 \boxed{\widehat H'=\mathbb U\widehat H\mathbb U^\dagger
                         +i\hbar\dot{\mathbb U}\mathbb U^\dagger.}
\]
Equivalently the connection term is
$-i\hbar\mathbb U\dot{\mathbb U}^\dagger$.
For an actually finite bosonic or constrained joint word this last
statement requires its already licensed lift and transported actual
projectors; a fixed modewise occupation cap is not presumed invariant.
No finite CCR squeezing or new canonical quantization is asserted.
\end{theorem}
\begin{proof}
Substitute $\psi=T^{-1}\psi'$ in the symmetrized kinetic action. Its
extra coefficient is $(A^\dagger G_p'-G_p'A)/2$, proving the displayed
parent rule. In the finite pullback, the extra parent term becomes
\[
 \tfrac12(J^\dagger\dot T^\dagger T^{-\dagger}G_pJ
                    -J^\dagger G_pT^{-1}\dot TJ).
\]
The two terms from differentiating $TJ$ in the embedding connection
are their negatives. Thus $B'=B$ exactly, while the $G,K$ equalities
are immediate. Interaction-slot substitution cancels $T$ with $T^{-1}$
in every slot. The same physical central allocation is unchanged.
Pointwise equality of these complete differentiable families proves
the jet statements, including all $T_a,T_{ab},\dot T_a$ terms.

With the additional frame $C$, the pulled-back coefficients are
\[
 G'=C^\dagger GC,\quad K'=C^\dagger KC,\quad
 B'=C^\dagger BC+\tfrac12(C^\dagger G\dot C-\dot C^\dagger GC).
\]
Consequently
$F^\dagger F=WC^{-\dagger}G'C^{-1}W=WGW=I$ and $q'=Fq$.
Substitute $q=F^\dagger q'$ in the normalized symmetrized action.
The extra connection is $F\dot F^\dagger=-\dot FF^\dagger$,
giving $\gamma'$ with the stated sign. Each inherited vertex and
actual projector transports by its licensed lift, and the exact CAR
Lie homomorphism lifts the one-particle connection. Finally
$i\hbar\partial_t(\mathbb Uu)=
(\mathbb U\widehat H\mathbb U^\dagger+
i\hbar\dot{\mathbb U}\mathbb U^\dagger)(\mathbb Uu)$.
This fixes the state convention independently of ladder notation.
\end{proof}

\begin{corollary}[Source-bound purely spatial Ward transfer]
\label{p4:cor:v4118-source-bound-spatial-ward}
Suppose an interior-supported, time-preserving spatial diffeomorphism
of the actual complete source family has licensed parent lift $T_\varphi$
and the actual embedding/mark diagram
\[
 J(x^\varphi)=T_\varphi J(x)C_\varphi
\]
with the foregoing complete tensor and projector transport. This is an
identity of the actual family before restriction to a scalar port or
gauge leaf. Then its canonical lift $\mathbb U_\varphi$ satisfies the
covariance antecedent of Theorem~\ref{p4:thm:v4106-canonical-stress-ward};
its actual $\mathfrak D_\xi$ and all existing first source variations
vanish. With $\mathbb U_{\epsilon\xi}=I-i\epsilon G_\xi/\hbar+
O(\epsilon^2)$,
\[
 \delta H[R_\xi]=\dot G_\xi-\frac{i}{\hbar}[G_\xi,H].
\]
The state is transported, not assumed invariant. The functional Hilbert
insertion retains $K[R_\xi]$, endpoint/initial data, central derivatives
and all external-force/mark terms. A force-free metric-only identity
follows only when the actual complete coupled equations omit no such
force. Differentiation includes $\delta R_\xi[v]$ and mixed parent jets.
Direct evaluation of the complete actual $\mathfrak D_\xi=0$ and its
first variation is an alternative to this sufficient diagram.
\end{corollary}

The theorem preserves the same actual carrier and supplies the stated
conditional spatial Ward transfer; it selects no physical moving cutoff.
That transfer consumes the actual complete source-band diagram or a direct
complete zero-defect and first-variation certificate, with every retained
mark force included. Compact isometries already natural to the identified
unmarked word require no extra symmetry ingredient or matrix printout.
Appendix~\ref{p4:app:v4118-transported-embedding-diagnostics} distinguishes
this exact presentation covariance from changing the carrier of a
metric-only source family and records the fixed-reference/comparison
scope without supplying a new physical band.

This result covers spatially dependent $T_\varphi$ and its time connection
when the foliation is preserved. It supplies no canonical lift for
spatially varying normal-time displacements, changes of clock/foliation,
or changed physical occupation marks. The passive homogeneous clock
Theorems~\ref{p4:thm:v4118-passive-clock-pullback} and
\ref{p4:thm:v4118-homogeneous-time-ward} retain their own domain and do
not enlarge this spatial result to all spacetime diffeomorphisms or close
the full local stress-Ward/projectability application.

% END INLINE SOURCE: Paper4-TransportedEmbeddingCovariance-Main-v1_0.tex



% BEGIN INLINE SOURCE: Paper4-SphericalDiracSource-Main-v1_0.tex
% Framework v2.166. Complete coordinate-source work and evaluated free band.


\subsection{Inductive finite-to-field bridge}
\label{p4:sec:part2-inductive-field-bridge}

The finite mode-space construction already contains the algebraic data needed for a controlled continuum interface.
The correct claim is not that the full continuum theory is primitive.
It is that compatible families of finite mode sectors define the local algebraic object from which the usual field-theoretic presentation is recovered.

\begin{lemma}[Unital full-matrix embedding obstruction]
\label{p4:lem:matrix-embedding-divisibility}
There exists a unital injective \(*\)-homomorphism
\[
 M_n(\mathbb C)\longrightarrow M_m(\mathbb C)
\]
if and only if \(m=kn\) for an integer \(k\ge1\).  Up to unitary conjugacy it
is \(A\mapsto A\otimes I_k\).  Thus a mere dimension inequality \(n\le m\)
does not supply a unital full-matrix embedding.
\end{lemma}

\begin{proof}
A unital representation of the simple algebra \(M_n\) on
\(\mathbb C^m\) is a direct sum of copies of its unique irreducible
\(n\)-dimensional representation.  Hence \(m=kn\), and the displayed
amplification gives the converse.
\end{proof}


\begin{theorem}[Inductive finite-to-field interface bridge]
\label{p4:thm:inductive-qft-bridge}
Let
\[
\mathcal K_1\subset \mathcal K_2\subset \cdots
\]
be a nested family of distinguished finite one-particle mode spaces with compatible embeddings of the corresponding finite many-body observable algebras
\[
\iota_{MN}:\mathcal A_N\hookrightarrow \mathcal A_M,
\qquad
M\ge N.
\]
The theorem-domain data have four distinct types:
\[
\begin{array}{c|l}
v_{MN}:\mathcal K_N\to\mathcal K_M&
\text{isometric one-particle/Hilbert maps},\\
\iota_{MN}:\mathcal A_N\to\mathcal A_M&
\text{injective \(*\)-homomorphisms},\\
\omega_M\circ\iota_{MN}=\omega_N&
\text{state compatibility, when a limit state is claimed},\\
\iota_{MN}\alpha_t^{(N)}
=\alpha_t^{(M)}\iota_{MN}&
\text{dynamics compatibility, when limit dynamics is claimed}.
\end{array}
\]
Canonical CAR embeddings and regular Weyl embeddings are induced by
isometric or symplectic \(v_{MN}\).  Tensor-factor refinements use
\(A\mapsto A\otimes I\).  Total-particle-number caps instead give natural
nonunital corner maps unless a different truncation is chosen, and arbitrary
full matrix algebras are subject to
Lemma~\ref{p4:lem:matrix-embedding-divisibility}.
Then the directed union
\[
\mathcal A_{\mathrm{loc}}:=\bigcup_N \iota_N(\mathcal A_N)
\]
is a $*$-algebra admitting a unique $C^*$-completion \(\mathcal A_\infty\).
Any compatible family of states \(\omega_M\circ \iota_{MN}=\omega_N\) defines a unique state on \(\mathcal A_\infty\), and any compatible family of internally supplied local finite dynamics defines a \(*\)-automorphism group on \(\mathcal A_{\mathrm{loc}}\).  It extends to a \(*\)-automorphism group on \(\mathcal A_\infty\) whenever the induced local action is norm-continuous, or more generally whenever the theorem-domain quasi-local dynamical-limit hypotheses are imposed.
\end{theorem}


\begin{proof}
\ProofPlan{We construct the finite-to-field bridge by putting a well-defined \(C^*\)-norm on the directed union, extending compatible states, and extending compatible finite dynamics to the completed algebra under the stated continuity hypotheses.}
\begin{proofsteps}

\item \textbf{Directed finite \(C^*\)-algebra norm and completion.}
Because each \(\mathcal A_N\) is a finite-dimensional $C^*$-algebra, every embedding \(\iota_{MN}\) is an injective isometric $*$-homomorphism.
Hence if \(A\in \mathcal A_N\) and \(\iota_N(A)=\iota_M(B)\) in the directed union, then after passing to a common level \(L\ge M,N\) we have
\[
\iota_{LN}(A)=\iota_{LM}(B),
\]
so
\[
\|A\|=\|\iota_{LN}(A)\|=\|\iota_{LM}(B)\|=\|B\|.
\]
Thus the formula \(\|\iota_N(A)\|:=\|A\|\) defines a well-posed norm on \(\mathcal A_{\mathrm{loc}}\).
Closure under multiplication and involution is immediate because any two elements of the directed union lie in some common algebra \(\mathcal A_L\), where the algebraic operations are already defined.
Since the $C^*$-identity holds on every \(\mathcal A_L\), it holds on the union and therefore on its completion.
This gives a $C^*$-completion, and uniqueness is the usual uniqueness of the completion of a normed space.


\item \textbf{Compatible state extension.}
For states, define on the union
\[
\omega\bigl(\iota_N(A)\bigr):=\omega_N(A).
\]
If \(\iota_N(A)=\iota_M(B)\), then at a common level \(L\)
\[
\omega_N(A)=\omega_L\bigl(\iota_{LN}(A)\bigr)=\omega_L\bigl(\iota_{LM}(B)\bigr)=\omega_M(B),
\]
so \(\omega\) is well defined.
It is positive and normalized because each \(\omega_N\) is, hence bounded with norm one and therefore extends uniquely by continuity to a state on \(\mathcal A_\infty\).

\item \textbf{Limit dynamics on the inductive system.}
Finally, let \(\alpha_t^{(N)}\) denote the local finite dynamics generated by the
compatible Hamiltonians.  Compatibility means
\[
\iota_{MN}\circ \alpha_t^{(N)} = \alpha_t^{(M)}\circ \iota_{MN}
\qquad (M\ge N),
\]
so
\[
\alpha_t\bigl(\iota_N(A)\bigr):=\iota_N\bigl(\alpha_t^{(N)}(A)\bigr)
\]
defines a well-posed \(*\)-automorphism of the directed union.  Suppose that
\(t\mapsto\alpha_t\) is norm-continuous on the directed union.  Then each
\(\alpha_t\) extends uniquely to a \(*\)-automorphism of \(\mathcal A_\infty\),
giving the claimed local dynamics.
\end{proofsteps}

\end{proof}

\begin{proposition}[Inheritance of reflection positivity under compatible inductive embeddings]
\label{p4:prop:inductive-reflection-positivity-inheritance}
Assume the inductive-system hypotheses of
Theorem~\ref{p4:thm:inductive-qft-bridge} together with compatible
positive-time subalgebras \(\mathcal A_{N,+}\subseteq\mathcal A_N\) and
conjugate-linear isometric reflections \(\Theta_N\) such that
\[
 \iota_{MN}(\mathcal A_{N,+})\subseteq\mathcal A_{M,+},
 \qquad
 \Theta_M\iota_{MN}=\iota_{MN}\Theta_N.
\]
Suppose the compatible states satisfy finite reflection positivity,
\[
 \omega_N\!\left(\Theta_N(F)F\right)\ge0,
 \qquad F\in\mathcal A_{N,+}.
\]
Then the algebraic positive-time union
\[
 \mathcal A_{\mathrm{loc},+}
 :=\bigcup_N\iota_N(\mathcal A_{N,+})
\]
has a well-defined conjugate-linear isometric reflection \(\Theta\), and the
inductive-limit state obeys
\begin{equation}
\label{p4:eq:inductive-os-inheritance}
 \boxed{
 \omega\!\left(\Theta(F)F\right)\ge0,
 \qquad F\in\mathcal A_{\mathrm{loc},+}.}
\end{equation}
The isometric reflection extends uniquely to the corresponding norm closure.
For a graded-fermionic branch, the same conclusion holds only when the graded
OS products/twists are also compatible with the embeddings, in which case
\(\Theta(F)F\) in the displayed formula is replaced by the compatible graded
product \(\Theta(F)\circ F\).

This proposition transfers a proved finite reflection-positivity property; it
does not supply the embeddings, compatible states/reflections, continuum
locality or covariance, distributional regularity, clustering, or convergence
needed for \((\mathrm{WICK}_{\rm int,\infty})\).
\end{proposition}

\begin{proof}
If \(F=\iota_N(F_N)=\iota_M(F_M)\) in the directed union, choose a common
level \(L\ge M,N\).  Reflection compatibility gives
\[
 \iota_{LN}(\Theta_NF_N)
 =\Theta_L\iota_{LN}(F_N)
 =\Theta_L\iota_{LM}(F_M)
 =\iota_{LM}(\Theta_MF_M),
\]
so \(\Theta\) is well defined on the union.  Because each \(\Theta_N\) is
isometric and the embeddings are isometric, the induced map is isometric and
therefore extends to the closure.  State compatibility gives, for
\(F=\iota_N(F_N)\),
\[
 \omega(\Theta(F)F)
 =\omega_N(\Theta_N(F_N)F_N)\ge0.
\]
The graded case is identical after replacing multiplication by the supplied
compatible graded OS product.  No limiting regularity beyond these exact
algebraic compatibilities is used.
\end{proof}


\begin{remark}[Embedding compatibility and the role of \texorpdfstring{$m=1$}{m=1}]
\label{p4:rem:truncation-compatibility}
Even after a one-particle assignment of a compatible Paper-II/interface
carrier or symbol ladder, those spaces do not, by themselves, determine
unital embeddings of every associated full matrix algebra.  The embeddings
\(\iota_{MN}\) in Theorem~\ref{p4:thm:inductive-qft-bridge} are therefore explicit
theorem-domain data and must be constructed for the chosen finite-field presentation.
For fermionic CAR systems, nested one-particle spaces give canonical CAR embeddings.
For modewise bosonic truncations with a fixed occupation cap, adjoining new mode
factors gives unital maps \(A\mapsto A\otimes\Id\).  For total-particle-number
truncations the natural maps are generally nonunital corner embeddings, so one must
use the corresponding nonunital inductive system or change the truncation scheme.
The protected datum \(m_p=1\) fixes the microscopic compact momentum seed; the
auxiliary ladder \(m_x=1,2,\ldots\), when separately constructed with the
displayed compatible maps, is a reciprocal-position/interface presentation
and does not enlarge that protected carrier.
\end{remark}


\begin{theorem}[Local or graded-local pre-net from compatible finite support subspaces]
\label{p4:thm:compatible-local-net}
On a local-support theorem domain, for each bounded region \(\mathcal O\) and each truncation level \(N\) let there be a finite mode subspace
\[
\mathcal K_N(\mathcal O)\subset \mathcal K_N
\]
and coherent isometric mode embeddings
\[
v_{MN}:\mathcal K_N\longrightarrow\mathcal K_M,
\qquad
v_{LM}v_{MN}=v_{LN}.
\]
Assume that the algebra embeddings \(\iota_{MN}\) are induced by these mode
maps on the chosen CAR or bosonic truncation scheme, so, for example,
\(\iota_{MN}(\mathfrak a_N(f))=\mathfrak a_M(v_{MN}f)\).  Require:
\begin{enumerate}[label=(\roman*), leftmargin=2em]
\item \emph{isotony at fixed level}: \(\mathcal O_1\subset \mathcal O_2\Rightarrow \mathcal K_N(\mathcal O_1)\subset \mathcal K_N(\mathcal O_2)\),
\item \emph{embedding compatibility}:
      \(v_{MN}(\mathcal K_N(\mathcal O))
      \subset \mathcal K_M(\mathcal O)\) for \(M\ge N\),
\item \emph{finite locality}: for disjoint regions, the corresponding finite supported generators commute in the bosonic case and graded-commute in the fermionic case.
\end{enumerate}
Let \(\mathcal A_N(\mathcal O)\subset \mathcal A_N\) be the finite local algebra
generated by smeared creation and annihilation operators and by
number-preserving polynomials built from \(\mathcal K_N(\mathcal O)\).
Then
\[
\mathcal A_{\mathrm{loc}}(\mathcal O):=
\bigcup_N \iota_N\bigl(\mathcal A_N(\mathcal O)\bigr)
\]
is an isotonic \(*\)-subalgebra of \(\mathcal A_{\mathrm{loc}}\), and its
norm closure
\[
\mathcal A_\infty(\mathcal O):=
\overline{\mathcal A_{\mathrm{loc}}(\mathcal O)}
\]
defines an algebraic local pre-net in the bosonic branch and a graded-local
field pre-net in the fermionic CAR branch.  The even fermionic observable
subnet is ordinarily local and hence Haag--Kastler-type.
In particular, the inductive bridge of
Theorem~\ref{p4:thm:inductive-qft-bridge} already yields the local algebraic
object customarily used in continuum QFT whenever the finite support calculus is
fixed.
\end{theorem}


\begin{proof}
For each fixed \(\mathcal O\), condition (ii) and the induced-generator
identity for \(\iota_{MN}\) imply
\(\iota_{MN}(\mathcal A_N(\mathcal O))
\subseteq\mathcal A_M(\mathcal O)\).
Thus the family \(\mathcal A_N(\mathcal O)\) is compatible under the algebra
embeddings, so its directed union is a \(*\)-subalgebra of
\(\mathcal A_{\mathrm{loc}}\).
If \(\mathcal O_1\subset \mathcal O_2\), condition (i) gives \(\mathcal A_N(\mathcal O_1)\subset \mathcal A_N(\mathcal O_2)\) for every \(N\), hence
\[
\mathcal A_{\mathrm{loc}}(\mathcal O_1)\subset \mathcal A_{\mathrm{loc}}(\mathcal O_2),
\]
so isotony holds on the directed union and therefore on the closures.
Because each embedding is isometric, the inherited norm from Theorem~\ref{p4:thm:inductive-qft-bridge} is well defined on every local union, and the completion gives \(\mathcal A_\infty(\mathcal O)\).
For disjoint regions, condition (iii) states that the finite generators commute or graded-commute at each level; since the embeddings preserve products and adjoints, the same relation holds in the directed union and survives norm closure.
Thus the completed family
\(\mathcal O\mapsto \mathcal A_\infty(\mathcal O)\) is isotonic and is local
or graded-local according to the branch.  In the CAR branch, even
observables in disjoint regions commute, so the even subnet is the ordinary
local observable pre-net. \qedhere
\end{proof}


\paragraph{Bandlimited localization domain.}
The exact-support spaces in
Theorem~\ref{p4:thm:compatible-local-net} are not identified with sharp
spectral subspaces.  On the bandlimited presentation, let \(E_N\) be the
spatial spectral projection on the smooth parent one-particle space.  For a
bounded region \(\mathcal O\), let
\[
J_{\mathcal O}:L^2(\mathcal O)\longrightarrow L^2(S^3_x)
\]
be extension by zero and define the contraction
\[
L_N(\mathcal O):=E_NJ_{\mathcal O}.
\]
Its range is only approximately localized.  Let
\(\mathcal A_N^{\rm sp}(\mathcal O)\) be generated by the fields
\[
\mathfrak a^\#(L_N(\mathcal O)f)
\]
on the chosen CAR, regular-CCR, or occupation-cutoff theorem domain.

\paragraph{Claim-specific dynamical-locality domains.}
For a selected compatible Hamiltonian family \(H_N\), distinguish
\[
\begin{array}{c|l}
(\mathrm{DLOC}_{\rm alg})
 & \text{bounded local-algebra propagation/approximant estimate},\\
(\mathrm{DLOC}_{\rm pkt})
 & \text{one-particle packet leakage and return estimate},\\
(\mathrm{DLOC}_{\rm LR})
 & \text{direct Lieb--Robinson or weighted-conjugation bound}.
\end{array}
\]
The legacy tag \((\mathrm{DLOC}_H)\) is incomplete unless one of these
subscripts and its norm/domain are stated.  These are dynamical properties of
\(H_N\).  They are not Paper~III-F's \(({\rm QLOC})\), which controls
equal-time compressed multiplication observables, and none follows from the
numerical speed ceiling alone.

\begin{theorem}[Weighted-conjugation locality for quadratic CAR dynamics]
\label{p4:thm:weighted-conjugation-car-locality}
Let \(h_N=h_N^*\) be a finite one-particle Hamiltonian with localization
projections \(\mathbf1_E,\mathbf1_F\).  Let \(\rho\) be a bounded
one-particle multiplication weight with \(\rho|_E=0\) and
\(\rho|_F\ge d(E,F)\).  Suppose that for \(0<\mu\le\mu_0\)
\[
 \bigl\|e^{\mu\rho}h_Ne^{-\mu\rho}-h_N\bigr\|
 \le\mu v_N.
\]
Then
\[
 \boxed{\|\mathbf1_Fe^{-it h_N}\mathbf1_E\|
 \le e^{-\mu(d(E,F)-v_N|t|)}.}
\]
For the number-conserving quadratic CAR Hamiltonian
\(\mathbb H_N=d\Gamma(h_N)\),
\[
 \alpha_t(c(f))=c(e^{-it h_N}f),
\]
so cross-region CAR anticommutators obey the same one-particle bound.
Finite even CAR polynomials satisfy an ordinary commutator bound obtained by
the graded Leibniz rule, with a constant depending only on their degrees and
norms.  Hence this hypothesis proves
\((\mathrm{DLOC}_{\rm LR})\) and the corresponding CAR instance of
\((\mathrm{DLOC}_{\rm alg})\).
\end{theorem}

\begin{proof}
Writing \(B_\mu=e^{\mu\rho}h_Ne^{-\mu\rho}-h_N\), the interaction-picture
Duhamel estimate gives
\[
 \|e^{\mu\rho}e^{-it h_N}e^{-\mu\rho}\|
 \le e^{|t|\|B_\mu\|}
 \le e^{\mu v_N|t|}.
\]
The weights on \(E\) and \(F\) contribute \(1\) and at most
\(e^{-\mu d(E,F)}\), respectively, proving the first bound.  The quadratic
CAR evolution formula follows from the CAR Heisenberg equation.  The
elementary anticommutator is the transported one-particle inner product, and
repeated graded Leibniz expansion proves the polynomial statement.
\end{proof}

\begin{theorem}[Approximate microcausality with spectral, occupation, and dynamical controls]
\label{p4:thm:approximate-microcausality}
Let \(A_N\in\mathcal A_N^{\rm sp}(\mathcal O_1)\) and
\(B_N\in\mathcal A_N^{\rm sp}(\mathcal O_2)\) be bounded finite
polynomials at a reference time \(s\) on the CAR or
occupation-truncated matrix branch.  On the regular-CCR branch, instead
take bounded finite linear combinations of Weyl words, with the number of
words and coefficient sums uniformly bounded in \(N\).  Let
\(A_N(t)=U_N(t,s)^*A_NU_N(t,s)\), and put
\[
\delta:=d(\mathcal O_1,\mathcal O_2)-c|t-s|>0.
\]
On the CAR and occupation-truncated matrix branches, assume
\((\mathrm{DLOC}_{\rm alg})\) in the following explicit form: there is
\(A_N^{\rm loc}(t)\in
\mathcal A_N^{\rm sp}(\mathcal O_1^{\,c|t-s|})\) such that
\[
\|A_N(t)-A_N^{\rm loc}(t)\|
\le\varepsilon_{{\rm dyn},N}(t,s).
\]
On the regular-CCR Weyl branch do not assume a norm-local Weyl approximant.
Instead assume the direct symplectic propagation estimate
\[
 |\sigma_N(U_N(t,s)f,g)|
 \le\varepsilon_{{\rm LR},N}(\delta)\|f\|\,\|g\|
\]
for the smeared vectors entering the two Weyl words.  This is the
\((\mathrm{DLOC}_{\rm LR})\) input for that branch.
Assume the spatial-localizer overlap estimate
\begin{equation}
\label{p4:eq:spectral-leakage-bound}
\bigl\|L_N(\mathcal O_1)^*L_N(\mathcal O_2)\bigr\|
\le
\varepsilon_{{\rm sp},N}\!\left(d(\mathcal O_1,\mathcal O_2)\right)
\le
C_{\rm loc}\bigl(K_Nd(\mathcal O_1,\mathcal O_2)\bigr)^{-2}
\end{equation}
uniformly for the enlarged region pairs used above, where \(K_N\) is the
largest resolved spatial momentum.  Let \(Q_N\) project onto an occupation
sector chosen away from the bosonic truncation boundary. On the truncated
bosonic branch let \(d_N\) bound the total number of elementary creation or
annihilation factors in each ordered word appearing in either product
\(A_N^{\rm loc}(t)B_N\) or \(B_NA_N^{\rm loc}(t)\). Assume either
\(Q_N\le\mathbf1_{\{n\le N-d_N\}}\) for a total-occupation
cutoff \(N\), or an explicit compressed \emph{word-level} bound, including
all intermediate occupation excursions, by
\(\varepsilon_{{\rm occ},N}(Q_N)\) for those expanded products.
The first alternative gives \(\varepsilon_{{\rm occ},N}=0\);
a small compressed elementary CCR defect alone is not sufficient.  In fermionic sectors take physical even
observables, or use the graded commutator for odd generators.  On the CAR
and regular-CCR branches set \(Q_N=\mathbf1\) and
\(\varepsilon_{{\rm occ},N}=0\).  Then
\begin{equation}
\label{p4:eq:approx-microcausality-bound}
\begin{aligned}
\bigl\|Q_N[A_N(t),B_N]_{\pm}Q_N\bigr\|
\le{}&
C(A_N,B_N)\varepsilon_{{\rm sp},N}(\delta)\\
&+
C'(A_N,B_N)\varepsilon_{{\rm occ},N}(Q_N)
+\mathcal E_{{\rm dyn},N}(A_N,B_N),
\end{aligned}
\end{equation}
where
\[
\mathcal E_{{\rm dyn},N}=
\begin{cases}
2\|B_N\|\varepsilon_{{\rm dyn},N}(t,s),
&\text{CAR or occupation-truncated matrix branch},\\
C_{\rm W}(A_N,B_N)\varepsilon_{{\rm LR},N}(\delta),
&\text{regular-CCR Weyl branch}.
\end{cases}
\]
In the Weyl branch the direct symplectic estimate may absorb the spatial
overlap completely, in which case take \(C(A_N,B_N)=0\).
Consequently, on any compatible refinement for which the polynomial
degrees and operator norms, or the Weyl-word coefficient bounds, keep
\(C(A_N,B_N)\), \(C'(A_N,B_N)\), and \(\|B_N\|\) uniformly bounded and
\(\varepsilon_{{\rm sp},N}(\delta)\to0\),
\(\varepsilon_{{\rm occ},N}(Q_N)\to0\), and
\(\mathcal E_{{\rm dyn},N}\to0\), the compressed finite observables
satisfy microcausality in norm outside the GP1 cone.  With
\(\ell_N=K_N^{-1}\), the
spatial term is
\begin{equation}
\label{p4:eq:framework-microcausality-bound}
O\bigl((\delta/\ell_N)^{-2}\bigr).
\end{equation}
\end{theorem}

\begin{remark}[Status of the three error controls]
\label{p4:rem:microcausality-leakage-status}
The spatial-localizer estimate and the occupation-boundary estimate are
theorem-domain hypotheses in Paper~IV.  Exact finite CAR is available in
fermionic sectors, whereas a finite bosonic carrier necessarily has a
boundary defect.  The continuous-QFT handoff may close these estimates for
a specified admissible refinement system; they are not hidden consequences
of finite dimensionality alone.  The dynamical leakage estimate is the
separate branch-specific \((\mathrm{DLOC}_{\rm alg})\) or
\((\mathrm{DLOC}_{\rm LR})\) statement: the GP1 speed bound defines the
operational cone but does not by itself prove that the chosen Hamiltonian
transports approximately local algebras inside that cone.  Paper~III-F's
static \(({\rm QLOC})\) estimate may be imported when the corresponding
compressed-observable presentation is used, but it neither replaces nor
implies either dynamical-locality tag.
\end{remark}

\begin{proof}
The enlarged region \(\mathcal O_1^{\,c|t-s|}\) is separated from
\(\mathcal O_2\) by \(\delta\).  In the CAR branch, the cross-region
elementary graded commutator is the one-particle overlap determined by
\(L_N(\mathcal O_1^{\,c|t-s|})^*L_N(\mathcal O_2)\), hence is bounded by
\eqref{p4:eq:spectral-leakage-bound}.  On the regular-CCR branch the direct
propagated symplectic estimate and the Weyl relation give
\[
\|[W(f),W(g)]\|
\le 2|\sin(\sigma(f,g)/2)|
\le |\sigma(f,g)|,
\]
so finite Weyl-word expansions give the branch-specific
\(C_{\rm W}\varepsilon_{{\rm LR},N}\) term directly.  No operator-norm
local Weyl approximant is used.  In a truncated bosonic sector one must expand each ordered word \emph{before}
compressing by \(Q_N\): intermediate factors can encounter the cutoff
when the elementary compressed CCR defect vanishes. Under the stated
\(d_N\)-occupation margin, every intermediate word acting on
\(Q_N\mathcal H\) stays below the cutoff, so the canonical word
relations hold exactly there. Under the alternative direct word-level
hypothesis their deviation is bounded by
\(\varepsilon_{{\rm occ},N}(Q_N)\) after the finite coefficient
and norm constants are included. Repeated graded Leibniz expansion, using
\[
[XY,Z]_{\pm}=X[Y,Z]_{\pm}\pm[X,Z]_{\pm}Y,
\]
then gives the displayed bound. Neither alternative follows merely
from \(\|Q_N([a,a^\dagger]-1)Q_N\|\) being small.  Finally,
On the CAR and occupation-truncated matrix branches,
\[
\|[A_N(t)-A_N^{\rm loc}(t),B_N]_\pm\|
\le2\|B_N\|\varepsilon_{{\rm dyn},N}(t,s),
\]
which gives their branch-specific final term.  The stated limit follows when all three
independent errors vanish.
\end{proof}

\begin{remark}[Exact-support locality versus bandlimited localization]
\label{p4:rem:microcausality-status}
Theorem~\ref{p4:thm:compatible-local-net} gives exact equal-time graded
locality only on its explicitly supplied exact-support domain.
Theorem~\ref{p4:thm:approximate-microcausality} is a different,
bandlimited-localizer domain; it is not an upgrade obtained from the
exact-support theorem alone.  Sharp support and sharp spectral bandlimit
cannot both hold for a nonzero smooth-parent mode, so the latter theorem
carries a spatial-overlap error.  A total-particle-number bosonic truncation
may also fail exact factor locality at its occupation boundary unless
condition~(iii) of Theorem~\ref{p4:thm:compatible-local-net} is separately
verified.  Spatial, occupation, and dynamical errors must all vanish before
continuum microcausality is claimed.
\end{remark}


\begin{theorem}[Exact CAR and regular-CCR quasi-free inductive reconstruction]
\label{p4:thm:quasifree-reconstruction}
Let \(\mathcal K_N\) be nested one-particle spaces with compatible
isometries \(v_{MN}\).  Consider either:
\begin{enumerate}[label=(\alph*), leftmargin=2em]
\item the canonical finite CAR algebras with their induced unital CAR
embeddings; or
\item the regular Weyl CCR algebras
\(\mathfrak W(\mathcal K_N,\sigma_N)\) with compatible symplectic
embeddings.
\end{enumerate}
Let \(\omega_N\) be compatible quasi-free states, determined by compatible
two-point forms \(C_N\), so that
\[
\omega_M\circ\iota_{MN}=\omega_N,
\qquad
C_M(v_{MN}f,v_{MN}g)=C_N(f,g)
\quad(M\ge N).
\]
In the Weyl branch, assume the usual CCR covariance positivity and
regularity at every level.  Then the corresponding CAR or Weyl inductive
limit carries a unique compatible quasi-free state \(\omega\).  Its
finite-support \(n\)-point functions are the stabilized Wick pairings.
Consequently the GNS triple
\((\pi_\omega,\mathcal H_\omega,\Omega_\omega)\) reconstructs the
corresponding quasi-free field presentation on the algebraic core
\[
\mathcal D_\omega:=
\mathrm{span}\,\pi_\omega(\mathcal A_{\mathrm{loc}})\Omega_\omega.
\]
The CAR branch is an inductive limit of finite matrix algebras.  The regular
CCR branch is an interface algebra and is not a finite
occupation-truncated carrier.
\end{theorem}

\begin{proof}
In the CAR branch,
Theorem~\ref{p4:thm:inductive-qft-bridge} gives the inductive algebra and
the unique state extending the compatible family.  In the Weyl branch,
symplectic isometries induce injective, hence isometric, unital
\(*\)-homomorphisms of Weyl algebras.  The same directed-union completion
and compatible-state argument gives the regular CCR inductive algebra and
state.

For every local CAR polynomial, or every bounded finite Weyl word, built
from finitely many smeared modes, choose \(N\) large enough to contain its
mode support.  Then
\[
\omega(P)=\omega_N(P),
\]
independently of the chosen sufficiently large level.  Quasi-free expansion
and covariance compatibility make the CAR pairings and Weyl characteristic
functions stabilize.  Field moments in the regular Weyl branch are the
derivatives at the identity on the common smooth GNS domain, not
\(C^*\)-algebra elements.  This proves that \(\omega\) is quasi-free.  Its
GNS construction gives the displayed dense algebraic core, with the smooth
Weyl-vector subcore used for unbounded field derivatives. \qedhere
\end{proof}

\begin{remark}[Occupation-truncated bosonic quasi-free boundary]
\label{p4:rem:truncated-boson-quasifree-boundary}
A matrix algebra with a finite occupation cutoff is not a regular CCR
algebra.  It does not carry an exact quasi-free state for all polynomials in
its truncated ladder operators.  The finite commutator defect of
Proposition~\ref{p4:prop:no-exact-ccr} prevents that identification.

For such a family the valid handoff is a degree-controlled moment limit.
For every fixed finite mode support and word degree \(r\), require the
discarded-tail and boundary-layer terms in
Theorem~\ref{p4:thm:finite-wick-quasifree-calculus} to vanish as
\(N_{\max}\to\infty\), together with compatibility of the limiting
two-point forms.  The limiting positive moments then define the quasi-free
state on the polynomial CCR algebra.  Regularity of the resulting Weyl
state is a further continuity condition.  A fixed occupation cap, even
along a large-\(m_x\) mode ladder, does not discharge this requirement.
\end{remark}


\begin{proposition}[Weak operator-valued-distribution bridge for smeared fields]
\label{p4:prop:weak-ovd-smeared-fields}
Let
\[
\mathfrak T:=\bigcup_N \mathcal K_N
\]
carry any locally convex test-profile topology for which the truncation
embeddings are continuous.  Work on one of the exact CAR/regular-CCR
inductive domains of
Theorem~\ref{p4:thm:quasifree-reconstruction}.  An occupation-truncated
family enters only after the moment-limit construction of
Remark~\ref{p4:rem:truncated-boson-quasifree-boundary}, through the resulting
algebraic polynomial-CCR moment functional rather than through a fictitious
exact embedding of the truncated ladder fields; denote that functional by
\(\omega_{\rm mom}\).
Let \(\omega\) be either a quasi-free state whose two-point form is
continuous on \(\mathfrak T\times\mathfrak T\), or more generally a state
for which all local matrix-element correlation functions needed on the
polynomial GNS core are continuous in the chosen test-profile topology.
For \(f\in\mathfrak T\) and local algebraic \(A,B\), choose a level \(N\)
large enough to contain \(f\) and the finite supports of \(A,B\).  On the
CAR branch let \(\Phi_N(f)\) be the bounded
smeared field.  On the regular-CCR branch write \(W_N(tf)\) for the Weyl
unitaries and require the local vectors used below to be smooth for their
one-parameter groups.  Define on \(\mathcal D_\omega\)
\[
\begin{aligned}
&\langle \pi_\omega(B)\Omega_\omega,\,
\Phi_\omega(f)\,\pi_\omega(A)\Omega_\omega\rangle\\
&\quad:=
\begin{cases}
\omega\!\bigl(B^*\,\iota_N(\Phi_N(f))\,A\bigr),
&\text{CAR branch},\\[1mm]
-i\,\displaystyle\left.\frac{d}{dt}\right|_{t=0}
\omega\!\bigl(B^*\,\iota_N(W_N(tf))\,A\bigr),
&\text{regular-CCR branch},\\[3mm]
\omega_{\rm mom}\!\bigl(B^*\Phi(f)A\bigr),
&\text{polynomial-CCR moment-limit branch},
\end{cases}
\end{aligned}
\]
for local algebraic \(A,B\), where the last line is evaluated by the
constructed moment functional.  Thus no unbounded Weyl field is treated as
an element of a \(C^*\)-algebra.
This is independent of the chosen sufficiently large level \(N\), linear in \(f\), and for every fixed core vectors \(\psi,\chi\in\mathcal D_\omega\) the map
\[
f\longmapsto \langle \chi,\Phi_\omega(f)\psi\rangle
\]
is continuous on \(\mathfrak T\).
Hence \(f\mapsto\Phi_\omega(f)\) is a weak sesquilinear-form-valued distribution on the algebraic GNS core.  When the additional boundedness or closability condition holds so that these forms are represented by operators on the core, this is the usual weak operator-valued-distribution statement.
\end{proposition}


\begin{proof}
If \(f\in\mathcal K_N\subset \mathcal K_M\) with \(M\ge N\), embedding
compatibility gives the corresponding field identity in the CAR branch and
\[
\iota_{MN}(W_N(tf))=W_M(tv_{MN}f)
\]
in the Weyl branch.  Differentiating the compatible Weyl matrix elements at
zero on the stated common smooth domain proves level independence there.
The moment-limit branch is level independent by construction.  Linearity in
\(f\) follows from the CAR/polynomial field relations and, in the Weyl
branch, from the regular infinitesimal representation of the Weyl
relations.
For fixed \(A,B\in\mathcal A_{\mathrm{loc}}\), the map
\[
f\mapsto
\langle\pi_\omega(B)\Omega_\omega,\Phi_\omega(f)
\pi_\omega(A)\Omega_\omega\rangle
\]
is built from the finite two-point function in the quasi-free case, or from the finite local higher correlations in the general hypothesis; by the strengthened hypothesis those matrix elements are continuous in the chosen test-profile topology.
Therefore each matrix element of \(\Phi_\omega(f)\) is continuous in \(f\), which is exactly the weak form-valued distribution statement; the operator-valued conclusion additionally uses the stated boundedness or closability condition. \qedhere
\end{proof}


\begin{theorem}[Typed finite coarse-graining and coefficient-transfer cocycles]
\label{p4:thm:renormalization-semigroup}
\leavevmode\par\noindent
Let \(\Phi^*_{M\leftarrow N}\) be compatible CPTP coarse-graining maps on states
from the finer finite resolution \(N\) to the coarser resolution \(M\le N\), and
let \(\Phi_{N\leftarrow M}\) denote the dual unital completely positive maps on
observables.  The state-picture maps satisfy
\[
\Phi^*_{L\leftarrow M}\circ \Phi^*_{M\leftarrow N}
=
\Phi^*_{L\leftarrow N},
\]
and the dual observable-picture maps compose contravariantly.
This state-map cocycle alone does not define an effective Hamiltonian or a
map on coupling constants.

On a threshold-free shell interval \(I_r\), let
\(G_r\cong\mathbb R^{d_r}\) be the coefficient space of a supplied closed
operator basis \(\{O_a^{(r)}\}\), and write
\[
H_r(g):=\sum_a g_aO_a^{(r)}.
\]
Assume that a stated effective matching prescription supplies maps
\[
R^{(r)}_{\ell_2,\ell_1}:G_r\longrightarrow G_r,
\qquad \ell_2\ge\ell_1,
\]
and is functorial under successive matching.  Then the typed coefficient
cocycle is
\[
R^{(r)}_{\ell_3,\ell_1}
=
R^{(r)}_{\ell_3,\ell_2}\circ R^{(r)}_{\ell_2,\ell_1}.
\]
If, locally uniformly in \(g\),
\[
R^{(r)}_{\ell+d\ell,\ell}(g)
=
g+d\ell\,\beta_r(\ell,g)+o(d\ell),
\]
then its trajectories obey
\[
\frac{dg}{d\ell}=\beta_r(\ell,g).
\]
The affine form
\(\beta_r(\ell,g)=A_r(\ell)g+B_r(\ell)\) follows only when the supplied
matching prescription preserves an affine finite basis.

Across a threshold, let the supplied typed matching map be
\[
M_r:G_{r,-}\longrightarrow G_{r,+}.
\]
When the retained Hilbert space and dimension are unchanged and the
threshold is only a basis conversion, \(M_r\) may be induced by
\(H\mapsto U_rHU_r^*\).  Channel opening, elimination, or dimension change
instead uses the appropriate Schur-complement, projection/embedding,
partial-isometry, or CP matching map.  The global flow is the ordered
composition of the interval maps \(R^{(r)}\) and threshold maps \(M_r\).
\end{theorem}


\begin{proof}
The first cocycle is the assumed compatibility of the supplied CPTP maps.
Closure of a chosen effective basis can define coefficient-transfer maps,
but their cocycle follows only from functoriality of the matching
prescription: projection back to a finite basis after each step can
otherwise destroy composition.  Differentiating the typed cocycle gives
the beta-vector field.  Threshold bookkeeping is ordinary composition of
the typed maps \(M_r\); unitary conjugation is only its
same-retained-space special case. \qedhere
\end{proof}

\begin{corollary}[Renormalization as an exact-continuation descendant]
\label{p4:cor:rg-universal-continuation-instance}
The state-map part of
Theorem~\ref{p4:thm:renormalization-semigroup} is an instance of
Theorem~\ref{p3a:thm:universal-generated-exact-continuation} with the
resolution preorder as \(\mathcal S_D\), the finite state spaces as its
objects, and CPTP maps as its arrow type.  It is generally a projective
cocycle rather than one autonomous time semigroup.

A closed coefficient flow is a descendant of that parent precisely when the
chosen coefficient readouts satisfy the projectability condition of
Theorem~\ref{p3a:thm:generated-process-projectability}.  On that domain the
functorial matching maps \(R^{(r)}\) are the descended continuation.  If they
are differentiable in a homogeneous scale coordinate, their infinitesimal
law is the displayed beta equation.  Thus finite nested resolution and exact
composition derive the RG architecture, while basis closure, the matching
prescription, differentiability, and the sector-specific numerical
\(\beta_r\) remain exactly the typed obligations stated in the theorem.
\end{corollary}

Thus algebraic local nets, quasi-free field reconstruction, smeared-field
distribution language, and state coarse-graining can be stated on
compatible finite or declared interface systems.  Running and matching
additionally require the typed closed-basis transfer and functoriality data
displayed above.  Fully interacting model-specific continuum
renormalization remains outside the present paper.



\subsection[Time-dependent mode spaces and curved-background bridge]
{Time-dependent modes and curved backgrounds}
\label{p4:sec:part2-curved-bridge}

A cosmological or otherwise time-dependent background is represented here not by postulating an infinite curved-spacetime field theory from the start, but by allowing the distinguished finite mode space itself to vary with the background data.



\begin{theorem}[One-mode nonstationary Hessian and explicit Bogoliubov source]
\label{p4:thm:v245-nonstationary-hessian-bogoliubov-source}
Let a real resolved mode have quadratic action
\[
 S^{(2)}=\frac12\int_I dt\,
 \left[m(t)\dot\eta^2-\ell(t)\eta^2\right],
 \qquad m(t)>0,
\]
with \(m\in C^3(I)\), \(\ell\in C^1(I)\), and the
resulting \(\Omega_{\rm eff}^2\) strictly positive on the interval
where its instantaneous derivative is used. These regularity conditions
make \(\Omega_{\rm eff}\in C^1\). Put \(q=\sqrt m\,\eta\).  Up to
an endpoint term,
\[
 S^{(2)}\doteq\frac12\int_I dt\,
 \left[\dot q^2-\Omega_{\rm eff}^2q^2\right],
\]
where
\[
 \boxed{
 \Omega_{\rm eff}^2
 =\frac{\ell}{m}
 -\frac12\frac{\ddot m}{m}
 +\frac14\left(\frac{\dot m}{m}\right)^2.}
\]
On any interval where \(\Omega:=\Omega_{\rm eff}>0\), define
\[
 a(t)=\frac1{\sqrt{2\hbar}}
 \left(\Omega^{1/2}q+i\Omega^{-1/2}p\right),
 \qquad p=\dot q.
\]
Then the exact instantaneous Nambu law is
\[
 \boxed{\dot a=-i\Omega a+\frac{\dot\Omega}{2\Omega}a^\dagger,}
 \qquad
 \boxed{r(t):=\frac{\dot\Omega}{2\Omega}.}
\]
Thus variation of the physical Hessian itself is the positive/negative
frequency pairing source.  The scalar formula applies to a decoupled
commuting mode.  Coupled matrix-valued modes require the full symplectic
connection/Nambu theorem below, and instantaneous particle language is not
asserted across intervals where \(\Omega^2\le0\).
\end{theorem}

\begin{proof}
Substitution of \(\eta=m^{-1/2}q\) gives a mixed term
\(-\tfrac12(\dot m/m)q\dot q\); integration by parts yields the displayed
\(\Omega_{\rm eff}^2\).  Differentiating \(a(t)\), using
\(\dot q=p\) and \(\dot p=-\Omega^2q\), and expressing \(q,p\) back in
terms of \(a,a^\dagger\) gives the stated two terms exactly.
\end{proof}



\begin{corollary}[Affine-defect form of the one-mode pairing source]
\label{p4:cor:v247-defect-bogoliubov-source}
Under the hypotheses of
Theorem~\ref{p4:thm:v245-nonstationary-hessian-bogoliubov-source}, suppose a
resolved mode is written relative to a fixed affine harmonic number \(\nu>0\)
as
\[
 \Omega^2(t)=\nu^2+\mathfrak D(t),
\]
where \(\mathfrak D\) is the affine spectral defect of
Paper~III-A, Definition~\ref{p3a:def:v247-affine-spectral-defect}, in the
corresponding dimensionless clock normalization.  Then
\[
 \boxed{\dot a=-i\Omega a
 +\frac{\dot{\mathfrak D}}{4(\nu^2+\mathfrak D)}a^\dagger.}
\]
Hence a constant nonzero defect is a stationary detuning but is not by itself
a particle-production source, whereas \(\dot{\mathfrak D}\ne0\) is the
single-mode positive/negative-frequency mixing source.  This formula retains
the same decoupled-mode and \(\Omega^2>0\) scope as the parent theorem.
\end{corollary}
\begin{proof}
Differentiate \(\Omega^2=\nu^2+\mathfrak D\) to obtain
\(2\Omega\dot\Omega=\dot{\mathfrak D}\), and substitute into the exact
coefficient \(\dot\Omega/(2\Omega)\).
\end{proof}

\begin{theorem}[Two-level fermionic mixing from a time-dependent mass owner]
\label{p4:thm:v247-two-level-fermion-mixing}
Let a resolved positive/negative-energy fermionic pair obey
\[
 i\frac{d\Psi}{d\tau}=\mathcal H(\tau)\Psi,
 \qquad
 \mathcal H(\tau)=\nu\sigma_3+M(\tau)\sigma_1,
 \qquad \nu>0,
\]
with real \(M\in C^1\).  Put
\[
 \Omega=\sqrt{\nu^2+M^2},
 \qquad
 \tan\theta=\frac{M}{\nu},
 \qquad
 R_\theta=e^{-i\theta\sigma_2/2}.
\]
For \(\Xi=R_\theta^\dagger\Psi\),
\[
 \boxed{
 i\frac{d\Xi}{d\tau}
 =\left(\Omega\sigma_3-\frac{\theta'}2\sigma_2\right)\Xi,}
\]
and the magnitude of the instantaneous positive/negative-energy mixing source
is
\[
 \boxed{r_f=\frac{|\theta'|}{2}
 =\frac{\nu|M'|}{2(\nu^2+M^2)}.}
\]
Thus a constant mass owner changes the instantaneous spectrum but does not
produce local pair mixing in this basis; a changing mass owner does.  This is
a local finite mode theorem and does not construct a fermion determinant,
Pfaffian, Majorana reality structure, or regulated quantum measure.
\end{theorem}
\begin{proof}
The rotation satisfies
\(R_\theta^\dagger\mathcal H R_\theta=\Omega\sigma_3\).  Differentiating
\(\Psi=R_\theta\Xi\) gives the connection term
\(-iR_\theta^\dagger R_\theta'=-(\theta'/2)\sigma_2\).  Finally,
\(\theta'=\nu M'/(\nu^2+M^2)\).
\end{proof}

\begin{theorem}[Primitive-turn conformal monodromy and exact finite Floquet reduction]
\label{p4:thm:v247-conformal-floquet}
Let \(\mathcal H_{\rm conf}\) be a finite free carrier assembled from the
clock-matched conformal scalar and transverse gauge towers and from the
massless Dirac tower of Paper~III-A,
Theorem~\ref{p3a:thm:v247-conformal-clock-hierarchy}.  On a selected finite
bosonic/fermionic occupation realization let
\[
 P_F:=(-1)^{\widehat N_F}
\]
be fermion parity and allow a scalar vacuum-energy offset \(E_0I\).  Then
\[
 \boxed{U_{\rm conf}(T_t)=e^{-iE_0T_t/\hbar}P_F.}
\]
Let \(V=V^*\) be any time-independent finite interaction with
\([V,P_F]=0\), and define
\[
 V_I(t)=U_{\rm conf}(t)^\dagger VU_{\rm conf}(t).
\]
Then
\[
 \boxed{V_I(t+T_t)=V_I(t).}
\]
Consequently the exact one-turn interaction-picture evolution is the finite
Floquet unitary
\[
 \boxed{\mathcal F
 =\mathcal T\exp\!\left[-\frac{i}{\hbar}
 \int_0^{T_t}V_I(t)\,dt\right].}
\]
If \(e^{-i\varepsilon_\alpha T_t/\hbar}\) is an eigenvalue of
\(\mathcal F\), its quasienergy is defined modulo
\(2\pi\hbar/T_t=\Eper\).  No small-interaction or Magnus-convergence
assumption is required for this exact finite statement.

If \(V_I(t)=\sum_{q\in\mathbb Z}V_qe^{-2\pi iqt/T_t}\), its one-period
average is exactly \(V_0\).  Thus the previously used zero-net-clock-number
selection rule is the first period/Fourier projection of the exact Floquet
problem; nonzero-\(q\) terms remain in the full time-ordered unitary and are
not declared absent.
\end{theorem}
\begin{proof}
Every bosonic conformal excitation has integer energy in units of \(\Eper\),
while every fermionic excitation has half-integer energy.  On an occupation
state the free one-turn phase is therefore the common vacuum phase times
\((-1)^{N_F}\), proving the first formula.  The vacuum phase cancels from
conjugation, and
\[
 V_I(t+T_t)=P_FV_I(t)P_F=V_I(t)
\]
by parity evenness.  Finite-dimensional time-dependent Schr\"odinger
existence gives the unitary \(\mathcal F\), and the quasienergy period follows
from \(T_t\Eper=2\pi\hbar\).  The final statement is ordinary Fourier
orthogonality on one period.
\end{proof}

\begin{theorem}[Finite interface for supplied time-dependent mode spaces]
\label{p4:thm:curved-background-finite-bridge}
Let \(\mathcal K\to I\) be a finite-rank \(C^1\) Hermitian vector bundle over
a time interval \(I\), equipped with a supplied metric connection
\(\nabla_t\).  Let \(h_t=h_t^*\) and \(\Delta_t\) be norm-continuous bundle
endomorphism data defining the fiberwise quadratic many-body Hamiltonian
\[
H_t=d\Gamma(h_t)+\frac12\bigl(\mathfrak a^\dagger \Delta_t \mathfrak a^\dagger + \mathfrak a\,\Delta_t^\dagger \mathfrak a\bigr).
\]
Use the induced connection \(\nabla_t^{\mathcal F}\) and the covariant
Schr\"odinger equation
\[
i\hbar\nabla_t^{\mathcal F}\Psi(t)=H_t\Psi(t).
\]
There are three distinct theorem domains:
\begin{enumerate}[label=(\alph*), leftmargin=2em]
\item the finite CAR Fock bundle;
\item a finite occupation-truncated bosonic Fock bundle; or
\item the regular bosonic CCR Fock bundle, provided the quadratic
Hamiltonians share a specified invariant dense domain, are self-adjoint
there, satisfy a standard common-domain condition ensuring a strongly
continuous unitary propagator, and their Bogoliubov evolution is unitarily
implementable.
\end{enumerate}
In cases (a)--(b), the many-body Hamiltonians are finite matrices and the
evolution is a unitary fiber-to-fiber propagator.  In case (c), the same
conclusion follows only from the displayed unbounded-operator and
implementability hypotheses; finite rank of the one-particle bundle does not
make the regular bosonic Fock space finite dimensional.

In finite CAR sectors, and in the regular bosonic CCR sector on domain (c),
the Heisenberg evolution of the mode operators induces the
Bogoliubov transformation
\[
\mathfrak a_t(f)=\mathfrak a_s(\alpha_{t,s}f)+\mathfrak a_s^\dagger(\beta_{t,s}\overline f).
\]
For finite occupation-truncated bosonic sectors the same formula holds up to
boundary-defect terms supported on the occupation cutoff.  Those terms
vanish only in specified state expectations or common-domain matrix
elements under an independent boundary-layer-weight estimate, or in an
explicit occupation-cutoff-removal topology.  No operator-norm convergence
of the truncated CCR defect is asserted, and enlarging the spatial window
alone does not control it.  Relative to a supplied instantaneous positive-frequency
splitting, nonzero \(\beta_{t,s}\) represents particle production,
squeezing, or vacuum mismatch on the supplied one-particle carrier.

If a separately supplied reciprocal-position/interface system provides
compatible one-particle embeddings, the induced CAR or Weyl algebra maps
(or explicitly controlled truncated-boson maps), compatible localizers,
states, and finite evolutions, then its observables fit into the
directed-system structure of
Theorem~\ref{p4:thm:inductive-qft-bridge}.  Paper~II supplies nested
retained-symbol spaces and their density statement; it does not by itself
supply these many-body, local-observable, state, or time-evolution maps.
The corresponding curved-background quantities arise as interface limits
only after that additional compatible system has been exhibited.
The bundle connection or an equivalent coherent trivialization is essential:
a \(C^1\) list of abstract mode spaces alone does not identify vectors at
different times and therefore does not define evolution.
\end{theorem}


\begin{proof}
\ProofPlan{We prove the finite curved-background bridge by choosing a finite moving frame, deriving the Nambu/Bogoliubov evolution, reading the boundary-defect and particle-production content, and then passing to the compatible truncation ladder.}
\begin{proofsteps}

\item \textbf{Finite-rank covariant time evolution.}
Parallel transport for the supplied metric connection gives a \(C^1\)
orthonormal trivialization of \(\mathcal K\) over the interval.  In a
parallel frame the connection coefficient vanishes, \(h_t\) and \(\Delta_t\)
are norm-continuous finite one-particle matrices.  On the finite CAR and
occupation-truncated bosonic carriers, \(H_t\) is consequently a
norm-continuous self-adjoint finite matrix, so the covariant Schr\"odinger
equation has a unique unitary propagator
\(U(t,s):\mathcal F_s\to\mathcal F_t\).  On the regular bosonic Fock carrier,
\(H_t\) is generally unbounded; its unitary propagator is supplied instead by
the common-domain/self-adjointness hypotheses in case (c).
In an arbitrary orthonormal frame
with \(\nabla_t=\partial_t+\Gamma_t\), the frame generator is
\(H_t-i\hbar\Gamma_t^{\mathcal F}\); metric compatibility makes
\(\Gamma_t^{\mathcal F}\) anti-Hermitian.  The resulting generator is again
self-adjoint on the finite carriers and on the common regular-CCR domain
assumed in (c).  Thus the result is independent of the chosen frame.


\item \textbf{Bogoliubov/Nambu evolution.}
Let
\[
\Xi_t:=\binom{\mathfrak a_t}{\mathfrak a_t^\dagger}
\]
denote the Nambu column of annihilation and creation operators in the Heisenberg picture.
In finite CAR sectors, and on the invariant regular-CCR domain of case (c),
the quadratic form of \(H_t\) implies that every commutator
\([H_t,\mathfrak a(f)]\) and
\([H_t,\mathfrak a^\dagger(f)]\) is again linear in creation and
annihilation operators.  In finite bosonic occupation truncations the
commutator acquires additional boundary-defect terms of the type in
Proposition~\ref{p4:prop:explicit-finite-ccr-boundary-defect}; these are
controlled only in the state-expectation or common-domain matrix-element
topology specified above, not in operator norm.
Thus, in the exact CAR/regular-CCR setting, \(\Xi_t\) satisfies a linear finite-dimensional ODE
\[
\frac{d}{dt}\Xi_t = M_t\,\Xi_t
\]
for a block matrix \(M_t\) built from \(h_t\) and \(\Delta_t\).
Its solution is a block propagator
\[
\Xi_t = \mathcal U_{t,s}\,\Xi_s,
\qquad
\mathcal U_{t,s}=
\begin{pmatrix}
\alpha_{t,s} & \beta_{t,s}\\
\overline{\beta}_{t,s} & \overline{\alpha}_{t,s}
\end{pmatrix},
\]
which is exactly the stated Bogoliubov form.  In the regular-CCR branch, the
implementability hypothesis ensures that this algebraic symplectic
transformation is realized by the unitary propagator:
\[
\mathfrak a_t(f)=\mathfrak a_s(\alpha_{t,s}f)+\mathfrak a_s^\dagger(\beta_{t,s}\overline f).
\]

\item \textbf{Boundary defects and particle-production readout.}
Because the Heisenberg evolution is implemented by the unitary \(U(t,s)\)
on the relevant theorem domain, it preserves the adjoint operation and the
exact CAR or regular CCR relations; in the occupation-truncated bosonic case
the same statement is read modulo the explicit boundary-defect contribution.
Thus \((\alpha_{t,s},\beta_{t,s})\) satisfies the corresponding Bogoliubov
constraints in the exact sectors and up to the stated
state-expectation/common-domain boundary error in the truncated bosonic
sectors.
The coefficient \(\beta_{t,s}\) vanishes precisely when no positive/negative frequency mixing occurs; nonzero \(\beta_{t,s}\) is the finite manifestation of particle production and squeezing.


\item \textbf{Directed limits require algebra and dynamics maps.}
Paper~II's retained Peter--Weyl symbol spaces and continuum-limit theorem
give a nested/dense approximation statement for symbols
\cite{Kugelmann2026b}; they do not assert literal inclusions of the
corresponding matrix algebras, nor do they construct Fock embeddings,
local-observable maps, compatible states, or time-dependent propagators.
If those maps are supplied and satisfy the hypotheses of
Theorem~\ref{p4:thm:inductive-qft-bridge}, applying the present construction
level by level gives the asserted directed system and interface limit.
Thus the symbol-density theorem is upstream input, whereas the many-body
and dynamical compatibility needed for this limit remains displayed
interface data.
\end{proofsteps}

\end{proof}

\begin{theorem}[Positive production density on a constituted final carrier]
\label{p4:thm:positive-production-density}
Work in an exact finite CAR sector or in the implementable regular-CCR sector
of Theorem~\ref{p4:thm:curved-background-finite-bridge}, and fix its supplied
initial/final positive-frequency splittings.  There is a unique positive
operator \(\Gamma_{t,s}\) on the final one-particle space \(\mathcal K_t\)
whose quadratic form is
\begin{equation}
 \label{p4:eq:production-density-quadratic-form}
 \langle f,\Gamma_{t,s}f\rangle_{\mathcal K_t}
 :=\|\beta_{t,s}\overline f\|_{\mathcal K_s}^{2},
 \qquad f\in\mathcal K_t.
\end{equation}
It is the type-correct pullback of \(\beta_{t,s}^*\beta_{t,s}\) from the
conjugate final space; in particular,
\begin{equation}
 \label{p4:eq:produced-number-and-energy}
 N_{t,s}:=\Tr_{\mathcal K_t}\Gamma_{t,s}
 =\|\beta_{t,s}\|_{\rm HS}^{2}\ge0.
\end{equation}
If the final one-particle Hamiltonian is nonnegative, the produced quadratic
energy is
\begin{equation}
 \label{p4:eq:produced-positive-energy}
 E_{t,s}:=\Tr_{\mathcal K_t}
 \bigl(h_t^{1/2}\Gamma_{t,s}h_t^{1/2}\bigr)
 =\Tr_{\mathcal K_t}(h_t\Gamma_{t,s})\ge0.
\end{equation}
For an orthogonal final-sector projector \(P\), its generated occupation is
\(N_{t,s}(P):=\Tr(P\Gamma_{t,s}P)\ge0\).  These statements use the exact
CAR/regular-CCR Bogoliubov transformation.  In a finite
occupation-truncated CCR sector they hold only with the boundary-defect
qualification and topology stated in the curved-background theorem.
\end{theorem}

\begin{proof}
The map \(f\mapsto\|\beta_{t,s}\overline f\|^2\) is a nonnegative Hermitian
quadratic form on the finite-dimensional space \(\mathcal K_t\).  Polarization
therefore represents it by a unique positive operator \(\Gamma_{t,s}\).
Taking its trace in an orthonormal basis gives the Hilbert--Schmidt norm of
\(\beta_{t,s}\), independently of basis.  Positivity of
\(h_t^{1/2}\Gamma_{t,s}h_t^{1/2}\) proves the energy statement; cyclicity of
the finite trace gives the second expression.  The compressed operator
\(P\Gamma_{t,s}P\) is positive, proving the sector statement.
\end{proof}

\begin{corollary}[Absent-target production theorem]
\label{p4:cor:absent-target-production}
Let \(P_C\) be the constituted final one-particle projector for a proposed
native carrier \(C\) in a supplied model.  If that construction yields
\(P_C=0\), then every quadratic production process internal to the supplied
one-particle space satisfies
\begin{equation}
 \label{p4:eq:absent-target-zero-production}
 \boxed{P_C\Gamma_{t,s}P_C=0,\qquad N_{t,s}(P_C)=0.}
\end{equation}
Thus nonadiabatic evolution can populate only constituted final carriers; it
cannot manufacture a missing one-particle target.  This is an internal
carrier statement.  It does not forbid extending the theory by an additional
external sector, and it does not decide whether such an extension occurs.
\end{corollary}

\begin{proof}
Both identities follow by substituting \(P_C=0\) in
Theorem~\ref{p4:thm:positive-production-density}.  The scope statement follows
because that theorem quantifies only over the supplied final one-particle
space and its constituted sector projectors.
\end{proof}

\begin{theorem}[Profile-independent sudden-layer squeeze law]
\label{p4:thm:profile-independent-thin-layer-squeeze}
Let a finite regular-CCR sector split into positive, mutually commuting
bosonic oscillator modes \(\lambda\), with one fixed mode basis throughout a
transition layer.  For a family of monotone profiles
\(s_\epsilon:\mathbb R\to[0,1]\) whose variation is supported in an interval
of width \(\epsilon\), assume
\[
 s_\epsilon(-\infty)=0,\qquad s_\epsilon(+\infty)=1,\qquad
 \omega_{\lambda,\epsilon}(t)
 =\omega_{\lambda,-}e^{-2d_\lambda s_\epsilon(t)},
 \qquad \omega_{\lambda,-}>0.
\]
If the mode Hamiltonians retain the canonical form
\(\tfrac12(p_\lambda^2+\omega_{\lambda,\epsilon}^2q_\lambda^2)\), then in
the zero-width limit
\begin{equation}
\label{p4:eq:thin-layer-bogoliubov-law}
 \boxed{\alpha_\lambda\longrightarrow\cosh d_\lambda,\qquad
 \beta_\lambda\longrightarrow-\sinh d_\lambda,\qquad
 |\beta_\lambda|^2\longrightarrow\sinh^2 d_\lambda.}
\end{equation}
In particular, if a same-action Hessian bridge yields
\(d_\lambda=w_\lambda\delta\), then
\begin{equation}
\label{p4:eq:delta-weighted-thin-layer-production}
 \boxed{|\beta_\lambda|^2\longrightarrow
 \sinh^2(w_\lambda\delta).}
\end{equation}
The limit depends only on the endpoint frequency ratio and not on the shape
of \(s_\epsilon\).
\end{theorem}

\begin{proof}
The canonical evolution across the shrinking layer is generated by the
bounded matrix
\[
 \frac{d}{dt}\binom{q_\lambda}{p_\lambda}
 =
 \begin{pmatrix}0&1\\-\omega_{\lambda,\epsilon}^2&0\end{pmatrix}
 \binom{q_\lambda}{p_\lambda}.
\]
Because the frequencies remain between two fixed positive endpoints, the
propagator over an interval of length \(\epsilon\) tends to the identity.
Thus \(q_\lambda,p_\lambda\) are continuous in the limit.  Expressing the
final annihilation variable
\[
 a_{\lambda,+}
 =\sqrt{\frac{\omega_{\lambda,+}}2}\,q_\lambda
  +\frac{i}{\sqrt{2\omega_{\lambda,+}}}\,p_\lambda
\]
in the initial basis gives
\[
 \alpha_\lambda
 =\frac12\left(
 \sqrt{\frac{\omega_{\lambda,+}}{\omega_{\lambda,-}}}
 +\sqrt{\frac{\omega_{\lambda,-}}{\omega_{\lambda,+}}}\right),
 \quad
 \beta_\lambda
 =\frac12\left(
 \sqrt{\frac{\omega_{\lambda,+}}{\omega_{\lambda,-}}}
 -\sqrt{\frac{\omega_{\lambda,-}}{\omega_{\lambda,+}}}\right).
\]
Since \(\omega_{\lambda,+}/\omega_{\lambda,-}=e^{-2d_\lambda}\),
these are \(\cosh d_\lambda\) and \(-\sinh d_\lambda\).
\end{proof}

\begin{remark}[What the thin-layer theorem does not close]
\label{p4:rem:thin-layer-squeeze-scope}
The theorem removes the profile-shape selector only on the zero-width branch.
A finite physical width, a rotating or noncommuting eigenbasis, the weights
\(w_\lambda\), the identification of \(\delta\) with a mode-frequency jump,
the initial state, backreaction, and any KMS or kinetic thermalization remain
separate data.  The output is a squeezed state, not automatically a thermal
state.
\end{remark}



\begin{theorem}[Zero-excess-energy ground-support theorem]
\label{p4:thm:v2105-zero-excess-ground-support}
Let \(\rho\) be a positive trace-one state on a finite-dimensional quantum
carrier and let \(H_{\rm exc}\succeq0\).  If
\begin{equation}
\label{p4:eq:v2105-zero-excess-ground-support}
 \boxed{\operatorname{Tr}(\rho H_{\rm exc})=0,}
\end{equation}
then
\begin{equation}
\label{p4:eq:v2105-ground-kernel-support}
 \boxed{\operatorname{supp}\rho\subseteq\ker H_{\rm exc}.}
\end{equation}
If \(\ker H_{\rm exc}=\mathbb C|0\rangle\) is one-dimensional, then
\begin{equation}
\label{p4:eq:v2105-unique-ground-state}
 \boxed{\rho=|0\rangle\langle0|.}
\end{equation}
The same support conclusion holds on a declared trace-class positive-energy
domain whenever the displayed expectation is finite.
\end{theorem}

\begin{proof}
Let \(H_{\rm exc}^{1/2}\) be the positive square root.  Positivity and
cyclicity give
\[
 0=\operatorname{Tr}(\rho H_{\rm exc})
  =\operatorname{Tr}(H_{\rm exc}^{1/2}\rho H_{\rm exc}^{1/2}).
\]
The operator under the final trace is positive, hence it vanishes.  Therefore
\(H_{\rm exc}^{1/2}\rho^{1/2}=0\), so the support of \(\rho\) lies in the
kernel of \(H_{\rm exc}\).  A one-dimensional kernel leaves only the displayed
rank-one normalized state.
\end{proof}

\begin{corollary}[Positive quadratic ground complex structure and covariance]
\label{p4:cor:v2105-quadratic-ground-complex-structure}
For one real canonical mode \(z=(q,p)\) with
\[
 H^{(2)}=\frac12 z^{\mathsf T}Gz,\qquad G=G^{\mathsf T}>0,
 \qquad
 \mathbb J=\begin{pmatrix}0&1\\-1&0\end{pmatrix},
\]
put \(\omega=\sqrt{\det G}\).  The Hamiltonian flow matrix satisfies
\((\mathbb JG)^2=-\omega^2I\), so the positive-energy complex structure and
pure ground covariance are fixed by
\begin{equation}
\label{p4:eq:v2105-ground-complex-structure}
 \boxed{J_G=\frac{\mathbb JG}{\omega}},
\end{equation}
\begin{equation}
\label{p4:eq:v2105-ground-covariance}
 \boxed{\Gamma_G=\frac{\hbar\omega}{2}G^{-1},
 \qquad \det\Gamma_G=\frac{\hbar^2}{4}.}
\end{equation}
For a finite positive multimode block, Williamson reduction gives the direct
sum of these canonical one-mode structures, transported back by the
symplectic diagonalizer.  Thus once a physical theorem supplies zero
excitation and a nondegenerate ground sector, no additional Gaussian vacuum
covariance or compatible positive-energy complex-structure datum remains.
This corollary determines the ground data; it does not actualize that ground
state physically.
\end{corollary}

\begin{proof}
For a positive real \(2\times2\) symmetric matrix, direct multiplication gives
\((\mathbb JG)^2=-\det(G)I\).  The displayed \(J_G\) therefore squares to
\(-I\) and is the Hamiltonian positive-frequency structure in the stated
orientation convention.  Symplectic diagonalization reduces the Hamiltonian
to \(\frac12(P^2+\omega^2Q^2)\), whose ground covariance is
\(\operatorname{diag}(\hbar/(2\omega),\hbar\omega/2)\); transport back gives
\(\Gamma_G=(\hbar\omega/2)G^{-1}\).  The determinant follows immediately.
The multimode statement is Williamson's theorem applied blockwise.
\end{proof}

\begin{theorem}[Finite field-theoretic support for supplied cosmological perturbation modes]
\label{p4:thm:finite-cosmological-perturbation-bridge}
Let Paper~VI-A supply a homogeneous background and a finite resolved
quadratic perturbation sector with scalar, tensor, or scalaron perturbation
variables collected in a finite mode space \(\mathcal K_t\).  Suppose the
gauge and constraint zero modes have been removed and, in real canonical
coordinates \(z=(q,p)\), its quadratic form
\[
H_t^{(2)}=\frac12 z^{\mathsf T}G_tz
\]
obeys one of two sufficient stability domains:
\begin{enumerate}[label=(\alph*), leftmargin=2em]
\item it is uniformly positive on the physical phase space,
\(G_t\ge\gamma_t I\) with
\(\inf_{t\in I}\gamma_t>0\); or
\item Paper~VI-A supplies a compatible positive-energy complex structure
and positive-frequency splitting whose positive one-particle energies stay
uniformly separated from zero.
\end{enumerate}
These alternatives are not asserted to be equivalent without the displayed
positive-energy compatibility.  Suppose the sector is represented either on
a finite occupation-truncated Fock sector, or on the regular-CCR branch
under the common-domain, self-adjointness, propagator, and implementability
hypotheses of
Theorem~\ref{p4:thm:curved-background-finite-bridge}.  Then:
\begin{enumerate}[label=(\roman*), leftmargin=2em]
\item each resolved perturbation mode has a stable oscillator
normal form on the stated physical phase space;
\item its evolution is unitary at finite resolution;
\item vacuum mismatch, squeezing, and particle production are described by the finite Bogoliubov coefficients of Theorem~\ref{p4:thm:curved-background-finite-bridge};
\item the usual scalar/tensor two-point functions are recovered as
large-\(m_x\) reciprocal-position/window limits only when Paper~VI-A
supplies compatible embeddings, states, dynamics, and convergence of the
finite correlators in a stated topology, in addition to its slow-background
and large-window estimates; on an occupation-truncated bosonic branch an
independent \(N_{\max}\)-tail estimate or occupation-cutoff-removal limit is
also required.
\end{enumerate}
Thus the finite quantum-field status of the Paper~VI-A perturbation and scalaron mode language is closed for the finite operator block, with the cosmological background and quadratic sector supplied by the Paper~VI-A interface.
\end{theorem}


\begin{proof}
\leavevmode\par\noindent
On the physical phase space, uniform positivity of the full quadratic form,
or the separately supplied compatible positive-energy complex structure
and splitting, gives the standard finite symplectic oscillator normal form
on its stated domain.  Positivity of the kinetic block
alone would not suffice: \(H=p^2-\omega^2q^2\) has positive kinetic form but
is an inverted, unstable oscillator.  On the occupation-truncated carrier
the corresponding Hamiltonian is a norm-continuous finite self-adjoint
matrix, so the finite Schr\"odinger equation has a unique unitary propagator.
On the regular-CCR carrier, the same conclusion instead uses the explicit
unbounded-operator hypotheses in
Theorem~\ref{p4:thm:curved-background-finite-bridge}.  Since the Hamiltonian
is quadratic, the Heisenberg evolution has the exact Bogoliubov form on the
regular-CCR domain and the boundary-corrected form on the finite
occupation-truncated domain.  Two-point functions are expectations of these
transported operators on the respective domain.

The finite calculation supplies resolved two-point functions.  Identifying
their limit with the standard scalar and tensor correlators additionally
uses the supplied compatible embeddings, states, dynamics, and convergence
topology, together with vanishing occupation-boundary error in the
truncated-boson branch.  Without that convergence theorem Paper~IV asserts
only the finite correlators.  It does not perform the subsequent CMB or
structure-growth propagation.  If one assumes only self-adjointness, the exact conclusion is
unitary finite quadratic evolution, not stability, a vacuum, or oscillator
normal form.
\end{proof}


\begin{theorem}[Finite reheating operator block and supplied particle-production interface]
\label{p4:thm:finite-reheating-interface}
On the Paper-VI-A cosmological interface domain, let the finite
scalaron/reheating and matter blocks be given on a fixed trivialized product
carrier \(\mathcal H_{\rm bg}\otimes\mathcal H_{\rm mat}\), or on a
finite-rank product bundle equipped with the metric connection of
Theorem~\ref{p4:thm:curved-background-finite-bridge}.  Suppose the supplied
coupled self-adjoint generator is
\[
H(t)=H_{\rm bg}(t)\otimes I+I\otimes H_{\rm mat}+V(t),
\qquad
H(t)=H(t)^*,
\]
and that the uncoupled part and \(V(t)\), in the fixed trivialization or
covariant frame, are norm-continuous (more generally, locally \(L^1\) in
operator norm).
Then finite unitary evolution gives transition probabilities
\[
P_{i\to f}(t)=\left|\langle f,U(t,t_0)i\rangle\right|^2.
\]
Writing
\[
H_0(t):=H_{\rm bg}(t)\otimes I+I\otimes H_{\rm mat}
\]
and denoting its propagator by \(U_0(t,t_0)\), the exact Duhamel estimate for
\(t\ge t_0\) is
\begin{equation}
\label{p4:eq:reheating-duhamel-bound}
\bigl\|U(t,t_0)-U_0(t,t_0)\bigr\|
\le
\frac1{\hbar}\int_{t_0}^{t}\|V(s)\|\,ds .
\end{equation}
Thus a weak-interaction approximation is controlled by the displayed
integral and is not an additional qualitative premise for the exact finite
claims.
For a fixed supplied matter state \(\sigma_{\rm mat}\), independent of the
variable background input \(\rho\), unresolved matter channels give the
exact CPTP map
\[
\Phi_{t,t_0}(\rho)
=
\Tr_{\rm mat}\!\left[
U(t,t_0)(\rho\otimes\sigma_{\rm mat})U(t,t_0)^*
\right].
\]
Selecting a matter effect \(0\le F\le I\) instead gives the CP
trace-nonincreasing instrument branch
\[
\Phi^F_{t,t_0}(\rho)
=
\Tr_{\rm mat}\!\left[
(I\otimes F^{1/2})U(t,t_0)(\rho\otimes\sigma_{\rm mat})U(t,t_0)^*
(I\otimes F^{1/2})
\right].
\]
Initial system--matter correlations do not, in general, define a
state-independent CP map on arbitrary background inputs; that conclusion
requires a CP assignment map or a restricted compatibility domain.  If the
reduced system also lies in one of the
mixing-reservoir, repeated-probe, or simultaneous large-bath/weak-coupling
domains of Theorem~\ref{p4:thm:microscopic-bath-gkls}, the limit has a finite
GKLS generator \(\mathcal L\).  Let
\(\{P_a\}_{a=1}^r\) be supplied orthogonal occupation projectors with
\(\sum_aP_a=I\), and put
\[
\mathcal N:=\operatorname{span}\{P_1,\ldots,P_r\}.
\]
The populations \(p_a(t):=\Tr(P_a\rho(t))\) obey one closed,
state-independent linear rate equation for every initial state if and only if
\begin{equation}
\label{p4:eq:reheating-population-invariance}
\mathcal L^*(\mathcal N)\subseteq\mathcal N,
\qquad\text{equivalently}\qquad
\mathcal L^*(P_a)=\sum_b W_{ab}P_b .
\end{equation}
On this invariant commutative branch,
\[
\dot p_a=\sum_bW_{ab}p_b,
\]
and \(W\) is a classical Markov generator: \(W_{ab}\ge0\) for
\(a\ne b\) and \(\sum_aW_{ab}=0\).  If
\eqref{p4:eq:reheating-population-invariance} fails, the GKLS equation remains
valid but the selected populations do not close; coherences or block-density
variables must be retained.
Hence reheating and particle production have an exact finite operator
interface, while irreversible thermalization is a separately displayed
open-system limit and scalar occupation kinetics is its computed invariant
commutative restriction.
\end{theorem}


\begin{proof}
Norm-continuity, or local operator-norm integrability, and finite dimension
give a unique absolutely continuous unitary propagator \(U(t,t_0)\) in the
fixed trivialization; the bundle version uses the supplied metric
connection.  The displayed transition probabilities are therefore ordinary
finite Born probabilities.  Expanding \(U(t,t_0)\) in the interaction picture
gives the finite Dyson series, which is norm convergent on bounded time
intervals.  The variation-of-constants identity
\[
U(t,t_0)-U_0(t,t_0)
=-\frac{i}{\hbar}\int_{t_0}^{t}
U(t,s)V(s)U_0(s,t_0)\,ds
\]
and unitarity prove \eqref{p4:eq:reheating-duhamel-bound}.  The fixed product assignment
\(\rho\mapsto\rho\otimes\sigma_{\rm mat}\), unitary conjugation, and partial
trace are CPTP; their composition gives \(\Phi_{t,t_0}\), as in
Theorem~\ref{p4:thm:exact-finite-system-environment-reduction}.  Inserting
the effect \(F\) gives the displayed CP trace-nonincreasing branch.  Without
the fixed product assignment (or an explicitly CP correlated assignment),
partial tracing a single correlated trajectory does not define a CP channel
on all possible inputs.  Only under the
reservoir/repeated-probe/large-bath hypotheses cited in the theorem does the
reduced family converge to a GKLS semigroup.  For that semigroup,
\[
\dot p_a
=\Tr\!\left(\mathcal L^*(P_a)\rho\right).
\]
If \(\mathcal L^*(P_a)=\sum_bW_{ab}P_b\), this is the displayed closed rate
equation.  Conversely, if the same closed equation holds for every density
matrix, finite trace duality forces the operator identity in
\eqref{p4:eq:reheating-population-invariance}; this proves necessity.  The
restricted Heisenberg semigroup is positive and unital on the commutative
algebra \(\mathcal N\).  In the projector basis its matrices are therefore
stochastic, so differentiation at zero gives the stated sign and
column-sum conditions on \(W\).

A finite unresolved block has finite unitary kinematics, but recurrence
is not automatic under an arbitrary time-dependent drive. An autonomous
finite-dimensional unitary group admits approximate returns by simultaneous
phase approximation; no such uniform recurrence is inferred for a general
nonautonomous propagator, and neither case establishes the kinetic limit.  Even after a GKLS limit is obtained, failure of
\eqref{p4:eq:reheating-population-invariance} forces a block-density rather
than scalar-population description.  The theorem supplies the finite
interface; the actual branching fractions require couplings and channel
projectors supplied by the sector under study, not by Paper~IV as primitive
inputs.
\end{proof}


\begin{theorem}[Finite transfer-kernel status of cosmological interface coefficients]
\label{p4:thm:finite-transfer-kernel-status}
Let
\[
\Phi_a:\End(\mathcal H_{\rm in})\longrightarrow\End(\mathcal H_a)
\]
be a finite CPTP Schr\"odinger-picture transfer map and let
\(\rho_{\rm in}\) be an input state.  A selected-output branch may instead
be a CP trace-nonincreasing map \(\Phi_a^F\); its conditional output is
\(\Phi_a^F(\rho_{\rm in})/\Tr\Phi_a^F(\rho_{\rm in})\) whenever the
denominator is nonzero.

For a supplied response functional
\(\ell_a:\End(\mathcal H_a)\to\mathbb C\), any Paper-VI-B transfer coefficient
of the schematic form
\[
\Xi_a
=
\frac{\ell_a\!\left(\Phi_a(\rho_{\rm in})\right)}
{\ell_a^{(0)}\!\left(\Phi_a^{(0)}(\rho_{\rm in})\right)},
\qquad
a\in\{d,*,g,\mathrm{BBN}\},
\]
when its denominator is nonzero, is a finite transfer-response ratio.
Finite dimensionality gives a unique operator \(O_a\) with
\(\ell_a(X)=\Tr(O_aX)\).  This pairing is real for self-adjoint \(O_a\),
nonnegative for \(O_a\ge0\), and a probability only for an effect
\(0\le O_a\le I\) (with the corresponding normalization).  A generic
response functional is not positive and therefore is not a Born
probability.

Paper~VI-B determinant objects such as \(\mathcal T_8(a)\),
\(\mathcal T_g\), and a derived growth coefficient \(\kappa_g\) are
separate finite determinant or transfer outputs once their resolved
kernels and extraction rules are supplied.  A determinant kernel is not
automatically an observable or a positive response functional.  Thus these
objects are computable finite interface outputs, not new microscopic
geometric constants, but their map, response, kernel, and input data are not
fixed by the finite carrier alone.
\end{theorem}


\begin{proof}
Complete positivity and trace preservation send every input state to a
state.  A CP trace-nonincreasing branch sends it to a subnormalized positive
operator and therefore defines the displayed conditional state precisely
when its occurrence probability is nonzero.  Finite-dimensional trace
duality represents every linear response functional by a unique operator;
the order properties of that operator give exactly the reality,
nonnegativity, and probability qualifications in the statement.  No such
order property follows for a generic response functional.

The ratio defining \(\Xi_a\) is therefore a well-typed transfer-response
quantity once both response functionals and the nonzero denominator are
supplied.  Determinants are finite algebraic functions of their supplied
kernels but do not become Born pairings without a separate positive-operator
representation.  Computing the numerical values requires
\(\Phi_a\), \(\ell_a\), the reference data, any determinant kernel and
extraction rule, and \(\rho_{\rm in}\); that is the cosmological propagation
problem left to Paper~VI-B or to a numerical validation pipeline.
\end{proof}




\subsection{Wavefunctions, kernels, and path sums as presentation limits}
\label{p4:sec:part2-wave-path}



\paragraph{Conditional kernel and commutative presentation limits}
\label{p4:sec:part2-kernel-limit}

Let \(\mathcal K_m\) be finite reciprocal-position presentation spaces with Hamiltonians
\(H_m\).  A Schr\"odinger presentation is obtained only on a supplied limit
domain consisting of a target Hilbert space
\(\mathcal K=L^2(X,\mu)\), compatible embeddings of the finite presentation spaces, and
self-adjoint extensions \(\widetilde H_m\) on \(\mathcal K\) that converge in
strong resolvent sense to a specified self-adjoint \(H\).  The
Trotter--Kato theorem then gives
\[
e^{-it\widetilde H_m/\hbar}\psi
\longrightarrow
e^{-itH/\hbar}\psi
\]
for every \(\psi\in\mathcal K\), uniformly for \(t\) in compact intervals.
Only when \(H\) has independently been identified as a Schr\"odinger operator
and its semigroup or propagator admits a kernel does ordinary wave-mechanical
kernel language follow.  Papers~II--III do not supply these embeddings,
operator convergence, or target identification merely from \(m_x\to\infty\).



\paragraph{Discrete path sums}
\label{p4:sec:part2-path-sums}

For a Hamiltonian split $H=T+V$, one has in finite dimension the Trotter product formula
\begin{equation}
\label{p4:eq:trotter}
e^{-itH/\hbar}
=\lim_{n\to\infty}
\bigl(e^{-itT/(n\hbar)}e^{-itV/(n\hbar)}\bigr)^n.
\end{equation}
Inserting resolutions of the identity in a chosen presentation basis yields a finite discrete path sum at fixed truncation.
For each time-slice number every label set is finite, and the time-slice limit
is exactly the finite matrix propagator.  This is the proved algebraic
``sum over histories'' statement.



\paragraph{Path-integral presentation limit}
\label{p4:sec:part2-path-integral-limit}

A continuum path-integral statement requires more than the finite Trotter
identity.  In addition to the operator-convergence domain above, one must
specify a topology on the path space and prove that the time-sliced kernels
converge either to a countably additive Euclidean measure or to a controlled
real-time oscillatory integral/distribution.  The required hypotheses depend
on the target configuration space, kinetic operator, potential, and boundary
conditions.  Real-time Feynman expressions are not in general probability
measures.  No such measure, tightness, oscillatory-integral, or interchange-of-
limits theorem follows from the finite matrix sequence alone.  Paper~IV
therefore proves the finite path sum \eqref{p4:eq:trotter}; continuum path
integrals and wavefunctions remain conditional presentation limits.


\begin{theorem}[No-independent-wave-postulate closure]
\label{p4:thm:v244-no-independent-wave-postulate}
Consider a constituted branch satisfying the following already separated
conditions: the compact reciprocal geometry supplies its harmonic carrier;
Paper~II-C supplies the licensed reciprocal harmonic/natural-bundle
presentation; Paper~III-F owns the local action, a stationary solution, the
required physical reduction, and a stable positive quadratic Hessian; and the
finite algebraic state/effect calculus of this paper applies.  Then no
additional wave-specific primitive is required on that branch:
\begin{enumerate}[label=(\roman*),leftmargin=2.2em]
\item harmonic mode decompositions are representation-theoretic consequences
of the compact carrier and its natural bundles;
\item the Lorentzian second-order propagation operator is the kinetically
normalized physical Hessian generated in
Theorem~\ref{p3f:thm:v244-stationary-physical-hessian-propagation};
\item the positive-frequency complex structure is generated by
Theorem~\ref{p4:thm:stationary-positive-frequency-mode-space}, and the
first-order Schr\"odinger equation/action are the complex-coordinate form of
the same real quadratic dynamics by
Theorem~\ref{p4:thm:v244-action-generated-schrodinger-form};
\item a finite wavefunction is a basis/MASA-dependent coordinate presentation
of a state vector representing the algebraic state, not a separate structural
state axiom;
\item within a simple coherent block, linear combinations are already vectors
of the constituted complex carrier and therefore define rays after
normalization; and
\item interference terms are the off-diagonal Born/effect contributions of
Section~\ref{p4:sec:part2-interference}, and disappear under the corresponding
dephasing map.
\end{enumerate}
Accordingly, on this theorem domain there is no independent wave ontology,
fundamental wavefunction postulate, second-order wave-equation postulate,
positive-frequency postulate, within-sector superposition postulate, or
interference postulate.

This closure does not select an actual state, \(\mathsf{DYN}\) member,
initial/history member, one-particle semantics, continuum configuration-space
wavefunction, nonstationary initial-vacuum/state selection, interacting nonlinear
completion, or a physical position--momentum calibration.  Those roles retain
their existing typed owners and frontiers.
\end{theorem}

\begin{remark}[Finite closure versus continuum wavefunction]
The theorem above concerns the finite harmonic and algebraic presentations.
The continuum statement \(\psi\in L^2(X,\mu)\), a continuum propagator kernel,
or a path-integral measure still requires the convergence hypotheses stated in
this section.  Thus removing an independent finite wave postulate does not
manufacture the continuum limit.
\end{remark}



\section{Three-level, threshold, perturbative and textbook normal forms}
\label{p4c:sec:v238-c5}
\noindent\textit{Source and claim boundary.} These are original extended calculations. Statements on physical sources, measurements or continuum limits retain the inherited antecedents of Papers IV-A/IV-B.

\subsection{Perturbation theory and adiabatic transport}
\label{p4:sec:part2-perturbation-adiabatic}

\paragraph{Role in textbook quantum mechanics.}
The finite Born framework already gives exact unitary evolution, spectral projectors, and the three standard dynamical pictures.
The next textbook layer is approximation theory: one studies a solvable reference Hamiltonian, weak deformations, slowly varying generators, and the induced transport of eigenspaces.
Because the carrier is finite-dimensional, all of these constructions are matrix-theoretic inside $\mathcal A_m^{\mathrm{env}}$; no domain subtleties or additional continuum axioms are needed.



\paragraph[Time-independent perturbation theory around isolated levels]
{Time-independent perturbation theory}
\label{p4:sec:part2-timeindep-perturbation}

Let
\begin{equation}
\label{p4:eq:H-lambda}
H(\lambda)=H_0+\lambda V,
\qquad
H_0=H_0^*,\quad V=V^*,
\end{equation}
with real perturbation parameter $\lambda$.
Suppose $E_a$ is an eigenvalue of $H_0$ with spectral projector $P_a$, and let
\begin{equation}
\label{p4:eq:Qa-def}
Q_a:=\mathbf 1-P_a.
\end{equation}
If $E_a$ is isolated from the rest of the spectrum by a nonzero gap, then for sufficiently small $|\lambda|$ the corresponding spectral block of $H(\lambda)$ deforms smoothly.
This is the finite matrix version of Rayleigh--Schr\"odinger perturbation theory.


\begin{theorem}[Finite nondegenerate perturbation formulas]
\label{p4:thm:nondeg-perturbation}
Let $E_a$ be a simple eigenvalue of $H_0$ with normalized eigenvector $\phi_a$, and let $\{\phi_b\}$ be an orthonormal eigenbasis of $H_0$.
Then for sufficiently small $|\lambda|$ there exist a unique nearby eigenvalue branch $E_a(\lambda)$ and normalized eigenvector branch $\psi_a(\lambda)$ such that
\begin{equation}
\label{p4:eq:Ea-expansion}
E_a(\lambda)=E_a+\lambda\,\langle \phi_a,V\phi_a\rangle+O(\lambda^2),
\end{equation}
and, after fixing the phase by requiring $\langle\phi_a,\psi_a(\lambda)\rangle\in\mathbb R_{>0}$,
\begin{equation}
\label{p4:eq:psi-first-order}
\psi_a(\lambda)
=
\phi_a
+\lambda\sum_{b\neq a}
\frac{\langle \phi_b,V\phi_a\rangle}{E_a-E_b}\,\phi_b
+O(\lambda^2).
\end{equation}
Equivalently, the first-order energy shift is the expectation value of $V$ in the unperturbed state, while the first-order state correction mixes in only the other unperturbed eigenstates with denominators given by spectral gaps.
\end{theorem}


\begin{proof}
Differentiate the eigenvalue equation
\[
H(\lambda)\psi_a(\lambda)=E_a(\lambda)\psi_a(\lambda)
\]
at $\lambda=0$ and project once onto $\phi_a$ and once onto the orthogonal complement spanned by the remaining eigenvectors.
Because $E_a$ is simple and isolated, the reduced resolvent on $Q_a\mathcal H_m$ is well-defined, yielding the stated formulas.
\end{proof}


\begin{theorem}[Finite degenerate perturbation theory]
\label{p4:thm:degenerate-perturbation}
Let \(P_a\) project onto a degenerate eigenspace of \(H_0\) with eigenvalue
\(E_a\).  To first order in \(\lambda\), the splitting of this level is
determined by diagonalizing the compressed perturbation
\begin{equation}
\label{p4:eq:deg-perturb-block}
V_a:=P_a V P_a
\end{equation}
on the finite-dimensional eigenspace \(P_a\mathcal H_m\).  More precisely, if
\(\mu\) is an eigenvalue of \(V_a\), then \(H(\lambda)\) has an eigenvalue
branch of the form
\begin{equation}
\label{p4:eq:deg-splitting}
E_{a,\mu}(\lambda)=E_a+\lambda\mu+O(\lambda^2).
\end{equation}
Thus the first-order degeneracy splitting is intrinsic to the finite corner
algebra
\[
  P_a\mathcal A_m^{\mathrm{env}}P_a .
\]
\end{theorem}


\begin{proof}
Restrict the eigenvalue problem for $H(\lambda)$ to the degenerate subspace $P_a\mathcal H_m$.
At order $\lambda$ one obtains the secular equation for $V_a$.
Since the carrier is finite-dimensional, this reduced diagonalization is exact and complete.
\end{proof}

\begin{corollary}[Second-order effective Hamiltonian on an isolated block]
\label{p4:cor:second-order-effective}
For an isolated spectral block $P_a$ of $H_0$, the effective Hamiltonian on that block takes the form
\begin{equation}
\label{p4:eq:heff-second-order}
H_{\mathrm{eff},a}(\lambda)
=
E_a P_a
+\lambda P_aVP_a
-\lambda^2 P_aVQ_a(H_0-E_a)^{-1}Q_aVP_a
+O(\lambda^3).
\end{equation}
Here \((H_0-E_a)^{-1}\) is the reduced inverse on \(Q_a\mathcal H_m\).  Hence second-order shifts are controlled by virtual transitions through the orthogonal complement $Q_a\mathcal H_m$.
\end{corollary}

\begin{proof}
Apply the finite Schur-complement/Feshbach expansion to the decomposition \(P_a\oplus Q_a\).  On \(Q_a\mathcal H_m\), the operator \(H_0-E_a\) is invertible because the block is isolated.  Expanding the resulting effective block in powers of \(\lambda\) gives the displayed second-order term; the finite-dimensional remainder is \(O(\lambda^3)\).
\end{proof}




\paragraph{Time-dependent perturbation theory in the interaction picture}
\label{p4:sec:part2-timedep-perturbation}

With the split $H=H_0+V(t)$ and interaction-picture quantities defined in \eqref{p4:eq:interaction-picture-def}--\eqref{p4:eq:UI-evolution}, the exact propagator is generated by the bounded operator-valued function
\begin{equation}
\label{p4:eq:VIt-reminder}
V_I(t)=e^{itH_0}V(t)e^{-itH_0}.
\end{equation}
Because the carrier is finite-dimensional, the Dyson expansion is the
convergent iterated-integral solution of a matrix ODE.  Put
\[
I_t:=[\min(0,t),\max(0,t)],
\qquad
L_t(V):=\int_{I_t}\|V_I(\tau)\|\,d\tau .
\]
The displayed ordered-simplex formula is for \(t\ge0\); for \(t<0\) one
uses the reverse-ordered simplex, equivalently the time-ordered exponential
\(U_I(t)=\mathcal T\exp(-i\int_0^tV_I(\tau)\,d\tau)\).


\begin{theorem}[Dyson series on the finite carrier]
\label{p4:thm:dyson-finite}
For \(t\ge0\), the interaction-picture propagator solving
\eqref{p4:eq:UI-evolution} admits the norm-convergent series
\begin{equation}
\label{p4:eq:dyson-series}
U_I(t)
=
\mathbf 1
+\sum_{n\ge 1}(-i)^n
\int_{0\le t_n\le\cdots\le t_1\le t}
V_I(t_1)\cdots V_I(t_n)
\,dt_1\cdots dt_n.
\end{equation}
For \(t<0\), the norm-convergent representation is the reverse-ordered
Picard series over \([t,0]\), equivalently
\(\mathcal T\exp(-i\int_0^tV_I(\tau)\,d\tau)\).
Moreover,
\begin{equation}
\label{p4:eq:dyson-bound}
\|U_I(t)-\mathbf 1\|
\le
\exp\bigl(L_t(V)\bigr)-1.
\end{equation}
Thus finite time-dependent perturbation theory is an exact convergent expansion rather than a merely formal series.
\end{theorem}


\begin{proof}
For \(t\ge0\), iterate the integral form of \eqref{p4:eq:UI-evolution},
\[
U_I(t)=\mathbf 1-i\int_0^t V_I(\tau)U_I(\tau)\,d\tau,
\]
and estimate the resulting simplex integrals in operator norm.  For \(t<0\)
the same iteration runs on the reverse-ordered simplex over \([t,0]\).
In either case the \(n\)-th term has norm at most
\(L_t(V)^n/n!\).  Summing proves absolute convergence and
\eqref{p4:eq:dyson-bound}.
\end{proof}

\begin{corollary}[First-order transition amplitudes]
\label{p4:cor:first-order-transition}
Let $H_0\phi_a=E_a\phi_a$ and $H_0\phi_b=E_b\phi_b$.
If the system starts in \(\phi_a\), then for \(b\ne a\) the transition
amplitude into \(\phi_b\) is
\begin{equation}
\label{p4:eq:first-order-transition}
\mathcal A_{b\leftarrow a}(t)
=
-i\int_0^t e^{i(E_b-E_a)\tau}
\langle \phi_b,V(\tau)\phi_a\rangle\,d\tau
+R_{ba}(t),
\end{equation}
where
\[
|R_{ba}(t)|
\le e^{L_t(V)}-1-L_t(V).
\]
If \(V=\lambda W\), this gives the uniform quadratic estimate
\[
|R_{ba}(t)|
\le
\frac{|\lambda|^2L_t(W)^2}{2}
e^{|\lambda|L_t(W)}.
\]
For a monochromatic drive, near-resonant transitions occur when the driving frequency matches the Bohr frequency $E_b-E_a$.
\end{corollary}

\begin{proof}
Retain the first nontrivial term of the Dyson expansion in
Theorem~\ref{p4:thm:dyson-finite}.  Taking the matrix element between
\(\phi_b\) and \(\phi_a\) in the interaction picture inserts the phase
\(e^{i(E_b-E_a)\tau}\), giving the displayed integral.  The sum of terms of
order \(n\ge2\) is bounded by
\(\sum_{n\ge2}L_t(V)^n/n!=e^{L_t(V)}-1-L_t(V)\).
For \(V=\lambda W\), the elementary bound
\(\sum_{n\ge2}x^n/n!\le\frac12x^2e^x\) gives the stated uniform
\(O(\lambda^2)\) estimate.
\end{proof}


Therefore resonant excitation, matrix-element selection, and weak transition probabilities are represented within the finite algebraic formalism.



\paragraph{Adiabatic transport of finite spectral subspaces}
\label{p4:sec:part2-adiabatic-transport}

Let \(s\mapsto H(s)\), \(s\in[0,1]\), be a \(C^2\) family of
self-adjoint Hamiltonians on the fixed finite carrier, and let \(P(s)\) be
the constant-rank spectral projector of a continuously tracked isolated
band whose gap from the complementary spectrum stays uniformly positive:
\begin{equation}
\label{p4:eq:adiabatic-gap}
g:=\inf_{s\in[0,1]}
\operatorname{dist}\!\left(
\Spec\!\left(H(s)|_{\operatorname{Ran}P(s)}\right),
\Spec\!\left(H(s)|_{\operatorname{Ran}(\mathbf1-P(s))}\right)
\right)>0.
\end{equation}
For a slow total runtime $T$, define the physical evolution by
\begin{equation}
\label{p4:eq:adiabatic-evolution}
i\frac{dU_T}{dt}=H(t/T)U_T(t),
\qquad
U_T(0)=\mathbf 1.
\end{equation}


\begin{theorem}[Finite adiabatic theorem]
\label{p4:thm:finite-adiabatic}
Under the gap condition \eqref{p4:eq:adiabatic-gap}, there exists a constant $C$, depending only on finitely many derivatives of $H(s)$ and on the inverse gap, such that
\begin{equation}
\label{p4:eq:adiabatic-bound}
\sup_{0\le t\le T}
\bigl\|(\mathbf 1-P(t/T))U_T(t)P(0)\bigr\|
\le
\frac{C}{T}.
\end{equation}
Hence every state initially supported in $P(0)\mathcal H_m$ remains in the transported spectral subspace $P(t/T)\mathcal H_m$ up to an $O(T^{-1})$ error.
In rank one, the transported state equals the instantaneous eigenvector multiplied by the dynamical phase and the Berry phase, up to the same order.
\end{theorem}


\begin{proof}
\ProofPlan{We prove the finite adiabatic theorem by controlling the spectral projection, constructing the adiabatic intertwiner, estimating leakage with the reduced resolvent, and reading the rank-one Berry-phase form.}
\begin{proofsteps}

\item \textbf{Local Riesz projections and gap-uniform bounds.}
Track the isolated spectral cluster with constant multiplicity.  For every
\(s_0\in[0,1]\), choose a contour \(\Gamma_{s_0}\) enclosing that cluster
at \(s_0\) and no complementary eigenvalue.  Spectral continuity gives a
neighborhood \(U_{s_0}\) on which the same contour separates the continued
cluster, and there
\[
P(s)=\frac{1}{2\pi i}
\oint_{\Gamma_{s_0}}(z-H(s))^{-1}\,dz .
\]
On overlaps the local formulas agree because both are spectral projections
onto the same continuously tracked cluster.  Compactness of \([0,1]\) gives
a finite subcover.  Differentiating the resolvent locally and taking the
maximum over this finite cover shows that \(P(s)\) is \(C^2\) with uniform
first- and second-derivative bounds depending on
\(g^{-1}\), \(\|\dot H\|_\infty\), and
\(\|\ddot H\|_\infty\).  No single contour separating the spectra for all
parameter values is assumed.

\item \textbf{Adiabatic Hamiltonian and intertwining.}
Define the adiabatic Hamiltonian
\[
H_{\mathrm a}(s):=H(s)+\frac{i}{T}[\dot P(s),P(s)],
\]
and let \(U_{\mathrm a}(t)\) be the propagator generated by \(H_{\mathrm a}(t/T)\).
By direct differentiation of \(U_{\mathrm a}(t)^*P(t/T)U_{\mathrm a}(t)\), using \(P^2=P\), one gets the exact intertwining relation
\[
U_{\mathrm a}(t)P(0)=P(t/T)U_{\mathrm a}(t).
\]


\item \textbf{Reduced-resolvent leakage estimate.}
Set \(\Omega(t):=U_{\mathrm a}(t)^*U_T(t)\).
Then
\[
\dot\Omega(t)=-\frac1T U_{\mathrm a}(t)^*[\dot P(t/T),P(t/T)]U_T(t).
\]
The commutator \([\dot P,P]\) is purely off-diagonal with respect to \(P\oplus(\mathbf1-P)\).
Using the reduced resolvent of the gapped band, one constructs a uniformly bounded off-diagonal operator \(X(s)\) such that
\[
[H(s),X(s)]=i[\dot P(s),P(s)],
\qquad
\|X(s)\|+\|\dot X(s)\|\le C.
\]
Substituting this identity into the Duhamel formula for \(\Omega(t)-\mathbf1\) and integrating by parts gives
\[
\sup_{0\le t\le T}\| (\mathbf1-P(t/T))U_T(t)P(0)\|
=
\sup_{0\le t\le T}\|(\mathbf1-P(t/T))(U_T(t)-U_{\mathrm a}(t))P(0)\|
\le \frac{C}{T},
\]
for a constant depending only on the stated bounds.
This is exactly \eqref{p4:eq:adiabatic-bound}.


\item \textbf{Rank-one Berry-phase form.}
\leavevmode\par\noindent
In rank one, write
\[
  P(s)=|\phi(s)\rangle\langle\phi(s)|,
  \qquad
  \|\phi(s)\|=1 .
\]
The exact intertwiner generated by \([\dot P,P]\) transports \(\phi(0)\) to
\[
  e^{i\gamma_B(s)}\phi(s),
\]
where
\[
\gamma_B(s)
=
i\int_0^s\langle \phi(\sigma),\dot\phi(\sigma)\rangle\,d\sigma .
\]
Multiplying by the dynamical phase from \(H(s)\) yields
\eqref{p4:eq:adiabatic-phase-form}, with the same \(O(T^{-1})\) error coming
from the leakage estimate.
\end{proofsteps}

\end{proof}

For an isolated nondegenerate branch $H(s)\phi(s)=E(s)\phi(s)$ with $\langle \phi(s),\phi(s)\rangle=1$, the adiabatically transported state has the standard form
\begin{equation}
\label{p4:eq:adiabatic-phase-form}
\psi_T(t)
=
\exp\Bigl(-iT\int_0^{t/T}E(\sigma)\,d\sigma\Bigr)
\exp\bigl(i\gamma_B(t/T)\bigr)
\phi(t/T)
+O(T^{-1}),
\end{equation}
where
\begin{equation}
\label{p4:eq:berry-phase-again}
\gamma_B(s)=i\int_0^s \langle \phi(\sigma),\dot\phi(\sigma)\rangle\,d\sigma.
\end{equation}
This is the same phase transport whose geometric meaning was already identified earlier through the Hopf-fiber description of Berry holonomy.
So adiabatic quantum transport and geometric phase are not separate add-ons: they are two presentations of the same finite projective geometry.

\begin{remark}[Status of the finite adiabatic theorem]
\label{p4:rem:adiabatic-status}
Theorem~\ref{p4:thm:finite-adiabatic} is the standard finite-dimensional adiabatic estimate specialized to the present carrier \cite{Kato1950,Teufel2003}.
The proof above records the only ingredients needed later, namely the gapped projector calculus and the $O(T^{-1})$ leakage control.
\end{remark}


\begin{theorem}[Transported-frame reduction of a finite spectral sector]
\label{p4:thm:transported-frame-sector}
Let $P(s)$ be the gapped spectral projector of Theorem~\ref{p4:thm:finite-adiabatic}, and let $W(s)$ be the Kato intertwiner solving
\begin{equation}
\label{p4:eq:kato-intertwiner}
\frac{dW}{ds}=[\dot P(s),P(s)]W(s),
\qquad
W(0)=\mathbf 1.
\end{equation}
Then $W(s)$ is unitary and intertwines the moving sector with the initial one,
\begin{equation}
\label{p4:eq:intertwiner-projector}
W(s)^*P(s)W(s)=P(0).
\end{equation}
In particular, if $Y(s)=P(s)Y(s)P(s)\ge 0$ is a $C^1$ family of positive observables supported on the moving rank-three sector, then
\begin{equation}
\label{p4:eq:transported-observable}
\widetilde Y(s):=W(s)^*Y(s)W(s)
\in
\mathcal A_{f,0}:=P(0)\mathcal A_m^{\mathrm{env}}P(0)\cong \mathrm{Mat}(3,\CC)
\end{equation}
is a $C^1$ family of positive observables on a fixed corner algebra, with the same spectrum as $Y(s)$ for every $s$.
If the exact physical evolution is adiabatic up to the $O(T^{-1})$ leakage bound \eqref{p4:eq:adiabatic-bound}, then the reduction to the fixed transported frame is valid up to the same order.
\end{theorem}


\begin{proof}
\ProofPlan{We reduce a moving finite spectral sector to a transported frame by proving unitarity of the Kato transport, the projection-intertwining identity, positivity and spectral preservation of transported observables, and the inherited adiabatic error control.}
\begin{proofsteps}

\item \textbf{Skew-adjointness of the Kato generator.}
Because \(P(s)=P(s)^*=P(s)^2\), differentiating \(P^2=P\) gives
\[
P\dot P P = 0,
\qquad
(\mathbf1-P)\dot P(\mathbf1-P)=0.
\]
Hence
\[
[\dot P,P]^* = P\dot P-\dot P P = -[\dot P,P],
\]
so the generator in \eqref{p4:eq:kato-intertwiner} is anti-self-adjoint.
Therefore
\[
\frac{d}{ds}\bigl(W(s)^*W(s)\bigr)=0,
\qquad W(0)^*W(0)=\mathbf1,
\]
which proves that \(W(s)\) is unitary.


\item \textbf{Intertwining of spectral projections.}
\leavevmode\par\noindent
Define
\[
  F(s):=W(s)^*P(s)W(s).
\]
Differentiating gives
\[
\dot F
=
W^*\bigl(-[\dot P,P]P+\dot P+P[\dot P,P]\bigr)W.
\]
Using \(P\dot P P=0\), differentiating \(P^2=P\) to get
\[
  \dot P=\dot P P+P\dot P,
\]
and using the identities
\[
  P[\dot P,P]=-P\dot P,
  \qquad
  [\dot P,P]P=\dot P P,
\]
the bracket simplifies to zero.  Thus \(\dot F=0\), and since \(F(0)=P(0)\) we
obtain
\[
W(s)^*P(s)W(s)=P(0).
\]

\item \textbf{Transported observable positivity and spectral preservation.}
If \(Y(s)=P(s)Y(s)P(s)\ge0\), then for every \(\psi\)
\[
\langle \psi,\widetilde Y(s)\psi\rangle
=
\langle W(s)\psi,Y(s)W(s)\psi\rangle\ge0,
\]
so \(\widetilde Y(s)\ge0\).
Because \(\widetilde Y(s)\) is unitarily conjugate to \(Y(s)\), they have the same spectrum; and since both \(Y\) and \(W\) are $C^1$, so is \(\widetilde Y\).

\item \textbf{Reduced-frame sector dynamics.}
Finally, Theorem~\ref{p4:thm:finite-adiabatic} says that the exact dynamics transports the physical state into \(P(s)\mathcal H_m\) up to an \(O(T^{-1})\) component orthogonal to that sector.
Conjugating by \(W(s)^*\) identifies \(P(s)\mathcal H_m\) with the fixed corner \(P(0)\mathcal H_m\), so the same \(O(T^{-1})\) control holds in the transported frame.
\end{proofsteps}

\end{proof}

The approximation toolkit of textbook quantum mechanics therefore survives in full finite form:
stationary perturbation theory resolves weak Hamiltonian deformations, Dyson expansions control driven transitions, and adiabatic transport governs slowly varying spectral sectors.
Only the infinite-dimensional analytic complications disappear.

\begin{remark}[Status of the transported-frame reduction]
\label{p4:rem:transported-frame-status}
Theorem~\ref{p4:thm:transported-frame-sector} is included because later papers use exactly this fixed-corner reformulation of moving finite sectors.
Its content is finite-dimensional and standard once the Kato intertwiner is introduced \cite{Kato1950,Teufel2003}.
\end{remark}



\paragraph{Three-Level Sectors, Mixing, and Threshold Algebra}
\label{p4:part:three-level-mixing-thresholds}

\paragraph{Orientation.}
This part isolates the finite rank-three, transported-block, mixing, block-elimination, and threshold-composition machinery used later by sector papers.  The purpose is to keep these exact finite spectral tools visible without treating downstream flavor, mass, or threshold assignments as inputs to Paper~IV.



\subsection{Spectral geometry of three-level sectors}
\label{p4:sec:part2-threelevel}

\paragraph{Status and motivation.}
This section is not part of the closed structural core of Part~I.
Its role is to show that whenever the finite quantum carrier contains a distinguished rank-three sector, the spectral data of positive observables on that sector has a canonical basis-free organization.
That organization is exactly the kind of quantum support used later in Paper~V-A for family triplets, shape variables, and relative mixing data.



\paragraph{Rank-three sectors and positive observables}
\label{p4:sec:part2-threelevel-sector}

Let $P_f\in\mathcal A_m^{\mathrm{env}}$ be a rank-three projector and define the associated corner algebra and carrier by
\[
\mathcal A_f:=P_f\,\mathcal A_m^{\mathrm{env}}\,P_f\cong \mathrm{Mat}(3,\CC),
\qquad
\mathcal H_f:=P_f\mathcal H_m\cong \CC^3.
\]
A positive observable on this sector is an element
\[
Y_f\in\mathcal A_f,
\qquad
Y_f\ge 0.
\]
In Paper~V-A such objects may be interpreted as family or Yukawa-type observables, but the present statement is purely quantum-mechanical: every strictly positive $Y_f$ on a rank-three sector is diagonalizable and admits canonical spectral invariants.


\begin{proposition}[Weyl basis for a distinguished rank-three sector]
\label{p4:prop:rank-three-weyl-basis}
Let \(\mathcal H_f\cong\CC^3\) be a distinguished rank-three sector with ordered orthonormal basis \(e_0,e_1,e_2\).
Define
\[
Z e_j=\omega^j e_j,
\qquad
X e_j=e_{j+1},
\qquad
\omega=e^{2\pi i/3},
\]
with indices in \(\ZZ_3\).
Then
\[
X^3=Z^3=\mathbf 1,
\qquad
ZX=\omega XZ,
\]
and the nine operators
\[
W_{ab}:=X^aZ^b,
\qquad a,b\in\ZZ_3,
\]
are trace-orthogonal:
\[
\operatorname{tr}(W_{ab}^*W_{cd})=3\delta_{ac}\delta_{bd}.
\]
Hence
\[
\End(\mathcal H_f)=\operatorname{span}_{\CC}\{X^aZ^b:a,b\in\ZZ_3\}.
\]
\end{proposition}


\begin{proof}
The relation \(ZX=\omega XZ\) follows by applying both sides to \(e_j\).
The trace-orthogonality follows from the usual finite Fourier sum \(\sum_{j=0}^{2}\omega^{j(b-d)}=3\delta_{bd}\) together with the fact that a nonzero cyclic shift has zero trace.
There are nine trace-orthogonal nonzero operators, equal to \(\dim\End(\CC^3)\), so they form a basis.
On the \(\mathsf{Supp}_{2T}\) domain, a Paper-I-A trimer representative gives
exactly the finite Weyl algebra of
Proposition~\ref{p1:prop:math-trimer-weyl-a2}; the present proposition records
the same statement for any separately supplied rank-three quantum carrier.
\end{proof}



\paragraph{Log-spectrum decomposition into scale and shape}
\label{p4:sec:part2-threelevel-logshape}

Let the eigenvalues of $Y_f$ on its support be $\{y_1,y_2,y_3\}$ with $y_i>0$ and define
\[
\nu_i:=\ln y_i.
\]
Then $\ln Y_f$ splits uniquely into trace and traceless parts,
\begin{equation}
\label{p4:eq:part2-logsplit}
\ln Y_f
=
I_f\,\mathbf 1_3 + H_f,
\qquad
I_f:=\frac13\Tr(\ln Y_f),
\qquad
\Tr(H_f)=0.
\end{equation}
In an eigenbasis one may write
\[
H_f=\mathrm{diag}(\delta_1,\delta_2,\delta_3),
\qquad
\delta_i:=\nu_i-I_f,
\qquad
\delta_1+\delta_2+\delta_3=0.
\]
Thus the traceless log-spectrum is two-dimensional.
Using the orthonormal basis
\[
e_1=\frac{(1,-1,0)}{\sqrt2},
\qquad
e_2=\frac{(1,1,-2)}{\sqrt6},
\]
of the diagonal traceless subspace, one obtains shape coordinates
\begin{equation}
\label{p4:eq:part2-xy}
(x_f,y_f):=(\delta\cdot e_1,\delta\cdot e_2)
=
\left(\frac{\delta_1-\delta_2}{\sqrt2},\ \frac{\delta_1+\delta_2-2\delta_3}{\sqrt6}\right).
\end{equation}
The scalar $I_f$ is the overall scale coordinate, while $(x_f,y_f)$ records the internal hierarchy of the triplet.



\paragraph{Weyl action and basis-free shape space}
\label{p4:sec:part2-threelevel-weyl}

The ordering of eigenvalues is not physical.
Conjugation by unitaries preserves the spectrum, and the residual ambiguity after diagonalization is exactly permutation of the three eigenvalues.
Hence the physical shape space is the Cartan plane modulo the Weyl group,
\[
\mathfrak h/W,
\qquad
W(\mathrm{SU}(3))\cong S_3.
\]
On the coordinates $(x_f,y_f)$ this acts as the dihedral symmetry of the regular hexagon.
A convenient generator pair is
\[
M_r=
\begin{pmatrix}
-\tfrac12 & -\tfrac{\sqrt3}{2}\\[0.2em]
 \tfrac{\sqrt3}{2} & -\tfrac12
\end{pmatrix},
\qquad
M_s=
\begin{pmatrix}
-1 & 0\\
0 & 1
\end{pmatrix},
\]
corresponding to a $120^\circ$ rotation and a reflection.
Thus any physically meaningful relation among triplet hierarchies should be expressible through symmetric functions of the log-spectrum, or equivalently through $S_3$-invariant functions on the shape plane.


\begin{corollary}[Resolved rank-one projector defects and the \(A_2\) plane]
\label{p4:cor:rank-three-a2-defect-plane}
Let \(P_f\) be the identity of a rank-three corner and let
\(P_0,P_1,P_2\) be rank-one mutually orthogonal projectors resolving that
corner:
\[
P_iP_j=\delta_{ij}P_i,
\qquad
\Tr P_i=1,
\qquad
P_0+P_1+P_2=P_f.
\]
Suppose this ordered projector triple is theorem-fixed by an upstream
carrier, by a spectral decomposition already present in the finite algebra,
or by an explicitly declared interface domain.  Define
\[
D_j:=P_j-\frac13P_f.
\]
Then
\[
D_0+D_1+D_2=0,
\qquad
\operatorname{tr}(D_j^2)=\frac23,
\qquad
\operatorname{tr}(D_iD_j)=-\frac13\quad(i\ne j).
\]
For the simple-root directions
\[
\alpha_1=P_0-P_1,
\qquad
\alpha_2=P_1-P_2,
\]
one has
\[
\operatorname{tr}(\alpha_i\alpha_j)=
\begin{pmatrix}
2&-1\\
-1&2
\end{pmatrix}_{ij}.
\]
Thus the centered diagonal projector plane of any such ordered rank-three theorem domain is the \(A_2\) root/weight plane.  The ordering is notation or downstream sector data, not a Paper-IV sector selector.
\end{corollary}


\begin{proof}
In a basis of the rank-three corner that diagonalizes the three rank-one
projectors,
\(D_0=\operatorname{diag}(2/3,-1/3,-1/3)\), with cyclic permutations for
\(D_1,D_2\), and the displayed trace identities follow by direct
multiplication.
The simple-root Gram matrix is computed from \(P_iP_j=\delta_{ij}P_i\).
On the \(\mathsf{Supp}_{2T}\) domain this specializes to the
Paper-I-A trimer \(A_2\) calculation in
Proposition~\ref{p1:prop:math-trimer-weyl-a2}; the displayed identity itself
is the general rank-three statement.
\end{proof}



\paragraph{Intrinsic invariants and affine shape flow}
\label{p4:sec:part2-threelevel-invariants}

Two canonical invariants of the traceless log-spectrum are
\begin{equation}
\label{p4:eq:part2-casimirs}
C_2(Y_f):=\Tr(H_f^2)=\delta_1^2+\delta_2^2+\delta_3^2=x_f^2+y_f^2,
\qquad
C_3(Y_f):=\Tr(H_f^3)=\delta_1^3+\delta_2^3+\delta_3^3.
\end{equation}
They are invariant under unitary conjugation and under permutations of eigenvalues.
This gives a basis-free language for triplet shape constraints.


\begin{theorem}[Intrinsic first jet and conditional affine linearization on a rank-three sector]
\label{p4:thm:local-affine-normal-form}
Let \(J\subset\mathbb R\) be an interval on which the spectrum is simple, and let
\(\widetilde Y_f(t)\) be a \(C^2\) family of strictly positive rank-three
observables.  Put
\[
I_f(t)=\frac13\sum_{i=1}^3\ln y_i(t),
\qquad
z_f(t)=\binom{x_f(t)}{y_f(t)},
\qquad
\nu_f(t)=\binom{I_f(t)}{z_f(t)}.
\]
Then \(\nu_f\in C^2(J)\), and for every \(t_0\in J\),
\[
\nu_f(t)=\nu_f(t_0)+\dot\nu_f(t_0)(t-t_0)+O(|t-t_0|^2).
\]
This first-jet statement alone does not define an autonomous affine evolution law.
If, on a neighborhood of \(\nu_f(t_0)\), the transported reduced dynamics induces an
autonomous \(C^2\) vector field \(\dot\nu=F(\nu)\), then its genuine first-order
linearization has the form
\begin{align}
\label{p4:eq:part2-affine-normal-form-shape}
\dot z_f
&=A_f(t_0)z_f+b_f(t_0)+R_f,\\
\label{p4:eq:part2-affine-normal-form-scale}
\dot I_f
&=\alpha_f(t_0)+\beta_f(t_0)I_f+\eta_f(t_0)\!\cdot z_f+r_f,
\end{align}
where the coefficients are the corresponding blocks of \(DF(\nu_f(t_0))\), with the
constant terms chosen so that the affine part equals
\(F(\nu_f(t_0))+DF(\nu_f(t_0))(\nu-\nu_f(t_0))\), and
\begin{equation}
\label{p4:eq:part2-affine-remainder}
\|R_f\|+|r_f|\le C_{t_0}\|\nu-\nu_f(t_0)\|^2
\end{equation}
locally.
\end{theorem}

\begin{proof}
Finite-dimensional Rellich regularity gives \(C^2\) ordered eigenvalues on a simple
spectral interval; strict positivity permits their logarithms, so \(I_f,z_f\), and
\(\nu_f\) are \(C^2\).  Taylor's theorem gives the intrinsic first jet in time.
No matrix \(A_f\) is determined by that curve alone.  Under the additional autonomous
hypothesis, Taylor expansion of the vector field in state space gives
\[
F(\nu)=F(\nu_0)+DF(\nu_0)(\nu-\nu_0)+O(\|\nu-\nu_0\|^2).
\]
Writing the Jacobian in scale--shape blocks yields
\eqref{p4:eq:part2-affine-normal-form-shape}--\eqref{p4:eq:part2-affine-normal-form-scale}
and the quadratic remainder bound.
\end{proof}

\begin{corollary}[Projectability criterion for autonomous spectral transport]
\label{p4:cor:intervalwise-affine-normal-form}
Let \(\mathcal L\) be a time-independent finite GKLS generator on the fixed
rank-three corner, and let
\(\pi(Y)=\nu(Y)=(I_f(Y),z_f(Y))\) be the ordered simple-spectrum chart.
An autonomous spectral equation exists exactly on a domain \(\mathcal U\)
of strictly positive, simple-spectrum matrices contained in one ordered
spectral chart, on which \(\mathcal L\) is projectable through \(\pi\):
there is a
well-defined vector field \(F\) such that
\begin{equation}
\label{p4:eq:gkls-spectral-projectability}
D\pi_Y[\mathcal L(Y)]=F(\pi(Y))
\qquad(Y\in\mathcal U).
\end{equation}
Equivalently, the left side must agree for every two matrices in
\(\mathcal U\) with the same spectral coordinates.  On this domain
\(\dot\nu=F(\nu)\), and Theorem~\ref{p4:thm:local-affine-normal-form}
gives the autonomous local affine linearization.

Two sufficient domains are:
\begin{enumerate}[label=(\roman*), leftmargin=2em]
\item \(\mathcal L\) is unitarily equivariant on \(\mathcal U\),
      \(\mathcal L(UYU^*)=U\mathcal L(Y)U^*\); or
\item a fixed maximal abelian algebra is invariant under \(\mathcal L\), the
      domain \(\mathcal U\) remains in one fixed diagonal-label/permutation
      chamber of that algebra, and a no-crossing eigenvalue labeling has
      been fixed.
\end{enumerate}
Approximate spectral closure requires a quantitative bound on the defect of
\eqref{p4:eq:gkls-spectral-projectability}; approximate diagonality in each
matrix's instantaneous eigenbasis is not sufficient.
\end{corollary}

\begin{proof}
The chain rule gives
\[
\dot\nu(t)=D\pi_{Y(t)}[\dot Y(t)]
=D\pi_{Y(t)}[\mathcal L(Y(t))].
\]
This depends only on \(\nu(t)=\pi(Y(t))\) exactly when
\eqref{p4:eq:gkls-spectral-projectability} holds, proving necessity and
sufficiency.  Under unitary equivariance, conjugate matrices have conjugate
velocities and therefore the same eigenvalue derivatives.  On an invariant
fixed maximal abelian algebra, the matrix entries are fixed functions of the
ordered eigenvalues, so the same conclusion holds.  Taylor expansion of the
resulting \(F\) is Theorem~\ref{p4:thm:local-affine-normal-form}.

The additional hypothesis is essential.  For qubit dephasing
\(\mathcal L(\rho)=\gamma(\sigma_z\rho\sigma_z-\rho)\), the isospectral
states
\[
\rho_z=\tfrac12(\mathbf1+r\sigma_z),
\qquad
\rho_x=\tfrac12(\mathbf1+r\sigma_x)
\]
have different spectral derivatives: \(\rho_z\) is stationary while the
Bloch radius of \(\rho_x\) decreases.  Thus time independence of a GKLS
generator alone does not define autonomous eigenvalue dynamics.
\end{proof}

\begin{corollary}[Transfer to any fixed linear symmetry-adapted chart]
\label{p4:cor:chart-transfer-affine}
Let $q_f\in\mathbb R^n$ be a fixed full-rank linear reparameterization of the controlled scale--shape variables $(I_f,z_f)$, for example $q_f=L(I_f,z_f)+q_0$ with $L$ injective on the three-dimensional scale--shape space and with a chosen inverse on its image.
If $(I_f,z_f)$ obeys the affine normal form of Theorem~\ref{p4:thm:local-affine-normal-form}, then so does $q_f$ after this linear change of coordinates:
\begin{equation}
\label{p4:eq:chart-transfer-affine}
\dot q_f(t)=M_f(t_0)q_f(t)+c_f(t_0)+\widetilde R_f(t;t_0).
\end{equation}
In particular, any fixed linear symmetry-adapted family chart that is a reparameterization of the scale--shape variables inherits the same intervalwise affine first-order transport law.
Nonlinear reparameterizations need not preserve this affine form.
\end{corollary}

\begin{proof}
Affine vector fields are preserved under invertible linear coordinate changes.
The remainder estimate transforms by the operator norms of the change-of-basis matrices.
The final statement follows because nonlinear coordinate changes introduce higher-order terms even when the original chart is affine.
\end{proof}



\paragraph{Relative diagonalization, overlap matrices, and mixing stability}
\label{p4:sec:part2-threelevel-mixing}

Let $A,B\in\mathcal A_f$ be self-adjoint observables with simple ordered spectra
\[
a_1<a_2<a_3,
\qquad
b_1<b_2<b_3,
\]
and write their spectral decompositions as
\[
A=\sum_{i=1}^3 a_i P_i,
\qquad
B=\sum_{j=1}^3 b_j Q_j,
\]
with rank-one spectral projectors $P_i,Q_j$.
Choosing normalized eigenvectors $u_i\in P_i\mathcal H_f$ and $v_j\in Q_j\mathcal H_f$ gives unitary diagonalizations
\[
A=U_A\,\diag(a_1,a_2,a_3)\,U_A^\dagger,
\qquad
B=U_B\,\diag(b_1,b_2,b_3)\,U_B^\dagger,
\]
where $U_A=(u_1\,u_2\,u_3)$ and $U_B=(v_1\,v_2\,v_3)$.
The relative unitary
\[
V:=U_A^\dagger U_B
\]
is defined only up to multiplication on the left and right by diagonal phase matrices, but its modulus data and the usual rephasing-invariant quartets are intrinsic.
Here the spectral projectors are not inserted selectors: they are generated by the already supplied observables \(A\) and \(B\) inside the finite algebra.  If the choice of \(A\) or \(B\) is later sector-specific, that sector choice remains downstream/interface data.


\begin{theorem}[Projector overlaps reconstruct the mixing moduli]
\label{p4:thm:projector-overlap-mixing}
With notation as above, define the overlap matrix
\begin{equation}
\label{p4:eq:mixing-overlap-matrix}
M_{ij}:=\Tr(P_iQ_j).
\end{equation}
Then
\begin{equation}
\label{p4:eq:mixing-overlap-moduli}
M_{ij}=|\langle u_i,v_j\rangle|^2 = |V_{ij}|^2.
\end{equation}
In particular:
\begin{enumerate}
\item $M$ depends only on the pair $(A,B)$ and not on the phase choices of the eigenvectors;
\item $M$ is doubly stochastic,
\[
\sum_{j=1}^3 M_{ij}=1,
\qquad
\sum_{i=1}^3 M_{ij}=1;
\]
\item the full modulus data of the relative mixing matrix is recovered basis-freely from the spectral projectors of $A$ and $B$.
\end{enumerate}
\end{theorem}


\begin{proof}
For rank-one spectral projectors,
$\operatorname{Tr}(|u_i\rangle\langle u_i|v_j\rangle\langle v_j|)
=|\langle u_i,v_j\rangle|^2$, which gives
\eqref{p4:eq:mixing-overlap-moduli} independently of eigenvector phases.
Summing either family of projectors to the identity gives the unit row
and column sums. This is the overlap data used in Paper~1B,
Theorem~\ref{p1b:thm:rank-three}. If the present Hermitian matrices are
not positive, independent scalar shifts make them positive without
changing any eigenprojector or relative mixing modulus.
\end{proof}


\begin{theorem}[Gap-controlled eigenframe mismatch for nearby rank-three observables]
\label{p4:thm:gap-controlled-eigenframe-mismatch}
Let $E:=B-A$.
For every pair $(i,j)$,
\begin{equation}
\label{p4:eq:exact-eigenframe-identity}
(a_i-b_j)\,\langle u_i,v_j\rangle
=
-\langle u_i,E v_j\rangle,
\end{equation}
and therefore
\begin{equation}
\label{p4:eq:exact-eigenframe-bound}
|V_{ij}| = |\langle u_i,v_j\rangle|
\le
\frac{\|E\|}{|a_i-b_j|}
\qquad\text{whenever }a_i\neq b_j.
\end{equation}
Define the local spectral gaps
\[
g_i^{(A)}:=\min_{k\neq i}|a_i-a_k|,
\qquad
g_j^{(B)}:=\min_{\ell\neq j}|b_j-b_\ell|.
\]
If $\|E\|<g_i^{(A)}$, then for every $j\neq i$,
\begin{equation}
\label{p4:eq:gap-controlled-offdiag-bound}
|V_{ij}|
\le
\frac{\|E\|}{g_i^{(A)}-\|E\|}.
\end{equation}
Symmetrically, if $\|E\|<g_j^{(B)}$, then for every $i\neq j$,
\begin{equation}
\label{p4:eq:gap-controlled-offdiag-bound-B}
|V_{ij}|
\le
\frac{\|E\|}{g_j^{(B)}-\|E\|}.
\end{equation}
\end{theorem}


\begin{proof}
\ProofPlan{We prove the gap-controlled mismatch estimate by deriving the exact eigenframe identity, bounding overlap entries by the perturbation norm, applying spectral gaps to off-diagonal entries, and then symmetrizing the estimate.}
\begin{proofsteps}

\item \textbf{Exact eigenframe identity.}
Using $B=A+E$, the eigenvalue equations $Au_i=a_i u_i$ and $Bv_j=b_j v_j$ give
\[
a_i\langle u_i,v_j\rangle
=
\langle u_i,Av_j\rangle
=
\langle u_i,(B-E)v_j\rangle
=
b_j\langle u_i,v_j\rangle - \langle u_i,Ev_j\rangle,
\]
which is exactly \eqref{p4:eq:exact-eigenframe-identity}.

\item \textbf{Norm bound for overlap entries.}
Taking absolute values yields
\[
|a_i-b_j|\,|V_{ij}|
\le
|\langle u_i,Ev_j\rangle|
\le
\|E\|.
\]
This proves \eqref{p4:eq:exact-eigenframe-bound}.


\item \textbf{Off-diagonal estimate from spectral gaps.}
To obtain the off-diagonal estimate, use the finite-dimensional Weyl bound
\[
|a_k-b_k|\le \|E\| \qquad (k=1,2,3)
\]
for the ordered eigenvalues.
If $j\neq i$, then
\[
|a_i-b_j|
\ge
|a_i-a_j|-|a_j-b_j|
\ge
g_i^{(A)}-\|E\|.
\]
Substituting this into \eqref{p4:eq:exact-eigenframe-bound} gives \eqref{p4:eq:gap-controlled-offdiag-bound}.

\item \textbf{Symmetric estimate with \(A\) and \(B\) interchanged.}
The proof of \eqref{p4:eq:gap-controlled-offdiag-bound-B} is identical after interchanging $A$ and $B$.
\end{proofsteps}

\end{proof}

\begin{corollary}[Near-identity modulus matrix under a small sector mismatch]
\label{p4:cor:near-identity-modulus-matrix}
For a fixed $i$ in the small-mismatch domain $\|E\|<g_i^{(A)}$, set
\[
\varepsilon_i^{(A)}:=\frac{\|E\|}{g_i^{(A)}-\|E\|}.
\]
Then
\begin{equation}
\label{p4:eq:near-identity-row-bound}
\sum_{j\neq i} M_{ij}
=
\sum_{j\neq i}|V_{ij}|^2
\le
2\bigl(\varepsilon_i^{(A)}\bigr)^2,
\qquad
M_{ii}\ge 1-2\bigl(\varepsilon_i^{(A)}\bigr)^2.
\end{equation}
So if two rank-three observables differ by a perturbation smaller than the protecting spectral gaps, the relative modulus matrix is quantitatively close to the identity.
\end{corollary}

\begin{proof}
For $j\neq i$, Theorem~\ref{p4:thm:gap-controlled-eigenframe-mismatch} gives
\[
M_{ij}=|V_{ij}|^2\le \bigl(\varepsilon_i^{(A)}\bigr)^2.
\]
Since there are exactly two off-diagonal terms in row $i$,
\[
\sum_{j\neq i}M_{ij}\le 2\bigl(\varepsilon_i^{(A)}\bigr)^2.
\]
Now use the row-sum identity $\sum_j M_{ij}=1$ from Theorem~\ref{p4:thm:projector-overlap-mixing}.
\end{proof}


\begin{theorem}[Transport stability of projector overlaps on a gapped interval]
\label{p4:thm:transport-stability-projector-overlaps}
Let $J\subset\mathbb R$ be an interval and let $t\mapsto A(t),B(t)\in\mathcal A_f$ be $C^1$ families of self-adjoint observables with simple spectra
\[
a_1(t)<a_2(t)<a_3(t),
\qquad
b_1(t)<b_2(t)<b_3(t)
\]
for all $t\in J$.
Let $P_i(t)$ and $Q_j(t)$ be the corresponding rank-one spectral projectors and define
\begin{equation}
\label{p4:eq:transport-overlap-matrix}
M_{ij}(t):=\Tr\!\bigl(P_i(t)Q_j(t)\bigr).
\end{equation}
Then each $M_{ij}$ is $C^1$ on $J$.
If
\[
g_i^{(A)}(t):=\min_{k\neq i}|a_i(t)-a_k(t)|,
\qquad
g_j^{(B)}(t):=\min_{\ell\neq j}|b_j(t)-b_\ell(t)|,
\]
then
\begin{equation}
\label{p4:eq:transport-overlap-derivative-bound}
|\dot M_{ij}(t)|
\le
2\frac{\|\dot A(t)\|}{g_i^{(A)}(t)}
+
2\frac{\|\dot B(t)\|}{g_j^{(B)}(t)}.
\end{equation}
Consequently, for any $s,t\in J$,
\begin{equation}
\label{p4:eq:transport-overlap-integrated-bound}
|M_{ij}(t)-M_{ij}(s)|
\le
2\int_{\min(s,t)}^{\max(s,t)}
\left(
\frac{\|\dot A(\tau)\|}{g_i^{(A)}(\tau)}
+
\frac{\|\dot B(\tau)\|}{g_j^{(B)}(\tau)}
\right)\,d\tau.
\end{equation}
\end{theorem}


\begin{proof}
\ProofPlan{We prove transport stability of projector overlaps by bounding spectral-projector derivatives with reduced resolvents, differentiating the overlap matrix, estimating the trace terms, and integrating the resulting Lipschitz bound.}
\begin{proofsteps}

\item \textbf{Projector derivative bounds from reduced resolvents.}
Because the spectra are simple and the families are $C^1$, the spectral projectors are $C^1$ on $J$.
For an isolated simple branch of $A(t)$, the reduced resolvent
\[
S_i^{(A)}(t):=\sum_{k\neq i}\frac{P_k(t)}{a_i(t)-a_k(t)}
\]
satisfies $\|S_i^{(A)}(t)\|=1/g_i^{(A)}(t)$, and the standard projector derivative formula gives
\begin{equation}
\label{p4:eq:projector-derivative-A}
\dot P_i
=
S_i^{(A)}\,\dot A\,P_i + P_i\,\dot A\,S_i^{(A)}.
\end{equation}
Hence
\begin{equation}
\label{p4:eq:projector-derivative-A-bound}
\|\dot P_i(t)\|
\le
2\frac{\|\dot A(t)\|}{g_i^{(A)}(t)}.
\end{equation}
The same argument for $B(t)$ yields
\begin{equation}
\label{p4:eq:projector-derivative-B-bound}
\|\dot Q_j(t)\|
\le
2\frac{\|\dot B(t)\|}{g_j^{(B)}(t)}.
\end{equation}


\item \textbf{Derivative of the overlap matrix.}
Differentiating \eqref{p4:eq:transport-overlap-matrix} gives
\[
\dot M_{ij}(t)
=
\Tr\!\bigl(\dot P_i(t)Q_j(t)\bigr)
+
\Tr\!\bigl(P_i(t)\dot Q_j(t)\bigr).
\]

\item \textbf{Trace-norm estimate for rank-one projectors.}
Because $P_i(t)$ and $Q_j(t)$ are rank-one projectors, $\|P_i(t)\|_1=\|Q_j(t)\|_1=1$.
Therefore
\[
|\Tr(\dot P_iQ_j)|
\le
\|\dot P_iQ_j\|_1
\le
\|\dot P_i\|\,\|Q_j\|_1
=
\|\dot P_i\|,
\]
and similarly $|\Tr(P_i\dot Q_j)|\le \|\dot Q_j\|$.

\item \textbf{Gap-controlled Lipschitz and integrated estimates.}
Combining these inequalities with \eqref{p4:eq:projector-derivative-A-bound} and \eqref{p4:eq:projector-derivative-B-bound} proves \eqref{p4:eq:transport-overlap-derivative-bound}.
The integrated estimate \eqref{p4:eq:transport-overlap-integrated-bound}
follows by integrating the absolute derivative over the interval with
endpoints \(s,t\); this also covers \(t<s\).
\end{proofsteps}

\end{proof}

\begin{proposition}[Finite cubic commutator invariant and perturbative smallness]
\label{p4:prop:cubic-commutator-invariant}
For self-adjoint $A,B\in\mathcal A_f$, define
\begin{equation}
\label{p4:eq:cubic-commutator-invariant}
\mathcal C_3(A,B):=\frac{1}{6i}\Tr\!\bigl([A,B]^3\bigr).
\end{equation}
Then $\mathcal C_3(A,B)\in\mathbb R$, it is invariant under simultaneous unitary conjugation
\[
\mathcal C_3(U A U^\dagger, U B U^\dagger)=\mathcal C_3(A,B),
\]
and it vanishes whenever $A$ and $B$ commute.
Moreover,
\begin{equation}
\label{p4:eq:cubic-commutator-bound}
|\mathcal C_3(A,B)|
\le
\frac12 \|[A,B]\|^3.
\end{equation}
If $B=A+E$, then
\begin{equation}
\label{p4:eq:cubic-commutator-perturbative-bound}
|\mathcal C_3(A,A+E)|
\le
4\|A\|^3\|E\|^3.
\end{equation}
\end{proposition}

\begin{proof}
\ProofPlan{We prove the cubic commutator invariant and its perturbative smallness by checking reality and unitary invariance, bounding the cubic trace in norm, and substituting \(B=A+E\).}
\begin{proofsteps}

\item \textbf{Reality and unitary invariance.}
Since $A$ and $B$ are self-adjoint, $[A,B]^*=-[A,B]$, so $[A,B]$ is anti-Hermitian and $[A,B]^3$ is also anti-Hermitian.
Therefore $\Tr([A,B]^3)$ is purely imaginary, which shows that $\mathcal C_3(A,B)$ is real.
Unitary invariance follows from
\[
[UAU^\dagger,UBU^\dagger]=U[A,B]U^\dagger
\]
and cyclicity of the trace.
If $[A,B]=0$, then clearly $\mathcal C_3(A,B)=0$.


\item \textbf{Cubic trace-norm bound.}
To prove \eqref{p4:eq:cubic-commutator-bound}, let $X:=[A,B]$.
In dimension three,
\[
|\Tr(X^3)|\le \sum_{r=1}^3 |\lambda_r(X^3)| \le 3\|X^3\| \le 3\|X\|^3,
\]
where $\lambda_r(X^3)$ are the eigenvalues of $X^3$.
Dividing by $6$ gives
\[
|\mathcal C_3(A,B)|\le \frac12 \|[A,B]\|^3.
\]


\item \textbf{Perturbative smallness for \(B=A+E\).}
If $B=A+E$, then $[A,B]=[A,E]$ and
\[
\|[A,E]\|
\le
\|AE\|+\|EA\|
\le
2\|A\|\,\|E\|.
\]
Substituting this into \eqref{p4:eq:cubic-commutator-bound} yields
\[
|\mathcal C_3(A,A+E)|
\le
\frac12 (2\|A\|\,\|E\|)^3
=
4\|A\|^3\|E\|^3.
\]
\end{proofsteps}

\end{proof}

\begin{proposition}[Rank-three normalized Jarlskog identity]
\label{p4:prop:rank-three-jarlskog-identity}
When $A$ and $B$ have simple spectra as above, define the cyclic Vandermonde factors
\[
\widetilde\Delta_A:=(a_1-a_2)(a_2-a_3)(a_3-a_1),
\qquad
\widetilde\Delta_B:=(b_1-b_2)(b_2-b_3)(b_3-b_1).
\]
Let
\[
J(V):=\Im\!\bigl(V_{11}V_{22}\overline{V_{12}V_{21}}\bigr),
\]
which is the standard rank-three rephasing-invariant quartet, unique up to the usual sign changes coming from permutations of rows or columns.
Then
\begin{equation}
\label{p4:eq:rank-three-jarlskog-identity}
\mathcal C_3(A,B)=\widetilde\Delta_A\,\widetilde\Delta_B\,J(V),
\end{equation}
equivalently,
\begin{equation}
\label{p4:eq:rank-three-jarlskog-normalized}
J(V)=\frac{1}{6i\,\widetilde\Delta_A\,\widetilde\Delta_B}\Tr\!\bigl([A,B]^3\bigr).
\end{equation}
If $B=A+E$, then
\begin{equation}
\label{p4:eq:rank-three-jarlskog-perturbative-bound}
|J(V)|
\le
\frac{4\|A\|^3\|E\|^3}{|\widetilde\Delta_A\,\widetilde\Delta_B|}.
\end{equation}
\end{proposition}

\begin{proof}
\ProofPlan{We derive the normalized Jarlskog identity by reducing to the eigenbasis of \(A\), rewriting the cubic commutator determinant, proving Vandermonde divisibility in the \(B\)-spectrum, identifying the Jarlskog factor, and then reading the normalized perturbative form.}
\begin{proofsteps}

\item \textbf{Reduction to the eigenbasis of \(A\).}
By Proposition~\ref{p4:prop:cubic-commutator-invariant}, $\mathcal C_3(A,B)$ is invariant under simultaneous unitary conjugation.
So we may work in the eigenbasis of $A$, where
\[
A=\diag(a_1,a_2,a_3),
\qquad
B=V\,\diag(b_1,b_2,b_3)\,V^\dagger.
\]

\item \textbf{Cubic commutator determinant reduction.}
Set $C:=[A,B]$.
Then $C$ is a traceless $3\times 3$ matrix with vanishing diagonal and off-diagonal entries
\[
C_{ij}=(a_i-a_j)B_{ij}
\qquad (i\neq j).
\]
For a traceless $3\times 3$ matrix one has $\Tr(C^3)=3\det C$, and for a matrix with zero diagonal
\[
\det C = C_{12}C_{23}C_{31}+C_{13}C_{21}C_{32}.
\]
Using $C^*=-C$ gives
\[
\det C
=
2i\,\widetilde\Delta_A\,\Im(B_{12}B_{23}B_{31}),
\]
hence
\begin{equation}
\label{p4:eq:jarlskog-proof-reduced}
\mathcal C_3(A,B)
=
\widetilde\Delta_A\,\Im(B_{12}B_{23}B_{31}).
\end{equation}


\item \textbf{Vandermonde divisibility in the \(B\)-spectrum.}
Now $B_{12}B_{23}B_{31}$ is a homogeneous cubic polynomial in $(b_1,b_2,b_3)$.
If two eigenvalues coincide, say $b_1=b_2$, then
\[
B=b_2\,\mathbf 1 +(b_3-b_2)\,|v_3\rangle\langle v_3|,
\]
so
\[
B_{12}B_{23}B_{31}
=
(b_3-b_2)^3\,|V_{13}|^2|V_{23}|^2|V_{33}|^2\in\mathbb R,
\]
and therefore $\Im(B_{12}B_{23}B_{31})=0$.
The same argument applies when any other pair of $b$'s coincide.
Thus $\Im(B_{12}B_{23}B_{31})$ is divisible by $\widetilde\Delta_B$, so there exists a scalar $K(V)$ depending only on the relative unitary such that
\[
\Im(B_{12}B_{23}B_{31})=K(V)\,\widetilde\Delta_B.
\]

\item \textbf{Identification of the Jarlskog factor.}
A direct expansion in the special case $(b_1,b_2,b_3)=(1,0,-1)$, using only row/column orthogonality of the unitary $V$, gives
\[
\Im(B_{12}B_{23}B_{31})
=
-2\,\Im\!\bigl(V_{11}V_{22}\overline{V_{12}V_{21}}\bigr).
\]
Since for this choice $\widetilde\Delta_B=-2$, one gets $K(V)=J(V)$.
Substituting into \eqref{p4:eq:jarlskog-proof-reduced} proves \eqref{p4:eq:rank-three-jarlskog-identity}, and \eqref{p4:eq:rank-three-jarlskog-normalized} is the same identity rewritten.


\item \textbf{Normalized and perturbative forms.}
Finally, \eqref{p4:eq:rank-three-jarlskog-perturbative-bound} follows immediately from
\eqref{p4:eq:rank-three-jarlskog-normalized} and the perturbative bound
\eqref{p4:eq:cubic-commutator-perturbative-bound}.
\end{proofsteps}

\end{proof}



\subsection{Canonical transported sectors, block elimination, and threshold composition}
\label{p4:sec:part2-canonical-transported-blocks}

\paragraph{Status and motivation.}
The preceding sections already show that rank-three overlap data, transported projector families, and finite resonance reductions are legitimate quantum-mechanical constructions.
For Papers~V-A--V-B one needs a slightly stronger statement: not only may one work in transported finite sectors, but the relevant overlap data and reduced block dynamics should be canonical up to the expected gauge freedoms of constant unitary change, diagonal rephasing, and permutation of spectral labels.
The present section isolates exactly that finite theorem package.
Nothing here assumes continuum QFT or an infinite-volume limit; all statements are finite-sector consequences of the quantum mechanics already developed in the representation/interface parts.



\paragraph{Transported frames and constant-unitary equivalence}
\label{p4:sec:part2-transported-frame-equivalence}

Let $P(s)$ be a $C^1$ family of rank-three projectors on a finite interval and let $F(s):\CC^3\to\mathcal H_m$ be a $C^1$ family of isometries such that
\begin{equation}
\label{p4:eq:transported-isometry-projector}
F(s)^*F(s)=\mathbf 1_3,
\qquad
F(s)F(s)^*=P(s).
\end{equation}
If moreover
\begin{equation}
\label{p4:eq:transported-isometry-kato}
\dot F(s)=[\dot P(s),P(s)]F(s),
\end{equation}
then $F(s)$ is a transported orthonormal frame of the moving rank-three sector.

\begin{theorem}[Canonical transported frames differ only by a constant unitary]
\label{p4:thm:transported-frame-constant-unitary}
Let $F_1(s)$ and $F_2(s)$ be two transported orthonormal frames for the same projector family $P(s)$, both satisfying \eqref{p4:eq:transported-isometry-projector}--\eqref{p4:eq:transported-isometry-kato}.
Then there exists a constant unitary $U_0\in U(3)$ such that
\begin{equation}
\label{p4:eq:transported-constant-unitary}
F_2(s)=F_1(s)U_0
\qquad\text{for all }s.
\end{equation}
Consequently, for any supported observable $Y(s)=P(s)Y(s)P(s)$ the reduced matrices
\begin{equation}
\label{p4:eq:transported-reduced-matrices-twoframes}
Y_j(s):=F_j(s)^*Y(s)F_j(s)\in\mathrm{Mat}(3,\CC)
\qquad (j=1,2)
\end{equation}
are related by a constant simultaneous conjugation,
\begin{equation}
\label{p4:eq:transported-reduced-constant-conjugation}
Y_2(s)=U_0^*Y_1(s)U_0.
\end{equation}
Hence all spectral-overlap moduli and normalized cubic/Jarlskog invariants extracted from a transported rank-three sector are independent of the transported-frame choice.
\end{theorem}

\begin{proof}
Set $U(s):=F_1(s)^*F_2(s)$.
Using \eqref{p4:eq:transported-isometry-kato} and the anti-self-adjointness of $[\dot P,P]$,
\[
\dot U
=
\dot F_1^*F_2+F_1^*\dot F_2
=
-F_1^*[\dot P,P]F_2+F_1^*[\dot P,P]F_2=0.
\]
Thus $U(s)=U_0$ is constant.
Since each $F_j(s)$ is an isometry, $U_0^*U_0=\mathbf 1_3$, so $U_0\in U(3)$.
The relation \eqref{p4:eq:transported-constant-unitary} follows from $F_2=F_1U$.
Equation~\eqref{p4:eq:transported-reduced-constant-conjugation} is then immediate:
\[
Y_2=F_2^*YF_2=U_0^*F_1^*YF_1U_0=U_0^*Y_1U_0.
\]
All spectral data and all simultaneous-conjugation invariants are therefore identical in the two transported frames.
\end{proof}

\begin{corollary}[Transport-gauge invariance of overlap and cubic data]
\label{p4:cor:transport-gauge-invariance}
In the setting of Theorem~\ref{p4:thm:transported-frame-constant-unitary}, any modulus matrix reconstructed from projector overlaps and any normalized cubic invariant of the form
\[
\frac{1}{6i\,\widetilde\Delta_A\widetilde\Delta_B}\Tr([A,B]^3)
\]
computed from reduced transported representatives is independent of the transported-frame gauge.
Only the usual rephasing and permutation freedoms of the final rank-three representative remain.
\end{corollary}

\begin{proof}
The projector-overlap data of Section~\ref{p4:sec:part2-threelevel-mixing} and the cubic invariant of Proposition~\ref{p4:prop:rank-three-jarlskog-identity} are invariant under simultaneous unitary conjugation.
Apply this to \eqref{p4:eq:transported-reduced-constant-conjugation}.
\end{proof}



\paragraph{Comparison representatives and quotient invariants}
\label{p4:sec:part2-comparison-quotient}

\begin{theorem}[Comparison-slice quotient invariance]
\label{p4:thm:comparison-slice-quotient}
Let $V\in U(3)$ be a relative mixing matrix between two simple-spectrum rank-three observables, and let
\begin{equation}
\label{p4:eq:comparison-quotient-transform}
V' = \Pi_L D_L\, V\, D_R \Pi_R,
\end{equation}
where $D_L,D_R$ are diagonal unitaries and $\Pi_L,\Pi_R$ are permutation matrices.
Then:
\begin{enumerate}
\item the modulus matrix of $V'$ is obtained from that of $V$ only by the corresponding row and column permutations;
\item every rephasing-invariant polynomial in the squared moduli is unchanged after relabeling;
\item the normalized cubic invariant satisfies
\begin{equation}
\label{p4:eq:comparison-jarlskog-sign}
J(V') = \det(\Pi_L)\det(\Pi_R)\,J(V).
\end{equation}
\end{enumerate}
Accordingly, the physical invariant object attached to a rank-three comparison sector is the equivalence class of $V$ modulo diagonal rephasing and spectral-label permutations.
\end{theorem}

\begin{proof}
Multiplication by $D_L$ and $D_R$ changes only row and column phases, so it does not affect squared moduli and leaves $J(V)$ unchanged.
Left and right multiplication by permutation matrices only relabel rows and columns.
For the quartet definition of $J(V)$, an odd row or column permutation reverses the orientation of the ordered label triple and changes the sign, while an even permutation leaves it unchanged.
This proves \eqref{p4:eq:comparison-jarlskog-sign}.
\end{proof}



\paragraph{Finite light/heavy block elimination}
\label{p4:sec:part2-block-elimination}

\begin{theorem}[Positive finite block elimination]
\label{p4:thm:positive-block-elimination}
Let
\begin{equation}
\label{p4:eq:block-elimination-matrix}
H=
\begin{pmatrix}
A & X\\
X^* & D
\end{pmatrix}
\ge 0
\end{equation}
on a finite decomposition $\mathcal H=\mathcal H_L\oplus\mathcal H_H$, with $D\ge \mu\,\mathbf 1$ on $\mathcal H_H$ for some $\mu>0$.
Define the Schur-complement reduction on the light block by
\begin{equation}
\label{p4:eq:block-elimination-schur}
H_{\mathrm{eff}}:=A-XD^{-1}X^*.
\end{equation}
Then $H_{\mathrm{eff}}\ge 0$ on $\mathcal H_L$.
Moreover, on the open set of triples $(A,X,D)$ with $D$ uniformly positive, the map
\begin{equation}
\label{p4:eq:block-elimination-map}
(A,X,D)\longmapsto H_{\mathrm{eff}}
\end{equation}
is continuous in operator norm.
Hence the eigenvalues and spectral projectors of $H_{\mathrm{eff}}$ vary continuously under finite perturbations of the heavy block and the coupling block, as long as the relevant spectral gaps stay open.
\end{theorem}

\begin{proof}
The stationary Schur calculation of Paper~1B,
Theorem~\ref{p1b:thm:schur-jets} (or its exact quadratic form in
Theorem~\ref{p1b:thm:all-order-schur}), applied to the underlying real
Hermitian form, gives $H_{\rm eff}=A-XD^{-1}X^*$. Evaluating the positive
form at $(x,-D^{-1}X^*x)$ proves $H_{\rm eff}\ge0$.
Matrix products and inversion are norm-continuous on the uniformly
positive branch; sufficiently small perturbations preserve, for example,
the bound $D\ge\mu I/2$. Weyl continuity gives the eigenvalue claim,
and the contour formula for an isolated spectral cluster gives continuity
of its projector while the relevant gap remains open.
\end{proof}



\paragraph{Palindromic composition and BCH cancellation}
\label{p4:sec:part2-palindromic-bch}

\begin{theorem}[Palindromic BCH cancellation]
\label{p4:thm:palindromic-bch}
Let $X,Y\in\End(\mathcal H)$ be finite operators and define, for sufficiently small real $t$,
\begin{equation}
\label{p4:eq:palindromic-composition}
S(t):=e^{tX/2}e^{tY}e^{tX/2}.
\end{equation}
Then the principal logarithm $Z(t):=\log S(t)$ exists for $|t|\ll1$ and is an odd power series in $t$.
In particular, there is no quadratic commutator term, and the expansion begins
\begin{equation}
\label{p4:eq:palindromic-bch-expansion}
Z(t)
=
t(X+Y)
+
\frac{t^3}{12}[Y,[Y,X]]
-
\frac{t^3}{24}[X,[X,Y]]
+
O(t^5).
\end{equation}
Hence the first correction to the naive sum $X+Y$ is nested-bracket order, not linear commutator order.
\end{theorem}

\begin{proof}
Because
\[
S(-t)=e^{-tX/2}e^{-tY}e^{-tX/2}=S(t)^{-1},
\]
and the principal logarithm of an inverse is minus the principal logarithm whenever the logarithm exists, one has
\[
Z(-t)=-Z(t).
\]
Thus all even powers of $t$ vanish in the Taylor expansion of $Z(t)$.
A direct Baker--Campbell--Hausdorff expansion to cubic order then yields \eqref{p4:eq:palindromic-bch-expansion}.
In particular the usual $\tfrac12[X,Y]$ term cancels by palindromic symmetry.
\end{proof}



\paragraph{Piecewise-affine composition across resolved thresholds}
\label{p4:sec:part2-piecewise-affine-thresholds}

Consider a finite control parameter $s\in[s_0,s_N]$ together with a partition
\[
s_0<s_1<\cdots<s_N,
\]
such that on each open interval $I_r:=(s_{r-1},s_r)$ a rank-three projector family $P_r(s)$ is gapped and admits a transported frame $F_r(s)$ in the sense of \eqref{p4:eq:transported-isometry-projector}--\eqref{p4:eq:transported-isometry-kato}.
On the affine-transport theorem domain, let the reduced observable of interest have the intervalwise form
\begin{equation}
\label{p4:eq:piecewise-affine-local-form}
Y_r(s):=F_r(s)^*Y(s)F_r(s)=A_r+(s-s_{r-1})B_r.
\end{equation}
When a named sector theorem, a continuous transported-frame domain, or an explicitly declared interface model supplies a unitary boundary identification between the incoming and outgoing rank-three sectors, the continuous-sector case means $P_{r+1}(s_r^+)=P_r(s_r^-)$, and one may define the boundary conversion unitary by
\begin{equation}
\label{p4:eq:threshold-boundary-unitary}
U_r:=F_{r+1}(s_r)^*F_r(s_r)\in U(3).
\end{equation}
More generally, $U_r$ is theorem-domain threshold matching data.  Under exact GP2 it is structurally admissible only when intrinsic to, generated by, or fixed by the already constituted finite carrier, or when imported from a named sector theorem; otherwise it is an explicitly declared interface representative with no sector-selecting power.  Without such a unitary matching condition, $F_{r+1}(s_r)^*F_r(s_r)$ is only an overlap matrix and the unitary-conjugation law below does not follow.

\begin{proposition}[Piecewise-affine same-rank threshold composition and jump defect]
\label{p4:prop:piecewise-affine-threshold-composition}
Under the preceding hypotheses, the reduced observable on the initial
reference corner is, on each interval \(I_r\), of the form
\begin{equation}
\label{p4:eq:global-piecewise-affine-law}
\widehat Y_r(s)
=
\mathcal U_r^*\bigl(A_r+(s-s_{r-1})B_r\bigr)\mathcal U_r,
\qquad
\mathcal U_r:=U_{r-1}\cdots U_1,
\end{equation}
with \(\mathcal U_1:=\mathbf 1\).  Put
\begin{equation}
\label{p4:eq:threshold-incoming-value}
Y_r^-:=A_r+(s_r-s_{r-1})B_r
\end{equation}
and define the threshold defect in the outgoing coordinates by
\begin{equation}
\label{p4:eq:threshold-jump-defect}
\boxed{\;
\Delta_r:=A_{r+1}-U_rY_r^-U_r^* .
\;}
\end{equation}
Then the exact one-sided jump on the initial reference corner is
\begin{equation}
\label{p4:eq:threshold-boundary-conjugation}
\boxed{\;
\widehat Y_{r+1}(s_r^+)-\widehat Y_r(s_r^-)
=\mathcal U_{r+1}^*\Delta_r\mathcal U_{r+1}.
\;}
\end{equation}
Hence
\[
\|\widehat Y_{r+1}(s_r^+)-\widehat Y_r(s_r^-)\|
=\|\Delta_r\|.
\]
The threshold is only a transported-frame conversion exactly when
\(\Delta_r=0\); otherwise \(\Delta_r\) records an independent observable
jump.

If \(P(s)\) and \(Y(s)\) have common norm-continuous boundary values, the
rank remains constant, the relevant gap stays open, and the two frames have
the stated one-sided limits, then
\[
A_{r+1}=U_rY_r^-U_r^*,
\qquad
\Delta_r=0.
\]
In particular, this zero-defect conclusion is derived for a continuously
transported observable on the open-gap domain of
Theorem~\ref{p4:thm:positive-block-elimination}; no independent endpoint
matching premise is needed there.  At a gap-closing or discontinuous
threshold, the affine interval laws and the unitary map alone do not force
\(\Delta_r\) to vanish.
\end{proposition}

\begin{proof}
The interval formula is
\eqref{p4:eq:piecewise-affine-local-form} pulled back by the cumulative
unitary \(\mathcal U_r\).  Since
\(\mathcal U_{r+1}=U_r\mathcal U_r\), its two one-sided endpoint values obey
\[
\begin{aligned}
\widehat Y_r(s_r^-)
&=\mathcal U_r^*Y_r^-\mathcal U_r,\\
\widehat Y_{r+1}(s_r^+)
&=\mathcal U_r^*U_r^*A_{r+1}U_r\mathcal U_r .
\end{aligned}
\]
Subtracting and using
\(U_r^*\Delta_rU_r=U_r^*A_{r+1}U_r-Y_r^-\) gives
\[
\widehat Y_{r+1}(s_r^+)-\widehat Y_r(s_r^-)
=\mathcal U_{r+1}^*\Delta_r\mathcal U_{r+1},
\]
and unitary invariance proves the norm identity.

On the continuous open-gap domain, the limiting frames are isometries onto
the same rank-three subspace.  The definition
\(U_r=F_{r+1}(s_r)^*F_r(s_r)\) therefore gives
\[
F_r(s_r)=F_{r+1}(s_r)U_r.
\]
Norm continuity of \(Y(s)\) then implies
\[
Y_r^-
=F_r^*Y(s_r)F_r
=U_r^*F_{r+1}^*Y(s_r)F_{r+1}U_r
=U_r^*A_{r+1}U_r.
\]
Equivalently \(A_{r+1}=U_rY_r^-U_r^*\), so
\(\Delta_r=0\).  Continuity of the Schur complement and its isolated
spectral projector on the stated domain follows from
Theorem~\ref{p4:thm:positive-block-elimination}.
\end{proof}

\begin{remark}[Interpretation]
\label{p4:rem:piecewise-affine-interpretation}
Proposition~\ref{p4:prop:piecewise-affine-threshold-composition} is the exact
finite statement behind later threshold-jump laws.  What later papers must
determine is which thresholds are physically distinguished and which affine
coefficients \((A_r,B_r)\) are selected by the geometry.  On a continuous
open-gap branch the defect vanishes as a theorem and the threshold is
coordinate transport only.  At a genuine gap closing or on a discontinuous
sector interface, later papers must compute or supply \(\Delta_r\); it cannot
be hidden inside the unitary frame conversion.
\end{remark}



\paragraph{Near-degenerate unstable pairs}
\label{p4:sec:part2-near-degenerate-unstable-pairs}

\begin{theorem}[Exact finite unstable-pair resonance formula]
\label{p4:thm:unstable-pair-resonance}
Let
\begin{equation}
\label{p4:eq:unstable-pair-matrix}
K=
\begin{pmatrix}
E_1-\frac{i}{2}\Gamma_1 & v\\
\overline v & E_2-\frac{i}{2}\Gamma_2
\end{pmatrix}
\end{equation}
with $E_1,E_2\in\RR$, $\Gamma_1,\Gamma_2\ge0$, and $v\in\CC$.
Set
\[
\bar E:=\frac{E_1+E_2}{2},
\qquad
\Delta_E:=E_1-E_2,
\qquad
\bar\Gamma:=\frac{\Gamma_1+\Gamma_2}{2},
\qquad
\Delta_\Gamma:=\Gamma_1-\Gamma_2.
\]
Then the two eigenvalues of $K$ are
\begin{equation}
\label{p4:eq:unstable-pair-eigs}
\lambda_\pm
=
\bar E-\frac{i}{2}\bar\Gamma
\pm
\sqrt{\left(\frac{\Delta_E-i\Delta_\Gamma/2}{2}\right)^2+|v|^2},
\end{equation}
and for every $z\notin\{\lambda_+,\lambda_-\}$ the off-diagonal resolvent element is
\begin{equation}
\label{p4:eq:unstable-pair-resolvent}
[(z-K)^{-1}]_{12}
=
\frac{v}{(z-\lambda_+)(z-\lambda_-)}.
\end{equation}
Put
\begin{equation}
\label{p4:eq:unstable-pair-complex-splitting}
\delta_{\mathbb C}:=\Delta_E-\frac{i}{2}\Delta_\Gamma .
\end{equation}
If \(\delta_{\mathbb C}\ne0\) and
\begin{equation}
\label{p4:eq:unstable-pair-perturbative-domain}
\frac{4|v|^2}{|\delta_{\mathbb C}|^2}<1,
\end{equation}
then the square root in \eqref{p4:eq:unstable-pair-eigs}, on the branch
continued from \(v=0\), has a convergent binomial expansion in
\(4|v|^2/\delta_{\mathbb C}^2\).  In particular, the eigenvalue corrections
from diagonal detuning are \(O(|v|^2/|\delta_{\mathbb C}|)\).  When
\(2|v|/|\delta_{\mathbb C}|\) is not small, that perturbative diagonal-level
description is not controlled.  This is a crossover criterion, not an
apparatus-independent maximization claim.  At a specified probe value \(z\),
the off-diagonal response is exactly the modulus of
\eqref{p4:eq:unstable-pair-resolvent}.
\end{theorem}

\begin{proof}
The characteristic polynomial of $K$ is
\[
\chi_K(\lambda)=\bigl(\lambda-E_1+\tfrac{i}{2}\Gamma_1\bigr)\bigl(\lambda-E_2+\tfrac{i}{2}\Gamma_2\bigr)-|v|^2,
\]
whose quadratic roots are exactly \eqref{p4:eq:unstable-pair-eigs}.
The inverse of the $2\times2$ matrix $z-K$ gives
\[
(z-K)^{-1}
=
\frac{1}{(z-\lambda_+)(z-\lambda_-)}
\begin{pmatrix}
 z-E_2+\frac{i}{2}\Gamma_2 & v\\
 \overline v & z-E_1+\frac{i}{2}\Gamma_1
\end{pmatrix},
\]
which yields \eqref{p4:eq:unstable-pair-resolvent}.
For \eqref{p4:eq:unstable-pair-perturbative-domain},
\[
\sqrt{\frac{\delta_{\mathbb C}^2}{4}+|v|^2}
=
\frac{\delta_{\mathbb C}}2
\sqrt{1+\frac{4|v|^2}{\delta_{\mathbb C}^2}},
\]
with the square-root branch fixed by continuation from \(v=0\).  The binomial
series converges because the modulus of its expansion parameter is strictly
less than one, and its first correction has order
\(|v|^2/|\delta_{\mathbb C}|\).  Formula
\eqref{p4:eq:unstable-pair-resolvent} itself shows that any statement about a
maximum response additionally depends on the specified \(z\); no such maximum
is asserted here.
\end{proof}



\subsection{Worked finite normal forms and textbook Hamiltonian representatives}
\label{p4:sec:part2-worked-models}

This section adds a compact model library in the corrected framework interpretation.
Exact finite-level Hamiltonians live on distinguished finite sectors of the structural carrier, with the microscopic compact momentum sector physically fixed at $m_p=1$.
By contrast, oscillator, well, and central-potential models belong to the
reciprocal position-side presentation.  Their large-\(m_x\) canonical
language applies after the exact \(\mathsf{SF}_{\Lambda_F}\) staircase or another
separately supplied compatible refinement has actually reached that regime.
Accordingly, the models below are labeled either \emph{exact finite} or \emph{large-position-side effective}.
The displayed parameters are coordinates of a supplied finite
geometric/apparatus Hamiltonian relative to a declared model subspace: matrix
elements on a sector, Hessian data near a stated minimum, Schur-complement
coefficients after a specified elimination, or spectral splittings of a
derived finite operator.  Finite-dimensional spectral theory does not by
itself put a generic compression into oscillator, well, Coulomb, Rabi, or
spin--orbit form.  A displayed textbook Hamiltonian is framework-native only
after the finite sector, apparatus channel, model subspace, and approximation
error have been specified.  Without those data it is a template and carries no
independent numerical prediction.

\begin{theorem}[Coefficient extraction in a supplied finite Hamiltonian representative]
\label{p4:thm:worked-hamiltonian-coefficients-derived}
Let \(H_{\rm geom}=H_{\rm geom}^*\) be a finite Hamiltonian supplied on a
specified framework/apparatus domain, let \(P\) be a resolved sector projector,
and put
\[
  H_P=P H_{\rm geom}P
\]
on \(P\mathcal H\).  Let \(T_1,\ldots,T_r\) be a declared linearly independent
self-adjoint family spanning the proposed model subspace and set
\[
G_{ab}:=\Tr(T_aT_b),
\qquad
c_a:=\sum_b(G^{-1})_{ab}\Tr(T_bH_P).
\]
Then
\[
H_{\rm model}:=\sum_a c_aT_a,
\qquad
R:=H_P-H_{\rm model},
\qquad
\Tr(T_aR)=0
\]
for every \(a\).  Thus \(H_{\rm model}\) is the unique Hilbert--Schmidt
projection of \(H_P\) onto the declared model subspace.  The proposed named
form is exact precisely when \(R=0\), and is an approximation only with a
stated bound on \(R\) in the norm relevant to the application.

For a rank-two sector, taking
\(\{T_a\}=\{\Id,\sigma_1,\sigma_2,\sigma_3\}\) gives the universal Pauli
expansion.  Oscillator, central-potential, tunnelling, and spin--orbit
subspaces require their additional coordinate, symmetry, kinetic, or channel
hypotheses.  If
\(\mathcal H=\mathcal H_L\oplus\mathcal H_H\) and
\(\mu\notin\Spec(H_{HH})\), the energy-dependent Schur complement
\[
H_{\rm eff}(\mu)
=H_{LL}+H_{LH}(\mu-H_{HH})^{-1}H_{HL}
\]
is an exact finite elimination identity; any subsequent named-form expansion
is again governed by the projection and residual above.
\end{theorem}

\begin{proof}
The Gram matrix is positive definite because the \(T_a\) are linearly
independent in the Hilbert--Schmidt inner product.  The displayed coefficients
solve the normal equations
\(\Tr(T_a(H_P-\sum_bc_bT_b))=0\), proving existence and uniqueness of the
orthogonal projection.  Exactness is equivalent to vanishing residual.  The
rank-two assertion is the fact that the Pauli family is a basis of the
self-adjoint part of \(M_2(\CC)\).  The Schur-complement identity follows by
solving the heavy component of
\((H-\mu)\binom{\psi_L}{\psi_H}=0\) and substituting it into the light
component.
\end{proof}



\paragraph{Exact finite two-level sectors and Rabi oscillation}
\label{p4:sec:part2-rabi}

Any isolated rank-two sector of the finite observable algebra carries the Pauli basis
\begin{equation}
\label{p4:eq:pauli-sector}
\mathcal A_{(2)}\cong M_2(\mathbb C)=\mathrm{span}\{\Id,\sigma_1,\sigma_2,\sigma_3\}.
\end{equation}
Thus a general self-adjoint two-level Hamiltonian is
\[
H(t)=h_0(t)\Id+\frac{\hbar}{2}\boldsymbol\Omega(t)\cdot\boldsymbol\sigma.
\]
After removal of the scalar phase, a supplied circular monochromatic transverse
drive with constant longitudinal splitting has the special representative
\begin{equation}
\label{p4:eq:rabi-hamiltonian}
H(t)=\frac{\hbar\Delta}{2}\sigma_3
+ \frac{\hbar\Omega}{2}\bigl(\cos(\omega t)\,\sigma_1+\sin(\omega t)\,\sigma_2\bigr),
\end{equation}
where \(\Delta,\Omega,\omega\) are angular frequencies fixed by that drive
block.  Because the carrier is finite, its time-ordered evolution is an exact
matrix evolution in \(U(2)\).

\begin{proposition}[Finite Rabi sector]
\label{p4:prop:rabi}
In the rotating frame $R(t)=\exp(-i\omega t\sigma_3/2)$, the Hamiltonian \eqref{p4:eq:rabi-hamiltonian} becomes the time-independent matrix
\begin{equation}
\label{p4:eq:rabi-effective}
H_{\mathrm{rot}}=\frac{\hbar\delta}{2}\sigma_3+\frac{\hbar\Omega}{2}\sigma_1,
\qquad
\delta:=\Delta-\omega.
\end{equation}
If the system is prepared in the lower eigenstate of $\sigma_3$, then the excited-state occupation probability is
\begin{equation}
\label{p4:eq:rabi-prob}
P_{\uparrow}(t)=\frac{\Omega^2}{\Omega_R^2}\sin^2\!\Bigl(\frac{\Omega_R t}{2}\Bigr),
\qquad
\Omega_R:=\sqrt{\delta^2+\Omega^2}.
\end{equation}
Thus the standard Rabi formula is realized exactly on a finite framework sector, with no continuum limit required.
\end{proposition}

\begin{proof}
A direct computation in the rotated state
\(\psi_{\rm rot}=R(t)^*\psi(t)\) gives
\[
R(t)^*H(t)R(t)-i\hbar R(t)^*\dot R(t)
=\frac{\hbar\delta}{2}\sigma_3+\frac{\hbar\Omega}{2}\sigma_1.
\]
The resulting constant \(2\times2\) Hamiltonian has eigenvalues
\(\pm\hbar\Omega_R/2\).
Exponentiating $H_{\mathrm{rot}}$ and transforming back yields \eqref{p4:eq:rabi-prob}.
\end{proof}

\begin{remark}[Microscopic carrier domain]
\label{p4:rem:rabi-framework}
The Pauli decomposition is exact on any supplied rank-two corner, and the
circular-drive solution is exact for the specified drive.  Neither the
physical rank-two sector nor that drive is selected merely by the existence of
an \(M_2(\CC)\) corner.  When such a sector and channel occur on the
\(m_p=1\) carrier, the calculation requires no enlargement of the microscopic
momentum truncation.
\end{remark}



\paragraph{Spin in a magnetic field and Zeeman/Larmor dynamics}
\label{p4:sec:part2-zeeman}

For a spin sector with generators $J_i$ and a specified finite magnetic/apparatus channel, the projected magnetic Hamiltonian has the normal form
\begin{equation}
\label{p4:eq:zeeman}
H_B=-\gamma\,\mathbf B\cdot \mathbf J.
\end{equation}
Here \(\gamma\mathbf B\) is the derived coefficient vector of that finite channel, not an unconstrained primitive parameter.  On the microscopic compact sector, the bounded angular-momentum theorem implies that only the finite multiplets allowed by $m_p=1$ occur physically.
For a spin-$\frac12$ sector one may write $J_i=\frac{\hbar}{2}\sigma_i$ and therefore
\begin{equation}
\label{p4:eq:zeeman-half}
H_B=-\frac{\gamma\hbar}{2}\,\mathbf B\cdot\boldsymbol\sigma.
\end{equation}

\begin{proposition}[Finite Larmor precession]
\label{p4:prop:larmor}
For constant $\mathbf B$, the Heisenberg evolution of the spin vector is
\begin{equation}
\label{p4:eq:larmor}
\frac{d\mathbf J_H}{dt}=\gamma\,\mathbf J_H\times \mathbf B.
\end{equation}
Hence the expectation value of spin precesses rigidly about $\mathbf B$ with angular frequency
\begin{equation}
\label{p4:eq:larmor-frequency}
\omega_L=|\gamma|\,|\mathbf B|.
\end{equation}
For $\mathbf B=B\hat z$, the two Zeeman levels are split by
\begin{equation}
\label{p4:eq:zeeman-splitting}
\Delta E = |\gamma|\hbar |\mathbf B|.
\end{equation}
The signed angular-velocity vector is
\(\boldsymbol\omega_L=-\gamma\mathbf B\); the sign of \(\gamma\) determines
the precession orientation and the level ordering.
\end{proposition}

\begin{proof}
\leavevmode\par\noindent
Use the commutation relations
\[
  [J_i,J_j]=i\hbar\epsilon_{ijk}J_k
\]
and the physical-units Heisenberg equation:
\[
\dot J_i
=
\frac{i}{\hbar}[H_B,J_i]
=
\gamma\epsilon_{ijk}J_jB_k.
\]
This is exactly the vector precession equation \eqref{p4:eq:larmor}.  For
\(\mathbf B=B\hat z\), the Hamiltonian is diagonal in the \(J_z\) basis,
yielding eigenvalues \(\mp\gamma\hbar B/2\) and hence the absolute splitting
\(|\gamma|\hbar|B|\).
\end{proof}



\paragraph{Effective harmonic oscillator on the reciprocal position side}
\label{p4:sec:part2-oscillator}

In a local reciprocal position-side regime, suppose a positive kinetic metric
\(G\), a positive-definite Hessian \(K\) at a specified nondegenerate local
minimum, and a canonical-pair approximation have been supplied.  Then the
normal-mode frequencies are determined by
\[
  \omega_i^2\in\operatorname{spec}(G^{-1}K).
\]
In one normal mode this gives the effective oscillator Hamiltonian
\begin{equation}
\label{p4:eq:oscillator-H}
H_{\mathrm{osc}}=\frac{P^2}{2M}+\frac12 M\omega^2 X^2.
\end{equation}
Define the ladder operators
\begin{equation}
\label{p4:eq:oscillator-a}
a:=\sqrt{\frac{M\omega}{2\hbar_{\mathrm{eff}}}}\,X + \frac{i}{\sqrt{2M\omega\hbar_{\mathrm{eff}}}}\,P,
\qquad
a^{\dagger}:=\sqrt{\frac{M\omega}{2\hbar_{\mathrm{eff}}}}\,X - \frac{i}{\sqrt{2M\omega\hbar_{\mathrm{eff}}}}\,P.
\end{equation}
For a general finite pair one has the exact identity
\begin{equation}
\label{p4:eq:oscillator-spectrum}
H_{\mathrm{osc}}
=\hbar_{\mathrm{eff}}\omega\Bigl(a^{\dagger}a+\frac12\Bigr)
-\frac{i\omega}{2}
\bigl([X,P]-i\hbar_{\mathrm{eff}}\Id\bigr).
\end{equation}
Thus a commutator-defect estimate controls this operator identity, but does not
by itself determine the spectrum of \(a^\dagger a\).

\begin{corollary}[Explicit finite oscillator ladder and comparison domain]
\label{p4:cor:truncated-oscillator}
Take the explicit truncated ladder operators of
Theorem~\ref{p4:thm:finite-heisenberg-contraction}, with \(\hbar\) replaced by
\(\hbar_{\mathrm{eff}}\) and \(\kappa=M\omega\), and define
\[
H_d:=\frac{P_d^2}{2M}+\frac12M\omega^2X_d^2
=\frac{\hbar_{\mathrm{eff}}\omega}{2}
\bigl(a_da_d^\dagger+a_d^\dagger a_d\bigr).
\]
Then
\begin{equation}
\label{p4:eq:oscillator-energies}
H_d|n\rangle
=
\hbar_{\mathrm{eff}}\omega\Bigl(n+\frac12\Bigr)|n\rangle,
\qquad 0\le n\le d-2,
\end{equation}
whereas
\[
H_d|d-1\rangle
=
\frac{\hbar_{\mathrm{eff}}\omega}{2}(d-1)|d-1\rangle .
\]
For a supplied local Hamiltonian \(H_{\rm loc}\), an approximate textbook
ladder follows only after an isometry \(W\) and an operator comparison
\[
\|W^*H_{\rm loc}W-H_d\|\le\varepsilon_{\rm loc}
\]
have been proved on the matched finite cluster.  Weyl's inequality then bounds
the corresponding eigenvalue errors by \(\varepsilon_{\rm loc}\).
\end{corollary}

\begin{proof}
Expanding \(X_d^2\) and \(P_d^2\) gives the symmetrized number-operator
formula.  For \(n<d-1\),
\(a_da_d^\dagger|n\rangle=(n+1)|n\rangle\) and
\(a_d^\dagger a_d|n\rangle=n|n\rangle\), which proves
\eqref{p4:eq:oscillator-energies}.  At the top state the first term vanishes
and the second equals \((d-1)|d-1\rangle\), giving the displayed boundary
energy.  The final assertion is the finite-dimensional Weyl eigenvalue bound.
\end{proof}


\begin{remark}[Reciprocal-position domain]
\label{p4:rem:oscillator-framework}
The harmonic oscillator is therefore not a fundamental microscopic momentum-space object.
It is an effective local model on the large reciprocal position side.
The explicit truncated model has the exact low ladder and the stated boundary
deformation.  A different local finite Hamiltonian is textbook-like only on a
proved comparison domain; compactness alone does not supply that comparison.
\end{remark}



\paragraph{Finite wells, barriers, and tunneling splittings}
\label{p4:sec:part2-well-barrier}

Let $V(X)$ be the derived effective potential on a local reciprocal position-side patch, obtained by projecting the finite geometric/material Hamiltonian to that patch, and define
\begin{equation}
\label{p4:eq:well-H}
H_{V}=\frac{P^2}{2M}+V(X)
\end{equation}
as an effective Hamiltonian on a finite resolved mode space.
Bound states are then ordinary eigenvectors of the finite matrix representing $H_V$.
For a symmetric double well, the low-energy sector is approximated by the two localized states $\ket{L},\ket{R}$ with effective Hamiltonian
\begin{equation}
\label{p4:eq:double-well}
H_{\mathrm{dw}}=
\begin{pmatrix}
E_0 & -\kappa \\
-\kappa & E_0
\end{pmatrix}
=E_0\Id-\kappa\sigma_1.
\end{equation}
Here \(E_0=\langle L|H_V|L\rangle=\langle R|H_V|R\rangle\).
Rephase \(\ket R\) so that
\(\langle L|H_V|R\rangle=-\kappa\) with
\(\kappa:=|\langle L|H_V|R\rangle|\ge0\).  In an asymmetric reduction the
diagonal entries and complex off-diagonal matrix element are read directly
from \(P H_VP\).

\begin{proposition}[Tunnel splitting as finite-mode transport]
\label{p4:prop:tunnel-splitting}
The double-well Hamiltonian \eqref{p4:eq:double-well} has eigenvalues
\begin{equation}
\label{p4:eq:tunnel-splitting}
E_{\pm}=E_0\pm \kappa,
\end{equation}
and, starting from $\ket{L}$, the probability of finding the system in $\ket{R}$ at time $t$ is
\begin{equation}
\label{p4:eq:tunnel-prob}
P_{R\leftarrow L}(t)=\sin^2\!\Bigl(\frac{\kappa t}{\hbar_{\mathrm{eff}}}\Bigr).
\end{equation}
Thus textbook barrier penetration appears as coherent finite-mode transport between nearly localized position-side sectors.
\end{proposition}

\begin{proof}
Diagonalizing \eqref{p4:eq:double-well} gives the symmetric and antisymmetric states with energies \eqref{p4:eq:tunnel-splitting}.
Evolving $\ket{L}=(\ket{+}+\ket{-})/\sqrt2$ and projecting onto $\ket{R}=(\ket{+}-\ket{-})/\sqrt2$ yields \eqref{p4:eq:tunnel-prob}.
\end{proof}



\paragraph[Central potentials, hydrogenic sectors, and spin--orbit splitting]
{Central potentials and spin--orbit splitting}
\label{p4:sec:part2-central}

In a local reciprocal position-side regime with approximate spherical symmetry, one may organize the effective Hamiltonian by partial waves,
\begin{equation}
\label{p4:eq:central-H}
H_{\mathrm{cent}}=\frac{P_r^2}{2M}+\frac{L^2}{2Mr^2}+V(r),
\end{equation}
where $L^2$ is the orbital angular-momentum operator of the large-position-side presentation.
The earlier bounded-angular-momentum analysis shows that each resolved sector still lives in a finite tower, even though the accessible orbital range may be large in an extended position regime.
\begin{proposition}[Conditional Coulomb spectral approximation]
\label{p4:prop:conditional-hydrogenic-spectrum}
Let
\[
H_C=-\frac{\hbar_{\mathrm{eff}}^2}{2M}\Delta-\frac{\kappa}{r}
\]
be the specified Coulomb Schr\"odinger operator with closed semibounded
quadratic form \(q_C\).  Suppose the reciprocal-position interface supplies
nested finite-dimensional form subspaces
\(\mathcal V_m\subset\operatorname{Dom}(q_C)\) whose union is dense in the
form norm, and let \(H_m\) be the corresponding Rayleigh--Ritz operators.
Then, for every fixed discrete Coulomb shell below the essential spectrum, the
finite eigenvalue cluster converges with multiplicity to
\begin{equation}
\label{p4:eq:hydrogenic}
E_n=-\frac{M\kappa^2}{2\hbar_{\mathrm{eff}}^2 n^2}.
\end{equation}
At every finite \(m\) only finitely many such clusters are represented.  A
quantitative error \(O(\varepsilon_m)\) requires a separately proved
form-approximation or Ritz estimate; neither its rate nor a finite
\(n_{\max}(m)\) follows from GP1--GP3 alone.
\end{proposition}

\begin{proof}
The min--max principle for semibounded self-adjoint operators and density of
\(\bigcup_m\mathcal V_m\) in the form norm imply convergence of every fixed
Ritz eigenvalue below the essential spectrum, including multiplicity.  The
target Coulomb eigenvalues are \eqref{p4:eq:hydrogenic}.  The theorem gives no
uniform claim over all \(n\) and no convergence rate without additional
approximation estimates.
\end{proof}

If an internal spin sector $S$ is present, then the projected effective spin--orbit term
\begin{equation}
\label{p4:eq:spin-orbit}
H_{\mathrm{so}}=\lambda\,\mathbf L\cdot\mathbf S
\end{equation}
is diagonal in the coupled basis of total angular momentum $\mathbf J=\mathbf L+\mathbf S$.
Its eigenvalue on a fixed $(\ell,s,j)$ channel is
\begin{equation}
\label{p4:eq:spin-orbit-eigs}
E_{\mathrm{so}}(j,\ell,s)=\frac{\lambda\hbar^2}{2}\Bigl(j(j+1)-\ell(\ell+1)-s(s+1)\Bigr),
\end{equation}
with the allowed values of \(j\) determined by the bounded addition rules of
Section~\ref{p4:sec:part2-addition}; here \(\lambda\) has units of
energy/action squared.  To extract \(\lambda\), choose a linearly independent
family of rotational scalars \(T_a\), including
\(T_{\rm so}=\mathbf L\cdot\mathbf S\), on a sector containing enough coupled
channels to distinguish them.  With
\[
G_{ab}=\Tr(T_aT_b),
\qquad
c_a=\sum_b(G^{-1})_{ab}\Tr(T_bPH_{\rm eff}P),
\]
the spin--orbit coefficient is the corresponding \(c_{\rm so}\).  A single
trace quotient is valid only after all other scalar components have been
subtracted or proved Hilbert--Schmidt orthogonal.  On one fixed \(j\)-channel,
\(\mathbf L\cdot\mathbf S\) is scalar and cannot be separated from an identity
energy shift.  Thus the coupled-spin algebra supplies the splitting pattern,
while its realization and coefficient remain conditional on the specified
reduced Hamiltonian and scalar decomposition.



\section{Interferometric protocols, control, monitoring and code calculations}
\label{p4c:sec:v238-c6}
\noindent\textit{Source and claim boundary.} These are original extended calculations. Statements on physical sources, measurements or continuum limits retain the inherited antecedents of Papers IV-A/IV-B.

\subsection{Finite protocol representations}
\label{p4:sec:part2-protocols}

The following finite examples use the envelope algebra $\mathcal A_m^{\mathrm{env}}$, states $\omega(a)=\Tr(\rho a)$, specified accessible subalgebras, and the POVM, instrument, and CPTP-map calculus of Part~I.



\paragraph{Interference, coherence, and superposition}
\label{p4:sec:part2-interference}

Relative to a supplied apparatus/interface family of mutually orthogonal projections $\{P_i\}$ with $\sum_iP_i=\mathbf 1$, the state is incoherent exactly when
\[
\rho=\sum_i P_i\rho P_i,
\]
that is, when the off-diagonal blocks $P_i\rho P_j$ for $i\neq j$ vanish.
Interference is the contribution of those off-diagonal coherences when alternatives recombine before final readout.



\paragraph{Double-slit as a two-path algebraic calculation}
\label{p4:sec:part2-double-slit}

Let $P_1,P_2\in\mathcal A_m^{\mathrm{env}}$ represent the two exclusive path alternatives at the aperture, and let $E_x$ be the POVM element for a detector pixel near $x$.
If the preparation is $\rho$ and only transmitted events are kept, then the post-selected state after the aperture is
\[
\rho'=
\frac{(P_1+P_2)\rho(P_1+P_2)}{\Tr\!\bigl(\rho(P_1+P_2)\bigr)}.
\]
Put
\(p_{\rm tr}:=\Tr(\rho(P_1+P_2))>0\).
The detection probability is therefore
\begin{align}
\label{p4:eq:part2-double-slit}
p(x)
&=\Tr(\rho' E_x) \\
&=\frac{1}{p_{\rm tr}}\Bigl(
\Tr(\rho P_1E_xP_1)+\Tr(\rho P_2E_xP_2)
+\Tr(\rho P_1E_xP_2)+\Tr(\rho P_2E_xP_1)
\Bigr).
\end{align}
The last two terms are the interference terms for the specified detector
effect \(E_x\). Dephasing the two path blocks is sufficient to remove
these terms for \emph{every} detector effect, but their disappearance for
one effect is not proof of dephasing: a path-diagonal \(E_x\) also makes
them vanish in a coherent state, as can a relative-phase cancellation.
Absence of the off-diagonal contribution for every possible detector effect
is equivalent to vanishing of the transmitted off-diagonal state blocks.
Which-path information may be represented by coupling to an ancilla and
tracing out the marker sector; its effect depends on the marker overlap.



\paragraph{Mach--Zehnder interferometry and relative projective phase}
\label{p4:sec:part2-mzi}

In a two-path interferometer with arm basis \(\ket{0},\ket{1}\), beam
splitters are unitaries that mix the two-dimensional sector and arm
propagation contributes a relative phase \(e^{i\varphi}\).  The output
probabilities depend on \(\varphi\) through recombination of the two
amplitudes.  Explicitly, with
\[
 B=\frac1{\sqrt2}
 \begin{pmatrix}1&1\\1&-1\end{pmatrix},
 \qquad
 D_\varphi=\begin{pmatrix}1&0\\0&e^{i\varphi}\end{pmatrix},
\]
one obtains
\[
 BD_\varphi B\ket0
 =
 e^{i\varphi/2}
 \left(
 \cos\frac\varphi2\,\ket0
 -i\sin\frac\varphi2\,\ket1
 \right).
\]
Thus the two sharp output probabilities are exactly
\[
 p_0=\cos^2\frac\varphi2,
 \qquad
 p_1=\sin^2\frac\varphi2.
\]
The prefactor is a global phase and cancels from both probabilities, whereas
the relative arm phase controls them.  This relative phase changes the projective state
\([\ket0+e^{i\varphi}\ket1]\); it is not the central global phase removed in
forming a ray.  It is generated by the supplied arm Hamiltonian or connection.
If an apparatus theorem identifies that transport with the degree-one
associated Hopf line, \(\chi_1\) supplies the comparison, but the
interferometer calculation itself needs no identification of the whole
carrier Hopf action with scalar phase.



\paragraph{Aharonov--Bohm phase as carrier holonomy}
\label{p4:sec:part2-aharonov-bohm}

The Aharonov--Bohm setup is the force-free interferometer version of the same phase mechanism \cite{AharonovBohm1959,WuYang1975}.
Along the accessible electron arms the local electromagnetic curvature may vanish, so a Lorentz-force account on either arm has no local field with which to produce the fringe shift.
What remains physical is the gauge-invariant closed-loop phase
\begin{equation}
\label{p4:eq:ab-holonomy}
\exp\!\left(\frac{iq}{\hbar}\oint_\gamma A\right)
=
\exp\!\left(\frac{iq}{\hbar}\int_\Sigma F\right),
\end{equation}
where \(\gamma\) is the loop formed by the two arms and \(\Sigma\) is read in the usual patched or relative sense when the enclosed flux is shielded from the electron paths.
When one gauge trivialization covers \(\Sigma\), the equality is ordinary
Stokes' theorem.  In the patched case, the right-hand side denotes the same
line-bundle holonomy computed from curvature integrals and transition
functions; it is not a claim that one global potential exists on \(\Sigma\).
Under a single-valued gauge change \(A\mapsto A+d\chi\),
\[
 \exp\!\left(\frac{iq}{\hbar}\oint_\gamma(A+d\chi)\right)
 =
 \exp\!\left(\frac{iq}{\hbar}\oint_\gamma A\right),
\]
and a patched large gauge change modifies the exponent only by
\(2\pi i\mathbb Z\).  This verifies gauge invariance of the displayed phase.
The phase in \eqref{p4:eq:ab-holonomy} is the holonomy of the supplied
electromagnetic line-bundle connection with charge representation \(q\).
Recombination reads that gauge-invariant relative holonomy through Born
probabilities.  The geometric Hopf circle is not automatically the
electromagnetic gauge bundle; relating them requires an explicit bundle map
or associated-character comparison.  Local curvature values on the two arms
therefore do not determine the interference phase by themselves.



\paragraph{Quantum eraser and delayed-choice conditioning}
\label{p4:sec:part2-eraser}

Suppose the path sector is entangled with a marker system $M$ so that
\[
\ket{\Psi}_{PM}=\tfrac{1}{\sqrt2}\bigl(\ket{0}\ket{m_0}+\ket{1}\ket{m_1}\bigr),
\qquad
\braket{m_0}{m_1}=0.
\]
Then
\[
 \rho_P=\Tr_M\ket\Psi\!\bra\Psi
 =\frac12\bigl(\ket0\!\bra0+\ket1\!\bra1\bigr),
\]
so the reduced path state is dephased and no unconditional interference
remains.  Conversely, measure the marker in the complementary basis
\[
 \ket{m_\pm}:=\frac{\ket{m_0}\pm\ket{m_1}}{\sqrt2}.
\]
Each result has probability \(1/2\), and the corresponding normalized path
state is
\[
 \ket{\psi_\pm}=\frac{\ket0\pm\ket1}{\sqrt2}.
\]
Thus conditioning on either complementary marker outcome restores a coherent
fringe or antifringe subensemble, while their equal mixture is the unchanged
dephased state \(\rho_P\).
This is conditioning on a joint state by an instrument of the type described in Section~\ref{p4:sec:part2-conditioning}; the conditional statistics do not modify the earlier reduced state.



\paragraph{Stern--Gerlach and spin from fuzzy \texorpdfstring{$\CP^1\simeq S^2$}{CP1=S2}}
\label{p4:sec:part2-sg}

Spin measurements are the simplest manifestation of noncommutativity in a distinguished two-level block.
Such blocks arise from the fuzzy $\CP^1\simeq S^2$ sector discussed in Section~\ref{p4:sec:part2-fuzzycp1}.
For a measurement along unit vector $\hat n$, the sharp outcomes are represented by
\[
P_{\pm}(\hat n)=\frac{\mathbb I\pm \hat n\cdot\vec\sigma}{2},
\]
and the probabilities are $p(\pm\mid \hat n)=\Tr(\rho P_{\pm}(\hat n))$.
The Pauli identity
\((\hat n\cdot\vec\sigma)^2=\mathbb I\) gives
\(P_\pm^2=P_\pm\), \(P_+P_-=0\), and \(P_++P_-=\mathbb I\).  For two
directions,
\[
 [P_+(\hat n),P_+(\hat m)]
 =\frac{i}{2}(\hat n\times\hat m)\cdot\vec\sigma,
\]
so the sharp tests fail to commute exactly when the directions are neither
parallel nor antiparallel.
Sequential Stern--Gerlach experiments probe the fact that different spin directions correspond to noncommuting projectors. They are the lowest-spin instance of the bounded angular-momentum and vector-selection structure derived in Sections~\ref{p4:sec:part2-angular} and~\ref{p4:sec:part2-selection}.



\paragraph{Entanglement and Bell tests}
\label{p4:sec:part2-bell}

Bell experiments combine the tensor-product structure of Section~\ref{p4:sec:part1-entanglement} with incompatible local measurement settings \cite{Bell1964,ClauserHorneShimonyHolt1969}.
For the singlet state
\[
|\Psi^-\rangle=\frac1{\sqrt2}(|01\rangle-|10\rangle),
\]
one obtains the correlation
\[
E(\hat a,\hat b)=\Tr\!\bigl(\rho_{\Psi^-}(\hat a\cdot\vec\sigma)\otimes(\hat b\cdot\vec\sigma)\bigr)=-\hat a\cdot\hat b,
\]
because
\(\langle\Psi^-|\sigma_i\otimes\sigma_j|\Psi^-\rangle=-\delta_{ij}\).
For the explicit settings
\[
 \hat a=\hat z,
 \qquad
 \hat a'=\hat x,
 \qquad
 \hat b=\frac{\hat z+\hat x}{\sqrt2},
 \qquad
 \hat b'=\frac{\hat z-\hat x}{\sqrt2},
\]
the CHSH combination
\[
 S:=E(\hat a,\hat b)+E(\hat a,\hat b')
   +E(\hat a',\hat b)-E(\hat a',\hat b')
\]
satisfies \(S=-2\sqrt2\), hence \(|S|=2\sqrt2>2\).
Non-signaling nevertheless remains automatic because the reduced state is independent of the remote measurement choice, as shown in Section~\ref{p4:sec:part1-nonsignaling}.



\paragraph{Accessible-algebra description of an observer}
\label{p4:sec:part2-observer}

An operational observer is specified by an accessible algebra.
The global description is a state on a global algebra $\mathcal A_{\mathrm{tot}}$, while an ``observer'' corresponds to an operational algebra $\mathcal A_{\mathrm{obs}}\subset \mathcal A_{\mathrm{tot}}$ determined by finite resolution and available controls.
The observer's effective state is the restriction or conditional reduction of the global state to that accessible algebra.
A definite measurement outcome is then a conditioning step on an event inside $\mathcal A_{\mathrm{obs}}$.
Conditioning changes the state restricted to the accessible algebra while leaving the stipulated closed global dynamics unitary.



\paragraph{Uncertainty and complementarity in operational form}
\label{p4:sec:part2-uncertainty-op}

Part~I already derived the uncertainty relation and the no-cloning obstruction from the noncommutative algebra itself.
Operationally, this means that incompatible devices correspond to noncommuting effects or observables, so no single sharp classical description can refine all experimental contexts simultaneously.
Finite truncation adds a second limitation: within a fixed plateau, there is only finite spectral resolution available for representing conjugate quantities.



\paragraph{Teleportation as an algebraic protocol}
\label{p4:sec:part2-teleportation}

Teleportation is a multipartite protocol built entirely from the structural ingredients of Part~I \cite{Bennett1993}:
\begin{enumerate}[leftmargin=2em]
\item Alice and Bob share an entangled resource state on $\mathcal A_A\otimes\mathcal A_B$.
\item Alice performs a Bell-type measurement on the unknown input and her share, obtaining a classical outcome label $i$.
\item The outcome $i$ is communicated classically.
\item Bob applies the corresponding unitary correction $U_i$.
\end{enumerate}
For completeness, let the input be
\(\ket\psi_C=\alpha\ket0+\beta\ket1\) and the resource be
\(\ket{\Phi^+}_{AB}=(\ket{00}+\ket{11})/\sqrt2\).  Direct expansion in the
Bell basis of \(C\otimes A\), with
\[
 \ket{\Phi^\pm}:=\frac{\ket{00}\pm\ket{11}}{\sqrt2},
 \qquad
 \ket{\Psi^\pm}:=\frac{\ket{01}\pm\ket{10}}{\sqrt2},
\]
gives the exact identity
\begin{align}
\label{p4:eq:teleportation-bell-identity}
 \ket\psi_C\ket{\Phi^+}_{AB}
 =\frac12\bigl(&
 \ket{\Phi^+}_{CA}\ket\psi_B
 +\ket{\Phi^-}_{CA}Z\ket\psi_B\nonumber\\
 &+\ket{\Psi^+}_{CA}X\ket\psi_B
 +\ket{\Psi^-}_{CA}XZ\ket\psi_B
 \bigr).
\end{align}
The four outcomes therefore have probability \(1/4\).  Conditional on the
announced outcome, Bob applies respectively
\(\mathbb I,Z,X,(XZ)^*\); in every branch his normalized state is
\(\ket\psi\) (an irrelevant overall phase may occur in the last branch).
Consequently the unconditional corrected reduced state on Bob's side is
also \(\ket\psi\!\bra\psi\), while before receipt of the classical label it
is \(\mathbb I/2\).
Teleportation is therefore a composition of entanglement, local measurement, classical communication, and conditional unitary control.



\paragraph{Berry phase from Hopf holonomy}
\label{p4:sec:part2-berry}

For a supplied \(C^1\) family of nondegenerate eigenlines
\(\mathbf R\mapsto[\psi(\mathbf R)]\) and an adiabatic closed loop in
parameter space, the geometric phase is the holonomy of the pulled-back
eigenline connection \cite{Berry1984},
\begin{equation}
\label{p4:eq:berry}
\gamma=i\oint \langle\psi(\mathbf R)\mid d\psi(\mathbf R)\rangle .
\end{equation}
For a normalized local eigenvector and a gauge change
\(\ket\psi\mapsto e^{i\chi}\ket\psi\),
\[
 i\langle e^{i\chi}\psi|d(e^{i\chi}\psi)\rangle
 =i\langle\psi|d\psi\rangle-d\chi.
\]
Hence a closed-loop representative changes \(\gamma\) by
\(-2\pi n\), \(n\in\mathbb Z\), and the observable holonomy
\(e^{i\gamma}\) is gauge invariant.  This calculation also shows why a
choice of local eigenvector changes the connection one-form but not its
closed-loop holonomy.
The sign depends on the Berry-connection convention and line orientation;
the Chern class fixes only the integral topological class, not the holonomy
of an arbitrary loop.  For a two-level ray space
\(\CP^1\), the universal eigenline is the Hopf line, so the framework's
degree-one associated character supplies the standard geometric model.
The actual Berry phase still requires the supplied eigenline map, adiabatic
transport, and loop; it is their pullback holonomy, not a consequence of
first Chern class one alone.



\paragraph{Tunneling as finite-mode transport}
\label{p4:sec:part2-tunneling}

Barrier penetration is described by ordinary unitary transport generated by the supplied effective finite generator together with the supplied detection algebra.
In the kernel and path-sum presentations of
Section~\ref{p4:sec:part2-wave-path}, tunneling is represented by a
transition amplitude between coarse regions that remains nonzero even when
the corresponding classical path would be forbidden.
The prediction follows from unitary evolution and Born evaluation.




\subsection{Entropy on declared coarse-graining domains}
\label{p4:sec:part3-entropy}

The von Neumann entropy
\begin{equation}
\label{p4:eq:entropy}
S(\rho)=-\Tr(\rho\ln\rho)
\end{equation}
is constant under closed unitary flow.  A partial trace or a generic
nonunital reduced channel need not increase it monotonically.

\begin{proposition}[Entropy monotonicity on declared unital coarse-graining domains]
\label{p4:prop:unital-coarse-entropy}
If \(\Phi\) is a unital CPTP map on a fixed finite matrix algebra, then
\[
S(\Phi(\rho))\ge S(\rho).
\]
If \(E_r\) is a nested family of trace-preserving unital conditional
expectations satisfying
\[
E_{r+1}E_r=E_{r+1},
\]
then
\[
S(E_{r+1}\rho)\ge S(E_r\rho).
\]
Likewise, for a unital CPTP semigroup \((\Phi_t)_{t\ge0}\),
\[
t\ge s\quad\Longrightarrow\quad
S(\Phi_t\rho)\ge S(\Phi_s\rho).
\]
The first two inequalities concern a supplied channel or coarse-graining
order.  A thermodynamic time arrow additionally requires an oriented
irreversible effective evolution and a suitable initial-state condition; it
does not follow from coarse description alone.
\end{proposition}

\begin{proof}
Because \(\Phi(I/d)=I/d\), monotonicity of relative entropy gives
\[
D(\Phi(\rho)\|I/d)
\le D(\rho\|I/d).
\]
Using \(D(\rho\|I/d)=\log d-S(\rho)\) proves the first claim.  The nesting
identity gives \(E_{r+1}\rho=E_{r+1}(E_r\rho)\), so the same result applied to
\(E_{r+1}\) proves the second.  For a semigroup,
\(\Phi_t=\Phi_{t-s}\Phi_s\), and applying the first claim to
\(\Phi_{t-s}\) proves the third.  Nonunital channels and finite subsystem
entropies can decrease, while exact finite unitary dynamics is recurrent;
hence none of these algebraic inequalities supplies the additional arrow
data.
\end{proof}



\subsection{Landauer principle}
\label{p4:sec:part3-landauer}

\begin{theorem}[Exact correlated finite Landauer balance]
\label{p4:thm:v2211-correlated-landauer}
Let \(M,B\) be the actual finite memory and bath of a licensed joint
unitary implementation, \(\rho'_{MB}=U\rho_{MB}U^\dagger\).  The initial
state may be correlated, but its bath marginal is the faithful thermal state
\[
 \rho_B=\tau_B:=Z_B^{-1}e^{-\beta H_B},\qquad
 \beta=(k_BT)^{-1}>0,
\]
for the fixed reference bath Hamiltonian \(H_B\).  Put
\[
 Q:=\operatorname{Tr}H_B(\rho'_B-\tau_B),\quad
 \Delta S_M:=S(\rho_M)-S(\rho'_M),\quad
 \Delta I:=I(M:B)_{\rho'}-I(M:B)_\rho,
\]
where \(S=-\operatorname{Tr}\rho\log\rho\) is dimensionless.  Then
\begin{equation}
 \boxed{\beta Q=\Delta S_M+\Delta I+
              D(\rho'_B\|\tau_B).}
 \label{p4:eq:v2211-landauer-correlation-balance}
\end{equation}
All other retained work systems, memories or correlated resources must be
included in the actual joint partition before using this identity.
\end{theorem}
\begin{proof}
Unitarity preserves \(S(\rho_{MB})\).  From
\(I(M:B)=S(\rho_M)+S(\rho_B)-S(\rho_{MB})\) it follows that
\[
 S(\rho'_B)-S(\tau_B)=\Delta S_M+\Delta I.
\]
The Gibbs logarithm gives the exact relative-entropy identity
\[
 D(\rho'_B\|\tau_B)
 =\beta\operatorname{Tr}H_B(\rho'_B-\tau_B)
       -[S(\rho'_B)-S(\tau_B)].
\]
Combining the two equalities proves the claim.  Faithfulness of \(\tau_B\)
makes the relative entropy finite; no Markov or weak-coupling limit is used.
\end{proof}

\begin{corollary}[Product-preparation erasure bound and finite-bath scope]
\label{p4:cor:v2211-product-landauer}
On the preceding domain, if the actual preparation is
\(\rho_{MB}=\rho_M\otimes\tau_B\), then
\[
 \beta Q=\Delta S_M+I(M:B)_{\rho'}+D(\rho'_B\|\tau_B)
 \ge\Delta S_M.
\]
For a maximally uncertain classical bit, an implementation whose final
memory entropy is at most \(\varepsilon\) therefore obeys
\begin{equation}
\label{p4:eq:landauer}
 \boxed{Q\ge k_BT(\ln2-\varepsilon).}
\end{equation}
The familiar \(k_BT\ln2\) value is its zero-error bound.  A mixed memory
and a faithful fixed finite bath cannot attain an exact pure-memory reset
by a joint unitary: exact reset in this product branch requires a further
limit or a different actual resource/domain.
\end{corollary}
\begin{proof}
The initial product has zero mutual information.  Nonnegativity of the
final mutual information and relative entropy proves the inequality
\cite{Landauer1961}.  For the final assertion, if
\(\operatorname{rank}\rho_M=r>1\) and \(\dim B=d_B\), the initial joint
rank is \(rd_B\).  A pure final memory forces the joint state to factor
as that pure state times a bath density and hence to have rank at most
\(d_B\).  Unitary rank preservation excludes such a reset at finite \(d_B\).
The heat inequality remains meaningful for finite-error operations.
\end{proof}

\begin{proposition}[Thermal marginal does not exclude a correlation resource]
\label{p4:prop:v2211-correlated-erasure-witness}
An initially thermal bath marginal, a cyclic memory Hamiltonian and a
unitary total implementation alone do not imply \(Q\ge k_BT\ln2\).
\end{proposition}
\begin{proof}
For two qubits ordered as \(|M,B\rangle\), take
\(H_M=H_B=0\) before and after the operation and
\[
 \rho_{MB}=\tfrac12|00\rangle\langle00|
                  +\tfrac12|11\rangle\langle11|.
\]
Both marginals equal \(I/2\), and the bath is thermal at every \(T>0\).
The bath-controlled NOT, which exchanges \(|01\rangle,|11\rangle\),
gives \(\rho'_{MB}=|0\rangle\langle0|\otimes I/2\).  Thus
\(\Delta S_M=\ln2,\ Q=0,\ \Delta I=-\ln2\), and the relative entropy in
the balance identity vanishes.  The Hermitian involution \(U\) is realized
for unit clock duration by \(H=\pi\hbar(I-U)/2\).  Thermal spectators may
be adjoined without changing the calculation.  This is erasure on the
stated correlated compatibility domain, not a reset map for every
independently prepared input.  It consumes correlation already present
in the complete initial owner.
\end{proof}

The product preparation above is the same assignment used by
Theorem~\ref{p4:thm:exact-finite-system-environment-reduction}; it is not
forced by the finite state space.  The heat statement is an interface
consequence of the actual preparation, resource partition and unitary
implementation, rather than a consequence of complete positivity or GP1
alone.



\subsection{Adiabatic behavior and speed limits}
\label{p4:sec:part3-adiabatic}

If the closed generator varies slowly compared with the relevant spectral gaps, the usual finite-dimensional adiabatic approximation follows.
Likewise, for a closed pure state evolving to an orthogonal state under a time-independent Hamiltonian, the Mandelstam--Tamm type speed limit
\begin{equation}
\label{p4:eq:speed-limit}
t\ge \frac{\pi\hbar}{2\Delta E}
\end{equation}
is a direct consequence of unitary generation and energy variance. Mixed-state and nonorthogonal versions use the corresponding Bures-angle form.
Adiabatic persistence and quantum speed limits therefore follow from the Part~I dynamics on their stated theorem domains.

\begin{corollary}[Conditional adiabatic realization of the midpoint connection]
\label{p4:cor:clifford-midpoint-adiabatic}
Let \(n:[0,1]\to S^3\) be a supplied \(C^2\) path on the protected
Clifford coefficient sphere, and suppose the actual finite Hamiltonian on
that block is
\[
 H(s)=E_+(s)P_+(n(s))+E_-(s)P_-(n(s)),\qquad
 P_\pm(n)=\tfrac12(I\pm\Gamma(n)),
\]
with \(E_\pm\) of class \(C^2\) and
\(\inf_s|E_+(s)-E_-(s)|>0\).  For physical duration \(T\), the
Schr\"odinger evolution driven by \(H(t/T)\) preserves each instantaneous
rank-two band up to \(O(T^{-1})\), with a constant controlled by the
actual derivatives and gap.  After removing its dynamical phase, either
band transports by the pullback of
\[
 \boxed{A=\tfrac12U^*dU,\qquad
 F=-\tfrac14(U^*dU)\wedge(U^*dU),\qquad
 U(n)=n_4I+i n_i\sigma_i.}
\]
Thus the Paper-I-B midpoint connection is the geometric transport of these
particular gapped protected eigenspaces when this Hamiltonian/history
domain is supplied.  The positive and negative bands carry equivalent
connections; choosing the complementary coordinate-sign projector alone
does not select a helicity imbalance.
\end{corollary}

\begin{proof}
Paper~II-C, Theorem~\ref{p2c:thm:protected-projector-midpoint}, gives
global orthonormal frames
\[
 V_\pm(n)=\frac1{\sqrt2}
 \begin{pmatrix}\pm U(n)\\I\end{pmatrix},\qquad
 V_\pm^*dV_\pm=\tfrac12U^*dU.
\]
The finite gapped adiabatic estimate applies to the stated \(C^2\)
Hamiltonian on its compact parameter interval.  In either frame, removal
of the scalar dynamical phase leaves
\(\dot W_\pm=-V_\pm^*\dot V_\pm W_\pm\); the two geometric
transport equations coincide.  The Maurer--Cartan identity gives the
displayed curvature.  This specialization constructs neither an actual
history nor an interaction Hamiltonian from the kinematic projectors alone.
\end{proof}




\subsection[Continuous monitoring, decoherence, and controlled inference]
{Monitoring, decoherence, and inference}
\label{p4:sec:part3-monitoring}

In the remaining finite-algebraic evolution formulas we set $\hbar=1$. Restoring physical units gives $U(t)=e^{-itH/\hbar}$ and $\dot\rho=-(i/\hbar)[H,\rho]+\cdots$.

\paragraph{Role in the framework.}
Part~I already supplied instruments, CP maps, and GKLS generators.
Weak monitoring, trajectory language, environment-induced decoherence, and finite estimation/control protocols arise by applying the finite channel calculus to repeated weak probes, reduced environments, and informationally complete measurements.



\paragraph{Feedback-conditioned evolution}
\label{p4:sec:part3-feedback}

On a control/interface theorem domain, if the observed current is fed back through a bounded self-adjoint control term,
\begin{equation}
\label{p4:eq:feedback-H}
H_{\mathrm{tot}}(t)=H_0+\sum_k u_k\!\bigl(Y_{[0,t]}\bigr)H_k,
\end{equation}
then the conditional state still evolves inside the same finite algebra by replacing $H$ in \eqref{p4:eq:sme} with $H_{\mathrm{tot}}(t)$.
So measurement-based feedback introduces no new kinematics: it is a theorem-domain choice of finite generator driven by an instrument record, not a constitutive Hamiltonian insertion in the exact-GP2 sense.



\subsection[Environment-induced decoherence and pointer sectors]
{Decoherence and pointer sectors}
\label{p4:sec:part3-decoherence}

\paragraph{Environment models.}
On an open-system interface domain, suppose a finite system--environment factorization and coupling have been supplied inside the already constituted finite algebra.  Exact GP2 forbids treating the split, bath, coupling observables, pointer projectors, or Hamiltonian generator as primitive sector selectors.  They are theorem-domain data for the reduced dynamics.  Write
\begin{equation}
\label{p4:eq:sys-env-H}
H = H_S\otimes \mathbf 1 + \mathbf 1\otimes H_E + \sum_{\alpha} A_{\alpha}\otimes B_{\alpha}.
\end{equation}
Before any Markov approximation, tracing out the finite environment gives an
exact CPTP reduced map.  A GKLS generator appears precisely when the reduced
maps form a norm-continuous Markov semigroup, or when a controlled
weak-coupling/large-bath interface theorem produces such a semigroup.  A
particularly transparent case is pure dephasing,
\begin{equation}
\label{p4:eq:pure-dephasing}
\dot\rho = -i[H_S,\rho]-\frac12\sum_{\alpha}\gamma_{\alpha}[A_{\alpha},[A_{\alpha},\rho]],
\end{equation}
with $\gamma_{\alpha}\ge0$.


\begin{theorem}[Exact finite system--environment reduction]
\label{p4:thm:exact-finite-system-environment-reduction}
Let \(\mathcal H_S\otimes\mathcal H_E\) be a finite supplied
system--environment factorization inside the already constituted algebra,
let \(U_t=e^{-itH_{\rm tot}}\) be the exact finite unitary generated by a
self-adjoint total Hamiltonian, and let \(\sigma_E\) be an environment
state.  Then
\[
\mathcal E_t(\rho_S)
:=
\Tr_E\!\bigl(U_t(\rho_S\otimes\sigma_E)U_t^*\bigr)
\]
is CPTP for every finite \(t\).  If \(\{\mathcal E_t\}_{t\ge0}\) is, on a
specified theorem domain, a norm-continuous semigroup, then its generator
has the GKLS form of Theorem~\ref{p4:thm:gkls}.  Thus exact finite
reduction gives CPTP maps unconditionally, while irreversibility and
Markovian GKLS evolution require the semigroup/interface hypotheses.
\end{theorem}

\begin{proof}
Choose a spectral decomposition
\(\sigma_E=\sum_a q_a\ket{a}\bra{a}\) and an orthonormal basis
\(\{f_b\}\) of \(\mathcal H_E\).  Define finite operators on the system by
\[
K_{ba}(t):=\sqrt{q_a}\,\bra{f_b}U_t\ket{a}.
\]
Then
\[
\mathcal E_t(\rho_S)=\sum_{a,b}K_{ba}(t)\rho_S K_{ba}(t)^*,
\]
so the map is completely positive.  Unitarity gives
\[
\sum_{a,b}K_{ba}(t)^*K_{ba}(t)
=
\sum_a q_a\,\bra{a}U_t^*U_t\ket{a}
=
\mathbf1_S,
\]
hence the map is trace preserving.  If the reduced maps additionally form
a norm-continuous semigroup, Theorem~\ref{p4:thm:gkls} applies directly.
The final distinction follows because finite unitaries generally have
recurrences, whereas a Markov semigroup is an extra theorem-domain regime.
\end{proof}

\begin{theorem}[Fixed finite-bath exact semigroup rigidity]
\label{p4:thm:fixed-finite-bath-semigroup-rigidity}
Under the hypotheses of
Theorem~\ref{p4:thm:exact-finite-system-environment-reduction}, suppose the
same fixed finite bath, fixed state \(\sigma_E\), and time-independent
\(H_{\rm tot}\) produce reduced maps satisfying
\[
 \mathcal E_{s+t}=\mathcal E_s\mathcal E_t
 \qquad(s,t\ge0)
\]
exactly for all times.  Then every \(\mathcal E_t\) has a CPTP inverse and is
a reversible channel.  On \(\End(\mathcal H_S)\) there is a unitary \(W_t\)
such that
\[
 \boxed{\mathcal E_t(\rho)=W_t\rho W_t^*.}
\]
Hence one fixed finite bath cannot produce a nonzero irreversible GKLS
dissipator as an exact all-time semigroup.
\end{theorem}

\begin{proof}
In fixed eigenbases of the finite total Hamiltonian, every matrix entry of
\(\mathcal E_t\) is a finite sum of exponentials
\(e^{it(\lambda_j-\lambda_k)}\), and therefore extends to an entire function
of complex \(t\).  For each fixed \(s>0\), the semigroup identity, initially
valid for \(t\ge0\), extends by the identity theorem to all complex \(t\).
Fixing such a real \(t\) and repeating in \(s\) extends the identity to all
real \(s,t\).  In particular,
\[
 \mathcal E_t\mathcal E_{-t}=\mathcal E_0
 =\mathcal E_{-t}\mathcal E_t.
\]
Both maps are CPTP because the exact finite reduction is defined for positive
and negative times.

Let \(T=\mathcal E_t^*\) be the unital CP Heisenberg dual and
\(S=\mathcal E_{-t}^*=T^{-1}\).  Kadison--Schwarz for \(T\), positivity of
\(S\), and Kadison--Schwarz for \(S\) give
\[
 x^*x
 =S(T(x^*x))
 \ge S(T(x)^*T(x))
 \ge S(T(x))^*S(T(x))
 =x^*x.
\]
Both inequalities are therefore equalities.  Since \(S=T^{-1}\) is
injective, the first equality gives
\(T(x^*x)=T(x)^*T(x)\); applying the same argument to \(xx^*\) gives the
second multiplicative-domain equality.  Every element consequently lies in
the multiplicative domain of \(T\), so \(T\) is a
\(*\)-automorphism.  Every automorphism of a full matrix algebra is inner;
duality gives the stated unitary channel.  Such a channel has no dissipative
GKLS part.
\end{proof}


\begin{theorem}[Reservoir or repeated-probe weak-coupling bridge to GKLS]
\label{p4:thm:microscopic-bath-gkls}
Let the system Hilbert space be finite dimensional.  Let the environment
interface be one of the following:
\begin{enumerate}[label=(\roman*), leftmargin=2em]
\item a stationary mixing reservoir satisfying the explicit Davies
correlation and cluster conditions below;
\item a repeated-interaction chain of freshly prepared finite probes whose
one-collision CPTP maps have a uniform continuous-time expansion; or
\item a family of finite baths \(\mathcal H_{E,n}\) in a simultaneous
large-bath/weak-coupling limit satisfying uniform versions of those
correlation and cluster conditions, for which the recurrence times satisfy
\(\tau_{{\rm rec},n}/\lambda_n^{-2}\to\infty\).
\end{enumerate}
In case (i), supply a \(C^*\)-dynamical reservoir
\((\mathfrak A_E,\tau_t,\omega_E)\), with stationary mixing state
\(\omega_E\), bounded self-adjoint \(B_\alpha\in\mathfrak A_E\), and
bounded interaction
\(\lambda\sum_\alpha A_\alpha\otimes B_\alpha\) in the standard GNS
weak-coupling representation.  Center by
\(\omega_E(B_\alpha)=0\), absorbing the means into \(H_S\), and put
\[
C_{\alpha\beta}(t)
:=
\omega_E\!\bigl(\tau_t(B_\alpha)B_\beta\bigr).
\]
No trace-class reservoir density is required.

In case (iii), start from the finite product assignment
\(\rho_S\otimes\sigma_{E,n}\), with \(\sigma_{E,n}\) stationary, and use
\[
H_{{\rm tot},n}
=
H_S\otimes\mathbf 1+\mathbf 1\otimes H_{E,n}
+
\lambda_n\sum_\alpha A_\alpha\otimes B_{\alpha,n}.
\]
Center \(\Tr(\sigma_{E,n}B_{\alpha,n})=0\), define the finite correlations
\[
C^{(n)}_{\alpha\beta}(t)
:=
\Tr\!\bigl(\sigma_{E,n}B_{\alpha,n}(t)B_{\beta,n}\bigr),
\]
and let \(C_{\alpha\beta}\) denote their integrable limit.  In cases (i)
and (iii), assume all of the following:
\begin{enumerate}[label=(\alph*), leftmargin=2em]
\item the coupling observables are bounded; in case (i) \(\omega_E\) is
stationary, while in case (iii) each displayed finite product assignment
and bath state is stationary;
\item
\[
\int_{\mathbb R}(1+|t|)\,
|C_{\alpha\beta}(t)|\,dt<\infty
\]
for every \(\alpha,\beta\), so the Fourier coefficients and one-sided
Lamb-shift integrals exist;
\item the stationary reservoir satisfies a mixing cluster condition at all
orders: the connected \(r\)-point correlation functions have integrable
time-difference majorants, uniformly in case (iii), such that the sum of
connected reduced Dyson cumulants of order \(r\ge3\) is \(o(1)\), uniformly
for the rescaled time \(\tau=\lambda^2t\) in case (i), respectively
\(\tau=\lambda_n^2t\) in case (iii), in every compact interval;
\item in case (iii), the two-point functions, their first weighted moments,
the connected-correlation majorants, and the resulting level-shift operators
converge uniformly to their limiting reservoir values on every compact
van Hove-time interval;
\item the finite system Bohr decomposition is kept fixed in the limit, and
the van Hove average retains all terms within each equal-frequency block,
including degeneracies, while averaging unequal Bohr frequencies.
\end{enumerate}
These are the theorem-domain hypotheses that justify the Davies
weak-coupling limit; two-point integrability or a long recurrence time alone
does not suffice for a general non-Gaussian bath.  Let
\(\mathcal E_\lambda(t)\) denote the actual Schr\"odinger-picture reduced
map in case (i), and use the corresponding index \(n\) in case (iii).
Convergence on the van Hove time \(\tau=\lambda^2t\) is in the
\emph{interaction picture}:
\[
 \widetilde{\mathcal E}_\lambda(\tau)
 :=\operatorname{Ad}(e^{iH_S\tau/\lambda^2})
                  \mathcal E_\lambda(\tau/\lambda^2),
 \qquad
 \sup_{0\le\tau\le\tau_0}
 \|\widetilde{\mathcal E}_\lambda(\tau)-e^{\tau\mathcal K}\|
 \longrightarrow0
 \quad(\tau_0<\infty),
\]
and analogously as \(n\to\infty\).  Its GKLS generator is
\[
\frac{d\widetilde\rho}{d\tau}=\mathcal K(\widetilde\rho)
=
-i[H_{\mathrm{LS}},\widetilde\rho]
+
\sum_{\omega,\alpha,\beta}
\gamma_{\alpha\beta}(\omega)
\Bigl(
A_\beta(\omega)\widetilde\rho A_\alpha(\omega)^*
-\frac12\{A_\alpha(\omega)^*A_\beta(\omega),\widetilde\rho\}
\Bigr),
\]
with positive coefficients
\[
\gamma_{\alpha\beta}(\omega)
=
\int_{-\infty}^{\infty} e^{i\omega t}C_{\alpha\beta}(t)\,dt.
\]
Writing \(\mathcal K=-i[H_{\rm LS},\,\cdot\,]+\mathcal D\), the equivalent
weak-coupling approximation in \emph{physical time} has generator
\[
 \frac{d\rho_{\rm app}}{dt}
 =-i[H_S+\lambda^2H_{\rm LS},\rho_{\rm app}]
                         +\lambda^2\mathcal D(\rho_{\rm app}).
\]
The secular construction makes \(\mathcal K\) commute with the free
system flow.  Consequently this physical-time approximation is uniform
on \(0\le t\le\tau_0/\lambda^2\) under the displayed interaction-picture
limit.  No unrenormalized Schr\"odinger-picture limit with finite
\(H_S\) on the \(\tau\)-scale is asserted.

In case (ii), let \(\Phi_{dt}\) be the exact one-probe reduced CPTP map and
assume, uniformly on the finite system algebra,
\[
\Phi_{dt}
=\operatorname{id}+dt\,\mathcal L+o(dt).
\]
Then \(\Phi_{t/n}^{\,n}\to e^{t\mathcal L}\) in norm, and
\(\mathcal L\) is GKLS.  This collision-model conclusion is separate from
the stationary-correlation Fourier formula above.
If, in addition, the secular components of the coupling are diagonal in a
common projector family \(\{P_a\}\) and \(H_S\) preserves these sectors, then
the generator has a pure-dephasing component with a nonnegative rate on each
off-diagonal block \(P_a\rho P_b\).  Blocks whose coupling difference lies
in the null space of the bath spectral matrix remain undamped; the
distinguishable \(P_a\)-classes are the corresponding pointer sectors.  For
general energy-exchange couplings, pointer sectors must be
identified from the actual Lindblad operators and need not coincide with the
eigenprojectors of the bare coupling observables.  If \(\omega_E\) is a
\(\beta\)-KMS state in case (i), or the limiting correlations in case (iii)
inherit the same KMS relation, the generator satisfies the associated
detailed-balance relation and the system Gibbs state is stationary.  A
repeated-probe limit has detailed balance only when its probe state and
scaled collision map satisfy the corresponding detailed-balance condition.

A single fixed finite bath is excluded from the dissipative convergence
claim: its correlations are quasiperiodic and it has recurrences.  It gives
the exact CPTP maps of
Theorem~\ref{p4:thm:exact-finite-system-environment-reduction}; if those maps
form an exact all-time semigroup,
Theorem~\ref{p4:thm:fixed-finite-bath-semigroup-rigidity} makes it reversible.
It may
approximate a Markov equation on a stated pre-recurrence window, but it does
not by itself yield a nontrivial irreversible semigroup in the van Hove
limit.
\end{theorem}


\begin{proof}
\ProofPlan{We separate exact finite reduction from the additional
reservoir/repeated-probe/large-bath limit, then apply interaction-picture
weak-coupling reduction, Bohr-frequency secular projection, positivity of the
limiting correlation transform, pointer-sector dephasing, and thermal detailed
balance.}
\begin{proofsteps}

\item \textbf{Interaction-picture Dyson expansion and bath trace for the
reservoir branches.}
In case (iii), pass to the interaction picture for the finite
system--bath Hamiltonian and write
\[
V_n(t):=\sum_\alpha A_\alpha(t)\otimes B_{\alpha,n}(t).
\]
Start from
\(\widetilde\rho_{SE,n}(0)=\rho_S\otimes\sigma_{E,n}\).
The centering condition removes the first-order bath mean.  The
interaction-picture Dyson expansion gives
\[
\widetilde\rho_{SE,n}(t)
=
\widetilde\rho_{SE,n}(0)
-\lambda_n i\int_0^t [V_n(s),\widetilde\rho_{SE,n}(s)]\,ds.
\]
Differentiating once more and tracing over the environment yields the second-order equation
\[
\dot{\widetilde\rho}_S(t)
=
-\lambda_n^2
\int_0^t
\Tr_E\!\Bigl(
[V_n(t),[V_n(s),\widetilde\rho_{SE,n}(s)]]
\Bigr)\,ds.
\]
In case (i), the same reduced Dyson-cumulant expansion is evaluated in the
GNS reservoir state by the conditional expectation
\(\operatorname{id}\otimes\omega_E\); it does not use a density matrix or
ordinary partial trace.

\item \textbf{Weak-coupling convergence of the limiting memory kernel.}
\leavevmode\par\noindent
In case (i), the displayed weighted correlation and all-order cluster
hypotheses are a Davies weak-coupling theorem domain.  They control the
Dyson-cumulant remainder uniformly for
\(\tau=\lambda^2t\) in compact intervals: the centered first-order term
vanishes, the second-order connected cumulant has a finite limit, and the
connected terms of order at least three are \(o(1)\) by hypothesis.  In
case (iii), the uniform versions of those bounds and convergence of the
level-shift operators give the same conclusion, while
\(\tau_{{\rm rec},n}\lambda_n^2\to\infty\) prevents recurrence on each fixed
rescaled-time window.  The Davies convergence theorem
\cite{Davies1974,Davies1976} therefore justifies replacing the memory state by
the instantaneous reduced state and extending the limiting second-order
correlation integral to \(+\infty\).  The resulting limiting kernel is
\[
\begin{aligned}
\dot{\widetilde\rho}_S(t)
={}&
-\lambda^2
\sum_{\alpha,\beta}
\int_0^\infty
\Bigl(
C_{\alpha\beta}(s)\,
[A_\alpha(t),A_\beta(t-s)\widetilde\rho_S(t)]
\\
&\hspace{5.5em}
+
\overline{C_{\alpha\beta}(s)}\,
[\widetilde\rho_S(t)A_\beta(t-s),A_\alpha(t)]
\Bigr)\,ds .
\end{aligned}
\]

\item \textbf{Separate repeated-probe/collision limit.}
In case (ii), each freshly prepared probe is discarded after one collision,
so successive system maps are independent compositions of the same exact
CPTP map \(\Phi_{dt}\).  The uniform expansion
\(\Phi_h=\operatorname{id}+h\mathcal L+r_h\), with
\(\|r_h\|/h\to0\), gives a direct norm estimate.  Put
\(A_h=\Phi_h\), \(B_h=\operatorname{id}+h\mathcal L\), and \(h=t/n\).
On the real space of self-adjoint trace-class matrices with its trace norm,
\(\|A_h\|=1\), \(\|B_h\|\le1+h\|\mathcal L\|\), and telescoping gives
\[
 \|A_h^n-B_h^n\|
 \le
 n\max(\|A_h\|,\|B_h\|)^{n-1}\|r_h\|
 \le
 t e^{t\|\mathcal L\|}\frac{\|r_h\|}{h}
 \longrightarrow0.
\]
The finite-dimensional Euler product \(B_{t/n}^n\to e^{t\mathcal L}\)
therefore gives
\[
\Phi_{t/n}^{\,n}\longrightarrow e^{t\mathcal L}
\]
in norm.  Every map in the limiting semigroup is CPTP, so
Theorem~\ref{p4:thm:gkls} makes \(\mathcal L\) a GKLS generator.  This
argument uses neither stationary bath correlations nor the Fourier
coefficients of the reservoir branches.


\item \textbf{Bohr-frequency decomposition and van Hove averaging in the
reservoir branches.}
Next decompose each system coupling into Bohr-frequency components for \(H_S\):
\[
A_\alpha(t)=\sum_\omega e^{-i\omega t}A_\alpha(\omega),
\qquad
[H_S,A_\alpha(\omega)]=-\omega A_\alpha(\omega).
\]
Substitution into the kernel above produces oscillatory terms of the form
\[
  e^{i(\omega'-\omega)t}.
\]
On the rescaled time, the Davies averaging map sends every term with
\(\omega\neq\omega'\) to zero, while terms in the same Bohr-frequency block
are retained, including degeneracies.  Thus this is a consequence of the
stated weak-coupling limit rather than an unquantified manual deletion.  The
limiting interaction-picture generator, with derivative in the rescaled
time \(\tau\), is
\[
\frac{d\widetilde\rho_S}{d\tau}
=
-i[H_{\mathrm{LS}},\widetilde\rho_S]
+
\sum_{\omega,\alpha,\beta}
\gamma_{\alpha\beta}(\omega)
\Bigl(
A_\beta(\omega)\widetilde\rho_S A_\alpha(\omega)^*
-\frac12\{A_\alpha(\omega)^*A_\beta(\omega),\widetilde\rho_S\}
\Bigr),
\]
where
\[
\Gamma_{\alpha\beta}(\omega)
:=
\int_0^\infty e^{i\omega s}C_{\alpha\beta}(s)\,ds
=
\frac12\gamma_{\alpha\beta}(\omega)+iS_{\alpha\beta}(\omega)
\]
and \(H_{\mathrm{LS}}=\sum_{\omega,\alpha,\beta}S_{\alpha\beta}(\omega)A_\alpha(\omega)^*A_\beta(\omega)\).
Each product \(A_\alpha(\omega)^*A_\beta(\omega)\) has zero Bohr
frequency, so \([H_S,H_{\rm LS}]=0\), and every dissipative term
intertwines the free flow.  Thus \(\mathcal K\) commutes with that flow.
Since \(d\tau/dt=\lambda^2\), returning to the Schr\"odinger picture in
physical time gives
\[
 \rho_{\rm app}(t)
 =e^{-itH_S}e^{\lambda^2t\mathcal K}(\rho_S(0))e^{itH_S},
 \qquad
 \dot\rho_{\rm app}
 =-i[H_S+\lambda^2H_{\rm LS},\rho_{\rm app}]
                         +\lambda^2\mathcal D(\rho_{\rm app}).
\]
Unitary conjugation is isometric in trace norm, so the interaction-picture
convergence estimate is exactly the corresponding physical-time
approximation estimate on \(0\le t\le\tau_0/\lambda^2\).
On the rescaled time itself, the Schr\"odinger free term would be
\(-i[H_S,\rho]/\lambda^2\).  Even a decoupled system with a coherent
initial state and noncentral \(H_S\) has rapid free oscillations and no
limit in that picture.  This establishes the required picture and
time-scaling distinction directly.


\item \textbf{Positivity of the limiting coefficient matrices and GKLS form.}
To see positivity of the coefficient matrices, note that the matrix-valued correlation function \(C_{\alpha\beta}(t)\) is of positive type:
for any finite families \(c_{j,\alpha}\) and times \(t_j\),
\[
\sum_{j,k,\alpha,\beta}
\overline{c_{j,\alpha}}\,C_{\alpha\beta}(t_j-t_k)\,c_{k,\beta}
=
\omega_E(Y^*Y)\ge 0
\quad\text{in case (i)},
\qquad
Y:=\sum_{k,\beta} c_{k,\beta}B_\beta(t_k).
\]
In case (iii), the corresponding finite expression is
\(\Tr(\sigma_{E,n}Y_n^*Y_n)\ge0\); positivity passes to the pointwise
correlation limit.
The matrix-valued Bochner theorem therefore makes the Fourier transform of
the integrable limiting correlation positive semidefinite, so the generator
is of GKLS form.  For a fixed finite bath this Fourier transform is instead
a pure-point distribution; the ordinary integrable transform used here is
not available.


\item \textbf{Pointer-sector dephasing.}
In the nondemolition pure-dephasing case, let the relevant secular coupling components share a common projector family \(\{P_a\}\), and let \(H_S\) preserve those sectors, so
\[
A_\alpha=\sum_a a_{\alpha,a}P_a.
\]
Then the dissipative part acts on an off-diagonal block \(P_a\rho P_b\) by
\begin{equation}
\label{p4:eq:general-complex-pointer-coefficient}
 \mathcal L_{\mathrm{diss}}(P_a\rho P_b)
 =\kappa_{ab}\,P_a\rho P_b,
\end{equation}
where the full, generally complex coefficient is
\begin{equation}
\label{p4:eq:general-complex-pointer-kappa}
 \kappa_{ab}
 :=
 \sum_{\alpha,\beta}\gamma_{\alpha\beta}(0)
 \left(
 a_{\beta,a}\overline{a_{\alpha,b}}
 -\frac12\overline{a_{\alpha,a}}a_{\beta,a}
 -\frac12\overline{a_{\alpha,b}}a_{\beta,b}
 \right).
\end{equation}
Let \(x_a=(a_{\alpha,a})_\alpha\) and
\(d_{ab}:=x_a-x_b\).  Since \(\gamma(0)\) is Hermitian positive
semidefinite,
\begin{align}
 \kappa_{ab}
 &=x_b^*\gamma(0)x_a
 -\frac12x_a^*\gamma(0)x_a
 -\frac12x_b^*\gamma(0)x_b,
 \label{p4:eq:general-complex-pointer-vector-form}\\
 -\operatorname{Re}\kappa_{ab}
 &=\frac12d_{ab}^*\gamma(0)d_{ab}\ge0.
 \label{p4:eq:general-complex-pointer-damping-rate}
\end{align}
Thus \eqref{p4:eq:general-complex-pointer-damping-rate}, rather than the
full coefficient \(\kappa_{ab}\), is the real quadratic-difference formula.
The imaginary part
\(\operatorname{Im}(x_b^*\gamma(0)x_a)\) is a possible dissipative
phase coefficient; for three or more sectors it need not be the difference
of two sector energies and therefore is not silently absorbed into one
diagonal Hamiltonian.  The block decays precisely when the quadratic form in
\eqref{p4:eq:general-complex-pointer-damping-rate} is positive and is
undamped precisely when \(d_{ab}\in\ker\gamma(0)\).  Equality of all
coupling data is sufficient for no damping, and is necessary only when
\(\gamma(0)\) is positive definite.  Thus the common eigensectors are
pointer candidates, with the actual distinguishable sectors determined by
the bath spectral form.


\item \textbf{Thermal detailed balance and stationary Gibbs state.}
Finally, in case (i), if \(\omega_E\) is a \(\beta\)-KMS state for
\(\tau_t\), the KMS relation implies
\[
\gamma_{\alpha\beta}(-\omega)=e^{-\beta\omega}\gamma_{\beta\alpha}(\omega).
\]
In case (iii), the same conclusion holds when the limiting correlations
inherit this KMS identity.  A finite Gibbs formula
\(\sigma_{E,n}\propto e^{-\beta H_{E,n}}\) is used only for those finite
approximants on which that density exists; it is not the definition of a
general mixing reservoir state.
This is the detailed-balance relation for the Davies generator, and it implies that the Gibbs state
\(
\rho_\beta\propto e^{-\beta H_S}
\)
is stationary.
The convergence on the van Hove scale is the Davies theorem on the explicit
weighted-correlation and all-order cluster domain displayed in the statement
\cite{Davies1974,Davies1976,BreuerPetruccione2002}, or the corresponding
uniform large-bath limit there.  The repeated-probe branch
instead has detailed balance only if its limiting collision generator is
assumed or proved to satisfy it.  The word
``finite'' applies to the system and each approximation, not to a claim that
one fixed finite bath is mixing.
\end{proofsteps}

\end{proof}

\begin{remark}[Status of the microscopic bath bridge]
\label{p4:rem:microbath-status}
Theorem~\ref{p4:thm:exact-finite-system-environment-reduction} is the
unconditional finite statement: a supplied fixed finite bath produces exact
CPTP reduced maps; if they obey an exact all-time semigroup law, they are
reversible by
Theorem~\ref{p4:thm:fixed-finite-bath-semigroup-rigidity}.
Theorem~\ref{p4:thm:microscopic-bath-gkls} is an
interface theorem whose dissipative limit requires a mixing reservoir,
fresh-probe process, or simultaneous large-bath/weak-coupling limit.
A merely pre-recurrence approximation may justify a finite-time effective
master equation, but not an exact semigroup limit on all \(t\ge0\).
Thus the rigorous hierarchy is: exact finite reduction gives CPTP maps;
an already supplied norm-continuous Markov semigroup has GKLS form by
Theorem~\ref{p4:thm:gkls}; and the stated non-finite limit interfaces can
produce such a semigroup.
\end{remark}



\subsection{Lie-algebraic controllability}
\label{p4:sec:part3-controllability}

For a controlled Hamiltonian family
\begin{equation}
\label{p4:eq:control-H}
H(t)=H_0+\sum_{k=1}^r u_k(t)H_k,
\end{equation}
let $\mathfrak g$ be the Lie algebra generated by $iH_0,\dots,iH_r$.


\begin{theorem}[Finite Lie-closure controllability]
\label{p4:thm:lie-closure-controllability}
Let $R\subset U(d)$ be the reachable set under the supplied piecewise-constant
bounded controls for Eq.~\ref{p4:eq:control-H}. Suppose the admissible controls
and drift permit the one-parameter groups of their real linear span in the
reachable closure, with the concatenations and inverses needed for the
Lie--Trotter and commutator constructions. Let $G$ be the connected analytic
Lie subgroup with Lie algebra
$\mathfrak g=\operatorname{Lie}(iH_0,\ldots,iH_r)$. Then
\[
 \overline R=\overline G.
\]
If $G$ is closed, this is $\overline R=G$. In particular
$\mathfrak g=\mathfrak{su}(d)$ gives dense reach of $SU(d)$ under the
stated admissibility conditions. No bounded control time is inferred.
\end{theorem}
\begin{proof}
Every admitted constant evolution is an exponential of a generator in
$\mathfrak g$, and every finite concatenation lies in $G$. Hence
$R\subset G$ and $\overline R\subset\overline G$.
By the reversibility hypothesis, $\overline R$ contains the one-parameter
groups of the real control span. Lie--Trotter limits give sums, and the
commutator product with both inverse flows gives the one-parameter groups
of every iterated bracket. Thus it contains the connected analytic subgroup
$G$ generated by $\mathfrak g$. Closedness of $\overline R$ then gives
$\overline G\subset\overline R$. The two inclusions prove the assertion.
If $G=SU(d)$ it is already closed and compact.
\end{proof}
\begin{remark}[Closure of an analytic control subgroup]
A connected analytic subgroup of a compact unitary group need not be closed.
For example the reversible group generated by
$\operatorname{diag}(1,\sqrt2)$ is dense in a two-torus. Its closure is the
closure of reachable operations. The compact native groups in
Corollary~\ref{p4:cor:v226-native-controls} are closed and retain the
corresponding stronger conclusion.
\end{remark}

Thus finite controllability in the framework reduces to a Lie-closure statement about admissible generators once the stated theorem-domain reversibility/accessibility conditions are supplied. Without such conditions the Lie algebra rank condition identifies the accessible Lie algebra but does not by itself prove $\overline R=\overline G$.



\subsection{Protected codes and correctable noise}
\label{p4:sec:part3-qec}

Let $\Phi(\rho)=\sum_a C_a\rho C_a^*$ be a noise channel on a finite sector and let $P$ project onto a candidate code subspace.
The Knill--Laflamme condition reads \cite{KnillLaflamme1997,NielsenChuang2010}
\begin{equation}
\label{p4:eq:knill-laflamme}
PC_a^*C_bP = \alpha_{ab}P
\end{equation}
for some Hermitian matrix $(\alpha_{ab})$.


\begin{theorem}[Finite Knill--Laflamme recovery]
\label{p4:thm:finite-kl-recovery}
If \eqref{p4:eq:knill-laflamme} holds, then there exists a CPTP recovery channel \(\mathcal R\) such that
\[
(\mathcal R\circ\Phi)(\rho)=\rho
\qquad
\text{for every code state }\rho=P\rho P.
\]
\end{theorem}


\begin{proof}
\ProofPlan{We prove finite Knill--Laflamme recovery by diagonalizing the error-correlation matrix, separating the noisy code sectors, constructing the recovery channel, and verifying exact recovery on code states.}
\begin{proofsteps}

\item \textbf{Diagonalization of the error-correlation matrix.}
Because the matrix \(\alpha=(\alpha_{ab})\) is Hermitian and positive semidefinite, there exists a unitary matrix \(u\) such that
\[
u\alpha u^*=\operatorname{diag}(\lambda_1,\dots,\lambda_N),
\qquad
\lambda_r\ge 0.
\]
Define new Kraus operators
\[
D_r:=\sum_a u_{ra}C_a.
\]
They represent the same channel as \(\{C_a\}\), and the Knill--Laflamme relation becomes
\[
PD_r^*D_sP=\lambda_r\delta_{rs}P.
\]

\item \textbf{Orthogonal noisy code sectors.}
Thus for \(r\neq s\) the images \(D_rP\mathcal H\) and \(D_sP\mathcal H\) are orthogonal, while for \(r=s\),
\[
(D_rP)^*(D_rP)=\lambda_r P.
\]
If \(\lambda_r>0\), the polar decomposition gives
\[
D_rP=\sqrt{\lambda_r}\,V_rP
\]
for a partial isometry \(V_r\) with initial space \(P\mathcal H\) and mutually orthogonal final spaces \(V_rP\mathcal H\).
If \(\lambda_r=0\), then \(D_rP=0\).


\item \textbf{Recovery channel construction.}
Define the recovery channel on the support of the noisy code space by
\[
\mathcal R(\sigma)
=
\sum_{\lambda_r>0} P\,V_r^*\,\sigma\,V_r\,P
+
\Tr\!\bigl((\mathbf 1-\Pi)\sigma\bigr)\tau,
\]
where
\[
\Pi:=\sum_{\lambda_r>0}V_rPV_r^*
\]
and \(\tau\) is any fixed code state.
This is CPTP because it is a Kraus sum on \(\Pi\mathcal H\) plus a trace-filling term on the orthogonal complement.


\item \textbf{Exact recovery on code states.}
Now let \(\rho=P\rho P\) be a code state.
Using \(D_rP=\sqrt{\lambda_r}V_rP\), one gets
\[
\Phi(\rho)=\sum_r D_r\rho D_r^*
=
\sum_{\lambda_r>0}\lambda_r\,V_r\rho V_r^*.
\]
Applying \(\mathcal R\) and using the orthogonality of the ranges,
\[
\mathcal R(\Phi(\rho))
=
\sum_{\lambda_r>0}\lambda_r\,P V_r^*V_r \rho V_r^*V_r P
=
\Bigl(\sum_{\lambda_r>0}\lambda_r\Bigr)\rho.
\]
Because \(\Phi\) is trace preserving on code states,
\[
1=\Tr(\Phi(\rho))=\sum_{\lambda_r>0}\lambda_r\,\Tr(\rho)=\sum_{\lambda_r>0}\lambda_r,
\]
so the prefactor is \(1\).
Hence \((\mathcal R\circ\Phi)(\rho)=\rho\), proving exact recovery.
\end{proofsteps}

\end{proof}

This is the finite-dimensional Knill--Laflamme sufficiency theorem on the finite carrier.  Fault-tolerant architectures and threshold theorems are outside the present domain.




\subsection[Conditional Reeb phase transfer interface]
{Conditional Reeb phase transfer interface}
\label{p4:sec:part3-signature}

\paragraph{Status.}
This section is not part of the finite quantum closure.  Paper~III's derived
common-rate hierarchy retains the fixed-shell family
\(\{\mathrm{REG}_\sigma\}_{\sigma_{\rm fin}\in\mathbb R}\).  Freeze one
member and write
\begin{equation}
\label{p4:eq:reeb-completed-step-import}
L_{\rm H}:=L_{\rm H}^{(\sigma)}
:=\Vlogcap+\delta_\sigma
=(2\pi)^2+\delta_\sigma,
\qquad
\delta_\sigma=\frac1{6\sqrt2}+\sigma_{\rm fin}.
\end{equation}
The source branch \(\mathrm{REG}_0\) is the specialization
\(\sigma_{\rm fin}=0\); it is not selected by the finite quantum closure.
Its rank valuation assigns \(L_{\rm H}\) to one completed rank-one atom and
\(7L_{\rm H}\) to a complete rank-seven descent.  It does not produce the quotient
\(L_{\rm H}/7\).  Neither equal rank increments nor the existence of the Reeb phase
direction alone implies that a laboratory residual is periodic in \(\ln E\).  A
periodic observable requires a separately derived theorem-domain transfer channel and
a discrete-scale covariance law.



\paragraph{Conditional discrete-scale covariance}
\label{p4:sec:part3-reeb-phase}

\begin{definition}[Reeb-covariant cyclic phase readout]
\label{p4:def:reeb-sensitive-readout}
Let \(y=\ln(E/E_{\rm geom})\) and fix \(s\in\mathbb N_{>0}\).  A specified
finite cyclic phase readout is called
\emph{Reeb-covariant with step \(sL_{\rm H}\)} when its residual after subtraction of
the standard Berry/adiabatic contribution is a real function
\[
R_{\mathcal E}\in L^2_{\rm loc}(\mathbb R;\mathbb R)
\]
and satisfies, almost everywhere,
\[
R_{\mathcal E}(y+sL_{\rm H})=R_{\mathcal E}(y),
\]
where \(s=1\) denotes a primitive-rank channel and \(s=7\) a channel proven to close
only after a full rank-seven descent; any other integer value requires its own rank-channel proof.  The covariance law and the value of \(s\) are
part of the specified apparatus/interface theorem domain; they are not consequences
of finite quantum mechanics alone.
\end{definition}

\begin{theorem}[Conditional Fourier form under supplied Reeb covariance]
\label{p4:thm:derived-reeb-phase-signature}
For a Reeb-covariant readout in the sense of
Definition~\ref{p4:def:reeb-sensitive-readout}, the residual descends to
\(\mathbb R/(sL_{\rm H}\mathbb Z)\) as a real \(L^2\) class and therefore has
the Fourier expansion, with convergence in \(L^2\),
\begin{equation}
\label{p4:eq:reeb-residual-periodic-class}
R_{\mathcal E}(y)
=
\sum_{n\in\mathbb Z}c_n^{(\mathcal E)}
\exp\!\left(\frac{2\pi iny}{sL_{\rm H}}\right),
\qquad c_{-n}^{(\mathcal E)}=\overline{c_n^{(\mathcal E)}}.
\end{equation}
Writing \(T=sL_{\rm H}\), the coefficients are
\[
c_n^{(\mathcal E)}
=
\frac1T\int_0^T
R_{\mathcal E}(y)e^{-2\pi iny/T}\,dy.
\]
The covariance hypothesis proves that \(T\) is a period:
\[
\operatorname{Per}(R_{\mathcal E})
:=
\{\tau>0:R_{\mathcal E}(y+\tau)=R_{\mathcal E}(y)
\text{ for almost every }y\}.
\]
\begin{equation}
\label{p4:eq:reeb-lab-log-period-derived}
\boxed{sL_{\rm H}\in\operatorname{Per}(R_{\mathcal E}).}
\end{equation}
It does not by itself prove that \(T\) is the least positive period.  If the
residual is nonconstant and
\[
d_{\mathcal E}:=
\gcd\bigl\{|n|:c_n^{(\mathcal E)}\ne0\bigr\},
\]
then its fundamental logarithmic period is
\[
\Delta\ln E_{\mathcal E}^{\rm fund}
=\frac{sL_{\rm H}}{d_{\mathcal E}}.
\]
In particular, the supplied step is fundamental exactly when
\(d_{\mathcal E}=1\).  A constant residual has no least positive period.
No universal first-harmonic amplitude, phase, or waveform follows from
\(\delta_\sigma\) alone.  These coefficients must be computed from the specified finite
transfer/readout functional.
\end{theorem}

\begin{proof}
The regularity and covariance hypotheses identify the residual with a real
\(L^2\) class on a circle of circumference \(T=sL_{\rm H}\).  The \(L^2\)
Fourier theorem gives \eqref{p4:eq:reeb-residual-periodic-class}, the coefficient
formula, and the conjugation relation.  A shift \(\tau\) is a period exactly
when
\[
c_n^{(\mathcal E)}
\bigl(e^{2\pi in\tau/T}-1\bigr)=0
\]
for every \(n\).  For a nonconstant residual, the least positive solution is
\(T/d_{\mathcal E}\); for a constant residual every shift is a period.  Thus
Paper~III fixes the conditional increment \(L_{\rm H}\), while the
theorem-domain covariance supplies a period and the channel's nonzero Fourier
support determines whether it is fundamental.
\end{proof}

\begin{corollary}[Specified readout fixes any observable harmonic]
\label{p4:cor:reeb-apparatus-overlap-interface}
If a specified apparatus/readout functional produces a real dimensionless transfer
function \(T_{\mathcal E}\in L^2(\mathbb R/\mathbb Z;\mathbb R)\), then its first harmonic is
\[
T_{\mathcal E}^{(1)}(x)
=a_{\mathcal E}\sin(2\pi x)+b_{\mathcal E}\cos(2\pi x),
\]
with
\[
a_{\mathcal E}=2\int_0^1T_{\mathcal E}(x)\sin(2\pi x)\,dx,
\qquad
b_{\mathcal E}=2\int_0^1T_{\mathcal E}(x)\cos(2\pi x)\,dx.
\]
Any physical conversion factor multiplying \(T_{\mathcal E}\), including a possible
appearance of \(\delta_\sigma/\pi\), must be derived from that apparatus channel and is not a
universal Paper-IV prediction.
\end{corollary}



\paragraph{Interface scope}
\label{p4:sec:part3-null-result}

Without an independently derived transfer coefficient and discrete-scale covariance law, the imported increment \(L_{\rm H}\) yields no observable log-periodic residual, universal waveform, amplitude, phase, or period \(L_{\rm H}/7\).  A null search therefore does not test the finite quantum core unless those interface hypotheses have first been established.



\section{Cofactor, moment, finite-readout countermodel and comparison tests}
\label{p4c:sec:v238-c7}
\noindent\textit{Source and claim boundary.} These are original extended calculations. Statements on physical sources, measurements or continuum limits retain the inherited antecedents of Papers IV-A/IV-B.

\subsection{Closed Hankel-pencil evaluation of source-visible poles}
\label{p4:subsec:v2188-source-moment-pencil}

\begin{theorem}[Exact source-cyclic moment determinant and positive weights]
\label{p4:thm:v2188-moment-pencil}
Let $(H,q,P)$ be an already licensed finite \emph{stationary} source chart
as in Theorem~\ref{p4:v2169:b10:thm:support} and
Lemma~\ref{p4:lem:v2186-exact-source-cyclic-action}. Let
$T=\operatorname{ad}_H$ on the complete support quotient and
$v=[q]\ne0$. To calculate a polynomial independent of a choice of
physical stationary density on \emph{this exact support}, use the
canonical \emph{calculation pairing} associated with
$\rho_P=P/\operatorname{rank}P$, and define
\[
 m_n=(v,T^n v)_{\rho_P},\qquad
 K_d=(m_{i+j})_{0\le i,j<d},\quad
 J_d=(m_{i+j+1})_{0\le i,j<d}.
\]
The first integer $d\ge1$ for which $K_d$ is positive definite and
$\det K_{d+1}=0$ equals the exact dimension of the source-cyclic space;
there $K_d>0$ and the \emph{monic} physical support polynomial is
\[
 \boxed{p_{q,P}(z)=\frac{\det(zK_d-J_d)}{\det K_d}.}
\]
Its $d$ roots are real, simple and exactly the source-visible frequencies,
with strictly positive canonical weights $w_a$ satisfying
\[
 \boxed{m_n=\sum_{a=1}^{d} w_a\lambda_a^n,\qquad
 w_a>0,\quad 0\le n\le 2d-1.}
\]
The weights are uniquely obtained from the first $d$ moment equations
by a Vandermonde solve after the roots are computed. Thus a claimed
terminal-to-pole equality at a pre-existing $\varepsilon$ has the
finite exact test
\[
 \boxed{p_{q,P}(\varepsilon)=0
 \ \Longleftrightarrow\ \det(\varepsilon K_d-J_d)=0.}
\]
The closing certificate uses moments through order $2d$; moments
through order $2d-1$ alone do not certify the end of the cyclic space.
For $v=0$ set $p_{q,P}=1$ separately.
\end{theorem}
\begin{proof}
Self-adjointness gives
$m_{i+j}=(T^i v,T^j v)_{\rho_P}$, so each $K_d$ is the
Gram matrix of the first $d$ Krylov vectors. These vectors remain
independent exactly until the first relation; after
$T^d v\in\operatorname{span}\{v,\ldots,T^{d-1}v\}$,
this span is $T$-invariant and contains every subsequent Krylov vector.
Hence $K_d>0$ and $\det K_{d+1}=0$ characterize its exact dimension.
In the Krylov basis, the matrix $A$ of $T$ satisfies $K_dA=J_d$;
therefore its characteristic polynomial is
$\det(zI-K_d^{-1}J_d)=\det(zK_d-J_d)/\det K_d$.
The restriction is self-adjoint and cyclic, so each eigenspace in this
restriction is one-dimensional and has nonzero source component.
The spectral theorem gives distinct real $\lambda_a$ and positive
$w_a=\|\Pi_{\lambda_a}v\|_{\rho_P}^2$, and hence the moment expansion.
The Vandermonde matrix for distinct roots is invertible. The scalar
polynomial is the first monic annihilator of $v$, and by the exact-support
theorem it is the physical source-visible polynomial. The terminal test
is simply substitution; it does not select the root.
\end{proof}


% BEGIN FRAMEWORK v2.239: VISIBLE SCHUR RESIDUE
\begin{proposition}[Regular Feshbach source visibility and positive pole residue]
\label{p4:prop:v239-visible-schur-residue}
For an already licensed finite self-adjoint source chart, write its
Hamiltonian on a fixed orthogonal source/memory split as
\[
 H=\begin{pmatrix}A&T^\dagger\\T&C\end{pmatrix},\qquad
 S(z)=zI-A-T^\dagger(z-C)^{-1}T.
\]
If $\lambda\in\mathbb R\setminus\operatorname{spec}(C)$ is a simple
eigenvalue of $H$, then it is visible to a specified source vector
$s$ in the upper block precisely when there is
$0\ne v\in\ker S(\lambda)$ with $s^\dagger v\ne0$.
The kernel is one-dimensional and the exact source-block inverse
residue is
\begin{equation}\label{p4:eq:v239-schur-positive-residue}
 \boxed{\operatorname{Res}_{z=\lambda}S(z)^{-1}
    =\frac{vv^\dagger}{v^\dagger S'(\lambda)v},\qquad
 S'(\lambda)=I+T^\dagger(\lambda-C)^{-2}T>0.}
\end{equation}
The scalar source residue is
$|s^\dagger v|^2/(v^\dagger S'(\lambda)v)$.
For a \emph{one-dimensional} bare source block $A=\varepsilon$,
$S(\varepsilon)=0$ with $T\ne0$ is impossible when
$\varepsilon-C$ is definite. An unshifted coupled regular root must
therefore receive complementary spectral contributions of both signs
(or involve a richer source/remainder block); nonzero coupling alone
does not exclude a protected endpoint.
\end{proposition}
\begin{proof}
The block eigenvalue equations give
$w=(\lambda-C)^{-1}Tv$ on the complement and
$S(\lambda)v=0$. The full eigenvector $(v,w)$ has squared norm
$\|v\|^2+\|(\lambda-C)^{-1}Tv\|^2=v^\dagger S'(\lambda)v$.
The simple-eigenvalue spectral projector therefore has the displayed
upper-left block. Differentiating the resolvent gives
$S'=I+T^\dagger(z-C)^{-2}T$. In the scalar case the endpoint
condition is $T^\dagger(\varepsilon-C)^{-1}T=0$; a definite inverse
has a strict sign on nonzero $T$. All statements retain the complete
actual $H$ and source domain; a kinematic finite constituent is not by
itself that physical chart.
\end{proof}
% END FRAMEWORK v2.239: VISIBLE SCHUR RESIDUE

\begin{corollary}[Finite common-moment comparison without a total-Hamiltonian map]
\label{p4:cor:v2188-common-moments}
Let independently licensed complete stationary source charts
$(H_\theta,q_\theta,P_\theta)$ each have the same, separately certified
source-cyclic degree $d$. If their canonical-support moments
$m_{n,\theta}$ agree for every $0\le n\le2d-1$, then the charts share
the same $p$ and the same \emph{canonical comparison} weights. The
Krylov-basis identification transports both the restricted source
operator and its canonical pairing isometrically. No map between their
full Hamiltonians or full occupation carriers is required.
If the same finite-moment equality is instead proved using each chart's
\emph{actual} stationary physical pairing, the resulting actual tagged
scalar resolvents and residues also agree. The canonical conclusion
alone says nothing about equality of residues at arbitrary distinct
physical densities with the same support.
\end{corollary}
\begin{proof}
The shared moments determine identical $K_d$ and $J_d$ and thus the
same determinant polynomial. They also determine the Gram matrix of
the entire Krylov bases and their $T$-action matrices, yielding the
isometric intertwiner. The spectral weights are the unique solution of
the shared Vandermonde system. The same proof works for actual pairings
when their moments coincide; support equality by itself gives only the
pole set, as in the existing exact-support theorem.
\end{proof}

\paragraph{Execution scope and no-regression rule.}
The canonical state $P/\operatorname{rank}P$ is a device for computing the
support-invariant pole set, \emph{not} a claim that the physical state is
maximally mixed. Every moment is generated from the \emph{complete}
$H$, including interaction, occupation, mixed source, constraint, metric,
contact and remainder returns, and from the \emph{actual} support $P$.
The pencil must not be formed on an arbitrary projected/truncated word
without the $m_{2d}$ closing test. For a changing support, a moving
connection, a nonstationary history, a physical rest/particle dictionary
or a continuum/renormalized readout, the corresponding antecedents remain.
In particular, C01/A05 is \emph{not} closed merely by this exact finite
algorithm. Appendix~\ref{p4:app:v2188-moment-pencil-test} contains a
native supplied-source replay including an invisible full-operator zero.

% END INLINE SOURCE: Paper4-SourceMomentPencil-Main-v1_0.tex



% BEGIN INLINE SOURCE: Paper4-SourceCofactor-Main-v1_0.tex
% Framework v2.191: exact finite execution formula; no physical tuple is selected.


\subsection{Direct source-visible cofactor extraction}
\label{p4:sec:v2191-source-cofactor}
The closed cyclic and moment constructions above have a direct finite-matrix
form. Paper V-B consumes it for the existing constant charged mass word.
All physical uses retain the already licensed complete generator, raw source
and exact support; this calculation supplies none of those inputs.

\begin{theorem}[Source-visible cofactor and rank certificate]
\label{p4:thm:v2191-source-cofactor}
Let $T=T^\dagger$ act on a finite complex Hilbert space and let $V$ be a
specified source map into it. Define
\[
 \chi_T(z)=\det(zI-T),\qquad
 \eta_{T,V}(z)=\operatorname{Tr}
       \bigl(V^\dagger\operatorname{adj}(zI-T)V\bigr).
\]
With a monic greatest common divisor over $\mathbb C[z]$, the first monic
polynomial annihilating the complete source image is
\[
 \boxed{p_{T,V}(z)=
   \frac{\chi_T(z)}{\gcd(\chi_T(z),\eta_{T,V}(z))}.}
\]
Its roots are real and simple. For every real $\varepsilon$,
\[
 \boxed{p_{T,V}(\varepsilon)=0
 \quad\Longleftrightarrow\quad
 \operatorname{rank}[\,T-\varepsilon I\ \ V\,]
       >\operatorname{rank}(T-\varepsilon I).}
\]
For $V=0$ use $p_{T,V}=1$, so there are no source-visible roots.
\end{theorem}
\begin{proof}
Writing $T=\sum_\lambda\lambda\Pi_\lambda$ gives the exact rational identity
\[
 \frac{\eta_{T,V}(z)}{\chi_T(z)}
 =\operatorname{Tr}\bigl(V^\dagger(zI-T)^{-1}V\bigr)
 =\sum_\lambda\frac{\|\Pi_\lambda V\|_{\rm HS}^2}{z-\lambda}.
\]
The nonnegative weight is nonzero exactly when that spectral space sees the
source. A nonzero weight leaves one factor $z-\lambda$ in the reduced
denominator; a zero weight cancels it completely, including all determinant
multiplicity. This reduced monic denominator is exactly the product of
the visible factors, and hence the first monic source-image annihilator.
Hermiticity gives
$\operatorname{ran}(T-\varepsilon I)=\ker(T-\varepsilon I)^\perp$.
Appending the columns of $V$ increases rank exactly when
$\Pi_\varepsilon V\ne0$. This proves the rank test, including degeneracy.
\end{proof}

\begin{corollary}[Actual adjacent-sector execution]
\label{p4:cor:v2191-adjacent-cofactor}
On the already licensed fixed-energy and fixed-charge support of
\eqref{p4:eq:v2169-b10-s9}, apply the theorem to
$(T_+,Q_+)$ and $(T_-,Q_-)$ separately. Their cofactor polynomials give
the same complete support polynomial through the LCM in
\eqref{p4:eq:v2169-b10-s10}. A null source leg contributes one.
On a global-ground support and at $\varepsilon>0$, the removal leg has no
pole at $\varepsilon$, and the displayed rank test on $(T_+,Q_+)$ is the
exact positive-frequency visibility test for densities with exact support.
\end{corollary}
\begin{proof}
Each $T_\pm$ is the complete Hermitian adjacent-sector matrix and each
$Q_\pm$ is its actual source map. The theorem supplies their source-image
annihilators. The inherited positive response weights prevent cancellation
between the two legs. For a global ground, $T_-\succeq0$, while the removal
term is $(z+T_-)^{-1}$, so it has no strictly positive root.
\end{proof}

A general stationary support still uses its full Liouvillian quotient or
closed moment pencil; fixed-energy adjacency is not imposed on it.
Nonzero source-image visibility on a degenerate support is also weaker than
visibility for every smaller-support state. Appendix~\ref{p4:app:v2191-cofactor-checks}
records the exact multiplicity, zero-mode and density-domain distinctions.

% END INLINE SOURCE: Paper4-SourceCofactor-Main-v1_0.tex


The supplied native constant-insertion comparison has a two-root singlet
cyclic polynomial with positive residues, while the other two roots of
its constant sector are source invisible. A determinant zero can therefore
miss this vector source. Appendix~\ref{p4:sec:native-source-comparison-support}
proves these exact outcomes, and separately evaluates the declared square;
the comparison does not evaluate the physical electron tagged polynomial.


% BEGIN INLINE SOURCE: Paper4-ActionStateCertificate-Main-v1_0.tex
% Framework v2.169: Batches 7--13, integrated at the owning source.


\subsection{A terminating all-orders nonvanishing test on an autonomous interval}
For a constant complete $H$, let $\mathcal L=\operatorname{ad}_H$ on the full finite
operator space, $v_1=\mathcal LQ$, and
$\ell(X)=\Tr(\rho_iX)$. Let $s$ be the dimension of the cyclic space
\[
 \mathcal K=\operatorname{span}_{\mathbb C}
 \{v_1,\mathcal Lv_1,\ldots\},\qquad s\leq d^2.
\]
For $v_1=0$ set $s=0$.

\begin{theorem}[Exact finite zero/nonzero charge certificate]\label{p4:v2171:thm:finite}
The charge change vanishes identically on the autonomous interval iff
\begin{equation}\label{p4:v2171:eq:moments}
 \ell(\mathcal L^nv_1)=0\quad(0\leq n<s).
\end{equation}
If one of these moments is nonzero, the charge change is a nonzero analytic
function of time. Its sign or value at a particular endpoint still requires
evaluation. For an aligned two-packet reversal diagnostic, replace
$\mathcal L$ by $\operatorname{ad}_H\oplus\operatorname{ad}_{H^\sharp}$, $Q$ by $(Q,Q^\sharp)$ and
$\ell$ by the corresponding signed state functional. The finite bound becomes
$d^2+(d^\sharp)^2$. Such a diagnostic has a physical CP interpretation only
under the actual inherited packet/charge convention; matrix conjugation
alone is not declared to be a physical CP implementation.
\end{theorem}
\begin{proof}
The exact Heisenberg charge is
$e^{i\mathcal L(t-t_i)/\hbar}Q$, so the derivative of its expectation is
\[
 \frac{i}{\hbar}\ell\left(
   e^{i\mathcal L(t-t_i)/\hbar}v_1\right).
\]
Finite cyclic closure supplies a recurrence expressing every later power
in terms of the first $s$. If \eqref{p4:v2171:eq:moments} holds, the entire
derivative vanishes, giving zero change. Conversely, a nonzero moment is a
nonzero Taylor derivative. Matrix exponentials are analytic, so the change
cannot vanish identically. Apply the same argument on the direct sum for
the signed diagnostic.
\end{proof}

The theorem does not posit a universal perturbative onset order. It tests
the actual complete generator at its fixed generated coupling strengths.
For a genuinely changing cosmological history use
Theorem~\ref{p4:v2171:thm:flow}; freezing that history merely to apply this theorem
would change the application domain.

% END INLINE SOURCE: Paper4-CompleteChargeTests-Main-v1_0.tex


% BEGIN INLINE SOURCE: Paper4-CompleteOccupationCommutator-Main-v1_0.tex
% Framework v2.171: integrated conversation proof on its stated domain.


\subsection{Cofactor multiplicities and exact-support visibility}
\label{p4:app:v2191-cofactor-checks}
These supplementary calculations explain the finite execution theorem;
they introduce no physical Hamiltonian, support, vacuum or source.
The complete Hermitian matrix and source map remain the theorem's inputs.

If an eigenvalue $\lambda$ has multiplicity $m$, write
$\chi_T(z)=(z-\lambda)^m c(z)$ with $c(\lambda)\ne0$.
For $w_\lambda=\|\Pi_\lambda V\|_{\rm HS}^2>0$, the spectral formula in
Theorem~\ref{p4:thm:v2191-source-cofactor} gives
\[
 \eta_{T,V}(z)
  =w_\lambda(z-\lambda)^{m-1}c(z)
     +(z-\lambda)^m a(z)
\]
locally, with $a$ analytic there. Exactly $m-1$ common factors cancel.
For an invisible eigenvalue the resolvent is analytic there, so all $m$
factors cancel. This is why a determinant root alone is not a source pole.
No inverse of a singular complementary block is used.

For the synthetic comparison $T=\operatorname{diag}(1,1,3)$ and $V=e_3$,
\[
 \chi_T=(z-1)^2(z-3),\qquad
 \eta_{T,V}=(z-1)^2,\qquad p_{T,V}=z-3.
\]
For $V=(e_1,e_3)$, the response is
$1/(z-1)+1/(z-3)$ and the reduced denominator is
$(z-1)(z-3)$. For $V=0$ it is one. These are comparison matrices, not
physical electron instances.

For the mass-source polynomial in
Corollary~\ref{p6:cor:v2191-mass-cofactor}, suppose $B\succeq0$ has a
source-visible zero eigenvalue of multiplicity $m$. Its determinant
$\chi_B(z^2)$ contains $z^{2m}$, while the visible term of
$zN_B(z^2)$ has order exactly $2m-1$. Cancellation therefore leaves
$z$, the correct simple zero-frequency factor. If that eigenspace is
source invisible, all its zero factors cancel. One must cancel the full
numerator $zN_B(z^2)$ rather than $N_B(z^2)$ alone.

For an actual density $\rho$ with exactly support $P$, positivity of
$\rho|_P$ makes
$\operatorname{Tr}(\rho V^\dagger\Pi_\lambda V)>0$ equivalent to
$\Pi_\lambda V\ne0$. The unweighted trace in the cofactor theorem
therefore computes the same root set without replacing the physical state.
If only containment $\operatorname{supp}\rho\le P$ is known, this
equivalence need not hold. Visibility for every state on $P$ instead
requires the corresponding positive effect to be positive definite on
$P\mathcal H$, as already stated in
\eqref{p4:eq:v2169-b10-s12}. The rank theorem does not discharge that
stronger family claim.

% END INLINE SOURCE: Paper4-SourceCofactorChecks-Appendix-v1_0.tex


% BEGIN INLINE SOURCE: Paper4-FiniteActionEvolution-Appendix-v1_0.tex
% Framework v2.190: implementation diagnostics, not added physical instances.


\subsection{Complete finite evolution diagnostics}
\label{p4:app:v2190-evolution-diagnostics}

The accompanying evidence checks identities on synthetic finite matrices.
Those fixtures are neither admitted physical source members nor alternative
actions or fitted parameters.  The analytic proofs in the main text control
the conclusions.

For \(S=G^{1/2}\), \(W=S^{-1}\), differentiating \(S^2=G\) gives
\(S\dot S+\dot SS=\dot G\).  The noncommuting term in
\(\gamma=WBW+(W\dot S-\dot SW)/2\) must remain.  The unnormalized
Euler equation is
\[
 i\hbar G\dot c+\tfrac{i\hbar}{2}\dot Gc=(K-i\hbar B)c.
\]
With \(q=Sc\) it is exactly \(i\hbar\dot q=(WKW-i\hbar\gamma)q\).
For \(G_p=I\), \(B_p=0\), \(Z=JG^{-1/2}\), the same generator is
\(Z^\dagger K_pZ-i\hbar Z^\dagger\dot Z\).  Dropping that connection
changes the law and its source jets.

The CAR derivative identities were tested by independent Jordan--Wigner
matrix multiplication on every even normal monomial in four modes,
including overlapping creation/annihilation subsets, pair creation and
annihilation, quartic and higher slots.  The proof uses the actual symbol
order; it does not treat the symbol quantization map as an algebra
homomorphism.  Time-dependent observable transport uses
\(U^\dagger[H(t),O]U\); replacing it by \([H(t),O_H]\) requires a
separate commutation/picture identification.

The moving-projector check tests the full equality
\(\dot P=-(i/\hbar)[H,P]\), not only \(P\dot P P=0\).
An arbitrary projected finite band need not satisfy it.  A jointly
finite bosonic word likewise retains its actual occupation-edge
commutation defects.  These limitations do not reopen an already
discharged finite action/quantization/evolution target, but prevent its
promotion to an unrestricted parent or continuum dynamics claim.

% END INLINE SOURCE: Paper4-FiniteActionEvolution-Appendix-v1_0.tex

% Framework v2.182: supplied native operator/source comparison.

% BEGIN INLINE SOURCE: Paper4-NativeComparisonSourceSupport-Appendix-v1_1.tex
% Framework v2.182 / v1.1: APPENDIX comparison; exact native source support with supplied constant insertion.


\subsection{An exact native source polynomial after complementary coupling}
\label{p4:sec:native-source-comparison-support}

This supplementary comparison supplies a constant insertion and singlet
source; neither is selected by the physical source-support theorem.
The finite support procedures of Paper IV consume their specified
operator and source. The following evaluates a comparison input from
the spatial geometry, making both that input and its source explicit.
It is a vector cyclic spectral polynomial, not a claimed evaluation of
the actual tagged electron commutator polynomial.

\begin{theorem}[Exact constant-insertion spectrum and cyclic source]
\label{p4:thm:native-source-constant-cyclic}
Let \(b\in\mathbb R^3\) be constant in the fixed left-invariant frame,
\(D_b=B_R+1_{\mathbb S}\otimes b\cdot\sigma_E\), and let the source
be the midpoint singlet \(s_0\). Its constant spatial sector reduces
\(D_b\). With \(\rho=|b|\), its four eigenvalues are
\begin{equation}
 \lambda_\pm=\frac1R\pm\sqrt{\frac1{R^2}+\rho^2},\qquad
 \tau_\pm=\frac2R\pm\rho.
 \label{p4:eq:native-source-four-eigenvalues}
\end{equation}
For \(\rho\ne0\) its source cyclic subspace has dimension two and
minimal polynomial
\begin{equation}
 p_{s_0,D_b}(z)=z^2-\frac2Rz-\rho^2,
 \qquad
 \langle s_0,(z-D_b)^{-1}s_0\rangle
       =\frac{z-2/R}{z^2-2z/R-\rho^2}.
 \label{p4:eq:native-source-cyclic-polynomial}
\end{equation}
At \(\rho=0\) the minimal polynomial is \(z\); the displayed
resolvent cancels to \(1/z\). For \(\rho\ne0\) the two spectral
weights are strictly positive,
\begin{equation}
 w_- =\tfrac12\left(1+\frac1{\sqrt{1+R^2\rho^2}}\right),\qquad
 w_+ =\tfrac12\left(1-\frac1{\sqrt{1+R^2\rho^2}}\right).
 \label{p4:eq:native-source-cyclic-weights}
\end{equation}
The two \(\tau\) eigenvalues have exactly zero source overlap.
\end{theorem}
\begin{proof}
A simultaneous constant spin and twist rotation fixes the singlet and
aligns \(b\) along the third axis. At the midpoint the constant sector
is \(2Q/R\), with \(Q\) the triplet projector. The two aligned triplet
vectors have eigenvalues \(2/R\pm\rho\). On the singlet and its
coupled triplet the matrix, after a fixed phase choice, is
\[
 \begin{pmatrix}0&\rho\\\rho&2/R\end{pmatrix}.
\]
Its eigenvalues, source resolvent and residues give every displayed
formula. Applying \(D_b\) once to the singlet generates the second
vector when \(\rho\ne0\); it cannot reach the other two. Constants
and their orthogonal complement are preserved by the left-invariant
operator, so this is the cyclic source subspace of the full spatial
operator, not merely a truncation of its source support.
\end{proof}

\begin{corollary}[Source-invisible determinant zeros]
\label{p4:cor:native-source-determinant-zero}
The constant-sector characteristic polynomial is
\[
 (z^2-2z/R-\rho^2)((z-2/R)^2-\rho^2).
\]
At \(\rho=2/R\) it has a zero eigenvalue, while
\(p_{s_0,D_b}(0)=-\rho^2\ne0\): that zero belongs to an invisible
triplet. Thus even an exact determinant zero of this supplied sector
does not substitute for the source-visible test.
\end{corollary}

\begin{corollary}[Exact energy polynomial for the declared square]
\label{p4:cor:native-source-square-polynomial}
For the declared comparison \(H_b=D_b^2\), \(\rho\ne0\),
\begin{equation}
 p_{s_0,H_b}(E)=E^2-(4/R^2+2\rho^2)E+\rho^4.
 \label{p4:eq:native-source-energy-polynomial}
\end{equation}
The roots are \(\lambda_-^2,\lambda_+^2\), with the same positive
weights above. If \(\rho<\sqrt3/(2R)\), the full spatial \(H_b\)
has a unique ground
\begin{equation}
 E_-(b)=\left(\sqrt{1/R^2+\rho^2}-1/R\right)^2
       =\frac{R^2\rho^4}{4}+O(\rho^6).
 \label{p4:eq:native-source-quartic-energy}
\end{equation}
The signed Dirac low branch is
\(\lambda_-=-R\rho^2/2+R^3\rho^4/8+O(\rho^6)\), consistently with
Theorem~\ref{p3f:thm:native-source-helicity-coefficient}.
\end{corollary}
\begin{proof}
Squaring the two-dimensional reducing matrix gives the stated trace
and determinant, and its two squares remain distinct. The full
midpoint nonzero gap is \(g=\sqrt3/R\). Bounded perturbation of norm
\(\rho<g/2\) leaves all other eigenvalues at absolute value at least
\(g-\rho>\rho\), whereas \(|\lambda_-|<\rho\). Its square is
therefore the full ground. The expansions follow from the exact roots.
\end{proof}

This closes the vector support polynomial, its visibility, and its
energy roots for the declared native comparison and supplied singlet
source. The actual C01 carrier, complete action, ground support,
electron insertion and tagged polynomial are different inputs. None
has been selected or replaced by this comparison calculation.

% END INLINE SOURCE: Paper4-NativeComparisonSourceSupport-Appendix-v1_1.tex


% BEGIN INLINE SOURCE: Paper4-MomentPencilDiagnostics-Appendix-v1_0.tex
% Framework v2.188: independent exact-rational replay of the pre-existing
% native constant-insertion comparison, not the physical electron action.


\subsection{Rational source-Hankel replay and invisible determinant zero}
\label{p4:app:v2188-moment-pencil-test}
In the \emph{supplied comparison} of
Theorem~\ref{p4:thm:native-source-constant-cyclic}, choose $R=1$ and
$\rho=2$ (the $\rho$ here is the constant insertion length, not a density).
The source-reducing matrix on the singlet/one-triplet span is
\[
 T=\begin{pmatrix}0&2\\2&2\end{pmatrix},\qquad v=(1,0)^{\mathsf T}.
\]
With its ordinary unit-source pairing, direct powers give
\[
 m_0=1,\quad m_1=0,\quad m_2=4,\quad m_3=8,\quad m_4=32,
\]
\[
 K_2=\begin{pmatrix}1&0\\0&4\end{pmatrix},\qquad
 J_2=\begin{pmatrix}0&4\\4&8\end{pmatrix},\qquad
 \det K_2=4,\quad \det K_3=0.
\]
The exactly closed pencil yields
\[
 \frac{\det(zK_2-J_2)}{\det K_2}=z^2-2z-4,
 \qquad \lambda_\pm=1\pm\sqrt5.
\]
Its source weights are
$w_\pm=\frac12(1\mp1/\sqrt5)>0$, respectively.
By contrast, the two invisible triplet frequencies of the \emph{full}
constant operator are $2\pm2=4,0$. The full determinant therefore
vanishes at $z=0$, but the exact source polynomial obeys $p(0)=-4$.
This tests the non-spurious-source-root rule of
Theorem~\ref{p4:thm:v2188-moment-pencil} on existing licensed
\emph{comparison} data, without selecting the physical electron
Hamiltonian, actual support, rest dictionary or terminal coefficient.

% END INLINE SOURCE: Paper4-MomentPencilDiagnostics-Appendix-v1_0.tex


% BEGIN C01 V2174 OCCUPATION COMPARISONS

% BEGIN INLINE SOURCE: Paper4-OccupationComparisons-Appendix-v1_0.tex
% v2.174: supplementary auxiliary and fixed-port comparisons.


\subsection{Auxiliary occupation and fixed-port comparison domains}
\label{p4:app:v2174-occupation-comparisons}
These comparisons protect the scope of the graded-port and action-moment
results. They assume only the explicitly stated auxiliary pair presentation
or standard complete sublevels. They supply no physical occupation selector,
Gibbs orientation, energy level or C01 quantum member.

\begin{lemma}[Exact auxiliary pair boundary]
\label{p4:lem:v2174-aux-pair-boundary}
On the auxiliary presentation of
Proposition~\ref{p2b:prop:v2174-geometric-pair-presentation}, let
$v_\mu$ ($1\le\mu\le10$) be the normalized pair basis. Then
\[
 B_\mu^\dagger=|v_\mu\rangle\langle0_{\rm aux}|,
 \qquad [B_\mu,B_\nu^\dagger]
     =\delta_{\mu\nu}P_{{\rm aux},0}
                         -|v_\nu\rangle\langle v_\mu|.
\]
On $\mathbf1\oplus\operatorname{Sym}^2S$ this is exactly the ten-mode
complete sublevel with pair occupation at most one; the separate spinor
block is dark. One pair quantum has auxiliary symmetric degree two.
The graded raising map of
$X_{{\rm aux},-}=\chi(B_{L_2L_2}^\dagger+B_{L_1L_1})$ has nilpotency
index two and pair cap one.
\end{lemma}
\begin{proof}
Direct multiplication of the displayed rank-one matrices gives the
commutator, and every product of two pair creators is zero.
For $N_{\rm pair}=P_{\operatorname{Sym}^2S}$, the extractor gives
$(X_{{\rm aux},-}+[N_{\rm pair},X_{{\rm aux},-}])/2
=\chi B_{L_2L_2}^\dagger\ne0$, whose square is zero.
\end{proof}
Its unit return is $X_{{\rm aux},-}^\dagger X_{{\rm aux},-}
|0_{\rm aux}\rangle=|\chi|^2|0_{\rm aux}\rangle$.
The actual standard physical coordinate instead has the full return
\eqref{p5c:eq:v2169-b13-o2}, including $|+,-\rangle$ when its actual cap
is at least two. Agreement with that one unit return at cap one does not
identify the physical field representation or the full Gibbs support.

\begin{proposition}[All-cap intertwiners of the same fixed graded ports]
\label{p4:prop:v2174-fixed-port-intertwiners}
Let $\mathcal F_{d,N}=\bigoplus_{n=0}^N\operatorname{Sym}^n\mathbb C^d$.
Its marked standard ladders generate $M_{D_N}(\mathbb C)$, where
$D_N=\binom{N+d}{d}$. For the two-charge-mode ports with the same
nonzero $\chi$ and their actual number grading,
\[
 \operatorname{Hom}_{N_B,X,X^\dagger}
       (\mathcal F_{2,N},\mathcal F_{2,M})
   =\begin{cases}\mathbb CI,&N=M,\\0,&N\ne M.\end{cases}
\]
There is no nonzero finite subspace of the unrestricted standard Fock
parent invariant under these fixed marked operations on their common domain.
\end{proposition}
\begin{proof}
The exact degree operator is $N_B=\sum_i a_i^\dagger a_i$.
Its unit projector is $P_{B,0}=\prod_{r=1}^N(I-N_B/r)$, using the empty
product at $N=0$. Thus
\[
 \frac{(a^\dagger)^{\mathbf n}P_{B,0}a^{\mathbf m}}
      {\sqrt{\mathbf n!\mathbf m!}}
        =|\mathbf n\rangle\langle\mathbf m|
 \qquad(|\mathbf n|,|\mathbf m|\le N)
\]
gives every matrix unit, with all intermediate states retained.
The graded extractor and its adjoint recover both creators and annihilators
from $N_B,X,X^\dagger$. An intertwiner sends the unique degree-zero unit
to $c$ times the target unit. If $N<M$, applying $(a_-^\dagger)^{N+1}$
gives zero in the source and $c\sqrt{(N+1)!}|0,N+1\rangle$ in the target,
so $c=0$; creation from the unit then makes the entire intertwiner zero.
For $N>M$ apply this argument to its adjoint. For equal caps creation from
the unit gives $cI$. Finally a finite invariant unrestricted subspace would
have a highest number eigenvector $v$; its creator is nonzero since
$\|a_-^\dagger v\|^2=\|a_-v\|^2+\|v\|^2$ and has larger degree,
a contradiction.
\end{proof}
This excludes precisely intertwiners of the same fixed marked ports and
proper quotients of their full matrix algebra. It does not exclude every
GP3 comparison, generated compression, different first-allocation candidate
or constrained quantization. Those routes consume their own generated
complete candidate and comparison domains.



\paragraph{Neutral amplitude versus projective direction and phase}
In the local quadratic chart of
Proposition~\ref{p4:prop:v2174-occupation-phase-moment}, take
$z_+=z_-=r>0$. Then $\mu_Q=0$, $\mu_{\rm occ}=2\hbar r^2$ and
$[z_+:z_-]=[1:1]$. For common frequency $\omega>0$, its phase orbit is
\[
 q_j=\sqrt{2\hbar/\omega}\,r\cos\theta,\qquad
 p_j=-\sqrt{2\hbar\omega}\,r\sin\theta,\qquad
 \oint\sum_jp_j\,dq_j=4\pi\hbar r^2=2\pi\mu_{\rm occ}.
\]
The last identity follows by integrating $\sin^2\theta$ over $[0,2\pi]$.
Thus the charge level, projective direction and phase period do not specify
the amplitude moment. This is a local parent calculation, not an admitted
global large-amplitude physical family or a new Bohr--Sommerfeld rule.



\paragraph{The consumed alternate-constructor domains}
The occupation theorem~\ref{p4:thm:gp1-bosonic-occupation-boundary} gives
existential finiteness; a complete initial sublevel additionally consumes
$\mathsf P_{\rm gbd}(\mathrm{occ},N)$. Mode K\"ahler geometry supplies
neither that boundary nor a complete interacting finite quantization.
II-E's normalized cotangent reduction consumes its moment level and
same-line identification
(Proposition~\ref{p2e:prop:normalized-cotangent-reduction} and
Corollary~\ref{p2e:cor:moment-prequantization-joint-locus}).
II-D's $\mathcal O(k)$ quantization consumes the integral line $k$
(Corollary~\ref{p2:cor:cp1-geometric-quantization}).
The scalar Born character lattice
(Lemma~\ref{p4:lem:scalar-born-character-lattice}) does not select a largest
oscillator degree. IV's reduced-source response
(Corollary~\ref{p4:cor:finite-reduced-source-response}) and III-F's Schur
construction (Corollary~\ref{p3f:cor:v1145-finite-quantum-source-schur})
consume an already complete finite joint source owner.
These route-specific domains are preserved; they are not a universal
quantization impossibility theorem.

% END INLINE SOURCE: Paper4-OccupationComparisons-Appendix-v1_0.tex

% END C01 V2174 OCCUPATION COMPARISONS



% BEGIN INLINE SOURCE: Paper4-ChargeTestComparisons-Appendix-v1_0.tex
% Framework v2.171: integrated conversation proof on its stated domain.


\subsection{Independent finite response comparisons and their limits}
The autonomous comparison uses
\[
 H=\begin{pmatrix}1&2+i\\2-i&-1\end{pmatrix},\quad
 H^\sharp=\overline H,\quad Q=\operatorname{diag}(0,-2),\quad
 \rho_i=\tfrac12\begin{pmatrix}1&1\\1&1\end{pmatrix}.
\]
The fixed-label reversal-odd charge comparison is
\[
 q_{\rm odd}(t)=\tfrac12\Tr\rho_i
 (e^{iHt}Qe^{-iHt}-e^{iH^\sharp t}Qe^{-iH^\sharp t})
 =\frac{\sin(\sqrt{24}\,t)}{\sqrt6}.
\]
Exact Gaussian-rational Krylov elimination gives degree two and
$\mathcal L^2v_1=24v_1$ for the doubled source vector. With the same
generators, either an empty comparison state or a tracial state gives an
identically zero reversal-odd difference. The displayed nonzero comparison
also has exact zero endpoints. These states, generators and endpoints are
comparison inputs, not native cyclic data or a physical CP dictionary.

The time-dependent comparison retains four CAR modes, a 16-dimensional
Fock carrier, complex pairing, two different normal generators and quartic
density contacts. Its generated charge-test space has dimension 128.
Time-ordered exponentiation of its real expectation matrices agrees with
independent full Fock propagation. Deleting the quartic contacts changes
the answer, and reversing the two intervals changes it. A separately declared
commuting late comparison owner preserves the calculated charge. This is a
positive check of the finite theorems and a negative control against omitted
terms, missing chronology and inferred freeze-out.



\subsection{Formula comparisons and their physical scope}
\label{p4:v2171:app:fullchainchecks}
These comparisons verify the new complete-state, export-coefficient and
readout identities. Their small matrices are not supplied native histories,
and their occupation labels do not choose a physical resolution or cap.

\paragraph{A mathematical trace cannot be the charge-production state.}
For a complete finite algebra, $\rho=I/d$ satisfies
$U\rho U^\dagger=\rho$ for every unitary $U$, and
$\operatorname{Tr}(\rho[H,Q])=0$. Therefore every time-ordered commutator
coefficient of its charge change vanishes, even if an individual marked
coefficient such as $\Im\Theta$ is nonzero. The native inherited root is
not replaced by this trace state. This identity explains why the local
coefficient trace is a test of retained algebra rather than a final
physical charge expectation.

\paragraph{Complete finite CAR evolution with signed conjugate charges.}
The independent occupation-matrix comparison uses six mathematical slots:
one quark, one antiquark, two leptons and two antileptons. Its three
noncommuting time stages retain both pair orientations, quark--antiquark
creation, normal mixing and nonlinear density contacts. Every stage
commutes with $Q_B$, while its lepton charge changes. The literal full
state propagation preserves $q_B=0$ and reproduces the Ward source with
both conjugate contributions. This verifies the theorem on its declared
domain. It does not compute the native K16 state, or introduce these
six slots as the physical mode content.

% END INLINE SOURCE: Paper4-ChargeTestComparisons-Appendix-v1_0.tex




% BEGIN INLINE SOURCE: Paper4-CutoffWard-Appendix-v1_0.tex
% Framework v2.162: diagnostics delimiting the conditional Ward theorem.


\subsection{Finite-window Ward defect and interacting guard leakage}
\label{p4:app:v4106-cutoff-ward-defect}
This appendix records scope diagnostics rather than a premise of the
functional response theorem. For a frozen many-body projector $P$,
let $Q_\perp=I-P$. Direct multiplication gives
\[
 [PHP,PGP]-P[H,G]P=PGQ_\perp HP-PHQ_\perp GP.
\]
When the parent canonical defect is zero and the window is frozen, its
compressed defect is
\[
 \mathfrak D_{\xi,F}=\frac{i}{\hbar}
                 (PHQ_\perp G_\xi P-PG_\xi Q_\perp HP).
\]
It is generally noncentral; finite Hermiticity does not make it vanish.
An actual covariant moving carrier can cancel this term with its own
projector/connection jets, but existence of mathematical transport alone
does not license replacing the physical carrier word.

An empty occupation guard is not an interacting invariance theorem.
For four CAR modes, take a genuinely inherited diagnostic quartic
$H=\hbar g(a_3^\dagger a_4^\dagger a_2a_1+
a_1^\dagger a_2^\dagger a_4a_3)$ and initial state $|12\rangle$.
On the invariant span of $|12\rangle,|34\rangle$ the Hamiltonian is
$\hbar g\sigma_1$, so
\[
 |\psi(t)\rangle=\cos(gt)|12\rangle-i\sin(gt)|34\rangle,\qquad
 \langle N_3+N_4\rangle=2\sin^2(gt).
\]
This diagnostic does not add a native four-fermion coupling. It proves
that a free empty-guard argument cannot be exported to generic interactions
without an actual leakage/invariance certificate. Likewise an interacting
quartic expectation need not factor into one-body covariances.
The main theorem uses the full finite density, so it needs no such
factorization. Numerical identities and exact moment checks are reproduced
in the packaged metric/Ward companion scripts; they test algebra, not
the occurrence of a physical source member.

% END INLINE SOURCE: Paper4-CutoffWard-Appendix-v1_0.tex


% BEGIN INLINE SOURCE: Paper4-TransportedEmbeddingDiagnostics-Appendix-v1_0.tex
% Framework v2.189 / Paper IV v4.118: presentation-versus-physical-band diagnostics.


\subsection{Dynamical fidelity, special source sectors and finite occupation diagnostics}
\label{p4:app:v2168-source-dynamics-diagnostics}
These conditional branches and synthetic witnesses test the exact finite
response constructions. They are not an evaluation of the native charged
Hamiltonian or a physical vacuum choice.
\begin{corollary}[A genuinely interacting but source-invisible finite branch]
\label{p4:cor:v2168-source-invisible-interaction}
Let $H_0$ be number preserving and suppose the admitted Fock unit $\Omega$
and the one-excitation source sectors are invariant. If $V$ is a self-adjoint
normal-ordered number-preserving fermion polynomial with at least two
annihilators in every monomial, then $V$ vanishes on the zero- and
one-occupation sectors. Consequently the $\Omega$ two-point function of
linear CAR ports is exactly the quadratic one, even if $V$ is nonzero on
higher occupation sectors.
\end{corollary}
\begin{proof}
Every displayed monomial kills both sectors. Its adjoint has the same property,
and number preservation prevents re-entry from higher sectors. The vacuum is
stationary and the one-occupation resolvent is the unaltered quadratic
resolvent. Its linear-source anticommutator response is therefore unchanged.
\end{proof}

This is a proved finite realization branch, not a statement that the actual
charged gauge/Yukawa action has this form. A fermion--boson interaction can
connect a one-fermion source to a fermion-plus-boson sector and must remain in
the complete memory block. Nor is the Fock unit automatically the ground
state or the occurring physical vacuum of the full charged theory.


\paragraph{The first commutator of a retained gauge-current vertex}
The native same-datum connection enters the fermion operator through
$c_g(\rho_*(A))$ in III-F's joint relation closure. In a fixed finite kinetic
frame, let $T_\alpha=T_\alpha^\dagger$ be its actual Hermitian one-body
vertex matrices, including the inherited coupling and spin/spatial/source
marks. Quantizing the retained bosonic mode ports $X_\alpha=X_\alpha^\dagger$
gives the current vertex
\[
 V_{\rm lin}=\sum_\alpha d\Gamma(T_\alpha)\otimes X_\alpha.
\]
The $T_\alpha$ are the coefficient matrices of the existing covariant
fermion word, not new adjustable gauge or mass matrices. Restricting to such
a finite mode chart still requires its actual licensed projection and
constraint domain. Its first source commutator is already evaluated by CAR:
\[
 \boxed{[V_{\rm lin},a^\dagger(j_r)\otimes I]
       =\sum_\alpha a^\dagger(T_\alpha j_r)\otimes X_\alpha.}
\]
It creates the fermion--boson partner columns that enter $B$ and the exact
memory. V-C's same-source electromagnetic derivative uses this same owned
connection/current, rather than an independent current normalization.

\begin{proposition}[Exact vertex contribution on a stationary product-unit branch]
\label{p4:prop:v2168-product-unit-vertex-moment}
Suppose the complete Hamiltonian on this declared branch is
\[
 H=d\Gamma(h)\otimes I+I\otimes H_B+V_{\rm lin},
\]
with $h=h^\dagger$, and let the bosonic unit satisfy
$H_B\Omega_B=E_B\Omega_B$ and
$\langle\Omega_B,X_\alpha\Omega_B\rangle=0$. The fermionic Fock unit
$\Omega_F$ gives a stationary product state
$\Omega_F\otimes\Omega_B$. For normalized source columns $J=(j_r)$ and
$q_r=a^\dagger(j_r)\otimes I$, define
\[
 We_r=\sum_\alpha T_\alpha j_r\otimes X_\alpha\Omega_B.
\]
Then its first two source moments are exactly
\[
 \boxed{\mu_1=J^\dagger hJ,\qquad
        \mu_2=J^\dagger h^2J+W^\dagger W.}
\]
In components the new Gram is
\[
 (W^\dagger W)_{rs}
 =\sum_{\alpha,\beta}j_r^\dagger T_\alpha T_\beta j_s\,
 \langle X_\alpha\Omega_B,X_\beta\Omega_B\rangle.
\]
\end{proposition}
\begin{proof}
The number-preserving one-body terms and $V_{\rm lin}$ kill $\Omega_F$,
so the product unit is stationary with energy $E_B$. The first commutator
is $a^\dagger(hj_r)\otimes I$ plus the displayed vertex columns. Its first
moment is the stated $J^\dagger hJ$ because the boson means vanish. The
second moment is the Gram of these commutator columns. On the fermion unit,
the annihilating side of the anticommutator contributes zero and the creation
side has the ordinary one-occupation inner product. The free and vertex
columns are orthogonal because each $X_\alpha\Omega_B$ is orthogonal to
$\Omega_B$. Their Grams therefore add as displayed.
\end{proof}

The formula evaluates a native vertex prefix and exhibits the next required
data: the actual mode matrices and boson-source covariance, both supplied by
the same action/quantization word on its admitted state domain. For a complete
action with further couplings, or an entangled interacting vacuum, retain all
additional commutator columns and cross expectations. The product-unit
formula cannot substitute for those terms.

If the free Gram is already $M^2$, a nonzero $W$ changes the whole-response
comparison even on this stationary branch. Thus exact preservation of the
entire free rational function is not the appropriate universal requirement
for interacting mass prediction. The next section supplies the weaker pole
test; Appendix~\ref{p4:app:v2168-source-dynamics-diagnostics} evaluates a finite single-channel illustration.



\paragraph{An exact witness for the fourth-moment test}
This appendix is a diagnostic of the moment criterion, not a second native
charged Hamiltonian or an action-selection prescription. On an abstract
source/complement calculation space take
\[
 L=\begin{pmatrix}0&M&0\\M&0&t\\0&t^\dagger&0\end{pmatrix},
 \qquad \mathcal Q=\begin{pmatrix}I\\0\\0\end{pmatrix},
 \qquad t=\begin{pmatrix}1\\0\\0\end{pmatrix},\quad
 M=\diag(2,3,5).
\]
Then $\mu_1=0$, $\mu_2=M^2$ and $\mu_3=0$, but
\[
 \boxed{\mu_4-M^4=Mtt^\dagger M=\diag(4,0,0).}
\]
The electron-tag response is
\[
 G_{ee}(z)=\frac{z^2-1}{z(z^2-5)},
\]
instead of $z/(z^2-4)$. Its extra zero-frequency response and shifted pair
are detected by the positive leakage; moments through order two do not see
them. The numerical energies here are synthetic diagnostic values.



\paragraph{A finite single-channel vertex illustration}
This is a synthetic instance of the preceding vertex proposition, not a
native electromagnetic mass prediction, source selection or physical-vacuum
assertion. Take two fermion modes with
$h=\begin{psmallmatrix}0&\varepsilon\\\varepsilon&0\end{psmallmatrix}$,
a finite two-state bosonic port
$H_B=\diag(0,\omega)$, $X=\begin{psmallmatrix}0&1\\1&0\end{psmallmatrix}$,
and $T=gI$. Its fermionic current vertex is $gN_F\otimes X$.
The product unit is stationary. For the first fermion creation source,
\[
 \mu_1=0,\qquad \mu_2=\varepsilon^2+g^2.
\]
On its one-occupation sector the exact Hamiltonian is
$h\otimes I+I\otimes(H_B+gX)$. Its source-visible frequencies are
\[
 \pm\varepsilon+\delta_\pm,\qquad
 \delta_\pm=\frac{\omega\pm\sqrt{\omega^2+4g^2}}2.
\]
Their residues are $\tfrac12$ times the unit's squared overlap with the two
bosonic eigenvectors. For $g\ne0$ neither $\delta_\pm$ is zero. This explicit
memory calculation explains why a finite gauge-current channel must be
included in a pole-only comparison rather than declared invisible because
its vacuum mean vanishes.

\paragraph{Why the faithful-family scope is necessary.}
For two fermionic modes with
$H=a n_0+b n_1+g n_0n_1$ and $q=a_0^\dagger$, the charged eigenoperators
are $a_0^\dagger(1-n_1)$ and $a_0^\dagger n_1$, with frequencies $a$ and
$a+g$. Their positive residue operators are $1-n_1$ and $n_1$.
Every faithful stationary state sees both frequencies. The Fock-unit state
sees only $a$. This illustrates the exact positivity/support boundary; it
does not select a physical charged vacuum or an interaction coefficient.

\paragraph{A finite occupation correction that must survive.}
If a joint occupation projection does not factor over bosonic coordinates,
two compressed coordinate matrices can fail to commute. Then the last
line of (C1) and the second sum of (C2) need not vanish. The verification
uses noncommuting Hermitian finite ports to check both terms against the
full Hamiltonian. Commuting classical coordinate names are insufficient
to discard these finite projected terms.

% END INLINE SOURCE: Paper4-ChargedResponseDiagnostics-Appendix-v1_0.tex



% v2.169: additional hypotheses and finite comparisons.

% BEGIN INLINE SOURCE: Paper4-ChargedPolynomialDiagnostics-Appendix-v1_0.tex
% Framework v2.169: Batches 7--13, integrated at the owning source.


\subsection{Finite detector and inclusive soft-sector identities}
\begin{proposition}[Finite channel probabilities and flux stability]
\label{p4:prop:v2210-lit-detector-stability}
Let $S_m,S_n$ be the coherently compared finite outgoing channel operators
of Section~\ref{p4:sec:part2-finite-window-scattering}, with
$\|S_m\|,\|S_n\|\le1$, common normalized incoming vector $\psi$,
$\|(S_m-S_n)\psi\|\le\eta$, and outgoing effects
$0\le Q_m,Q_n\le I$, $\|Q_m-Q_n\|\le\delta$. Then
\[
 |p_m-p_n|\le2\eta+\delta,\qquad
 p_j=\langle S_j\psi,Q_jS_j\psi\rangle.
\]
If the actually calibrated incident fluxes satisfy
$f_m,f_n\ge f_*>0$ and $|f_m-f_n|\le\delta_f$, then
\[
 \left|\frac{p_m}{f_m}-\frac{p_n}{f_n}\right|
 \le\frac{2\eta+\delta}{f_*}+\frac{\delta_f}{f_*^2}.
\]
\end{proposition}
\begin{proof}
Telescope the two quadratic forms in their vectors and effects. Cauchy--Schwarz
and the contraction bounds give $2\eta+\delta$. Telescope the two flux
ratios and use $p_n\le1$ and the positive flux floor.
\end{proof}
The M\o ller criterion supplies $\eta\to0$ only on its actual locality,
leakage, return, tail and embedding domain. Detector/flux comparison is
explicitly retained, while asymptotic completeness is not required for the
finite probabilities. Quasi-local Haag--Ruelle and finite-volume decay
methods \cite{LitQFTBachmannScattering2016,LitQFTLellouchLuscher2001}
are possible stronger comparisons; their mass-shell or elastic-volume
conditions must be proved for the actual native channel and geometry.

\begin{corollary}[Inclusive probabilities on a licensed finite soft factor]
\label{p4:cor:v2210-lit-inclusive-soft}
On an already licensed hard/soft tensor decomposition, let $\rho$ be the
actual outgoing joint density, $\{Q_h\}$ the hard detector POVM and
$\{F_s\}$ any unresolved soft POVM with $\sum_sF_s=I$. Then
\[
 \sum_s\operatorname{Tr}\rho(Q_h\otimes F_s)
 =\operatorname{Tr}\rho(Q_h\otimes I)
 =\operatorname{Tr}\rho_{\rm hard}Q_h.
\]
\end{corollary}
\begin{proof}
Sum the effects and use the already derived partial trace.
\end{proof}
Entanglement with the unresolved soft factor does not invalidate this
identity. It is independent of a basis or finer unresolved POVM describing
the same detector question. Changing the physical soft-energy cut changes
the question. It proves neither an infrared limit nor a sharp charged
Fock pole. Gauge-compatible dressings and memory-sensitive in/out algebra
scattering \cite{LitQFTHiraiSugishita2021,LitQFTPrabhuSatishchandran2024}
remain claim-scoped routes requiring actual soft/current history, charge
identities, resolution and in/out comparisons. They do not insert a new
particle ontology.



\subsection{Exact finite Gaussian determinant jets}
\begin{theorem}[Owned determinant contribution to source jets]
\label{p4:thm:v2210-lit-determinant-jets}
Let $Q_a(J)$ be the complete constraint-reduced nonsingular finite
quadratic determinant blocks of one occurring action/source family on a regular
scalar source interval, with its admitted Gaussian factorization or actual
graded Schur reduction already performed. Assume $C^3$ dependence and use the local determinant
branches and weights $w_a$ fixed by
Corollary~\ref{p4:cor:v2140-relative-determinant-weights}; no weight is
independently selected. The exact Gaussian determinant contribution is
\[
 \mathcal G(J)=\sum_aw_a\bigl[\log\det Q_a(J)-\log\det Q_a(0)\bigr].
\]
Writing $G_a=Q_a^{-1}$ and primes for complete derivatives on that same
source chart gives
\begin{align*}
 \mathcal G'&=\sum_aw_a\operatorname{Tr}(G_aQ_a'),\\
 \mathcal G''&=\sum_aw_a\operatorname{Tr}
       (G_aQ_a''-G_aQ_a'G_aQ_a'),\\
 \mathcal G'''&=\sum_aw_a\operatorname{Tr}
       (G_aQ_a'''-3G_aQ_a'G_aQ_a''+2(G_aQ_a')^3).
\end{align*}
These are exact for the admitted Gaussian integral. They are the corresponding
one-loop determinant term only on an independently admitted perturbative
effective-action domain, not the complete interacting effective action.
\end{theorem}
\begin{proof}
The native finite Gaussian/Berezin determinant identities fix the powers
and signs. Jacobi's formula gives $d\log\det Q=\operatorname{Tr}(Q^{-1}dQ)$,
and differentiating $Q^{-1}Q=I$ gives $G'=-GQ'G$. Differentiate twice more;
trace cyclicity combines the three mixed third-derivative terms. Local
nonsingularity licenses each logarithm and every differentiation.
\end{proof}
All source-dependent frames, constraints, measure factors and boundary
terms entering the actual quadratic owner must enter $Q_a$ or their
separately owned determinant contributions. For any protected bare
source/terminal jet, protection after this Gaussian correction is the
explicit equality of the corrected jet with the licensed prediction after
already owned source transport. The displayed trace sums provide its finite
test; symmetry alone does not make them zero. The exact finite graded flow
of Theorem~\ref{p4:thm:v4100-exact-graded-eaa} remains separate from
truncating it to this determinant. The spectral-action counterterm result
\cite{LitQFTvanSuijlekom2012} concerns a flat four-dimensional
higher-derivative/asymptotic presentation and permits a spectral-function
shift there. It does not establish either a shift or radiative protection
of the native exact heat profile. That conclusion requires the complete
finite source calculation above, or a stronger interacting calculation.
% END LITERATURE-CHECKED QUANTUM CERTIFICATES v2.210



% Framework v2.251: exact arbitrary-state evaluation of the quadratic finite source.
\subsection{Exact arbitrary-state evaluation of the finite source influence}
\label{p4c:sec:v247-arbitrary-state-influence}
\begin{theorem}[Complete correlation polynomial for a quadratic CAR source]
\label{p4c:thm:v247-source-correlation-polynomial}
Let $\mathfrak h\cong\mathbb C^d$ be the actual complete retained
one-particle band. Suppose the complete source Hamiltonian has exactly
$H_q(t)=d\Gamma(h_q(t))+c_q(t)I$, with Hermitian $h_q$ and the
inherited clock and frame. Let $\rho_i$ be any supplied positive
trace-one density on $\bigwedge\mathfrak h$; neither a quasifree
preparation nor particle-number diagonality is assumed. Write
$W=u_-^\dagger(t_f)u_+(t_f)$ for the actual one-particle history return
and $\Delta c=\int(c_{q_+}-c_{q_-})dt$. For increasing $k$-tuples
$I=(i_1,\ldots,i_k)$ and $J=(j_1,\ldots,j_k)$ define the occupied
correlation operator on $\bigwedge^k\mathfrak h$ by
\[
 R^{(k)}_{I,J}
 =\operatorname{Tr}\!\left(\rho_i
 a_{j_1}^\dagger\cdots a_{j_k}^\dagger
 a_{i_k}\cdots a_{i_1}\right),\qquad R^{(0)}=1.
\]
Then the exact finite source influence is the degree-at-most-$d$
polynomial in the complete one-particle return:
\begin{equation}
 \boxed{Z[q_+,q_-]=e^{-i\Delta c/\hbar}
 \sum_{k=0}^d\operatorname{tr}_{\wedge^k\mathfrak h}
       \left(R^{(k)}\bigwedge\!{}^k(W-I)\right).}
 \label{p4c:eq:v247-source-correlation-polynomial}
\end{equation}
The matrix of $\bigwedge^k(W-I)$ consists of its $k$-minors.
All coefficients in this polynomial are contractions of the actual
supplied state; no new state, weight or preparation selector enters.
\end{theorem}
\begin{proof}
The exact exterior evolution gives
$U_q=e^{-i\int c_qdt/\hbar}\bigwedge u_q$.
Cyclicity therefore reduces the influence to
$e^{-i\Delta c/\hbar}\operatorname{Tr}(\rho_i\bigwedge W)$.
Unitarily diagonalize the actual unitary $W$, with eigenvalues $w_j$.
In that basis its exterior lift is
\[
 \bigwedge W=\prod_{j=1}^d(1+(w_j-1)a_j^\dagger a_j).
\]
Expanding the product gives the principal diagonal entries of the
claimed trace sum. Under a unitary change of one-particle basis,
$R^{(k)}$ and $\bigwedge^k(W-I)$ transform by the same exterior
conjugation, so their trace pairing is invariant. This returns the
formula to the initial frame, including the off-diagonal minors.
Particle-number coherences of $\rho_i$ make no contribution to this
number-preserving exterior operator, consistently with the stated
correlation contractions. The formula is an exact finite identity,
not a truncation of a cumulant or loop expansion.
\end{proof}

\begin{corollary}[Quasifree compression and the exact limit of covariance data]
\label{p4c:cor:v247-correlation-determinant-compression}
For a supplied gauge-invariant quasifree state with covariance
$C_{ij}=\operatorname{Tr}(\rho_i a_j^\dagger a_i)$,
$R^{(k)}=\bigwedge^k C$ and the sum is
\[
 \sum_{k=0}^d\operatorname{tr}\bigwedge\!{}^k(C(W-I))
 =\det(I-C+CW).
\]
For a general state its covariance alone is insufficient, even though
its one-particle mean and causal commutator jets are fixed by that
covariance on this quadratic source branch.
\end{corollary}
\begin{proof}
The supplied Wick rule gives $R^{(k)}=\bigwedge^k C$.
Multiplicativity of exterior powers and their characteristic-polynomial
trace identity give the determinant. For the insufficiency, on two
modes compare
\[
 \rho_{\rm corr}=\tfrac12(|00\rangle\langle00|+|11\rangle\langle11|),
 \qquad \rho_{\rm qf}=I_4/4.
\]
Both have $C=I_2/2$, but $R^{(2)}=1/2$ and $1/4$, respectively.
At $W=e^{i\theta}I_2$ their influence factors are
$(1+e^{2i\theta})/2$ and $(1+e^{i\theta})^2/4$.
At $\theta=\pi/2$ the former is zero and the latter is $i/2$.
The total-number variances are $1$ and $1/2$.
The mean and commutator statement follows from
$\langle d\Gamma(A)\rangle=\operatorname{tr}(CA)$ and
$[d\Gamma(A),d\Gamma(B)]=d\Gamma([A,B])$.
\end{proof}

Thus the complete finite quadratic source can be evaluated for its
actual preparation, including a paired or non-Gaussian preparation,
without replacing it by a state having the same two-point data.
The determinant is an exact compression only on its proved preparation
domain. In contrast, an actual quartic interaction in $H$ is not an
exterior lift of a one-particle generator and still requires the full
interacting propagator. The correlation polynomial does not remove
such an interaction or compute an interacting charged particle pole.


% Proposed positive closure after v2.247. No new physical interaction or state.
\subsection{A controlled determinant evaluation of the interacting finite source}
\label{p4:sec:v248-interacting-determinant-parent}

\begin{theorem}[Interacting finite-current influence as a convergent determinant sum]
\label{p4:thm:v248-current-determinant-sum}
On one actually licensed complete finite CAR/boson source chart let
$\mathfrak h=\mathbb C^d$, let $\mathcal B$ be the actual finite bosonic
carrier, and retain the complete self-adjoint Hamiltonian
\[
 H[q](t)=d\Gamma(h[q](t))\otimes I+I\otimes H_B[q](t)
 +\sum_{\alpha=1}^m d\Gamma(T_\alpha[q](t))\otimes X_\alpha[q](t)
 +c[q](t)I.
\]
All matrices, couplings, cutoffs, constraints and marks in this expression
are inherited from that word; every $h,H_B,T_\alpha,X_\alpha$ is Hermitian.
Assume the supplied initial state on this branch is
$\rho_i=\rho_F(n)\otimes\rho_B$, where $\rho_F(n)$ is gauge-invariant
quasifree, $0\le n\le I$, and $\rho_B$ is any positive trace-one density
on $\mathcal B$. This is a preparation hypothesis, not a reset during
evolution. For two admitted histories assume $H_\pm(t)$ is $C^1$ on
$[t_i,t_f]$.

Write the noncentral summands of $H_\sigma(t)$ as $H_{\sigma,a}(t)$,
$a=1,\ldots,m+2$, in the fixed order displayed above. For $T=t_f-t_i$,
$\delta=T/N$ and $t_j=t_i+j\delta$, form the unitary slice
\[
\begin{aligned}
 S_{\sigma,j}&=e^{-i\delta c_\sigma(t_j)/\hbar}
 \bigl(\bigwedge e^{-i\delta h_\sigma(t_j)/\hbar}
                        \otimes e^{-i\delta H_{B,\sigma}(t_j)/\hbar}\bigr)\\
 &\qquad{}\cdot\prod_{\alpha=1}^m
 e^{-i\delta d\Gamma(T_{\alpha,\sigma}(t_j))
                         \otimes X_{\alpha,\sigma}(t_j)/\hbar},\\
 U_{\sigma,N}&=S_{\sigma,N-1}\cdots S_{\sigma,0}.
\end{aligned}
\]
Spectrally resolve each finite bosonic port at its slice,
$X_{\alpha,\sigma}(t_j)=\sum_r x_{\alpha,\sigma,j,r}
P_{\alpha,\sigma,j,r}$. Expanding these finite sums gives exactly
\[
 U_{\sigma,N}=e^{-iC_{\sigma,N}/\hbar}
 \sum_{\mathbf r}\bigwedge w_{\sigma,\mathbf r}
                                  \otimes M_{\sigma,\mathbf r},
 \qquad C_{\sigma,N}=\delta\sum_jc_\sigma(t_j),
\]
where each $w_{\sigma,\mathbf r}$ is the chronological product of the
actual one-particle factors $e^{-i\delta h/\hbar}$ and
$e^{-i\delta x_rT_\alpha/\hbar}$, chronological across slices
and in the fixed printed product order within each slice, and each
$M_{\sigma,\mathbf r}$ is
the identically ordered bosonic product of $e^{-i\delta H_B/\hbar}$
and the spectral projectors. The sliced influence is the finite sum
\begin{equation}
\boxed{\begin{aligned}
 Z_N={}&e^{-i(C_{+,N}-C_{-,N})/\hbar}
 \sum_{\mathbf r,\mathbf s}
 \operatorname{Tr}_{\mathcal B}
       (M_{+,\mathbf r}\rho_B M_{-,\mathbf s}^\dagger)\\[-2pt]
 &\hspace{32mm}\cdot
 \det\bigl(I-n+n w_{-,\mathbf s}^\dagger w_{+,\mathbf r}\bigr).
\end{aligned}}
\label{p4:eq:v248-interacting-determinant-sum}
\end{equation}
Let $U_\sigma$ be the exact complete interacting propagator, and put
\[
\begin{aligned}
 L_\sigma&=\sup_t\|\dot H_\sigma(t)\|,\\
 C_\sigma^{\rm split}&=\sup_t\sum_{a<b}
                         \|[H_{\sigma,a}(t),H_{\sigma,b}(t)]\|,\\
 \epsilon_{\sigma,N}&=\frac{T\delta}{2}
       \left(\frac{L_\sigma}{\hbar}
                       +\frac{C_\sigma^{\rm split}}{\hbar^2}\right).
\end{aligned}
\]
Then
\[
 \boxed{\|U_\sigma-U_{\sigma,N}\|\le\epsilon_{\sigma,N},\qquad
 |Z-Z_N|\le\epsilon_{+,N}+\epsilon_{-,N},\qquad
 Z=\lim_{N\to\infty}Z_N.}
\]
Thus the full interacting finite influence is evaluated by the convergent
sum of one-particle determinants and actual bosonic matrix traces.
No quasifree factorization of the evolved interacting state is assumed.
\end{theorem}
\begin{proof}
For each port, spectral functional calculus and the exact exterior lift give
\[
 e^{-i\delta d\Gamma(T)\otimes X/\hbar}
 =\sum_r\bigwedge e^{-i\delta x_rT/\hbar}\otimes P_r.
\]
Multiplying these finite identities proves the branch expansion, since the
exterior lift is multiplicative. Substitution in
$\operatorname{Tr}(U_{+,N}\rho_iU_{-,N}^\dagger)$ and the v2.247 exact
quasifree determinant identity prove \eqref{p4:eq:v248-interacting-determinant-sum}.
Noncommuting bosonic projectors and cross paths remain inside their stated
ordered matrix trace; no intermediate projection or state re-preparation
has been performed.

On each slice, Duhamel's formula and unitarity bound the error from freezing
$H_\sigma(t)$ at $t_j$ by $L_\sigma\delta^2/(2\hbar)$.
The finite Lie--Trotter commutator estimate bounds the error of splitting
that frozen Hermitian sum by
$\delta^2\sum_{a<b}\|[H_{\sigma,a},H_{\sigma,b}]\|/(2\hbar^2)$.
This estimate follows for two terms by integrating their conjugation
commutator over the time triangle; induction on the number of terms and
$[A,B+C]=[A,B]+[A,C]$ give the displayed sum. All conjugating factors
are unitary, so no exponential growth factor is needed. Telescoping the
$N$ unitary slices gives $\epsilon_{\sigma,N}$.
Finally $\|\rho_i\|_1=1$, so replacing either propagator in the influence
trace costs at most its operator-norm error. This proves the influence
bound and convergence.
\end{proof}

\begin{corollary}[Normalized interacting local parent with certified error]
\label{p4:cor:v248-interacting-parent-error}
The sliced and exact influences satisfy $Z_N[q,q]=Z[q,q]=1$ and
$|Z_N|,|Z|\le1$. Their matrices on every finite list of histories are
positive semidefinite. On a sufficiently small simply connected history neighborhood of the
diagonal where the normalized logarithm exists (for example
$|Z-1|<1/2$), suppose $|Z|\ge\zeta>0$ and
$\epsilon_N:=\epsilon_{+,N}+\epsilon_{-,N}<\zeta/2$, the logarithms
normalized at equal histories obey
\[
 \boxed{\Gamma_N=-i\hbar\log Z_N,\quad
 \Gamma=-i\hbar\log Z,\quad
 |\Gamma_N-\Gamma|\le\frac{2\hbar\epsilon_N}{\zeta}.}
\]
The limit is uniform on any such history set with uniform $C^1$ bounds and
uniform split-commutator bounds. The finite sum has complex coherent
weights; it is not an independently chosen stochastic auxiliary field.
\end{corollary}
\begin{proof}
Each $U_{\sigma,N}$ is unitary, giving normalization and the trace
Cauchy--Schwarz bound. For histories $q_a$, the entries
$\operatorname{Tr}(U[q_a]\rho_iU[q_b]^\dagger)$ form a Gram matrix;
the same argument applies to $U_N$. The segment from $Z$ to $Z_N$ stays
at distance at least $\zeta/2$ from zero, so integrating $dz/z$ bounds
the normalized local logarithm difference by $2\epsilon_N/\zeta$.
The stated common neighborhood and normalization fix the same branch.
Uniform bounds in the theorem give uniform convergence.
\end{proof}

This theorem extends the exact determinant evaluation to the actually
retained finite gauge-current interaction. Its spectral labels are algebraic
summation indices, with all interference retained. It neither selects the
physical mode/state nor supplies a continuum measure, an independently
missing current coefficient, a source-to-metric incidence, or an interacting
mass pole. The number of path terms can grow with the number of slices;
no computational-efficiency claim is made. Splitting factors are
mathematical approximants, not separately occurring physical operations.
If only the full Hamiltonian preserves a constraint, no extra projection
is inserted between its factors. For an actual non-quasifree
initial fermionic state, replace each determinant by its v2.247 finite
occupied-correlation polynomial. An initially entangled CAR/boson state
requires its retained joint correlations instead of the product trace used
here.

\begin{lemma}[Finite discrete decoupling of an actually retained quartic square]
\label{p4:lem:v248-quartic-discrete-decoupling}
Let $S=S^\dagger$ on the actual $d$-mode band, $F=d\Gamma(S)$,
$\lambda\in\mathbb R$, and let
$x_1,\ldots,x_\ell$ be the distinct real subset-sum eigenvalues of $F$.
Choose real $s\ne0$ so the numbers $z_r=e^{isx_r}$ are distinct, and
solve the finite Vandermonde system
\[
 \sum_{k=0}^{\ell-1}a_k z_r^k=e^{-i\delta\lambda x_r^2/\hbar}
 \qquad(r=1,\ldots,\ell).
\]
Then the exact operator identity is
\[
 \boxed{e^{-i\delta\lambda[d\Gamma(S)]^2/\hbar}
       =\sum_{k=0}^{\ell-1}a_k\bigwedge e^{iksS}.}
\]
Any supplied self-adjoint number-preserving polynomial through quartic order
can be expressed as a sum of a central term, a one-body term and finitely
many real multiples of such Hermitian squares. Consequently the determinant
sum construction and its unitary splitting error also apply when that
actual quartic word is retained among the Hamiltonian summands.
\end{lemma}
\begin{proof}
Choose $s$ sufficiently small that the finitely many distinct $x_r$ remain
distinct modulo $2\pi/s$. The Vandermonde system is invertible; the spectral
theorem then gives the polynomial interpolation identity in $e^{isF}$.
Since $e^{iksF}=\bigwedge e^{iksS}$, this is the displayed decoupling.
For the final assertion write $E_{ij}=a_i^\dagger a_j$. CAR give
\[
 a_i^\dagger a_j^\dagger a_l a_k=E_{ik}E_{jl}-\delta_{kj}E_{il}.
\]
Expand the $E_{ij}$ in a real Hermitian one-body basis. In a self-adjoint
quadratic polynomial in its Hermitian lifts, the imaginary antisymmetric
coefficient part is a real linear combination of $i[F_a,F_b]$, hence
one-body by the CAR Lie identity. The remaining real symmetric quadratic
coefficient matrix diagonalizes orthogonally, giving real multiples of
Hermitian squares. This is a rewriting of the supplied interaction, not a
new quartic coupling. Splitting uses the original unitary square
exponential; its finite coherent expansion need not have positive weights.
\end{proof}


\begin{corollary}[Exact interacting evaluator for an arbitrary supplied joint state]
\label{p4:cor:v248-joint-state-interacting-determinants}
Retain the actual Hamiltonian, finite carriers and unitary slices of
Theorem~\ref{p4:thm:v248-current-determinant-sum}, but replace its product
quasifree preparation hypothesis by an arbitrary supplied positive
trace-one density $\rho_i$ on
$\bigwedge\mathfrak h\otimes\mathcal B$, fixed under the admitted source
variations. For increasing $k$-tuples $I=(i_1<\cdots<i_k)$ and
$J=(j_1<\cdots<j_k)$ define the boson-valued occupied correlations
\[
 \mathcal R^{(k)}_{I,J}
 :=\operatorname{Tr}_{\wedge\mathfrak h}
 \!\left[\left(a_{j_1}^\dagger\cdots a_{j_k}^\dagger
                 a_{i_k}\cdots a_{i_1}\otimes I\right)\rho_i\right],
 \qquad
 \mathcal R^{(0)}=\operatorname{Tr}_{\wedge\mathfrak h}\rho_i.
\]
For every pair of branch paths set
$W_{\mathbf s,\mathbf r}=w_{-,\mathbf s}^\dagger w_{+,\mathbf r}$ and
$B_{\mathbf s,\mathbf r}=M_{-,\mathbf s}^\dagger M_{+,\mathbf r}$.
Then the exact sliced influence is
\begin{equation}
\boxed{\begin{aligned}
 Z_N={}&e^{-i(C_{+,N}-C_{-,N})/\hbar}
 \sum_{\mathbf r,\mathbf s}\sum_{k=0}^d\sum_{|I|=|J|=k}
 \det\!\left((W_{\mathbf s,\mathbf r}-I)_{I,J}\right)\\[-2pt]
 &\hspace{25mm}\cdot
 \operatorname{Tr}_{\mathcal B}
       \left(\mathcal R^{(k)}_{J,I}B_{\mathbf s,\mathbf r}\right).
\end{aligned}}
\label{p4:eq:v248-joint-state-interacting-evaluator}
\end{equation}
The empty minor is $1$. The same bounds
$\|U_\sigma-U_{\sigma,N}\|\le\epsilon_{\sigma,N}$ and
$|Z-Z_N|\le\epsilon_{+,N}+\epsilon_{-,N}$ hold, as do the normalized
local-parent limit and its error bound. Thus initial CAR/boson entanglement,
non-quasifree occupation correlations and paired preparations do not
obstruct this finite evaluator.

On the supplied product gauge-invariant quasifree branch the hierarchy
factorizes as
$\mathcal R^{(k)}_{I,J}=(\bigwedge^k n)_{I,J}\rho_B$;
\eqref{p4:eq:v248-joint-state-interacting-evaluator} then reduces exactly
to \eqref{p4:eq:v248-interacting-determinant-sum}.
\end{corollary}
\begin{proof}
For any one-particle matrix $A$, the finite exterior-lift identity is
\[
 \bigwedge(I+A)
 =\sum_{k=0}^d\sum_{|I|=|J|=k}
 \det(A_{I,J})\,
 a_{i_1}^\dagger\cdots a_{i_k}^\dagger
                       a_{j_k}\cdots a_{j_1}.
\]
Apply it with $A=W_{\mathbf s,\mathbf r}-I$ to each branch pair after
cyclically moving the negative-branch factors in the full trace. By the
definition of the partial trace,
\[
 \operatorname{Tr}\!\left[\rho_i
 \left(a_{i_1}^\dagger\cdots a_{i_k}^\dagger
       a_{j_k}\cdots a_{j_1}\otimes B\right)\right]
 =\operatorname{Tr}_{\mathcal B}(\mathcal R^{(k)}_{J,I}B).
\]
This fixes both the reversed creation/annihilation ordering and the
$J,I$ indices in the formula. No commutation of $\mathcal R^{(k)}_{J,I}$
with $B$ is assumed. The propagation error proof used only
$\|\rho_i\|_1=1$, so it is unchanged. On the product quasifree branch,
the finite CAR Wick determinant gives the stated factorization and the
exterior-minor determinant sum recovers the earlier expression.
\end{proof}

All $\mathcal R^{(k)}$ are extracted from the same supplied initial joint
state; they are not independently selectable source parameters. Their
finite hierarchy is an exact state representation for this number-preserving
source evaluator. It does not prepare that state, license an additional
interaction, or identify an uncomputed physical tagged pole.







% BEGIN INTEGRATED V2.249 PROOF PACKAGE
% BEGIN FRAMEWORK v2.249: CERTIFIED LOG DOMAIN AND NATIVE INFLUENCE ZERO
\subsection{An explicit uniform logarithm domain and its finite obstruction}
\label{p4:sec:v249-certified-log-domain}

\begin{proposition}[A certified tube around the complete history diagonal]
\label{p4:prop:v249-uniform-log-tube}
Let the supplied complete finite Hamiltonian word act on one fixed Hilbert
space on a time interval of length $T$, and let its admitted histories satisfy,
in their declared history norm $\|\cdot\|_X$,
\[
 \sup_t\|H[q](t)-H[\widetilde q](t)\|
 \le L_H\|q-\widetilde q\|_X.
\]
The constant includes every retained source, central term and moving-frame
contribution on the declared domain. For coordinate-only words $X$ may be
the uniform norm; retained moving-frame terms involving $\dot q$ generally
require a norm such as $C^1$ which also controls that derivative.
For any fixed positive trace-one joint
state, put $Z[q_+,q_-]=\operatorname{Tr}(U_+\rho_i U_-^\dagger)$ and
$r=T L_H\|q_+-q_-\|_X/\hbar$. Then
\begin{equation}
 \boxed{|Z-1|\le r.}
 \label{p4:eq:v249-log-tube}
\end{equation}
Consequently, throughout any admitted tube with $r\le r_*<1$, the principal
logarithm gives one normalized influence action
$\Gamma=-i\hbar\operatorname{Log}Z$, with
$|Z|\ge1-r_*$ and $\Gamma[q,q]=0$.
If a computed amplitude obeys $|Z_N-Z|\le\varepsilon_N$ uniformly and
$r_*+\varepsilon_N<1$, its principal logarithm lies on that same normalized
branch and satisfies
\begin{equation}
 |\Gamma_N-\Gamma|
 \le \frac{\hbar\varepsilon_N}{1-r_*-\varepsilon_N}.
 \label{p4:eq:v249-log-tube-error}
\end{equation}
When the exact histories and word have $C^2$ dependence on the retained
control parameters, this action is $C^2$ on the tube. This last assertion
concerns the exact action; convergence of approximate source derivatives
requires the separate derivative estimates.
\end{proposition}
\begin{proof}
Duhamel's difference identity between the two full unitary evolutions gives
\[
 \|U_+-U_-\|\le\hbar^{-1}\int
       \|H[q_+](t)-H[q_-](t)\|\,dt\le r.
\]
Cyclicity and $\|\rho_i\|_1=1$ imply
$|Z-1|=|\operatorname{Tr}(\rho_i U_-^\dagger(U_+-U_-))|\le r$.
The disk $|z-1|<1$ avoids zero and the principal-log cut, fixing a unique
normalized logarithm without any assumption about the topology of the
history tube. The segment from $Z$ to $Z_N$ stays in
$|z-1|\le r_*+\varepsilon_N$. Integrating the derivative $1/z$ along that
segment proves \eqref{p4:eq:v249-log-tube-error}. Finite-dimensional
parameter-dependent ODE differentiation gives $C^2$ dependence of the exact
propagators; composition with the logarithm proves the regularity assertion.
\end{proof}

\begin{proposition}[An exact influence zero on the retained paired restriction]
\label{p4:prop:v249-paired-influence-zero}
Use the two-spin restriction of
Proposition~\ref{p5c:prop:v248-paired-history-compiler}, with fixed real
$v\ne0$, real $g>0$, $\hbar$ explicit, and the admitted supplied state
$\rho_i=I_4/4$. Compare the two constant histories
\[
 H_-=-\tfrac v2(n_0+n_1)+v n_0n_1,
 \qquad H_+=H_-+g(a_0^\dagger a_1^\dagger+a_1a_0).
\]
Their complete finite interacting influence is
\begin{equation}
 \boxed{Z(T)=\cos^2\!\left(\frac{gT}{2\hbar}\right).}
 \label{p4:eq:v249-paired-influence-zero}
\end{equation}
It vanishes at $T=(2k+1)\pi\hbar/g$, $k\ge0$. Thus the retained class of
finite paired words does not admit one finite continuous logarithmic action
on its entire history-duration domain. This is a supplied-word comparison,
not a claim that the physical electron preparation or coefficients have
been selected.
\end{proposition}
\begin{proof}
On the even basis $(|0\rangle,|01\rangle)$, $H_-=0$ and
$H_+=g\sigma_1$. On each odd line the two Hamiltonians equal $-v/2$.
The relative unitary therefore has even trace $2\cos(gT/\hbar)$ and odd
trace $2$. Dividing by four gives \eqref{p4:eq:v249-paired-influence-zero}.
The quartic coefficient is retained and nonzero in both histories.
If a finite continuous $\Gamma$ with
$Z=\exp(i\Gamma/\hbar)$ existed through one of these zeros, its exponential
would vanish there, which is impossible.
\end{proof}

The amplitude $Z$ itself remains globally defined for every supplied finite
history in its ODE domain. On overlaps of normalized nonzero logarithm
charts, logarithms differ by a locally constant $2\pi\hbar$ integer, so
their derivatives agree. Along the physical diagonal $Z[q,q]=1$, the tube
proposition supplies the local action needed for mean-source variations.
A global logarithm away from that diagonal is a separate, and here
explicitly obstructed, assertion.
% END FRAMEWORK v2.249: CERTIFIED LOG DOMAIN AND NATIVE INFLUENCE ZERO


% v2.249: controlled source-jet convergence of the actual interacting evaluator.
\subsection{Certified first and second source jets of the interacting parent}
\label{p4:sec:v249-interacting-jet-convergence}

The value error in Theorem~\ref{p4:thm:v248-current-determinant-sum}
does not by itself authorize differentiation of its limit. The following
estimate supplies that missing certificate for the complete finite word,
including the arbitrary joint-state extension of
Corollary~\ref{p4:cor:v248-joint-state-interacting-determinants}.

\begin{theorem}[Uniform two-jet convergence of the interacting slicing word]
\label{p4:thm:v249-interacting-two-jet-convergence}
Let $E$ be a declared Banach space of real source histories, for example
$C^1([t_i,t_f],\mathbb R^p)$ with its fixed physical norm. On an open
$\mathcal U\subset E$, let the actual finite-dimensional self-adjoint
split terms satisfy
\[
 H_a:\mathcal U\longrightarrow
       C^1([t_i,t_f],\operatorname{Herm}(\mathcal H)),\qquad
 H_a\in C^2,\qquad H=\sum_{a=1}^mH_a.
\]
The history space must be strong enough to satisfy this displayed target
regularity. For example, retained $\dot q$ moving-frame terms generally
require $C^2$ histories for a $C^1$ time-dependent Hamiltonian; smoothness
of coefficients alone does not imply this property on a $C^1$ history ball.
Include the actual central term among the $H_a$. The fixed carrier
$\mathcal H$ may be the full CAR/boson carrier, and the retained terms
may include its actual current vertices and quartic squares. Let
$\mathcal K\subset\mathcal U$ be a source family with the following
finite, explicitly bounded constants. All source derivative norms are
multilinear operator norms on $E$, and all matrix norms are spectral norms:
\[
 M_{a,j}=\sup_{q\in\mathcal K,t}\|D^jH_a(q,t)\|\quad(j=0,1,2),
 \qquad L_j=\sup_{q\in\mathcal K,t}
                      \|\partial_tD^jH(q,t)\|,
\]
\[
 C_{ab,j}=\sup_{q\in\mathcal K,t}
                    \|D^j[H_a,H_b](q,t)\|\quad(a<b,\ j=0,1,2).
\]
These are assumptions on a specified coefficient family, not a compactness
claim for an infinite-dimensional history ball. Define
\[
 R=\sum_a(2M_{a,1}+M_{a,2}),\qquad
 \widehat L=L_0+2L_1+L_2,\qquad
 \widehat C=\sum_{a<b}(C_{ab,0}+2C_{ab,1}+C_{ab,2}).
\]
For $T=t_f-t_i$, $\delta=T/N$, let $U(q)$ be the exact propagator and
let $U_N(q)$ be the same frozen-time unitary Lie--Trotter word used in
v2.248, in its fixed complete factor order. Then for $j=0,1,2$,
\begin{equation}
 \boxed{\sup_{q\in\mathcal K}\|D^jU_N(q)-D^jU(q)\|
 \le\varepsilon_N:=\frac{T\delta}{2}e^{TR/\hbar}
              \left(\frac{\widehat L}{\hbar}
                           +\frac{\widehat C}{\hbar^2}\right).}
 \label{p4:eq:v249-propagator-two-jet-error}
\end{equation}
For the commutator constants one may use the directly computable bounds
\begin{align*}
 C_{ab,0}&\le2M_{a,0}M_{b,0},\\
 C_{ab,1}&\le2(M_{a,1}M_{b,0}+M_{a,0}M_{b,1}),\\
 C_{ab,2}&\le2(M_{a,2}M_{b,0}+2M_{a,1}M_{b,1}
                                      +M_{a,0}M_{b,2}),
\end{align*}
or smaller bounds evaluated on the actual commutator word.
\end{theorem}
\begin{proof}
For unit source directions $u,v\in E$ define the finite jet lift
\[
 \mathsf J(A)=
 \begin{pmatrix}
 A&A_u&A_v&A_{uv}\\
 0&A&0&A_v\\
 0&0&A&A_u\\
 0&0&0&A
 \end{pmatrix}.
\]
The ordinary product rule gives exactly
$\mathsf J(AB)=\mathsf J(A)\mathsf J(B)$, and therefore
$\mathsf J(e^A)=e^{\mathsf J(A)}$ by the convergent finite matrix power
series. Differentiable dependence of the finite ODE, or the same uniformly
convergent Duhamel series, gives
\[
 i\hbar\partial_t\mathsf J(U(t,t_i))
       =\mathsf J(H(t))\mathsf J(U(t,t_i)),
 \qquad \mathsf J(U(t_i,t_i))=I.
\]
All four displayed jets, including the mixed contact $H_{uv}$, are thus
blocks of one finite propagator. Likewise the product of exponentials of
$\mathsf J(H_a(t_j))$ is exactly $\mathsf J(U_N)$.

The diagonal of $-i\mathsf J(H_a)/\hbar$ is skew-Hermitian.
Its off-diagonal part has norm at most
$(\|(H_a)_u\|+\|(H_a)_v\|+\|(H_a)_{uv}\|)/\hbar$:
the three corresponding fixed nilpotent matrices each have norm one.
Duhamel--Gronwall therefore bounds every lifted exact or split subpropagator
by $e^{\ell R/\hbar}$ over accumulated slice length $\ell$, with the
split lengths weighted by their own $2M_{a,1}+M_{a,2}$. In particular a
full product over total physical time $T$ is bounded by $e^{TR/\hbar}$;
no norm of the diagonal Hermitian $H$ is inserted into this growth factor.

The time derivative of the lifted total generator has norm at most
$\widehat L$. The lifted commutator obeys
$[\mathsf J(H_a),\mathsf J(H_b)]=\mathsf J([H_a,H_b])$, whose norm is
at most $C_{ab,0}+2C_{ab,1}+C_{ab,2}$. The v2.248 Duhamel freezing proof
therefore applies to the lifted matrices with a per-slice bound
\[
 \frac{\delta^2}{2}e^{\delta R/\hbar}
          (\widehat L/\hbar+\widehat C/\hbar^2).
\]
For clarity, the splitting estimate does not require unitarity of the lifted
matrices. For generators $A,B$, differentiate
$e^{(\delta-s)(A+B)}e^{sA}e^{sB}$ and insert
\[
 e^{sA}B-Be^{sA}
       =\int_0^s e^{(s-r)A}[A,B]e^{rA}dr.
\]
Every product in this time triangle has the same total growth bound
$e^{\delta R/\hbar}$. Its area is $\delta^2/2$; induction gives the sum
of pairwise commutator bounds for all factors. Telescoping $N$ slices
then gives \eqref{p4:eq:v249-propagator-two-jet-error} for the lifted
propagators. Each value, first derivative and mixed second derivative
is a compression of their difference, so the same bound holds for each.
Taking the supremum over unit $u,v$ proves the asserted operator derivative
norms. The elementary displayed commutator bounds are the product rule and
$\|[A,B]\|\le2\|A\|\|B\|$.
\end{proof}

\begin{corollary}[Complete interacting influence-jet error for the actual joint state]
\label{p4:cor:v249-joint-state-influence-jet-error}
Keep any supplied positive trace-one joint density $\rho_i$ fixed under
the source variations, and define
\[
 Z=\operatorname{Tr}(U_+\rho_iU_-^\dagger),\qquad
 Z_N=\operatorname{Tr}(U_{+,N}\rho_iU_{-,N}^\dagger).
\]
Use uniform constants of the theorem for both source branches, and give
$E\times E$ the maximum norm. Put
\[
 a_j=\sum_aM_{a,j},\qquad b_1=\frac{Ta_1}{\hbar},\qquad
 b_2=\frac{Ta_2}{\hbar}+b_1^2,
\]
\[
 \eta_0=2\varepsilon_N,\qquad
 \eta_1=2(1+b_1)\varepsilon_N,\qquad
 \eta_2=2(1+2b_1+b_2)\varepsilon_N.
\]
Then, uniformly on the stated two-history source family,
\[
 \boxed{\begin{aligned}
       \|D^jZ_N-D^jZ\|&\le\eta_j\quad(j=0,1,2),\\
       \|DZ\|,\|DZ_N\|&\le z_1:=2b_1,\\
       \|D^2Z\|,\|D^2Z_N\|&\le z_2:=2b_2+2b_1^2.
       \end{aligned}}
\]
Thus the exact finite path/correlation evaluator converges with its complete
first and second source variations, without assuming that the joint state
is a product or remains quasifree during interaction.
\end{corollary}
\begin{proof}
Exact Duhamel variation and unitarity give
$\|DU\|\le b_1$ and
$\|D^2U\|\le Ta_2/\hbar+b_1^2=b_2$. The same bounds hold for $U_N$:
apply Duhamel to each unitary factor and sum its integrated derivative
norms; the sum of its second derivative contacts is $Ta_2/\hbar$, and
the two ordered first-derivative insertions are bounded by $b_1^2$.
Differentiate the trace product. At order zero its error is at most
$2\varepsilon_N$. At first order the two product terms cost at most
$2(1+b_1)\varepsilon_N$. At second order the two second-derivative
terms and two ordered cross terms cost at most
$2(1+2b_1+b_2)\varepsilon_N$. Every estimate uses only
$\|\rho_i\|_1=1$. The same product rule gives $z_1,z_2$.
The determinant or occupied-correlation path sum equals this entire sliced
trace algebraically, so its summed jets satisfy these estimates as well.
\end{proof}

\begin{corollary}[Certified logarithmic action two-jet on a nonzero chart]
\label{p4:cor:v249-interacting-log-jet-error}
Use the certified diagonal disk tube of
Proposition~\ref{p4:prop:v249-uniform-log-tube}, or another connected
chart on which the normalized branches
$\Gamma=-i\hbar\log Z$ and $\Gamma_N=-i\hbar\log Z_N$ already
exist and agree at an equal-history point. Suppose
$|Z|\ge\zeta>0$ and $\eta_0<\zeta/2$. Mere connectedness and
nonvanishing are not being used as a global logarithm-existence theorem.
Define the computable quantities
\[
 r_0={\eta_0\over\zeta},\qquad
 r_1={\eta_1\over\zeta}+{\eta_0z_1\over\zeta^2},\qquad
 r_2={\eta_2\over\zeta}
       +{2\eta_1z_1+\eta_0z_2\over\zeta^2}
       +{2\eta_0z_1^2\over\zeta^3}.
\]
Then
\begin{equation}
 \boxed{\begin{aligned}
 |\Gamma_N-\Gamma|&\le2\hbar r_0,\\
 \|D\Gamma_N-D\Gamma\|&\le2\hbar r_1,\\
 \|D^2\Gamma_N-D^2\Gamma\|&\le\hbar(2r_2+4r_1^2).
 \end{aligned}}
 \label{p4:eq:v249-log-two-jet-error}
\end{equation}
All three errors tend to zero at the certified $O(N^{-1})$ rate with
fixed coefficient bounds and fixed $\zeta$.
\end{corollary}
\begin{proof}
Put $e=Z_N-Z$ and $r=e/Z$. The quotient rule gives
$\|r\|\le r_0<1/2$, $\|Dr\|\le r_1$ and $\|D^2r\|\le r_2$.
The normalized local branches obey
$\log Z_N-\log Z=\log(1+r)$; their difference is fixed by equality
on the diagonal. Since $|1+r|\ge1/2$,
\[
 |\log(1+r)|\le2|r|,\quad
 \|D\log(1+r)\|\le2\|Dr\|,\quad
 \|D^2\log(1+r)\|\le2\|D^2r\|+4\|Dr\|^2.
\]
Multiplying by $\hbar$ proves the bounds.
\end{proof}

\paragraph{Contacts, retarded blocks and the topology of the certificate.}
The second propagator variation used above is exactly
\begin{align*}
 D^2U[u,v]={}&-\frac{i}{\hbar}\int
 U(t_f,t)H_{uv}(t)U(t,t_i)dt\\
 &-\frac{1}{\hbar^2}\int_{s<t}U(t_f,t)
 \bigl[H_u(t)U(t,s)H_v(s)+H_v(t)U(t,s)H_u(s)\bigr]
 U(s,t_i)dsdt.
\end{align*}
Thus neither the direct second variation nor either noncommuting ordered
insertion is discarded. For example, on a varying current vertex its
$H_{uv}$ includes
\[
 d\Gamma(T_{uv})\otimes X+d\Gamma(T_u)\otimes X_v
 +d\Gamma(T_v)\otimes X_u+d\Gamma(T)\otimes X_{uv}.
\]
For $F=d\Gamma(S)$ the second variation of an actual quartic square
$F^2$ includes $F_{uv}F+FF_{uv}+F_uF_v+F_vF_u$, as well as any
actual coefficient derivatives. A central term retains its full first
and second source contacts even though it has no commutator memory.

For a word depending pointwise on the control values $q(t)$, without time
derivatives of those controls, put $B_a=H_{,a}$, $Q_{ab}=H_{,ab}$.
The limit then has the already owned complete mixed block
\[
 \Gamma_{\Delta a,c b}(t,s)
 =-\langle Q_{ab}(t)\rangle\delta(t-s)
 +\frac{i}{\hbar}\Theta(t-s)
          \langle[B_a^H(t),B_b^H(s)]\rangle.
\]
The new result is convergence to this full interacting block with the
stated error after testing against the declared history space. In
particular sampling deltas converge here as continuous bilinear forms
on the $C^1$ test-history domain; no pointwise convergence of delta
kernels, or operator-norm convergence against arbitrary $C^0$ tests,
is asserted. The exact and sliced full two-branch Hessians remain
symmetric with their inherited branch transposition, while the physical
mixed block has retarded support.
For derivative-dependent words, the general $D^2U$ formula above retains
the additional differential contacts; the displayed simple delta-contact
formula is not asserted for such words without their functional variations.

All source-dependent projectors, frames, density factors, coefficients and
boundary contacts consumed by the actual finite word must first be retained
in the $H_a$ and their derivatives on this common fixed carrier. The theorem
does not supply any missing physical incidence or select a state. Individual
spectral projectors or Vandermonde labels in the path representation can
change presentation at degeneracies. The certificate differentiates their
complete smooth operator sum, using Duhamel or the jet lift; it does not
license differentiating an arbitrary discontinuous eigenvector choice.


\subsection{Certified stationary source derivatives and visible-pole clusters}
\label{p4:sec:v249-certified-source-poles}
% v2.249: quantitative bridge from complete finite coefficients to source-visible poles.
\begin{theorem}[Fixed-space realization of the actual stationary source response]
\label{p4:thm:v249-source-hs-realization}
Let an actual licensed finite carrier $\mathcal K$ carry
$H=H^\dagger$, $\rho\succeq0$, $\operatorname{Tr}\rho=1$ and
$[H,\rho]=0$, with its complete inherited source columns $Q_j$.
Equip $\operatorname{End}(\mathcal K)$ with its Hilbert--Schmidt
inner product and put
\[
 T=\operatorname{ad}_H,\qquad \mathbb T=T\oplus T,\qquad
 Je_j=(Q_j\sqrt\rho,\sqrt\rho Q_j).
\]
The exact anticommutator response has the positive realization
\[
 \boxed{G_{rs}(z)=\operatorname{Tr}\rho
 \{Q_r^\dagger,(z-\operatorname{ad}_H)^{-1}Q_s\}
 =[J^\dagger(z-\mathbb T)^{-1}J]_{rs}.}
\]
Its complete residue at a frequency $\omega$ is
$R(\omega)=J^\dagger(\Pi_\omega\oplus\Pi_\omega)J\succeq0$.
In a simultaneous eigenbasis $H|n\rangle=E_n|n\rangle$,
$\rho|n\rangle=p_n|n\rangle$,
\[
 \boxed{R_{rs}(\omega)=
 \sum_{E_m-E_n=\omega}(p_m+p_n)
 \overline{(Q_r)_{mn}}(Q_s)_{mn}.}
\]
Thus a source tag $c$ sees precisely those frequencies with
$c^\dagger R(\omega)c>0$, including degenerate frequencies and
nonfaithful physical states. If the columns are canonical CAR
creation sources, $J^\dagger J=I$.
\end{theorem}
\begin{proof}
Trace cyclicity makes $T$ self-adjoint in the Hilbert--Schmidt
inner product. Stationarity implies $[H,\sqrt\rho]=0$, so
$J(TA)=(T\oplus T)J(A)$ for the map
$J(A)=(A\sqrt\rho,\sqrt\rho A)$. Its inner product is
\[
 \langle J(A),J(B)\rangle
 =\operatorname{Tr}\rho(A^\dagger B+BA^\dagger).
\]
Intertwining the resolvents proves the realization. In the simultaneous
eigenbasis, the two Hilbert--Schmidt components of $J(Q_s)$ have
entries $(Q_s)_{mn}\sqrt{p_n}$ and
$\sqrt{p_m}(Q_s)_{mn}$, while $T$ acts at $E_m-E_n$.
Their Gram gives the printed residue. It is positive and includes a
transition precisely when its populated endpoint and source amplitude
are nonzero. The canonical CAR identity gives $J^\dagger J=I$.
\end{proof}
This constructs the actual positive pairing without replacing $\rho$
by a canonical calculation density. Exact finite termination and the
Hankel-pencil extraction then follow from the already proved
Theorem~\ref{p4:thm:v2188-moment-pencil} on this actual pairing.

\begin{lemma}[Computable source-state perturbation bound]
\label{p4:lem:v249-source-state-error}
For two positive trace-one densities on the same finite carrier,
\[
 \boxed{\|\sqrt\rho-\sqrt{\widetilde\rho}\|_{\rm HS}^2
 \le\|\rho-\widetilde\rho\|_1.}
\]
For one source, the realization columns of the theorem therefore obey
\[
 \boxed{\|J-\widetilde J\|
 \le\sqrt2\|Q-\widetilde Q\|
 +\sqrt2\|\widetilde Q\|
                   \sqrt{\|\rho-\widetilde\rho\|_1}.}
\]
For several columns the same inequality holds with operator norms of
the coefficient maps $c\mapsto\sum_jc_jQ_j$; alternatively
$\|J-\widetilde J\|$ is calculated directly from its finite Gram.
\end{lemma}
\begin{proof}
Put $A=\sqrt\rho$, $B=\sqrt{\widetilde\rho}$,
$D=A-B$ and $S=\operatorname{sign}D$. Since
$A^2-B^2=AD+DB$, cyclicity gives
\[
 \operatorname{Tr}S(A^2-B^2)
 =\operatorname{Tr}|D|(A+B)\ge\operatorname{Tr}D^2.
\]
For the inequality, diagonalize $D$: its diagonal entries satisfy
$A_{ii}-B_{ii}=d_i$ and $A_{ii},B_{ii}\ge0$, so
$A_{ii}+B_{ii}\ge|d_i|$. Trace duality and $\|S\|\le1$
prove the square-root bound. In each of the two realization components,
add and subtract $\widetilde Q\sqrt\rho$ or
$\sqrt\rho\widetilde Q$. The Hilbert--Schmidt product bound,
$\|\sqrt\rho\|_{\rm HS}=1$, the triangle inequality in the doubled
space and the first result give the source bound. Apply it to every
unit coefficient vector for the multicolumn statement.
\end{proof}

\begin{theorem}[Certified spectral derivatives and visible-pole clusters]
\label{p4:thm:v249-certified-visible-poles}
Compare two complete stationary tuples of
Theorem~\ref{p4:thm:v249-source-hs-realization} on the same inherited
finite carrier and frame. Put
\[
 \epsilon_H=\|H-\widetilde H\|,\quad
 \delta=2\epsilon_H,\quad
 \eta=\|J-\widetilde J\|,\quad
 j=\|J\|,\quad\widetilde j=\|\widetilde J\|.
\]
At every $z$ of distance at least $d>0$ from both full Liouvillian
spectra, and for every $k\ge0$,
\[
 \boxed{\|G^{(k)}(z)-\widetilde G^{(k)}(z)\|
 \le k!\left[
 \frac{(j+\widetilde j)\eta}{d^{k+1}}
 +\frac{\widetilde j^2(k+1)\delta}{d^{k+2}}
 \right].}
\]
For nonreal $z$, $d=|\operatorname{Im}z|$ is always valid.
Now let $\widetilde\omega$ be an eigenvalue of the full original
$\widetilde T=\operatorname{ad}_{\widetilde H}$, separated from all
other distinct frequencies by $g>0$. Choose $0<r\le g/2$ and
$\delta<r$. Let $\widetilde\Pi$ be its full spectral projector and
$\Pi$ the actual projector in $|z-\widetilde\omega|<r$. Then
\[
 \boxed{\|\Pi-\widetilde\Pi\|
 \le\frac{\delta}{r-\delta},\quad
 \operatorname{rank}\Pi=\operatorname{rank}\widetilde\Pi,\quad
 |\omega-\widetilde\omega|\le\delta
 \quad(\omega\hbox{ in the actual cluster}).}
\]
For the complete cluster weights
$R=J^\dagger(\Pi\oplus\Pi)J$ and
$\widetilde R=\widetilde J^\dagger
 (\widetilde\Pi\oplus\widetilde\Pi)\widetilde J$,
\[
 \boxed{\|R-\widetilde R\|\le\epsilon_R
 =(j+\widetilde j)\eta+
       \widetilde j^2\frac{\delta}{r-\delta}.}
\]
Hence $c^\dagger\widetilde R c>\epsilon_R\|c\|^2$
certifies at least one actual tag-visible pole within $\delta$ of
$\widetilde\omega$, with total cluster tag weight at least
$c^\dagger\widetilde R c-\epsilon_R\|c\|^2>0$.
If $\operatorname{rank}\widetilde\Pi=1$ on the original
Hilbert--Schmidt carrier, the cluster has one actual frequency and
the same bound applies to its individual residue. A splitting
degenerate cluster is certified by its total weight.
\end{theorem}
\begin{proof}
For a Hilbert--Schmidt matrix $A$,
$\|[H-\widetilde H,A]\|_{\rm HS}
 \le2\epsilon_H\|A\|_{\rm HS}$, so
$\|T-\widetilde T\|\le\delta$, also on the doubled carrier.
Write $R_z=(z-\mathbb T)^{-1}$ and
$\widetilde R_z=(z-\widetilde{\mathbb T})^{-1}$.
The resolvent identity gives
$\|R_z-\widetilde R_z\|\le\delta/d^2$.
The identity
\[
 R_z^n-\widetilde R_z^n
 =\sum_{a=0}^{n-1}R_z^{n-1-a}
 (R_z-\widetilde R_z)\widetilde R_z^a
\]
gives norm at most $n\delta/d^{n+1}$.
Since $G^{(k)}=(-1)^kk!J^\dagger R_z^{k+1}J$,
add and subtract $\widetilde J^\dagger R_z^{k+1}\widetilde J$
and use the triangle inequality to obtain the first bound.

On the circle $|z-\widetilde\omega|=r$, the reference resolvent
of $\widetilde T$ has norm at most $1/r$. The Neumann series for
$z-\widetilde T-s(T-\widetilde T)$, $0\le s\le1$, shows that
the circle is disjoint from every intervening spectrum and gives the
actual resolvent bound $1/(r-\delta)$. Integrating their resolvent
difference around the circle gives
$\|\Pi-\widetilde\Pi\|\le\delta/(r-\delta)$.
The Riesz projectors vary continuously on the homotopy, so their
integer ranks agree. The same Neumann argument gives spectral
inclusion in the $\delta$-neighbourhood of the reference spectrum.
Inside this circle the other reference frequencies are farther than
$\delta$, so every actual cluster frequency lies within $\delta$
of $\widetilde\omega$.

Finally, add and subtract
$\widetilde J^\dagger(\Pi\oplus\Pi)\widetilde J$ to bound the
cluster-weight difference. Positivity of the exact residues converts
a strictly positive lower cluster weight into existence of a visible
pole. Rank one on the original, undoubled carrier gives one actual
frequency; otherwise only the stated cluster assertion is needed.
\end{proof}
The derivative estimate controls the spectral variable $z$. Physical
source-control or metric derivatives retain their complete coefficient,
preparation, frame and endpoint jets. Value-only finite-time evolution
errors do not establish exact pole equality or a visibility margin.

\begin{corollary}[Finite correlation derivatives and the Laplace bridge]
\label{p4:cor:v249-correlation-resolvent-bridge}
In the same source units $\hbar=1$, set
$C(t)=J^\dagger e^{-it\mathbb T}J$ and similarly
$\widetilde C(t)$. Let
$M=\max(\|T\|,\|\widetilde T\|)$. For every integer $k\ge0$,
\[
 \boxed{\|\partial_t^kC(t)-\partial_t^k\widetilde C(t)\|
 \le (j+\widetilde j)\eta M^k
 +\widetilde j^2\delta
       [kM^{k-1}+|t|M^k].}
\]
The first bracketed term is zero when $k=0$; $M^0=1$ also when
$M=0$. For $\operatorname{Im}z>0$, the exact all-time transform is
\[
 \boxed{G(z)=-i\int_0^\infty e^{izt}C(t)\,dt.}
\]
Differentiation under this convergent integral gives the spectral
$z$-derivative certificate of the theorem with
$d=\operatorname{Im}z$.
\end{corollary}
\begin{proof}
Duhamel's identity and unitary norm one give
$\|e^{-it\mathbb T}-e^{-it\widetilde{\mathbb T}}\|
 \le|t|\delta$.
The telescoping identity for powers gives
$\|\mathbb T^k-\widetilde{\mathbb T}^k\|
 \le kM^{k-1}\delta$ for $k\ge1$.
Differentiate the two finite exponentials, add and subtract the
corresponding source and propagator terms, and apply these bounds.
Finite self-adjoint spectral resolution gives the displayed transform
since $-i\int_0^\infty e^{i(z-\omega)t}\,dt=(z-\omega)^{-1}$.
The difference at $k=0$ is bounded by
$(j+\widetilde j)\eta+\widetilde j^2|t|\delta$.
After $k$ derivatives in $z$, absolute integration uses
$\int_0^\infty t^ke^{-yt}\,dt=k!/y^{k+1}$ and yields exactly
the stated resolvent-derivative bound. This justification uses the
complete coefficient-norm bound for all time, rather than only a
value-error bound on a finite time window.
\end{proof}

\begin{corollary}[Explicit certification data for the interacting two-spin source]
\label{p4:cor:v249-paired-coefficient-certification}
For two instances of the complete interacting restriction in
Theorem~\ref{p5c:thm:v248-paired-quartic-source-poles}, put
$\Delta E=\Delta\epsilon_0+\Delta\epsilon_1+\Delta v$.
Their full Hamiltonian error is exactly
\[
 \boxed{\epsilon_H=
 \max\left\{|\Delta\epsilon_0|,|\Delta\epsilon_1|,
 \frac{|\Delta E|}{2}
 +\sqrt{\frac{(\Delta E)^2}{4}+|\Delta g|^2}\right\}.}
\]
For the unchanged canonical source $Q=a_0^\dagger$ and their exact
full Gibbs states at the same finite $\beta\ge0$,
\[
 \boxed{\|\rho_\beta(H)-\rho_\beta(\widetilde H)\|_1
 \le\min(2,2\beta\epsilon_H),\quad
 j=\widetilde j=1,\quad
 \eta\le\min(2,2\sqrt{\beta\epsilon_H}).}
\]
Together with the computed full-Liouvillian gap and the positive
residues of that theorem, these are complete numerical inputs to the
preceding pole and spectral-derivative certificates on this restriction.
\end{corollary}
\begin{proof}
The Hamiltonian error has odd scalar blocks
$\Delta\epsilon_0,\Delta\epsilon_1$ and even block
$\bigl(\begin{smallmatrix}0&\overline{\Delta g}\\
 \Delta g&\Delta E\end{smallmatrix}\bigr)$.
The latter has eigenvalues $\Delta E/2\pm
\sqrt{(\Delta E)^2/4+|\Delta g|^2}$, proving its exact norm.
For $H_s=\widetilde H+s(H-\widetilde H)$, differentiate the finite
matrix exponential and its trace normalization:
\[
 \partial_s\rho_s=-\beta\int_0^1
 \rho_s^{1-u}(H-\widetilde H)\rho_s^u\,du
 +\beta\rho_s\operatorname{Tr}(\rho_s(H-\widetilde H)).
\]
Finite Schatten H\"older gives norm at most $\epsilon_H$ for
$\rho_s^{1-u}(H-\widetilde H)\rho_s^u$ in trace norm,
because the two powers have conjugate Schatten norms one. The
central expectation term has the same bound. Integration in $s$
gives $2\beta\epsilon_H$; trace-one positivity also gives the bound
$2$. The canonical source has norm one and $J^\dagger J=1$.
Lemma~\ref{p4:lem:v249-source-state-error} then gives
$\eta\le2\sqrt{\beta\epsilon_H}$, while the triangle inequality
between two unit realization columns gives $\eta\le2$.
\end{proof}
These statements certify the observable of the complete supplied finite
restriction. They do not supply an unexhibited physical Hamiltonian,
state, source marking or particle/rest dictionary.






% BEGIN INTEGRATED V2.251 FORWARD PACKAGE
% Forward common-packet integration: finite joint responses and exact jet closure.
% Standalone fragment. Uses only standard theorem/proof environments and amsmath.
\subsection{Joint source response and its exact finite jet carrier}
\label{p4:sec:v251-forward-joint-response}

This construction consumes the actual complete finite stationary chart;
it does not select a Hamiltonian, preparation or source member. The fixed
Hilbert--Schmidt realization of
Theorem~\ref{p4:thm:v249-source-hs-realization} supplies a convenient common
carrier. For its actual finite $H(q)=H(q)^\dagger$, stationary density
$\rho(q)\succeq0$, $\operatorname{Tr}\rho(q)=1$, and every consumed raw tag
$Q_a(q)$, put
\[
 \mathbb T(q)=\operatorname{ad}_{H(q)}\oplus\operatorname{ad}_{H(q)},
 \qquad J(q)e_a=(Q_a(q)\sqrt{\rho(q)},\sqrt{\rho(q)}Q_a(q)).
\]
Retain all actual kinetic, field, state, support, frame and boundary
transport. On a fixed finite carrier and a smooth fixed-rank stationary
chart the displayed matrices and their admitted jets are defined. No
maximally mixed calculation state is substituted. The complete matrix
anticommutator response is
\[
 G(q,z)=J(q)^\dagger(zI-\mathbb T(q))^{-1}J(q),\qquad z\notin\mathbb R.
\]
If commutator rather than anticommutator ports are consumed, retain the
two ordered source columns separately,
$J_+e_a=(Q_a\sqrt\rho,0)$ and $J_-e_a=(0,\sqrt\rho Q_a)$.
The appropriate sum or difference of their response blocks then gives
the actual port without replacing its population weights.

For the following finite theorem write $T=\mathbb T$ or the corresponding
actual self-adjoint operator in an already supplied fixed physical pairing.
A moving positive pairing $M(q)$ can instead be transported by
$\widetilde T=M^{1/2}TM^{-1/2}$, $\widetilde J=M^{1/2}J$; every derivative
of that transport must be retained, as in the native finite-metric
constructor. A rank-changing or nonsmooth chart is outside the theorem.

\begin{theorem}[Exact joint cyclic response]
\label{p4:thm:v251-joint-cyclic-response}
Let $T_0=T(q_0)$ be self-adjoint on a $d$-dimensional Hilbert space and
let $J_0$ collect all consumed source columns. Set
\[
 \mathscr C_0=\operatorname{span}\{T_0^nJ_0u:0\le n<d,\ u\in\mathscr S\},
 \qquad P_0=\operatorname{proj}_{\mathscr C_0}.
\]
Then $[P_0,T_0]=0$ and the complete matrix response equals the response
of $T_0|_{\mathscr C_0}$ with the same columns and pairing. Its full
residue matrices, including all cross-tag entries, are preserved.
\end{theorem}
\begin{proof}
Cayley--Hamilton makes the displayed finite span invariant under $T_0$.
Self-adjointness makes its orthogonal complement invariant as well.
The resolvent is therefore block diagonal, while $J_0$ takes values in
the first block. The identity follows by compression; contour integration
gives the same residue matrices. This is the joint-column extension of
Lemma~\ref{p4:lem:v2186-exact-source-cyclic-action}, not a replacement of
its complete-action antecedent.
\end{proof}

The fixed cyclic carrier preserves the static response. It need not
preserve a source or metric jet. The next theorem constructs the finite
carrier that does preserve those jets.

For a multi-index $\alpha$ let $T_\alpha=\partial^\alpha T(q_0)$ and
$J_\alpha=\partial^\alpha J(q_0)$. Put
$\operatorname{cyc}_{T_0}V=\operatorname{span}\{T_0^nv:0\le n<d,v\in V\}$.
Recursively define, for $r\ge1$,
\[
 \mathscr C_r=\operatorname{cyc}_{T_0}\left(
 \mathscr C_{r-1}
 +\sum_{1\le|\alpha|\le r}T_\alpha\mathscr C_{r-|\alpha|}
 +\sum_{|\beta|=r}\operatorname{ran}J_\beta\right).
\]
Equivalently it is the span of all words formed from powers of $T_0$,
derivative insertions and a differentiated source column whose total
derivative order is at most $r$. It is increasing, reducing for $T_0$,
and has dimension at most $d$. Its construction is finite linear algebra
from the supplied finite jets.

\begin{theorem}[Finite carrier preserving every consumed response jet]
\label{p4:thm:v251-finite-jet-cyclic-carrier}
Suppose the actual fixed-pairing matrices $T(q)=T(q)^\dagger,J(q)$ are
$C^k$ near $q_0$. Let $r=\lfloor k/2\rfloor$ and let $P_r$ project onto
$\mathscr C_r$. Define the fixed compressed family
\[
 T^{[r]}(q)=P_rT(q)P_r|_{\mathscr C_r},\quad
 J^{[r]}(q)=P_rJ(q),\quad
 G^{[r]}(q,z)=J^{[r]}(q)^\dagger(zI-T^{[r]}(q))^{-1}J^{[r]}(q).
\]
Then for every $|\alpha|\le k$ and nonreal $z$,
\[
 \boxed{\partial^\alpha G(q_0,z)=\partial^\alpha G^{[r]}(q_0,z).}
\]
In particular $\mathscr C_0$ preserves the value and first jet, and
\[
 \boxed{\mathscr C_1=\operatorname{cyc}_{T_0}
   \left(\mathscr C_0+\sum_iT_i\mathscr C_0
                         +\sum_i\operatorname{ran}J_i\right)}
\]
preserves the entire second and third response jets. The actual higher
contact matrices $T_{ij},J_{ij}$ are still compressed and differentiated;
they are not inferred from first derivatives.
\end{theorem}
\begin{proof}
Let $R_0=(zI-T_0)^{-1}$. Repeated differentiation of
$\partial_iR=RT_iR$ and the product rule expresses an order-$\ell$
response derivative as a finite linear combination of chains
\[
 J_\beta^\dagger R_0T_{\alpha_1}R_0\cdots
                T_{\alpha_m}R_0J_\gamma,
 \qquad |\beta|+\sum_j|\alpha_j|+|\gamma|=\ell,
 \quad |\alpha_j|\ge1.
\]
Every derivative matrix $T_\alpha$ is self-adjoint. Since $P_r$ reduces
$T_0$, each resolvent splits into its $P_r$ and $I-P_r$ blocks. The chain
with only $P_r$ resolvents is precisely the corresponding compressed
derivative chain. Consider any other, potentially nonzero block path.
At its rightmost complementary resolvent, the derivative order accumulated
from the right source must be at least $r+1$: otherwise the recursive
construction puts that vector in $\mathscr C_r$, and the complementary
projection kills it. Reversing adjoints gives the same lower bound
$r+1$ from the left source to its leftmost complementary resolvent.
Those two portions have disjoint derivative insertions. Thus a nonzero
complementary path needs $\ell\ge2r+2$. This is impossible for
$\ell\le k\le2r+1$. All complementary paths vanish, proving the identity.
The argument uses only a finite jet at $q_0$, not constancy of the cyclic
rank at neighboring $q$ or equality of the responses away from $q_0$.
\end{proof}

\begin{corollary}[The derivative-depth bound is sharp in general]
\label{paper4c:corollary:atlas-the-derivative-depth-bound-is-sharp-in-general-36af441}% stable atlas label; provenance v2_277

For each $r\ge0$ there is a finite positive one-particle CAR comparison
whose static source is the first basis vector, with
$\mathscr C_r=\operatorname{span}(e_1,\ldots,e_{r+1})$, and whose first
missed return after that compression occurs at order $2r+2$.
\end{corollary}
\begin{proof}
Take a distinct positive diagonal $T_0$ on $r+2$ sites and let the sole
first derivative be a nonzero nearest-neighbor Hermitian path matrix.
Small parameter values keep the matrix positive. To reach site $r+2$
from the first source and return requires $r+1$ edges in each direction.
The first complementary resolvent chain therefore has order $2r+2$;
its coefficient is nonzero for real energy above the spectrum. The
recursive spaces are the stated initial path segments.
\end{proof}

\begin{corollary}[Covariance on the inherited flat source connection]
\label{p4:cor:v251-jet-carrier-frame-covariance}
Suppose the actual fixed-rank source bundle is supplied with a flat
metric-compatible connection on the local parameter chart. Compute
$\mathscr C_r$ from its covariant jets in a parallel orthonormal frame.
Under a smooth unitary frame change $T'=UTU^\dagger$, $J'=UJ$, and
$A_i'=UA_iU^\dagger-(\partial_iU)U^\dagger$, the background carrier
transforms as $\mathscr C_r'=U(q_0)\mathscr C_r$. Its parallel projector
family and complete compressed response retain the same $k$-jet.
\end{corollary}
\begin{proof}
The covariant derivatives $\nabla_iT=\partial_iT+[A_i,T]$ and
$\nabla_iJ=\partial_iJ+A_iJ$ transform by $U$; the same is true of
their iterates. Every defining finite word for $\mathscr C_r$ therefore
transforms by $U(q_0)$. In the parallel frame the preceding theorem
applies directly. Transport its projector and all induced connection and
source terms to the new frame; equality of the scalar source response
then preserves its ordinary parameter derivatives.
\end{proof}

Without a specified transport, $\mathscr C_r$ is a sufficient carrier
in the declared fixed frame, not a uniquely selected frame-free
subbundle. Freezing a projector after a parameter-dependent frame change
would omit its moving-frame contacts. A general nonflat source bundle
retains its own curvature/transport data; flatness is not inferred here.

\begin{proposition}[Exact second-jet loss and a necessary/sufficient block test]
\label{p4:prop:v251-static-cyclic-jet-defect}
Let $P=P_0$, $Q=I-P$, $A=T_0|_{\mathscr C_0}$ and
$D=T_0|_{\mathscr C_0^\perp}$. Compare the actual response with the family
obtained by keeping only the fixed $P$ block. Define
\[
 V_i(z)=QJ_i+QT_iP(zI-A)^{-1}J_0.
\]
Their second-derivative difference is exactly
\[
 \boxed{\Delta G_{ij}(z)=
 V_i(\bar z)^\dagger(zI-D)^{-1}V_j(z)
 +V_j(\bar z)^\dagger(zI-D)^{-1}V_i(z).}
\]
Consequently the fixed static cyclic compression preserves all second
response derivatives in the consumed directions, as functions of energy,
if and only if
\[
 \boxed{QJ_i=0,\qquad QT_iP=0\quad\hbox{for every consumed direction }i.}
\]
\end{proposition}
\begin{proof}
Differentiate $G=J^\dagger RJ$ twice. Terms involving $J_{ij}$ or one
$T_{ij}$ have unperturbed cyclic source vectors on their other ends and
therefore agree after compression. The remaining complementary terms are
the two $J_i^\dagger R_DJ_j$ terms, four source/interaction cross terms,
and two $T_iR_DT_j$ returns. Expanding the displayed factorization gives
exactly these eight terms with their inherited signs and pairings.

For real $z>\|T_0\|$ the complementary resolvent is positive and
\[
 \Delta G_{ii}(z)=2V_i(z)^\dagger(zI-D)^{-1}V_i(z)\succeq0.
\]
If every pure second jet agrees as a meromorphic function, this positive
matrix vanishes for every such $z$, hence $V_i(z)=0$. Letting $z\to\infty$
gives $QJ_i=0$. The remaining Laurent coefficients give
$QT_iPA^nJ_0=0$ for every $n$. Cyclicity of $J_0$ for $A$ gives
$QT_iP=0$. Conversely those block identities give every $V_i=0$ and
therefore the entire second jet. Equality on nonreal energies transfers
to the regular real interval by rational continuation.
\end{proof}

For reference, the full first two derivatives on the transported fixed
pairing are
\begin{align*}
 G_i={}&J_i^\dagger RJ+J^\dagger RJ_i+J^\dagger RT_iRJ,\\
 G_{ij}={}&J_{ij}^\dagger RJ+J^\dagger RJ_{ij}
   +J_i^\dagger RJ_j+J_j^\dagger RJ_i\\
 &+J_i^\dagger RT_jRJ+J_j^\dagger RT_iRJ
   +J^\dagger RT_iRJ_j+J^\dagger RT_jRJ_i\\
 &+J^\dagger RT_{ij}RJ
   +J^\dagger RT_iRT_jRJ+J^\dagger RT_jRT_iRJ.
\end{align*}
These preserve the actual moving tags, state/pairing and direct contacts.
They are stationary response derivatives on the admitted family. They
are not a substitute for the separately generated nonstationary causal
influence functional and its actual preparation convention.

\begin{proposition}[A dark sector returns in an exact finite CAR example]
\label{p4:prop:v251-dark-sector-second-return}
On a supplied two-mode finite CAR chart take
\[
 H(q)=d\Gamma\begin{pmatrix}2&3q\\3q&5\end{pmatrix},\quad |q|<1,
 \qquad \rho=|0\rangle\langle0|,\qquad Q_1=a_1^\dagger.
\]
Its vacuum addition response is
\[
 G(q,z)=\frac{z-5}{(z-2)(z-5)-9q^2}.
\]
At $q=0$ its exact static cyclic carrier is one-dimensional and sees
only the pole $2$. Nevertheless
\[
 G_q(0,z)=0,\qquad
 \boxed{G_{qq}(0,z)=\frac{18}{(z-2)^2(z-5)}\ne0.}
\]
The static compression falsely gives zero second derivative; the new
carrier $\mathscr C_1$ includes the second mode and is exact. A joint
packet containing both actual addition tags already has the full two-mode
static cyclic carrier. This is a comparison example within the declared
finite CAR class, not an identification of the physical electron.
\end{proposition}
\begin{proof}
The one-particle matrix is positive for $|q|<1$, so the vacuum is an actual
stationary ground of this supplied CAR example. The resolvent inverse and
two derivatives give the displayed formulas. At $q=0$, $T_1e_1=3e_2$;
the recursive definition therefore includes both eigenlines at depth one.
\end{proof}

\begin{remark}[Absolute energy remains in the common parent]
For fixed tags and physical state, $H(q)\mapsto H(q)+c(q)I$ leaves
$\operatorname{ad}_{H(q)}$, and hence all the displayed matter response
jets, unchanged. The actual metric/source energy and contacts instead
shift by $c,c_i,c_{ij}$. For example a central $c(q)=q^2$ is invisible
to this complete stationary pole family while contributing $2$ to its
absolute energy second derivative. The action, central allocation and
metric/source ports must therefore remain upstairs in the common packet;
the response carrier is an exact observable-specific reduction only.
\end{remark}

% END INTEGRATED V2.251 FORWARD PACKAGE


% BEGIN V2.254 PASTERSKI-TECHNICAL-OWNERSHIP
\section{Finite detector charge covariance, retained static noise and angular inversion}
\label{p4c:sec:v254-detector-technical}
The physical target literature is
\cite{MoultNarayananPasterski2026MemoryCorrelators,OertelMoultPasterski2026Detectors,PasterskiStromingerZhiboedov2016Memory,JorstadPasterski2026SoftOmissions,Pasterski2023NoetherCurrents}.
We use no holographic premise. The following identities are independent
finite operator and round-sphere derivations; where physical angular
charges have not been constituted, they remain \emph{conditional
applications}, not physical closure claims.

\begin{proposition}[Nonsharp incoming-charge identity and source-residual bound]
\label{p4c:prop:v254-nonsharp-ward-covariance}
On the finite operator domain of Paper~IV-B,
Theorem~\ref{p4b:thm:v254-in-in-charge-detector-covariance},
let $B_f=Q^-_f+R_f$ and use its centered, symmetrized covariance.
Without any sharp-charge assumption, one has exactly
\begin{align*}
\operatorname{Cov}_\rho(Q^{S,+}_f,Q^{S,+}_g)
-\operatorname{Cov}_\rho(Q^{H,+}_f,Q^{H,+}_g)
={}&\operatorname{Cov}_\rho(B_f,B_g)
-\operatorname{Cov}_\rho(Q^{H,+}_f,B_g)\\
&-\operatorname{Cov}_\rho(B_f,Q^{H,+}_g).
\end{align*}
If both incoming charges are sharp, all their covariance and cross
terms vanish, recovering Equation~\eqref{p4b:eq:v254-charge-covariance}.
If, moreover, the true total endpoint charge is the canonical Ward
charge, its Heisenberg residual is bounded by the integral of the
actual source-history variation and the full Ward defect as in
Corollary~\ref{p4b:cor:v254-ward-operator-residual}.
\end{proposition}
\begin{proof}
The charge residual identity gives
$Q^{S,+}_f=-Q^{H,+}_f+B_f$. Expand the bilinear covariance
without commuting the Hermitian charges; its symmetrization removes
operator-order ambiguity. Positivity gives
$|\operatorname{Cov}_\rho(A,B)|\le
\sqrt{\operatorname{Cov}_\rho(A,A)\operatorname{Cov}_\rho(B,B)}$.
For a sharp charge $Q^-_f$, the centered operator annihilates the
support of $\rho$ in the $\rho$-weighted seminorm. The derivative
$\frac d{dt}(U^\dagger GU)=U^\dagger\mathfrak d_tG U$
then establishes the charge-transport norm estimate in IV-B;
if $\delta H[R^{\rm sym}]$ does not vanish, it must be counted.
\end{proof}

\begin{theorem}[Finite stationary static-noise retention and retarded-only no-go]
\label{p4c:thm:v254-zero-mode-no-omission}
Let $H_L$ be a time-independent finite self-adjoint operator and
$\rho_L\ge0$ a density matrix satisfying $[\rho_L,H_L]=0$.
Write $P_E$ for its distinct-energy spectral projections, and for
an \emph{action-owned} Hermitian insertion $J_f$ let
$\widetilde J_f=J_f-\operatorname{Tr}(\rho_LJ_f)I$ and
$\mathsf P_0(J)=\sum_E P_EJP_E$.
Then with $J_f(t)=e^{iH_Lt/\hbar}J_fe^{-iH_Lt/\hbar}$,
\[
\boxed{\ 
\lim_{T\to\infty}\frac1{2T}\int_{-T}^T
\frac12\operatorname{Tr}\rho_L
\{\widetilde J_f(t),\widetilde J_f(0)\}\,dt
=\operatorname{Tr}\rho_L[\mathsf P_0(\widetilde J_f)]^2\ge0.
\ }
\]
Polarization gives the positive-semidefinite two-test form
\[
C^{(0)}_{fg}=\tfrac12\operatorname{Tr}\rho_L
\{\mathsf P_0(\widetilde J_f),\mathsf P_0(\widetilde J_g)\}.
\]
The retarded commutator kernel alone does not fix this form:
$H_L=0$, $\rho_L=I/2$, $J=\sigma_z$ yield a zero retarded
commutator for all times, but $C^{(0)}_{JJ}=1$.
\end{theorem}
\begin{proof}
Expand $\widetilde J_f(t)$ in finite spectral matrix blocks
$P_E\widetilde J_fP_{E'}$ with phases
$e^{it(E-E')/\hbar}$. Symmetric Ces\`aro averaging annihilates
all terms with $E\neq E'$. Since $[\rho_L,H_L]=0$, the
remaining diagonal term equals
$\operatorname{Tr}\rho_L[\mathsf P_0(\widetilde J_f)]^2$.
This is nonnegative by positivity of $\rho_L$; polarization yields
the bilinear form. In the displayed two-dimensional example
$J(t)=J$, so $[J(t),J]=0$ while
$\operatorname{Tr}(I\sigma_z^2/2)=1$.
\end{proof}

\begin{remark}[Zero energy at finite resolution is not a soft theorem]
\cite{JorstadPasterski2026SoftOmissions} demonstrates that an
asymptotic boundary-correlator construction may miss a soft
contribution when only ordinary scattering data are retained.
The proven $\mathsf P_0$ statement above concerns a \emph{finite
stationary} spectral diagonal and is neither a derivation of their
noncontact angular boundary term nor an interchange of
$T\to\infty$, $L\to\infty$, detector-radius $r\to\infty$ and
$\omega\to0$. Uniform limit and action-generated radiative-operator
hypotheses must be proved separately; see Paper~VII-C,
Remark~\ref{p12:rem:v254-static-noise-limit-gate}.
\end{remark}

\begin{proposition}[Round angular tracefree-Hessian inversion for $\ell\ge2$]
\label{p4c:prop:v254-angular-spin-two-inverse}
On the local projective null-direction sphere from Papers~II-E and
III-E, suppose the \emph{unit-round} angular metric $q_{AB}$ is
specified as a comparison/readout geometry. With its covariant
derivative $D_A$ and Laplacian $\Delta=D^AD_A$, let
\[
\mathcal C_{AB}(c)=D_AD_Bc-\tfrac12q_{AB}\Delta c.
\]
Then the exact identity and eigenvalues are
\[
\boxed{\ D^AD^B\mathcal C_{AB}(c)
=\tfrac12\Delta(\Delta+2)c,\qquad
\lambda_\ell=\tfrac12(\ell-1)\ell(\ell+1)(\ell+2).\ }
\]
For every smooth source $F$ with no $\ell=0,1$ components,
$D^AD^B\mathcal C_{AB}(c)=F$ has the unique solution
modulo those components,
\[
\boxed{\ c_{\ell m}=
\frac{2F_{\ell m}}{(\ell-1)\ell(\ell+1)(\ell+2)},
\qquad\ell\ge2.\ }
\]
No physical gravitational memory normalization is assumed.
\end{proposition}
\begin{proof}
On the unit two-sphere,
$\operatorname{Ric}_{AB}=q_{AB}$, so commuting the
covariant derivatives on the one-form $D_Ac$ yields
$D^AD^BD_AD_Bc=\Delta^2c+\Delta c$. Subtract
$\tfrac12\Delta^2c$ and expand in harmonics with
$\Delta Y_{\ell m}=-\ell(\ell+1)Y_{\ell m}$.
The eigenvalues vanish only for $\ell=0,1$ and are strictly
positive for $\ell\ge2$.
\end{proof}

\begin{remark}[The remaining physical-memory firewall]
The projective $S^2$ does not supply a globally asymptotically flat
null boundary, gravitational soft charge, angle-resolved total energy
conservation, or the normalization of $f(R)$ spin-two and scalaron
responses. The gravitational memory and spin-memory detector
interpretations in \cite{PasterskiStromingerZhiboedov2016Memory}
are concrete targets; neither follows just from the angular inverse.
The $\ell=0,1$ kernel must \emph{not} be interpreted as an absence
of physical momentum/charge information. The hard and soft operators
in IV-B must be co-consumed by one complete action, not defined by
requiring their covariance equality.
\end{remark}
% END V2.254 PASTERSKI-TECHNICAL-OWNERSHIP

% BEGIN FRAMEWORK v2.320: PAPER IV-C PETERSEN-HOPF EXACT CERTIFICATES
\section{Petersen--Hopf quantum comparison certificates}
\label{p4c:sec:v320-petersen-certificates}
\noindent\textit{Technical scope.} This section supplies explicit
coordinate and operator certificates for the conditional Paper~IV-A
and IV-B Petersen comparison results. The input ten Hopf axes and
all finite maps are imported without alteration from Papers~I-A--III.
The packaged executable script
\texttt{Paper4C-PetersenQuantum-v1\_1-verify.py} uses exact
$\mathbb Q(\sqrt5)$ arithmetic and finite graph enumeration. Its
success establishes stated finite identities, not physical action
ownership or an apparatus realization.

\begin{proposition}[Finite signed $S_5$ parity-twist certificate]
\label{p4c:prop:v320-s5-signed-lift}
Identify the ten vertices with unordered pairs
$v=\{i,j\}\subset\{1,2,3,4,5\}$, and set
$S_{vw}=\operatorname{sgn}(u_v\cdot u_w)$ for disjoint pairs,
zero otherwise. For every $\pi\in S_5$ there exists a sign function
$\epsilon_{\pi}:V\to\{\pm1\}$ for which the signed permutation
$M_\pi e_v=\epsilon_\pi(v)e_{\pi v}$ obeys
\[
 \boxed{M_\pi^*SM_\pi=\operatorname{sgn}(\pi)S.}
\]
In particular, $M_\pi^*P_+M_\pi=P_-$ for odd $\pi$, while
$M_\pi^*P_0M_\pi=P_0$. The sign-lift existence holds for all 120
permutations; it makes no claim that an arbitrary choice of the lifts
forms a strictly linear $S_5$ representation, nor that odd lifts
are physical space rotations.
\end{proposition}
\begin{proof}
Choose a vertex $v_0$ and prescribe $\epsilon(v_0)=1$.
For each tree edge $v\sim w$, propagate
\[
 \epsilon(w)=\operatorname{sgn}(\pi)
 S_{vw}S_{\pi v,\pi w}\epsilon(v).
\]
The tree fixes all ten signs and leaves six nontree edges. Direct
substitution of the ten unit vectors
\[
 \begin{gathered}
 (1,1,1),\ (1,1,-1),\ (1,-1,1),\ (1,-1,-1),\\
 (\phi^{-1},-\phi,0),\ (\phi,0,-\phi^{-1}),\
 (0,\phi^{-1},-\phi),\\
 (0,\phi^{-1},\phi),\ (\phi,0,\phi^{-1}),\
 (\phi^{-1},\phi,0)
 \end{gathered}
 \qquad\text{each divided by }\sqrt3,
\]
into the six remaining sign equations gives no contradiction for
any of the 60 even or 60 odd permutations. Since $M_\pi$ is
orthogonal, substituting $M_\pi^*SM_\pi=\pm S$ into the polynomial
expressions for $P_\pm,P_0$ proves the projector transformation.
This is finite exhaustive enumeration, recorded independently in the
verification script.
\end{proof}

\begin{proposition}[Exact transfer instrument and pole/resolvent
certificate]
\label{p4c:prop:v320-transfer-proof}
For every one of the ten two-core selections, the map
$T_{ij}=J_i+J_j$ satisfies
$T_{ij}^*T_{ij}=4P_a+\tfrac14P_t$,
$T_{S^c}=-T_{ij}$. Consequently
\[
 \left(\frac{T_{ij}}2\right)^*\left(\frac{T_{ij}}2\right)
 =P_a+\frac1{16}P_t,
 \qquad
 \sqrt{I-\tfrac14T_{ij}^*T_{ij}}
 =\frac{\sqrt{15}}4P_t.
\]
The six-dimensional auxiliary block
$H_T=\begin{psmallmatrix}0&T^*\\T&0\end{psmallmatrix}$
has spectral polynomial
$(z^2-4)(z^2-1/4)^2$. At each regular $z$ the diagonal resolvent
is
\[
 [(z-H_T)^{-1}]_{11}
 =\frac{z}{z^2-4}P_a+\frac{z}{z^2-1/4}P_t.
\]
In particular the subset complement changes off-diagonal
amplitudes by a sign without changing the characteristic polynomial,
the diagonal resolvent or the CP success operation.
\end{proposition}
\begin{proof}
Functional calculus of the two orthogonal projections reduces each
formula to elementary scalar identities. The eigenvalues of
$T^*T$ are $4,1/4,1/4$, so the singular values of $T$ are
$2,1/2,1/2$. For the off-diagonal matrix, the Schur complement
or $H_T^2=\operatorname{diag}(T^*T,TT^*)$ gives the characteristic
polynomial and resolvent. The sign transformation is conjugation by
$\operatorname{diag}(I_3,-I_3)$. All ten choices are checked exactly.
\end{proof}

\begin{proposition}[Two explicit and independent readout-obstruction
certificates]
\label{p4c:prop:v320-quantum-readout-no-go}
(i) For the ten axis effects $F_v=\frac3{10}u_vu_v^{\mathsf T}$,
$\operatorname{span}_{\mathbb R}\{F_v\}
=\operatorname{Sym}(3,\mathbb R)$, and
$\operatorname{Tr}(F_v iA)=0$ for every real antisymmetric $A$.
Every real-Kraus history preserves this orthogonality.

(ii) For two copies, the 100 product effects span dimension 36 and
annihilate the nine-dimensional real symmetric space
$\Lambda^2\mathbb R^3\otimes\Lambda^2\mathbb R^3$.
The explicit positive states
$R_\pm=I_9/9\pm(A\otimes A)/18$ for
$A_{12}=1=-A_{21}$ exhibit nonseparation by all products, while
$\operatorname{Tr}[(R_+-R_-)(A\otimes A)]=4/9$.
The full complex Hermitian nine-state comparison space has dimension
81, of which 45 directions are orthogonal to the product-effect span.
\end{proposition}
\begin{proof}
The $Q_v=u_vu_v^{\mathsf T}-I/3$ span the traceless real symmetric
subspace by the Paper~II-C tight-frame identity. Adding their sum
recovers $I_3$, giving the six-dimensional span. Traces of products
of symmetric and antisymmetric real matrices vanish.
For real matrices the decomposition
\[
 \operatorname{Sym}^2(V\otimes W)
 \simeq(\operatorname{Sym}^2V\otimes\operatorname{Sym}^2W)
 \oplus(\Lambda^2V\otimes\Lambda^2W)
\]
has respective dimensions 36 and 9 when $\dim V=\dim W=3$.
$A\otimes A$ is symmetric with spectrum in $\{-1,0,1\}$;
thus $R_\pm\succeq0$. Finally
$\operatorname{Tr}(A^2)=-2$ gives the displayed distinction.
\end{proof}

\begin{remark}[Claim ownership and downstream interface tests]
\label{p4c:rem:v320-quantum-audit}
The signed five-cycle holonomies are nontrivial as a $\mathbb Z_2$
line-transport class. Their comparison requires coherent path
transport; it is not automatically the physical Berry connection.
Similarly the auxiliary $H_T$ is not the complete action-generated
Hamiltonian, and the ten-axis and reciprocal product POVMs are not
physical instruments until their test-family and detector occurrence
are established. An action-visible Petersen mass or energy requires
an explicit source-owning injection and source-cyclic residue test
in Paper~IV-B; none is furnished by finite graph eigenvalues alone.
The existing $m_p=1$ carrier and
$\mathcal A_{\rm kin}\cong\mathbb C\oplus M_2\oplus M_2\oplus M_4
\oplus M_3\oplus M_3$ are unchanged by these auxiliary comparisons.
\end{remark}
% END FRAMEWORK v2.320: PAPER IV-C PETERSEN-HOPF EXACT CERTIFICATES



% BEGIN FRAMEWORK v2.347: FINITE CUTOFF MOMENTS AND MULTIMODE COVARIANCE
\section{Exact vacuum-moment cutoff threshold and multimode covariance}
\label{p4c:sec:v347-cutoff-covariance}
These results separate the occupation-truncated matrix representation
from the regular-CCR Gaussian oscillator.  They do not infer that a
physical Born--Hopf branch has selected a specific occupation cap, a
quadrature detector, or a state preparation.

\begin{theorem}[First cutoff-sensitive vacuum moment and cumulant]
\label{p4c:thm:v347-first-cutoff-cumulant}
For a single oscillator let \(N\ge1\), and let \(a_N\) be the
occupation-truncated ladder operator on
\(\operatorname{span}\{|0\rangle,\ldots,|N\rangle\}\) with
\(a_N^\dagger|N\rangle=0\).  Put
\[
 q_N=\sigma(a_N+a_N^\dagger),\qquad
 \sigma^2=\frac{\hbar}{2\omega},\qquad \omega>0.
\]
For every \(1\le k\le N\), the exact vacuum moments equal those of a
centred regular-CCR Gaussian of variance \(\sigma^2\):
\begin{equation}\label{p4c:eq:v347-low-cutoff-moments}
 \boxed{\langle0|q_N^{2k}|0\rangle=(2k-1)!!\,\sigma^{2k}.}
\end{equation}
The first affected even moment is
\begin{equation}\label{p4c:eq:v347-first-cutoff-moment}
 \boxed{\langle0|q_N^{2N+2}|0\rangle
 =\big[(2N+1)!!-(N+1)!\big]\sigma^{2N+2}.}
\end{equation}
All nonzero lower connected cumulants beyond order two vanish, while
the first non-Gaussian cumulant is exactly
\begin{equation}\label{p4c:eq:v347-first-cutoff-cumulant}
 \boxed{\kappa_{2N+2}(q_N)=-(N+1)!\,\sigma^{2N+2}.}
\end{equation}
The identities are about the vacuum of the stated ladder system, not
about the ground state of an arbitrary quadratic polynomial in the
\emph{truncated} matrices.
\end{theorem}
\begin{proof}
Expand \((a_N+a_N^\dagger)^{2k}\) into words.  A nonzero vacuum
matrix element is a walk of \(2k\) unit raising/lowering steps on
nonnegative occupation numbers, beginning and ending at zero; each
step has its exact ladder weight.  Such a walk has maximum height
at most \(k\).  If \(k\le N\), every contributing word survives the
cutoff, so the Gaussian Wick count \((2k-1)!!\) is unchanged.
For \(k=N+1\), exactly one closed walk reaches height \(N+1\):
\(N+1\) upward steps followed by \(N+1\) downward steps.  Its
weight is \((N+1)!\).  Removing it gives the displayed deficit.
Odd moments vanish by occupation parity.  The moment--cumulant
triangular relations and equality of all lower moments with the
Gaussian reference imply that the first changed cumulant equals
the difference of the \((2N+2)\)-th moments.
\end{proof}

\begin{proposition}[Projection and polynomial multiplication do not commute]
\label{p4c:prop:v347-cutoff-ordering}
The fixed finite occupation representation is not a faithful
homomorphism of the regular CCR polynomial algebra.  Already for
\(N=1\), in the basis \(|0\rangle,|1\rangle\),
\[
 q_1=\sqrt{\frac{\hbar}{2\omega}}\,\sigma_x,
 \qquad p_1=\sqrt{\frac{\hbar\omega}{2}}\,\sigma_y,
 \qquad q_1^2=\frac{\hbar}{2\omega}I,\quad
 p_1^2=\frac{\hbar\omega}{2}I.
\]
Hence the explicitly substituted matrix Hamiltonian
\(H_1=\frac12(g_q q_1^2+g_p p_1^2)\), \(g_q,g_p>0\), equals
\[
 \boxed{H_1=\frac\hbar4\Big(\frac{g_q}{\omega}+g_p\omega\Big)I,}
\]
and has a degenerate ground space.  By contrast the positive
regular-CCR quadratic Hamiltonian has a unique Gaussian ground state.
Thus a positive classical quadratic form and Williamson's theorem do
not imply an identical squeezed ground covariance for arbitrary
finite occupation-truncated matrix substitutions.
\end{proposition}
\begin{proof}
The two truncated ladder matrices are \(a_1=|0\rangle\langle1|\)
and \(a_1^\dagger=|1\rangle\langle0|\).  Substitution gives the
Pauli matrices and their scalar squares.  The regular-CCR positive
oscillator has nonzero spacing, whereas \(H_1\) is proportional to
the identity.  In particular, compression of a polynomial and
multiplication of compressed variables cannot be identified.
\end{proof}

\begin{theorem}[Basis-independent multimode regular-CCR ground covariance]
\label{p4c:thm:v347-multimode-covariance}
Let \(G=G^{\mathsf T}>0\) on \(\mathbb R^{2m}\) with canonical
symplectic matrix \(\mathbb J\), \(\mathbb J^2=-I\), and set
\[
 D=\mathbb JG,\quad
 R=(-D^2)^{1/2},\quad
 \mathcal J_G=DR^{-1},\quad
 \Gamma_G=-\frac\hbar2\mathcal J_G\mathbb J.
\]
Here \(R\) denotes the \emph{principal matrix square root}:
\(-D^2\) is similar to a real positive-definite matrix, so this
functional calculus is well defined even though \(-D^2\) need not
be symmetric in the original Euclidean coordinates.  Then
\begin{equation}\label{p4c:eq:v347-multimode-covariance}
 \boxed{\mathcal J_G^2=-I,\qquad
 \Gamma_G=\Gamma_G^{\mathsf T}>0,\qquad
 \Gamma_G\mathbb J\Gamma_G=\frac{\hbar^2}{4}\mathbb J.}
\end{equation}
These are the compatible positive-frequency complex structure and
pure centred Gaussian ground covariance in the \emph{regular-CCR}
representation.  At \(m=1\), they reproduce
\(\Gamma_G=\hbar\sqrt{\det G}\,G^{-1}/2\).
\end{theorem}
\begin{proof}
By Williamson's theorem there is a real symplectic \(S\) such that
\(G=S^{\mathsf T}G_0S\),
\(G_0=\operatorname{diag}(\Omega,\Omega)\),
\(\Omega=\operatorname{diag}(\omega_1,\ldots,\omega_m)>0\).
The symplectic identity \(\mathbb JS^{\mathsf T}=S^{-1}\mathbb J\)
gives \(D=S^{-1}(\mathbb JG_0)S\).  Since \(G_0\) commutes with
\(\mathbb J\), its principal square root yields
\(R=S^{-1}G_0S\).  Consequently
\[
 \mathcal J_G=S^{-1}\mathbb JS,\qquad
 \Gamma_G=\frac\hbar2S^{-1}S^{-\mathsf T}.
\]
The latter matrix is symmetric positive.  Its symplectic purity
identity follows from \(S^{-1}\mathbb JS^{-\mathsf T}=\mathbb J\).
For one mode use \((\mathbb JG)^2=-(\det G)I\) and the identity
\(-\mathbb JG\mathbb J=(\det G)G^{-1}\) in dimension two.
The result is not asserted for a truncated CCR representation.
\end{proof}

\begin{remark}[Physical readout boundary]
\label{p4c:rem:v347-cumulant-license}
A high-order vacuum cumulant is a potential cutoff diagnostic only after
an action-selected bosonic occupation cap, the appropriate ground or
vacuum preparation, a physically specified quadrature source, and an
actual higher-moment readout are certified.  Finite moments by
themselves do not identify that cap with a physical length, particle
mass, cosmological power spectrum, or experimental threshold.
\end{remark}
% END FRAMEWORK v2.347

% BEGIN FRAMEWORK v2.349: EXACT FINITE HERMITE VACUUM LAW
\clearpage
\section{Complete truncated-oscillator vacuum law and multimode cutoff diagnosis}
\label{p4c:sec:v349-hermite-vacuum-law}
The next two statements concern exactly the sharply
occupation-truncated oscillator of the preceding section;
no physical occupation cap or detector is added.

\begin{theorem}[All-order finite vacuum quadrature distribution and resolvent]
\label{p4c:thm:v349-hermite-vacuum-law}
Let $a_N|n\rangle=\sqrt n\,|n-1\rangle$ on
$0\leq n\leq N$, $N\ge1$, and put
$q_N=\sigma(a_N+a_N^\dagger)$,
$\sigma^2=\hbar/(2\omega)>0$.
Let $\mathrm{He}_0=1$, $\mathrm{He}_1=x$, and
$\mathrm{He}_{k+1}=x\mathrm{He}_k-k\mathrm{He}_{k-1}$
be the \emph{probabilists'} Hermite polynomials.  Then
\[
 \det(xI-q_N/\sigma)=\mathrm{He}_{N+1}(x).
\]
If $x_1,\ldots,x_{N+1}$ are the simple real roots of
$\mathrm{He}_{N+1}$, the complete vacuum spectral measure is
\[
 d\mu_N(x)=\sum_{j=1}^{N+1}
 \frac{N!}{(N+1)\mathrm{He}_N(x_j)^2}
 \delta_{\sigma x_j}(dx).
\]
Every vacuum Borel observable and its characteristic function
are obtained from this positive atomic measure.  For the
\emph{associated} polynomials $H_0^{[1]}=1$, $H_1^{[1]}=x$,
$H_{k+1}^{[1]}=xH_k^{[1]}-(k+1)H_{k-1}^{[1]}$,
\[
 \langle0|(z-q_N)^{-1}|0\rangle
 =\frac1\sigma\frac{H_N^{[1]}(z/\sigma)}
 {\mathrm{He}_{N+1}(z/\sigma)}.
\]
The numerator is the trailing principal Jacobi minor;
for $N\ge2$ it is \emph{not} the ordinary $\mathrm{He}_N$.
The measure reproduces all Gaussian moments through degree
$2N+1$ and tends weakly to the normal law of variance
$\sigma^2$ as $N\to\infty$.
\end{theorem}
\begin{proof}
The vacuum Jacobi matrix of $q_N/\sigma$ has zero diagonal
and off-diagonal entries $\sqrt1,\ldots,\sqrt N$.
Laplace expansion of its characteristic determinant gives the
Hermite recursion.  The Christoffel--Darboux identity
normalizes the first-coordinate eigenvector weight as
$N!/[(N+1)\mathrm{He}_N(x_j)^2]$.
Cramer's rule expresses the $(0,0)$ resolvent entry as the
trailing minor divided by the full determinant, establishing
the associated-Hermite recurrence.  These are precisely the
$(N+1)$-point Gauss--Hermite quadrature weights, exact for
polynomial degree at most $2N+1$.  The Gaussian moment problem
is determinate and the second moments are bounded, giving
weak convergence.
\end{proof}

\begin{theorem}[First independent multimode non-Gaussian signal cannot cancel]
\label{p4c:thm:v349-multimode-first-cumulant}
Take finitely many independent tensor-factor truncated modes
in their product vacuum, real coefficients $c_a$ not all zero, widths
$\sigma_a>0$ and occupation caps $N_a\ge1$.  For
$Q=\sum_a c_aq_{N_a}^{(a)}$ and
$m=\min_{c_a\ne0}(N_a+1)$, one has
\[
 \phi_Q(t)=\prod_a\phi_{N_a}(c_at),\qquad
 \boxed{\kappa_{2m}(Q)=
 -m!\sum_{a:N_a+1=m}(c_a\sigma_a)^{2m}<0}.
\]
Every non-Gaussian cumulant of order below $2m$ vanishes.
Thus no relative sign among real mode coefficients can
cancel the first sharp-cutoff signature in this independent
finite product-vacuum family. In the excluded all-zero-coefficient case
\(Q=0\), \(\phi_Q(t)\equiv1\), and all positive-order cumulants
vanish; no finite minimum \(m\) is defined.
\end{theorem}
\begin{proof}
The commuting mode observables act on independent factors,
so their characteristic functions multiply and their
cumulants add.  The preceding single-mode cutoff theorem
gives its first nonzero non-Gaussian cumulant
$-(N_a+1)!\sigma_a^{2N_a+2}$ at order $2N_a+2$.
Rescaling by $c_a$ and summing only the smallest contributing
orders gives the formula.  At least one nonzero real
coefficient realizes the minimum, so the even-power sum is
strictly positive before the minus sign.
\end{proof}

\begin{remark}[Finite quadrature is not automatically a physical measurement]
\label{p4c:rem:v349-hermite-physical-guard}
The probability law is the \emph{spectral measure of a chosen
finite operator in a specified vacuum}; it is not an assertion
that the universe has that cutoff, or that a detector measures
$q_N$.  Different caps, preparations and action-selected
observables can have different response laws.  These results
do not replace the regular CCR Gaussian theorem on its own domain.
\end{remark}
% END FRAMEWORK v2.349

\printbibliography
\end{document}
